\documentclass{situationist-preprint}
\usepackage{longtable}
\usepackage{array}
\usepackage{needspace}
\usepackage{graphicx}

\title{Commentary on On the Meaning of Life\\[0.4em]\large Volume I: Text, Commentary and Concordance}
\msid{EXI-01-b}
\msfamily{EXI-01}
\msdiscipline{Meaning, Death and the Absurd}
\mstype{Companion manuscript, Volume I of III}
\mskeywords{meaning in life; objective worth; engagement; Wolf; fitting fulfilment; availability of meaning; the absurd; Nagel; commentary; concordance}

% Generated by gen/generate.py; do not edit.
\newcommand{\cnLemmata}{40}
\newcommand{\cnSentLemmata}{12}
\newcommand{\cnPhraseLemmata}{28}
\newcommand{\cnMss}{37}
\newcommand{\cnOtherMss}{36}
\newcommand{\cnTokens}{200}
\newcommand{\cnTypes}{88}
\newcommand{\cnContent}{54}
\newcommand{\cnCorpusWords}{51396}
\newcommand{\cnKwic}{2307}
\newcommand{\cnStatements}{249}
\newcommand{\cnBearing}{83}
\newcommand{\cnNotBearing}{166}
\newcommand{\cnTestimonia}{20}
\newcommand{\cnTestimoniaMss}{11}
\newcommand{\cnAbsurd}{29}
\newcommand{\cnAssTwo}{11}
\newcommand{\cnAssThree}{11}
\newcommand{\cnThmFour}{16}
\newcommand{\cnDefOne}{14}
\newcommand{\cnUnshared}{5}
\newcommand{\cnUnsharedList}{\emph{EXI-01-b}, \emph{aeternitatis}, \emph{commentary}, \emph{specie}, \emph{sub}}
\newcommand{\cnTopKwic}{\emph{life} (450); \emph{meaning} (229); \emph{engagement} (148); \emph{agent} (104); \emph{value} (93)}
\expandafter\newcommand\csname cnFammean\endcsname{348}
\expandafter\newcommand\csname cnFamlife\endcsname{578}
\expandafter\newcommand\csname cnFamengag\endcsname{232}
\expandafter\newcommand\csname cnFamactiv\endcsname{135}
\expandafter\newcommand\csname cnFamproject\endcsname{138}
\expandafter\newcommand\csname cnFamworth\endcsname{177}
\expandafter\newcommand\csname cnFamobjectiv\endcsname{25}
\expandafter\newcommand\csname cnFamvalu\endcsname{135}
\expandafter\newcommand\csname cnFamattitud\endcsname{63}
\expandafter\newcommand\csname cnFamagent\endcsname{130}
\expandafter\newcommand\csname cnFamavailab\endcsname{83}
\expandafter\newcommand\csname cnFamdegree\endcsname{64}
\expandafter\newcommand\csname cnFammatter\endcsname{66}
\expandafter\newcommand\csname cnFamabsurd\endcsname{55}
\expandafter\newcommand\csname cnFamsceptic\endcsname{4}

\emergencystretch=3em

% A lemma of EXI-01-a, quoted exactly; gen/generate.py checks every \lem
% against the source text and against gen/lemmata.py.
\newcommand{\lem}[2]{\par\needspace{4\baselineskip}\medskip\noindent\textbf{[#1]}\quad\enquote{#2}\par\nopagebreak\smallskip}
\newcommand{\kwichead}[2]{\par\needspace{4\baselineskip}\medskip\noindent{\normalsize\textbf{#1}}\quad(#2)\par\nopagebreak\smallskip}
\newcommand{\kwicnone}{\noindent No occurrences in the other manuscripts.\par}
\newcommand{\Lm}[1]{[#1]}

\begin{document}
\maketitle

\begin{abstract}
EXI-01-a answers the question of what the meaning of life is in one page: a definition, two assumptions, a theorem with a three-sentence proof, and a remark. This Commentary explains that page sentence by sentence and, for its central sentences, phrase by phrase. Volume I divides the text into \cnSentLemmata{} sentence lemmata and \cnPhraseLemmata{} phrase lemmata and comments on each: its claim, its sources and parallels in Wolf, Metz, Nagel, Taylor, Nozick, Frankfurt, Williams, Camus, Tolstoy and Aristotle, the rival readings it admits, the objections it faces, and its relation to the rest of the collection. We argue that the text is more precise than its brevity suggests. Its notion of objective worth is restricted, requiring independence only from the attitudes of those engaged; its proof does not use Assumption~2, whose role is to fix the interpretation of the theorem; its availability claim rests on Assumption~3 together with a bridge principle that the Lean module makes explicit; and its remark on the absurd identifies exactly the premise of Nagel's challenge that the assumption of objective worth denies. The volume contains the defence of Assumption~2 and the discussion of the absurd that EXI-01-a assigns to the Commentary, a commentary on every declaration of the Lean module, and a computed apparatus: a concordance of the \cnTypes{} word types of the text across \cnOtherMss{} manuscripts of the discipline, a table of the \cnStatements{} numbered statements of those manuscripts with the lemmata they bear on, and the testimonia.
\end{abstract}

\section*{Plan of the work}

\begin{center}\small
\begin{tabular}{@{}rlr@{}}
\toprule
Vol. & Title & Pages \\
\midrule
I & Text, Commentary and Concordance & 1--114 \\
II & Objective Worth and the Absurd & 115--612 \\
III & Availability, Degree and Death & 613--1104 \\
\bottomrule
\end{tabular}
\end{center}

\paragraph{Volume I} (the present volume) gives the text of EXI-01-a divided into numbered lemmata, a commentary on every lemma and on every declaration of the Lean module that formalizes Theorem~4, and three computed appendices: a concordance of the words of the text across the manuscripts of the discipline, a table of the numbered statements of those manuscripts with the lemmata they bear on, and the testimonia, that is, every sentence of those manuscripts that cites EXI-01. It contains the defence of Assumption~2 (Section~\ref{s:defence}) and the discussion of the absurd (Section~\ref{s:absurd}) to which Remark~5 of EXI-01-a refers.

\paragraph{Volume II, Objective Worth and the Absurd,} extends the defence of Assumption~2 given in Section~\ref{s:defence}. Part~A states the theory of value that the restricted notion of objective worth requires and shows that it is supplied by the account of the good in ETH-01 and by the naturalism of MTE-01, and that the arguments from queerness, from evolutionary debunking and from disagreement do not reach it (MTE-05, MTE-09, MTE-12). Part~B examines subjectivist theories of meaning, from Taylor's Sisyphus to Frankfurt's account of caring, and the hybrid views that answer them, and argues that the dilemma usually urged against hybrids is resolved by the reading of Definition~1 given in Section~\ref{s:def}. Part~C is a history of the absurd from Ecclesiastes to Nagel, with chapters on Tolstoy's \emph{Confession}, on Camus, on Nagel's two treatments of 1971 and 1986 and on recent work on cosmic insignificance; it develops the analysis of the vindication premise given in Section~\ref{s:absurd} and relates it to EXI-04 and EXI-06. Part~D takes up theism: it relates the additional meaning that EXI-11 assigns to a theistic framework to the argument of REL-01, and shows that the conclusion of EXI-01 is independent of both.

\paragraph{Volume III, Availability, Degree and Death,} extends the commentary on Assumption~3, on the word \enquote{extent} in Definition~1 and on Remark~5. Part~A examines the scope of the availability claim: the reading of \enquote{everyone} as every agent capable of intentional action, the lives that fall outside that scope, and the hard cases of deprivation, illness and depression that EXI-01-c answers. Part~B gives a measure of the degree of meaning: it states Definition~1 with an explicit measure on the parts of a life, proves that the propositions of Volume~I hold under it, and examines the alternatives, including the shape of a life (EXI-19) and the role of achievement (EXI-17). Part~C relates meaning to death: the badness of death (EXI-02), the undesirability of immortality (EXI-03), the fittingness of fear and regret (EXI-05) and the limit that gives a life its shape (EXI-19). Part~D contains a formal appendix, the bibliography of the three volumes and the indices.

\section*{Preface}

EXI-01-a is the shortest principal manuscript of its discipline. It occupies one page and contains, under its title, eleven sentences. It is also the manuscript to which the other manuscripts of the discipline most often refer when they need an account of what makes a life meaningful: EXI-03-a uses it to restrict the categorical desires that give life its point, EXI-04-a accepts it and argues that it leaves the absurd in place, EXI-06-b uses it to show that cosmic indifference takes nothing from meaning, and EXI-11-a and EXI-19-b take its account as their starting point. A text so short and so much used is in the condition of a classical text: each phrase carries weight, and each phrase can be read in more than one way. That is the reason for a commentary.

The form of the Commentary is the old one. The text is divided into lemmata, each quoted exactly and numbered; each lemma is followed by comment on its claim, its sources, its parallels, the readings it admits, the objections it faces and its relation to other texts. Where the text is a definition or an assumption, the comment is phrase by phrase, because in a definition every phrase constrains what is defined. Where the text is a proof, the comment is step by step. Where the text refers outward, to Wolf, Metz or Nagel, or to the Commentary itself, the comment follows the reference. We have tried throughout to explain the text and not to replace it: where we propose a reading that the text does not compel, we say so and give the alternatives, and where we find the text deficient, we say what would repair it and whether the repair is consistent with the conclusion the text defends.

That conclusion is the conclusion of its family: life is meaningful to the extent that one engages with projects of objective worth, and meaning in this sense is available to everyone. The Commentary defends it. A commentary that did not would be a different kind of work, and the reader who wants the case against the conclusion will find its strongest forms stated in the objections that follow each lemma.

\paragraph{Conventions.} Lemmata of EXI-01-a are cited by number in square brackets: \Lm{2} is the sentence of Definition~1, \Lm{2.5} the phrase \enquote{active engagement} within it. The list of lemmata is given in Section~\ref{s:text}; a program (Appendix~\ref{app:conc}) checks that every quotation in the Commentary agrees with the text. Numbered statements of EXI-01-a are cited as they are numbered there (Definition~1, Assumptions~2 and~3, Theorem~4, Remark~5); numbered statements of the Commentary are cited by section (Proposition~\ref{p:redundant}). Other manuscripts are cited by identifier, and their statements by the numbers they bear in those manuscripts. Declarations of the Lean module are cited by name in typewriter type. The text of EXI-01-a uses the spellings \enquote{defence} and \enquote{sceptical}; we keep them in quotation and use American spelling elsewhere.

\section{Introduction}\label{s:intro}

\subsection{The question and the kind of answer}

The question to which EXI-01-a is addressed is the oldest question of its discipline and the one most often thought to have no answer: what is the meaning of life? It has been asked in three ways, and an answer to it must say which way it takes it.

On the first, the question asks for the point of human existence as such: what human beings, or the universe that contains them, are for. Answers of this kind assign a purpose to the species or to the whole, and they are typically theological. The individual life is meaningful, on such answers, to the extent that it plays its part in that purpose. On the second, the question asks what makes an individual life meaningful: what a life must contain, or how it must be lived, for its subject to have lived a meaningful life rather than a wasted, empty or trivial one. \citet{Metz2013} calls this the question of meaning \emph{in} life, as against the meaning \emph{of} life, and most analytic work since the 1970s addresses it. On the third, the question is conceptual: what we are asking when we ask whether a life has meaning, and what distinguishes that question from the questions whether the life was happy or whether it was morally good. \citet{Metz2001} argued that the concept is a cluster, united by a family of ideas: purposes worth pursuing, transcending one's animal nature, being an apt object of esteem or admiration, and making a narrative whole of one's life.

EXI-01-a takes the question in the second way, and its title keeps the wording of the first. That choice is examined at \Lm{1} and again at \Lm{9.1}. The answer it gives belongs to a family of answers that \citet{Metz2013} classifies as naturalist, since they do not make meaning depend on a spiritual realm, and within that family to the subclass of hybrid views, which make meaning depend both on the subject's attitudes and on something that does not depend on them. Its source, acknowledged in the text, is the fitting-fulfilment view of \citet{Wolf1997,Wolf2010}, on which meaning arises when the subject's active engagement meets an object worthy of it. Its distinctive additions are three: it makes the degree of meaning explicit, it isolates the assumption that there are objects of the required kind, and it claims that meaning so understood is available to every agent.

\subsection{The debate before EXI-01-a}

A commentary must place its text in the debate from which it comes. We give the history in outline here; the lemmata at which each strand becomes relevant return to it in detail.

\paragraph{Aristotle.} The \emph{Nicomachean Ethics} does not ask what the meaning of life is, but it asks the closest ancient question: what the chief good of a human life is, the end for whose sake everything else is done. Its answer, that the human good is activity of the soul in accordance with excellence over a complete life (\emph{NE} I.7, 1097b22--1098a20), contains three ideas that recur in EXI-01-a: that the good of a life lies in an activity (\emph{energeia}) rather than a state or a possession, that the activity must be of something worth doing, since the excellence in question is the excellence of the function proper to human beings, and that the assessment is of a life as a whole and not of a moment (\emph{NE} I.7, 1098a18--20; I.10, 1100a10--1101a21). The first idea is the ancestor of \enquote{active engagement} in \Lm{2.5}, the second of \enquote{objective worth} in \Lm{2.7}, and the third of \enquote{a life} in \Lm{2.1}. Aristotle also shows how far apart the ancient and modern questions are: his question concerns happiness (\emph{eudaimonia}), and EXI-07 argues that happiness is not the goal of life and is best understood as a by-product of engagement in worthwhile activities. EXI-01-a separates meaning from happiness in a way Aristotle does not.

\paragraph{Tolstoy and Camus.} The modern question is the question of the crisis. Tolstoy, in the \emph{Confession} \citep{Tolstoy1987}, describes how, at the height of his success, his life came to a halt before questions he could not answer: what would come of what he was doing, and whether there was any meaning in his life that would not be destroyed by the death that awaited him. He compared his condition to that of a traveller in an Eastern fable who, fleeing a beast, hangs from a branch over a well in which a dragon waits, while two mice gnaw at the branch, and who licks drops of honey from the leaves. The honey, the goods of his life, no longer tasted sweet, because he could see the dragon and the mice. \citet{Camus1955} opened \emph{The Myth of Sisyphus} with the claim that whether life is worth living is the one serious philosophical problem, and located the absurd in the confrontation between the human demand for meaning and the silence of the world. Both texts present the question of meaning as a question that arises when the subject steps back from her life; both take the answer to depend on whether something outside the life could vindicate it. EXI-01-a's Remark~5 engages that structure directly, and Section~\ref{s:absurd} examines it.

\paragraph{Taylor, Nagel and Wiggins.} The analytic discussion begins with three texts of the 1970s. \citet[ch.~18]{Taylor1970} took the myth of Sisyphus as his model of a meaningless existence and asked what would make it meaningful. Giving Sisyphus a purpose outside his labor, a temple to be built from the stones, would not do, since the completed temple would only be a scene of boredom; implanting in him a desire to roll stones would, since he would then be doing what he most wanted to do. Taylor concluded that meaning comes from within, from the will to do what one does. \citet{Nagel1971} argued that the sense of the absurd arises from a collision within the reflective subject between the seriousness with which we take our lives and the perpetual possibility of regarding everything about which we are serious as arbitrary, and that the appropriate response is irony. \citet{Wiggins1976} argued against Taylor that a purely subjective account of meaning cannot be lived from the inside, since the agent who finds her life meaningful takes its objects to be worthy of her concern, and that the values on which meaning depends are neither simply invented nor simply discovered.

\paragraph{Nozick.} \citet[ch.~6]{Nozick1981} connected meaning with transcending limits: a life acquires meaning by connecting with value beyond itself, and the question of the meaning of that larger value recurs until it ends in something unlimited. He also gave, in the experience machine \citep[42--45]{Nozick1974}, the case that ETH-05 uses to show that well-being includes contact with reality and that EXI-17-a uses to show that meaning requires engagement with a real object (its Principle~P2).

\paragraph{Wolf.} \citet{Wolf1997} argued that meaning is a value distinct from happiness and from morality, and that a life is meaningful insofar as its subject is actively engaged in projects of worth. Her Tanner Lectures \citep{Wolf2010} developed the view, defended it against subjectivism and against purely objective accounts, and argued that the objective element should be understood as worth that does not derive from the subject's attitudes, while declining to give a general theory of that worth. They were published with comments by John Koethe, Robert M. Adams, Nomy Arpaly and Jonathan Haidt, and Wolf's replies. Wolf's view is the source of Definition~1.

\paragraph{Metz.} \citet{Metz2013} surveyed and classified the theories, argued against supernaturalism and subjectivism, and defended an objective naturalist theory, the fundamentality theory, on which a life is more meaningful the more its subject positively orients her rational nature toward conditions that are responsible for much else about human life. Metz is cited in EXI-01-a for the claim that there are projects of objective worth (\Lm{3.5}).

\paragraph{Later work.} The last fifteen years have added work on the temporal structure of meaning \citep{Kauppinen2012}, on the relation between meaning and consequences \citep{Smuts2013}, on cosmic insignificance \citep{Kahane2014}, on the pessimist's case that meaning is limited at best \citep{Benatar2017}, on the dependence of the meaning of our lives on the survival of humanity after us \citep{Scheffler2013}, on meaning in imperfect lives \citep{Landau2017} and on the variety of meaningful lives \citep{May2015}. Each enters the Commentary at the lemma to which it is relevant.

\subsection{The form of EXI-01-a}

EXI-01-a is written as a small axiomatic theory. It contains a definition (Definition~1), two assumptions (Assumptions~2 and~3), a theorem (Theorem~4) with a proof, and a remark (Remark~5) that says what the theorem answers and what it does not. The numbering is continuous from 1 to 5 and is not divided by section, as the numbering of longer manuscripts of the discipline is. Three features of the form matter for interpretation.

First, the form separates what is stipulated from what is asserted. Definition~1 is not true or false in the way an assumption is; it says what the text means by \enquote{meaningful to degree $d$}. Whether it captures what we ordinarily mean by a meaningful life is a question about the definition, and the definition is cited to Wolf because Wolf argued that it does. Assumptions~2 and~3 are assertions, and the text does not argue for them: the first is cited to Metz, and the second is left to its consequences and to the other manuscripts of its family (EXI-01-c, EXI-17-a). The theorem follows from the definition and the assumptions by a proof that occupies three sentences. Everything philosophically contested is therefore located in one of four places: the adequacy of the definition, the truth of either assumption, or the bridge between the assumptions and the theorem. The Commentary takes these in order.

Second, the form makes the text unusually easy to formalize, and it is formalized: Theorem~4 is proved in a Lean module whose axioms are listed in Section~\ref{s:lean}. The formalization does what formalization does in this discipline: it shows that the conclusion follows, and it shows exactly which premises it follows from. The second service is the more interesting here. It shows that Assumption~2 is not among the premises, and that a premise not stated in the text is (Propositions~\ref{p:redundant} and~\ref{p:bridge}).

Third, the brevity of the form leaves much unsaid that a longer text would say: the defence of the assumptions, the replies to objections, the relation to rival views. Remark~5 acknowledges this and refers the reader to the Commentary for two of the missing parts. The other missing parts are supplied in part by the other manuscripts of the family. EXI-01-c gives a second proof of Theorem~4 that does not use Assumption~3. EXI-01-d draws out what the availability claim implies and what it does not. The Commentary draws on both.

\subsection{The text in its discipline}

The discipline to which EXI-01-a belongs has twenty families, and most of them bear on it. We list the relations here so that later comments can refer to them briefly; Appendix~\ref{app:xref} gives the full table, computed from the numbered statements of every manuscript.

EXI-17 holds that a meaningful life requires active engagement with something worthwhile and that achievement is not necessary. Its principal, EXI-17-a, defines active engagement (Definition~2.1) and worth (Definition~2.2) in terms that agree with EXI-01-a, states the fitting-fulfilment view as its Principle~P1, adds a principle of contact with reality (Principle~P2), and derives that meaning is not hostage to luck in outcomes (Corollary~3.6), which supports the availability claim of EXI-01. It is the manuscript closest to EXI-01-a, and the two should be read together.

EXI-04 holds that life is absurd in Nagel's sense but that the appropriate response is irony rather than despair. EXI-04-a accepts the conclusion of EXI-01 and argues that the absurd survives it (its Objection~2 and reply), and it cites Remark~5 of EXI-01-a. EXI-04-b reconsiders Camus's Sisyphus. EXI-06 holds that the universe is indifferent and that indifference is not hostility, and EXI-06-b argues that this indifference takes nothing from meaning as EXI-01 understands it. EXI-11 holds that meaning does not require God, though a theistic framework would add to it; its principal defines meaning as EXI-01-a does (its Definition~2.1) and states as a principle that worth is not conferred by assignment (its Principle~P1). EXI-03 holds that an endless life would exhaust the categorical desires that give life its point, and EXI-03-a identifies those desires, in line with EXI-01, with desires for projects the person takes to be worthwhile. EXI-19 holds that death has no meaning in itself but gives life a shape by limiting it, and EXI-19-b draws out what this implies for meaning as EXI-01 and EXI-17 understand it. EXI-07 holds that happiness is a by-product of other pursuits; EXI-09 that existence precedes essence for human beings; EXI-20 that whether producing a body of work was meaningful depends on whether it was engagement with something worthwhile. Each of these is relevant at one or more lemmata.

Outside the discipline, the text depends on ETH-01, which holds that the good is what idealized reflective endorsement converges on and that it has seven irreducible components, and on the metaethical families MTE-01, MTE-05, MTE-09 and MTE-12, which hold that there are moral facts and that the arguments from queerness, evolution and disagreement do not show otherwise. These bear on Assumption~2. It touches ETH-05, which holds that well-being includes contact with reality, and REL-01, which holds that a maximally great being exists; the second bears on the additional meaning that EXI-11 assigns to a theistic framework, and so on the degree but not the availability of meaning.

\subsection{Note on the text}

The text of EXI-01-a consists of a title and eleven sentences, which contain \cnTokens{} words when citations are not counted, followed by a bibliography of three items. Its statements are introduced by bold labels (\enquote{Definition~1.}, \enquote{Assumption~2.} and so on) rather than by the numbered environments of the longer manuscripts, and the proof of Theorem~4 is introduced by an italic label and closed by a square. The Lean module refers to the statements by the same numbers: its docstrings name \enquote{Definition~1, positive case}, \enquote{Assumption~3} and \enquote{Theorem~4}, and they agree with the text. The text uses two British spellings, \enquote{defence} and \enquote{sceptical}, against the American spelling of the rest of the discipline; we quote them as they stand. The last sentence of Remark~5 refers to the present work by its identifier. No variant of the text is known to us.

\subsection{The commentary form and its precedent}

The commentary on an axiomatic text has a long precedent, and the precedent illuminates EXI-01-a. Proclus's commentary on the first book of Euclid's \emph{Elements} \citep{Proclus1970} comments on the definitions, the postulates, the common notions and the propositions in order, and distinguishes their roles: a definition says what is meant, a postulate demands something to be granted that is neither self-evident nor proved, a common notion states what is evident to all, and a proposition is demonstrated. Much of what he says about the postulates concerns their status. He held that the fifth postulate, on parallels, did not belong among the postulates, since it was not evident and its converse was proved as a proposition, and that it should be proved as a theorem; and he offered a proof, which later readers found to rely on an assumption equivalent to the postulate itself.

EXI-01-a has the same parts in miniature. Definition~1 says what is meant. Assumption~2 is closer to a common notion in the intention of its author, who cites rather than argues for it, and the Commentary has argued that its role is interpretive rather than demonstrative (Proposition~\ref{p:redundant}). Assumption~3 is a postulate in Proclus's sense: it carries the weight of the theorem, it is not self-evident, and it is the premise a critic will press. The family's response to it is Proclus's response to the fifth postulate. EXI-01-c proves it as a theorem from a premise thought to be better supported, its Assumption~2.2 (Reach), and anticipates the objection that Proclus's proof met, that the new premise is the old one in other words; its reply is that the new premise is supported by the ubiquity of persons and of the world as an object of understanding, and that the cases that threatened Assumption~3 satisfy it. Whether the reply succeeds is discussed at \Lm{4}. The parallel shows at least that the question it raises is the right one: in an axiomatic text, the dialectic concentrates on the premise that is neither definitional nor evident, and the commentary's task is to say whether it can be derived, from what, and at what cost.

\subsection{Method of this volume}

The text is divided in Section~\ref{s:text} into \cnSentLemmata{} sentence lemmata, numbered \Lm{1} (the title) to \Lm{12}, and \cnPhraseLemmata{} phrase lemmata, numbered by sentence and phrase, for the sentences that carry the argument: the definition, the two assumptions, the theorem, the central step of the proof and the remark on the absurd. Every word of each of those sentences belongs to at least one phrase lemma, except connectives whose role is discussed under the sentence lemma.

Sections~\ref{s:title} to~\ref{s:remark} comment on the lemmata in order. Each section opens with the sentence lemma, comments on the sentence as a whole, and then takes its phrases in turn. Where a phrase raises a question that the Commentary can settle, we state the result as a numbered proposition and prove it; where it raises a question that belongs to a later volume, we say so. Section~\ref{s:bib} comments on the bibliography of EXI-01-a, Section~\ref{s:lean} on the Lean module, declaration by declaration, and Section~\ref{s:whole} on the argument as a whole: the readings it admits, the objections to it that are not objections to any one lemma, and what it does not establish.

The appendices are computed. Appendix~\ref{app:conc} is a concordance: every word type of EXI-01-a, with the lemmata in which it occurs and its occurrences in each of the other \cnOtherMss{} manuscripts of the discipline, a table of fifteen term families by manuscript, and a keyword-in-context listing of every occurrence of the content words of the text in those manuscripts. Appendix~\ref{app:xref} lists every numbered statement of those manuscripts, read from their compiled \texttt{.aux} files, with its concept classes and the lemmata of EXI-01-a it bears on by a stated rule. Appendix~\ref{app:test} gives the testimonia. Each appendix has a preface stating what was computed and how.

\section{The text}\label{s:text}

The text of EXI-01-a is given below divided into sentence lemmata. The phrase lemmata are listed at the head of the comment on each sentence. Citations are reproduced as they appear in the text, and the bold labels of the statements are given in the first column. The table is generated from the list of lemmata used by the program of Appendix~\ref{app:conc}, which checks every entry against the source.

{\small% Generated by gen/generate.py; do not edit.
\begin{longtable}{@{}p{0.13\textwidth}p{0.07\textwidth}p{0.74\textwidth}@{}}
\toprule Statement & Lemma & Text \\ \midrule \endhead
\bottomrule \endfoot
Title & [1] & On the Meaning of Life \\[2pt]
Definition 1 & [2] & A life $L$ is \emph{meaningful to degree $d$} if $d$ is the extent to which $L$ consists in active engagement with projects of objective worth \citep{Wolf2010}. \\[2pt]
Assumption 2 & [3] & There exist projects of objective worth: projects whose value does not depend on the attitudes of those engaged in them \citep{Metz2013}. \\[2pt]
Assumption 3 & [4] & For every agent capable of intentional action, there is at least one project of objective worth in which she is able to engage. \\[2pt]
Theorem 4 & [5] & For every such agent, a life of positive meaning is available. \\[2pt]
Proof & [6] & By Assumption 3, there is a project $p$ of objective worth in which the agent can engage. \\[2pt]
 & [7] & A life that includes engagement with $p$ consists, to some positive extent, in active engagement with a project of objective worth. \\[2pt]
 & [8] & By Definition 1, it is meaningful to a positive degree. \\[2pt]
Remark 5 & [9] & The theorem answers the question of what the meaning of life is, and of whether meaning is available to everyone. \\[2pt]
 & [10] & It does not answer the question of whether any particular life is meaningful, which depends on what is done with it. \\[2pt]
 & [11] & Nor does it answer the sceptical challenge that, \emph{sub specie aeternitatis}, nothing matters \citep{Nagel1971}: Assumption 2 denies the premise of that challenge rather than refuting it. \\[2pt]
 & [12] & A defence of Assumption 2, and a discussion of the absurd, are given in the commentary (EXI-01-b). \\[2pt]
\end{longtable}
}

The bibliography of EXI-01-a, on which Section~\ref{s:bib} comments, lists \citet{Metz2013}, \citet{Nagel1971} and \citet{Wolf2010}.

\section{The title}\label{s:title}

\lem{1}{On the Meaning of Life}

\paragraph{The form of the title.} The title has the form of the classical treatise title, \emph{De} with the ablative: \emph{De brevitate vitae}, \emph{De senectute}, \emph{De finibus bonorum et malorum}. Such a title announces a subject rather than a thesis, and the text that follows it is expected to say what is true about the subject. The catalogue question to which EXI-01-a is the principal answer is phrased differently: \enquote{What is the meaning of life?} The title takes the subject of that question and drops its interrogative form; the text supplies the answer in its definition, and the remark at its end says that the answer has been given (\Lm{9}).

\paragraph{\enquote{the Meaning}.} The definite article presupposes that there is one meaning of life to be stated. On the first of the three readings distinguished in Section~\ref{s:intro}, the presupposition is substantive and contested: it presupposes that human life as such has a point, and the text would owe an account of what that point is. On the second reading, the presupposition is harmless: \enquote{the meaning of life} denotes the property in virtue of which lives are meaningful, and the text says what that property is. On the third, \enquote{the meaning} denotes the content of the concept, and the text would owe an analysis of what we mean.

The text takes the second reading, and its definition shows this. Definition~1 does not say what human life is for. It says what a life's being meaningful to a degree consists in. A theory that answered the first question would have the form \enquote{human life exists in order to $X$}, and the meaning of an individual life would derive from its contribution to $X$. Definition~1 has no such form: the projects of objective worth with which a life may be engaged are many, they need not be related to one another, and nothing in the definition requires that they serve a single end. The text therefore uses the definite article in the way in which one speaks of \emph{the} nature of knowledge or \emph{the} source of normativity: as naming whatever it is that the theory identifies, without presupposing that it is a purpose.

The second reading has a cost that the text accepts. Many people who ask what the meaning of life is are asking the first question, and an answer to the second may seem to change the subject. \citet{Metz2013} records the distinction and treats the second question as the one that analytic philosophy can address; \citet{Wolf2010} addresses the question of meaning in individual lives and does not make her answer depend on whether human life as such has a point. EXI-01-a follows both. It does not deny that the first question has a sense. It holds, and EXI-11 argues, that the answer to the second does not depend on the answer to the first: a life can be meaningful though human life as such has no cosmic purpose (Theorem~3.3 of EXI-11-a). The first question is not dismissed but separated, and EXI-06 and EXI-11 take it up: the universe does not care about us (EXI-06-a), and a cosmic purpose, if there is one, would add to the meaning available to a life without being required for it (Proposition~4.1 of EXI-11-a).

\paragraph{\enquote{of Life}.} \enquote{Life} in the title is used without an article, as a mass term, and Definition~1 replaces it with a count noun: \enquote{A life $L$}. The shift is significant. It locates the bearer of meaning in individual lives rather than in life as such, or in humanity, or in the biological phenomenon. On the first reading of the title, the bearer would be human life as a whole, and individual lives would be meaningful derivatively. On the reading the text adopts, individual lives are the primary bearers, and \enquote{the meaning of life} names a property that individual lives have to various degrees. The comment on \Lm{2.1} examines what a life is for the purposes of the definition.

\paragraph{The subject of the discipline.} The heading under which EXI-01-a stands, \enquote{Meaning, Death and the Absurd}, groups the question of meaning with two others. The grouping reflects a historical connection: in Tolstoy and in Camus, the question of meaning arises from the thought of death and is answered, or left unanswered, under the heading of the absurd. EXI-01-a severs the first connection and reshapes the second. It does not mention death, and Definition~1 makes the meaning of a life depend on what the life consists in rather than on whether it ends; EXI-19 adds that death gives a life its shape, which bears on one kind of meaning, the meaning of a life as a whole, without being a source of meaning (Proposition~1.1 of EXI-19-b). It mentions the absurd only in Remark~5, and there it treats the absurd as a challenge to one of its premises rather than as a feature of the condition its definition describes (\Lm{11}).

\paragraph{Parallels.} The title is shared, nearly or exactly, with several works in the literature: \citet{Cottingham2003}, \emph{On the Meaning of Life}, a defence of the view that meaning requires a religious framework of hope; Taylor's chapter \enquote{The Meaning of Life} \citep[ch.~18]{Taylor1970}; and the anthologies that collect the debate under that heading. EXI-01-a's answer is opposed to Cottingham's, since it makes meaning available without any religious framework, and it agrees with EXI-11 that such a framework would add to meaning without being required for it. Its answer is opposed to Taylor's, since it requires that the objects of engagement have worth that the engagement does not confer.

\paragraph{Rival readings.} Two readings of the title should be recorded and set aside. On a deflationary reading, the title is ironic: it promises an answer to the great question and delivers a definition, and the brevity of the text signals that the great question has no answer worth more than a page. Nothing in the text supports this reading. Remark~5 says plainly that the theorem answers the question of what the meaning of life is, and the register of the text is that of the rest of the discipline. On an ambitious reading, the title claims that the text answers all three questions. Remark~5 also refutes this reading, since it states two questions that the theorem does not answer (\Lm{10}, \Lm{11}). The correct reading is the one the text itself gives: the title names the subject; Definition~1 says what the meaning of life, in the sense of the second question, consists in; Theorem~4 says that it is available to every agent; and Remark~5 marks the boundaries of both claims.

\section{Definition 1}\label{s:def}

\lem{2}{A life $L$ is \emph{meaningful to degree $d$} if $d$ is the extent to which $L$ consists in active engagement with projects of objective worth \citep{Wolf2010}.}

\paragraph{The sentence as a whole.} Definition~1 defines a two-place predicate, \enquote{meaningful to degree $d$}, true of a life and a degree. The definiendum is italicized, as is customary, and the definiens follows \enquote{if}. The definiens has three components: a bearer (a life), a measure (the extent to which the life consists in something) and a condition (active engagement with projects of objective worth). Each is the subject of a phrase lemma below.

Three features of the sentence as a whole should be noted before the phrases are examined. First, the definition mentions no subject. It does not say that a life is meaningful to the extent that \emph{its subject} is engaged, as EXI-17-a (Principle~P1) and EXI-11-a (Definition~2.1) do; it says that a life is meaningful to the extent that it \emph{consists in} engagement. The subject is implied, since engagement is someone's engagement and a life is someone's life, but the grammatical subject of the condition is the life. The choice is deliberate and has a consequence examined at \Lm{2.4}: it makes meaning a matter of what the life is made of, not of what the subject brings about or of how she regards her life.

Second, the definition is a real definition, not a stipulation. It is cited to \citet{Wolf2010}, whose lectures argue that the concept we use when we call lives meaningful or meaningless is the concept of a life engaged with worthwhile things. A stipulative definition would need no citation and could not be wrong. Definition~1 can be wrong: if there were lives that reflection judged meaningful and that contained no active engagement with anything of objective worth, or lives so engaged that reflection judged meaningless, the definition would fail. The comment on \Lm{2.7} considers the cases that have been offered.

Third, the definition is general: it applies to every life, actual or possible. This matters for Theorem~4, which concerns the lives available to an agent and not only the life she actually leads. A definition that applied only to actual lives could not support a claim about availability, which is a claim about possible lives (\Lm{5.3}).

\paragraph{Parallels in the collection.} The definition is stated three times in the discipline, in slightly different words. EXI-17-a states it as a principle: \enquote{A life is meaningful to the extent that its subject is actively engaged with what is worthwhile} (its Principle~P1). EXI-11-a states it as a definition with the same words (its Definition~2.1). EXI-01-c restates Definition~1 and takes it over without change. The differences between EXI-01-a and the later statements are three: the later statements mention the subject, they replace \enquote{projects} by \enquote{what is worthwhile}, and they drop the explicit degree variable while keeping \enquote{to the extent that}. The second difference is the most important, and the comment on \Lm{2.6} argues that it is a clarification of EXI-01-a rather than a departure from it. The concordance (Appendix~\ref{app:conc}, Table~A.2) shows how the vocabulary of the definition is distributed across the discipline: the families of \emph{engage-} and \emph{worth-} occur \cnFamengag{} and \cnFamworth{} times in the other manuscripts, against \cnFamproject{} occurrences of \emph{project}, which is the measure of how far the later manuscripts have moved from the word \enquote{project} to the more general vocabulary of objects of engagement.

\lem{2.1}{A life $L$}

The bearer of meaning is a life. Three alternatives are excluded by this choice.

The first alternative is the person. One speaks of meaningful lives, not of meaningful persons, and the usage reflects a real distinction: a person's character or worth is assessed at a time, while a life is assessed as a temporally extended course. A person whose life has been empty may become someone whose life is meaningful; we do not say that she has become a meaningful person. Definition~1 follows usage. The second alternative is the action or the episode. Actions and episodes can be meaningful in a sense: a gesture can be meaningful, and a conversation can be the most meaningful hour of a week. Definition~1 does not deny this, but it does not define it. It defines the meaning of a life, and the meaning of the parts of a life enters only through the measure (\Lm{2.3}). The third alternative is humanity or life as such, which the comment on \Lm{1} has set aside.

\paragraph{Whole and part.} Metz distinguishes the meaning of a life as a whole from the meaning in a life, the meaning of its parts, and asks whether the first is the sum of the second or depends also on the relations among the parts \citep{Metz2013}. \citet{Kauppinen2012} argues that the meaning of a life depends on its narrative structure, on whether its later parts build on and redeem its earlier ones. \citet{MacIntyre1981} had earlier argued that a human life has the unity of a narrative and that its goods are intelligible only within that unity. Definition~1 takes a position on this question only through the phrase \enquote{the extent to which $L$ consists in}: the degree of meaning of the whole is the extent to which the whole consists in a certain kind of activity. On its most natural reading, this makes whole-life meaning depend on the meaning of the parts, without saying how. It is compatible with the view that the parts contribute additively and with the view that their arrangement matters. EXI-19 holds that the shape of a life, the order of its goods and evils, bears on its value as a whole (Principle~P1 of EXI-19-a), and EXI-19-b argues that finitude is a condition of the meaning of a life as a whole (its Proposition~1.2). These claims concern a kind of meaning that Definition~1 neither includes nor excludes; Volume~III examines whether they require a revision of the measure.

\paragraph{Lives as possibilities.} The variable $L$ ranges over lives, and Theorem~4 quantifies over lives available to an agent. A life in this sense is a possible course of an agent's existence, individuated by its subject and its content. The Lean module makes this explicit: \texttt{Life} is a family of types indexed by agents, so that each agent has her own space of possible lives (Section~\ref{s:lean}). Two features of this conception should be noted. A life is someone's life; there are no lives without subjects. And a life is a whole; the life of an agent who dies young is a complete life, not a fragment of the longer life she might have had.

\lem{2.2}{is \emph{meaningful to degree $d$}}

\paragraph{Gradualism.} The definiendum makes meaning a matter of degree. This is the common view in the literature: \citet{Wolf1997} speaks of lives as meaningful insofar as they are engaged with worthwhile projects, and \citet{Metz2013} treats meaning as a gradable property throughout. The alternative, that meaning is all or nothing, has few defenders, and Definition~1 does not argue against it. It is nevertheless worth stating why gradualism is right. Reflection compares lives in respect of meaning: the life of someone who devoted thirty years to the care of her disabled brother and the life of someone who gave him an afternoon a month are both meaningful in respect of that care, and the first more so. A binary concept could not register the comparison.

\paragraph{\enquote{Meaningful} simpliciter.} The definition defines \enquote{meaningful to degree $d$} and leaves \enquote{meaningful} without qualification undefined. The text uses \enquote{positive meaning} in Theorem~4 and \enquote{meaningful to a positive degree} in its proof, and Remark~5 uses \enquote{meaningful} without qualification (\Lm{10}). The natural reading is that a life is meaningful simpliciter if and only if it is meaningful to some positive degree. Ordinary usage is more demanding: to call a life \enquote{meaningful} is usually to say that it has a good deal of meaning, as to call a person \enquote{tall} is to say that she is taller than most. On the ordinary usage, Theorem~4 would claim much less than \enquote{a meaningful life is available}, and the family conclusion that meaning is available to everyone would be weaker than it sounds. The text avoids the ambiguity by stating Theorem~4 in terms of positive meaning, and the reader should hold it to that. EXI-01-d makes the point explicit in its reply to the objection that availability so secured is trivial: availability is not possession, and possession comes in degrees.

\paragraph{The range of $d$.} The text does not say what the degrees are. Three possibilities should be distinguished. On the cardinal reading, $d$ is a real number, and differences between degrees are meaningful: one life can be twice as meaningful as another. On the ordinal reading, degrees are ranks in a complete ordering of lives. On the partial reading, degrees are elements of a partial order, so that some pairs of lives are incomparable in meaning. The text is neutral. It uses only one feature of the degrees: that there is a distinguished degree, zero or none, and that degrees above it are positive. Every one of the three readings supplies this.

We think the partial reading is the most plausible, and we record the reason, though nothing in Volume~I depends on it. The meaning of a life spent in scientific research and the meaning of a life spent in the care of others are both substantial, and reflection does not rank them; it does not judge them equal either, since a small improvement in one would not make it more meaningful than the other. This is the mark of incomparability in the sense discussed by \citet{Raz1986}. A cardinal or complete ordinal measure would force a ranking that reflection does not make.

\begin{proposition}[No negative degrees]\label{p:noneg}
On Definition~1, no life is meaningful to a negative degree.
\end{proposition}

\begin{proof}
The degree of a life is the extent to which it consists in active engagement with projects of objective worth. An extent is a measure of how much of something there is, and is not negative. A life that consists in no such engagement is meaningful to the degree of zero, whatever else it consists in.
\end{proof}

Proposition~\ref{p:noneg} has a consequence that some will find objectionable. A life devoted to cruelty, or to the destruction of what is valuable, is on Definition~1 no less meaningful than a life devoted to nothing: both have zero meaning in respect of those activities. One might think that the destroyer's life has negative meaning, that it is not merely empty but anti-meaningful. Definition~1 does not allow this, and we think it right not to. The destroyer's life is bad, and morally worse than the idle life, but the respect in which it is worse is a moral respect. Wolf's argument that meaning is a value distinct from morality \citep{Wolf1997} implies that a life can be bad without being anti-meaningful, and Definition~1 keeps the distinction clean. A theory that allowed negative meaning would have to say what it is about destruction, as opposed to its moral wrongness, that subtracts meaning, and no account of this that does not collapse into moral assessment has been given.

\lem{2.3}{if $d$ is the extent to which}

\paragraph{\enquote{if}.} As is usual in definitions, \enquote{if} is to be read as \enquote{if and only if}. The proof of Theorem~4 uses the definition only in the direction from right to left (from engagement to meaning), and the Lean module formalizes only that direction (Section~\ref{s:lean}). But the definition is a definition, and Remark~5 uses the other direction when it says that a particular life may fail to be meaningful because of what is done with it (\Lm{10}): a life with no engagement with projects of objective worth has no meaning, and that is the left-to-right direction.

\paragraph{\enquote{the extent}.} The definite article implies that each life has exactly one degree of meaning. On the partial reading of the degrees, this is unproblematic: a life has one place in a partial order. What is problematic is the measure. \enquote{The extent to which $L$ consists in} engagement admits two readings that diverge in important cases.

On the \emph{proportional} reading, the extent is the share of the life that consists in engagement: a life of which half is spent in engagement with worthwhile projects has a higher degree than a life of which a quarter is. On the \emph{absolute} reading, the extent is the amount of engagement the life contains: a life with thirty years of engagement has a higher degree than a life with twenty, whatever the length of each.

The readings diverge in cases of the following kind. Two people spend thirty years in devoted engagement with worthwhile work. One dies at sixty; the other lives to ninety, and spends the last thirty years in idleness. On the proportional reading, the first life is more meaningful, since a larger share of it consists in engagement. On the absolute reading, the two lives are equally meaningful, since they contain the same engagement. Reflection is uncertain here, and its uncertainty is instructive. There is something to the thought that the idle years diminish the second life, and something to the thought that they subtract nothing from what it contained.

\begin{proposition}[Measure neutrality]\label{p:neutral}
Theorem~4 holds on every reading of \enquote{the extent to which $L$ consists in} on which a life that contains active engagement with a project of objective worth, sustained over some non-null part of the life, consists in such engagement to a positive extent.
\end{proposition}

\begin{proof}
The proof of Theorem~4 uses Definition~1 only to infer positive meaning from positive extent (\Lm{8}), and the inference from inclusion of engagement to positive extent (\Lm{7.2}). Both the proportional and the absolute readings satisfy the stated condition: a non-null part of a life is a positive share of it, and a positive amount of engagement. Any reading that satisfies the condition supports the same inference.
\end{proof}

Proposition~\ref{p:neutral} explains why the text can leave the measure unspecified: the availability claim does not depend on it. The measure matters for comparisons of degree, and for the claims of EXI-01-d about what affects degree (its Proposition~2.4). Volume~III states a measure and examines its consequences. One feature of the measure can be fixed now: it cannot be purely temporal. Engagement differs in depth as well as in duration, and the worth of its objects differs too, and EXI-11-a states as a principle that engagement with what is of greater worth makes a life more meaningful, other things equal (its Principle~P2). The extent to which a life consists in engagement with projects of objective worth is therefore at least a function of three variables: how much of the life is so engaged, how active the engagement is, and how great the worth of its objects. Definition~1 compresses these into one word. The compression is legitimate in a definition that is meant to say what meaning consists in, and it would be a defect in a theory that claimed to measure it.

\lem{2.4}{$L$ consists in}

\paragraph{Constitution, not causation.} \enquote{Consists in} is the language of constitution. On Definition~1, the meaning of a life is not an effect of engagement, as the satisfactions of a successful career are an effect of the work; it is the engagement, measured. This sets Definition~1 against two families of rival views.

The first family makes meaning depend on consequences. On the good cause account of \citet{Smuts2013}, a life is meaningful to the extent that it causes good outcomes, whether or not its subject is engaged with them and whether or not she knows of them. On the transcendence view of \citet{Nozick1981}, a life gains meaning by making a difference to, and leaving traces in, something of value beyond itself. Both views make the meaning of a life depend on facts outside the life: on what happens as a result of it. Definition~1 makes it depend on what the life consists in.

The difference shows in two kinds of case. In the first, a life produces great goods without engagement: a person invents, idly and as a joke, a device that later saves many lives. On the good cause account, her life is very meaningful. On Definition~1, the invention adds to the meaning of her life only to the extent that her life consisted in engagement with it, which by hypothesis it did not. We think Definition~1 gives the right verdict: her life was fortunate in its effects, and the world owes her gratitude, but her life was not given meaning by an outcome she did not care about. In the second, a life is engaged with a worthwhile project that produces nothing: the unsuccessful mathematician of EXI-17-a, who spends her career on a problem she never solves (its Lemma~3.1). On the good cause account, her life has little meaning in respect of that work. On Definition~1, it has much. EXI-17-a argues that reflection sides with Definition~1, and the argument depends on the constitutive reading of \enquote{consists in}.

\paragraph{Posthumous events.} A consequence of the constitutive reading is that the meaning of a life is fixed by its course and cannot be altered after the life has ended. A painter whose work is discovered after her death did not live a more meaningful life than she would have if the work had been lost; her life consisted in the same engagement in either case. This must be distinguished from the claim of EXI-13 that the dead can be harmed: a setback to a surviving interest is bad for the ante-mortem person (Principle~P1 of EXI-13-a), and the painter's interest in being recognized can be thwarted or satisfied after her death. Harm and meaning come apart here, as they do elsewhere. The painter's life went better for her, perhaps, if her work was recognized; it was not more meaningful.

\paragraph{Meaning without awareness.} A further consequence is that a subject can be mistaken about the meaning of her own life. Whether her life consists in active engagement with projects of objective worth is a fact about her activity and about the worth of its objects, and she can be wrong about the second. Someone who believes that her life is meaningful because she is passionately engaged in a project that has no worth is mistaken; someone who believes that her life is meaningless, because she does not see the worth of the care she gives, is mistaken in the other direction. The possibility of this second kind of mistake is part of what gives the question of meaning its urgency, and it is discussed further under Assumption~2 (Section~\ref{s:defence}).

\paragraph{Meaning and happiness.} EXI-07 holds that happiness is a by-product of other pursuits, so that aiming directly at it is self-defeating; EXI-07-a argues that a large part of happiness consists in or depends on wholehearted engagement in activities and relationships whose objects are valued for their own sake (its Principle~P1). Meaning, on Definition~1, consists in engagement of a certain kind. The two claims fit: engagement with worthwhile objects constitutes meaning and tends to cause happiness. They also explain why meaning and happiness come apart. A life of engagement with worthwhile but painful work, the care of the dying for instance, may be meaningful and unhappy; a life of pleasant engagement with trivial pursuits may be happy and nearly meaningless. Wolf's case for treating meaning as a distinct dimension of the good life rests on such pairs \citep{Wolf1997}.

\lem{2.5}{active engagement}

\paragraph{Wolf's notion.} \enquote{Active engagement} is Wolf's phrase. She describes the engaged subject as gripped, excited, involved, absorbed in what she is doing, as against the subject who goes through the motions, and she contrasts it with the passive enjoyment of comfort or diversion \citep{Wolf2010}. EXI-17-a makes the notion precise: an agent is actively engaged with an object if she acts in ways directed at it, cares how it fares, and is gripped or absorbed by it, so that her activity expresses her concern rather than being merely performed (its Definition~2.1). The definition has three clauses, a practical clause (directed action), a conative clause (care) and an affective clause (absorption), and the commentary on \Lm{2.5} must say how they are related.

\paragraph{Two readings.} On the \emph{affective} reading, engagement is fundamentally a felt state: the experience of being gripped. Action and care matter because they are what the feeling is about. On the \emph{conative} reading, engagement is fundamentally caring about an object in a way that issues in action directed at it. The feeling of absorption is the normal manifestation of such caring, but it is not its essence.

The readings diverge in the case of someone who acts with care but without feeling. A nurse who is clinically depressed continues to work devotedly: she attends to her patients, worries about them, does what she can for them, but she feels nothing in doing so; the work does not grip her, and she takes no satisfaction in it. On the affective reading, she is not actively engaged, and her work adds nothing to the meaning of her life during her illness. On the conative reading, she is engaged, and her depression diminishes her well-being without diminishing the meaning of her life.

We think the conative reading is the better reading of the text and the more plausible view. Three considerations support it. First, the text's adjective is \enquote{active}, not \enquote{felt}: it contrasts engagement with passivity, not with numbness. Second, EXI-17-a's definition ends with the clause \enquote{so that her activity expresses her concern rather than being merely performed}, which makes concern, not feeling, the decisive element. Third, the affective reading would make the meaning of a life depend on a feature of its subject's affective life that is not under her control and that bears no relation to the worth of what she does. The depressed nurse is doing exactly what she would be doing if she were well, for the same reasons and with the same concern; it would be strange if her illness, which takes so much from her, also took the meaning from her work.

The conative reading draws on Frankfurt's account of caring. \citet{Frankfurt1982} argued that what a person cares about is distinct from what she desires or values, that caring consists in a person's identifying herself with what she cares about and making herself vulnerable to its fortunes, and that it gives her life a structure of concern that persists over time. Active engagement on the conative reading is caring in Frankfurt's sense that issues in activity. Frankfurt himself held that the importance of what we care about derives in large part from our caring about it, a view that Definition~1 rejects through the qualifier \enquote{of objective worth} (\Lm{2.7}); but his account of what caring is can be separated from his account of where importance comes from, and Definition~1 takes the first without the second.

\paragraph{Activity and process.} \enquote{Active} also recalls Aristotle's distinction between activity (\emph{energeia}) and process (\emph{kinesis}) (\emph{Metaphysics} $\Theta$.6, 1048b18--35). A process is directed at an end outside itself and is incomplete until the end is reached: building is incomplete until the house is built. An activity is complete at every moment: one who sees has seen, one who contemplates has contemplated. Engagement as Definition~1 understands it has the structure of an activity. The mathematician engaged with her problem is engaged throughout, and her engagement is not incomplete because the problem remains unsolved. This is the structural reason why, as EXI-17-a argues, meaning does not require achievement: achievement belongs to processes, which are completed or not, while the engagement that constitutes meaning is an activity, which is complete in being exercised. The activity can be directed at an end outside itself, as the mathematician's is, without being incomplete until the end is reached.

\paragraph{Appreciation.} \enquote{Active} might seem to exclude receptive engagement: listening to music, contemplating a landscape, reading. It does not. EXI-17-a argues that the listener who spends much of her life studying music she loves, with growing understanding, is actively engaged, since attentive listening and study are activities directed at their object (its Lemma~3.2). Activity in the relevant sense is not productive activity but directed attention, and receptive engagement involves it as fully as production does. What the adjective excludes is passive consumption: the music played in the background of other activity, the diversion that fills time without engaging attention.

\paragraph{Ground projects.} \citet{Williams1981a} argued that a person's ground projects, the commitments that give her reason to go on living, cannot be required of her to be abandoned by an impartial morality without loss of her integrity. The projects that give life its point in Williams's sense are, in the terms of Definition~1, the objects of the deepest engagement, and EXI-03-a uses Williams's related notion of categorical desire, a desire that gives one reason to go on living, and identifies the categorical desires that give life its point with desires for projects one takes to be worthwhile (its Definition~2.1). Engagement, on the conative reading, is the practical form of such desires. Williams's argument against the desirability of immortality \citep{Williams1973}, which EXI-03-a reconstructs, then becomes an argument about the persistence of engagement: an endless life would eventually exhaust the projects with which a person of a given character can engage (Theorem~3.3 of EXI-03-a). EXI-19-b draws the consequence that finitude is a condition of a life's remaining meaningful throughout (its Proposition~1.3).

\lem{2.6}{with projects}

\paragraph{What a project is.} \enquote{Project} has a narrow and a broad sense. In the narrow sense, a project is an extended undertaking with an aim and a completion condition: writing a book, building a house, raising money for a cause. In the broad sense, a project is anything to which a person devotes sustained, directed activity: a relationship, a practice, an inquiry, the care of a garden. Williams's ground projects are projects in the broad sense, since they include personal relationships, and so are Wolf's examples, which include relationships of love and the pursuit of understanding \citep{Wolf2010}.

The text does not say which sense it intends, and the difference matters. On the narrow reading, Definition~1 excludes meaning that comes from love of a person, since a person is not an undertaking, and Assumption~3 becomes doubtful for agents whose circumstances leave no room for undertakings. On the broad reading, Definition~1 includes relationships, and Assumption~3 becomes much easier to defend.

\begin{proposition}[The broad reading]\label{p:broad}
\enquote{Projects} in Definition~1 is to be read in the broad sense: any object of sustained, caring activity is a project in the sense of the definition.
\end{proposition}

\begin{proof}[Argument]
Three considerations require the broad reading. First, the definition is cited to Wolf, whose central cases of meaning include relationships of love, and a definition that excluded them would not be the definition Wolf defends. Second, the narrow reading would make Definition~1 false: lives devoted to the love and care of particular persons are among the paradigms of meaningful lives, and on the narrow reading they would have no meaning in respect of that love. Third, the family of EXI-01 adopts the broad reading. EXI-01-c proves that sustained, caring activity directed at a worthwhile object of engagement is a project of objective worth in the sense of Definition~1 (its Lemma~2.3), and EXI-17-a counts persons, practices and inquiries among the objects of engagement (its Definition~2.2). A reading of a text that makes it false and inconsistent with its family, where another reading makes it true and consistent, is not the correct reading.
\end{proof}

On the broad reading, the later formulations of EXI-17-a and EXI-11-a, which speak of engagement \enquote{with what is worthwhile}, say the same as Definition~1. The word \enquote{projects} contributes something nonetheless: it marks that the engagement must be sustained and directed, not momentary. A single encounter with a worthwhile object, attended to once, is not a project; a relationship of care sustained over time is. EXI-01-c makes the same point in reply to the objection that its lemma trivializes the notion of a project.

\paragraph{The plural.} The definition speaks of projects in the plural and imposes no condition on how they are related. A life may be meaningful through several unrelated engagements, and its meaning is not reduced by their failure to form a whole. This distinguishes Definition~1 from views that require a unifying purpose. \citet{Cottingham2003} argues that a meaningful life requires that one's projects be integrated within a framework that gives them a common point, and that such a framework is in the end religious. EXI-11-a grants that a theistic framework would add meaning, in part by unifying a person's projects under a single purpose that is itself good (its Proposition~4.1), but denies that unification is required. Definition~1 is consistent with both claims: it does not require unity, and its measure is open to the view that a unified set of engagements constitutes a greater extent of engagement than a scattered one.

\paragraph{Projects and defining projects.} \enquote{Project} is used in a different sense in EXI-09, which holds with Sartre that for human beings existence precedes essence. EXI-09-a defines a being's defining project as the end that it adopts and in the light of which it understands itself (its Definition~2.1). A defining project in this sense need not be worthwhile, and a life can have one without being meaningful; EXI-09-b notes that a life can be less meaningful, in the sense of EXI-01, than it might have been, because it engages little with what is worthwhile, without being a defective specimen of anything. The projects of Definition~1 are the objects of a person's engagement; the defining project of EXI-09 is the self-understanding that organizes her ends. The two may coincide, as when a person understands herself through her commitment to a worthwhile cause, but they are different notions.

\lem{2.7}{of objective worth}

\paragraph{Why a qualifier is needed.} The qualifier \enquote{of objective worth} is what makes Definition~1 a hybrid view rather than a subjective one. Without it, the definition would say that a life is meaningful to the extent that it consists in active engagement with projects, whatever their worth, and it would be a version of the view of \citet{Taylor1970}: meaning comes from within, from the will to do what one does. Three cases show that the qualifier is needed.

The first is Taylor's own. Sisyphus is condemned to roll a stone to the top of a hill, from which it rolls down again, forever. Taylor imagines that the gods, in a perverse act of mercy, implant in Sisyphus a compulsive desire to roll stones, so that he does what he most wants to do. Taylor concluded that Sisyphus's life then has all the meaning a life can have. The second is Rawls's grass counter, a person whose only pleasure is to count blades of grass in various geometrically shaped areas such as park squares and well-trimmed lawns \citep[\S65]{Rawls1971}. Rawls introduced him to show that the good of a rational person may consist in odd activities; he has served the literature on meaning as the paradigm of engagement with a worthless object. The third is a case of the kind \citet{Wolf2010} discusses: a person who devotes herself with passion and skill to a pursuit that has no value beyond her attachment to it.

Reflection judges that these lives are not meaningful, or not meaningful in respect of the activities described, although their subjects are actively engaged. The best explanation of the judgment is that engagement confers meaning only when its object has worth that the engagement does not confer. That is what the qualifier says. Whether the explanation is correct, and whether there are any objects with worth of the required kind, is the question of Assumption~2, and the defence in Section~\ref{s:defence} returns to these cases.

\paragraph{\enquote{Worth}.} The text says \enquote{worth}, not \enquote{value}, and Assumption~2 glosses the first in terms of the second (\Lm{3.2}). The words are close, but \enquote{worth} carries an implication that \enquote{value} lacks: that what has it is worthy of something, here of engagement. A project of objective worth is one that merits the engagement it receives. This links Definition~1 to the fitting-attitude tradition, in which values are understood in terms of the attitudes that are fitting toward their bearers, and to EXI-04-a's principle that an attitude is fitting to the extent that the evaluative features under which it presents its object are features the object has (its Principle~P1). Active engagement presents its object as worth engaging with. It is fitting when the object is, and Definition~1 says that only fitting engagement confers meaning. This is why Wolf calls her view the fitting-fulfilment view.

\paragraph{The kinds of worth.} Definition~1 does not say what has objective worth, and the text gives no list. ETH-01 holds that the good has seven irreducible components: well-being, knowledge, achievement, relationships, fairness, beauty and autonomy. EXI-17-a takes these as examples of the worthwhile (its comment on Definition~2.2), and the natural reading of Definition~1 within the collection does the same: projects of objective worth are projects that serve or realize one or more of these goods, the well-being of others, understanding, the making or appreciation of beautiful things, relationships of love and friendship, the advancement of fairness. Two features of this list matter for the definition. It is plural, so that meaning can be found in many different kinds of project, which is what the availability claim needs. And it is not moral: worthwhile projects include aesthetic and epistemic projects that are morally neutral, so that meaning is not the same as moral goodness. A life devoted to mathematics can be meaningful without being morally admirable, and a morally decent life can lack meaning if it engages with nothing.

\paragraph{The degree of worth.} \enquote{Of objective worth} is grammatically a threshold: projects either have objective worth or lack it. But worth comes in degrees, and EXI-11-a states as a principle that engagement with what is of greater worth makes a life more meaningful, other things equal (its Principle~P2). The comment on \Lm{2.3} has argued that the measure in Definition~1 must be sensitive to the worth of the objects of engagement. The threshold reading and the degree reading are therefore both correct, at different points: a project contributes meaning only if it has some objective worth, and it contributes more meaning the greater its worth.

\lem{2.8}{\citep{Wolf2010}}

\paragraph{What is taken from Wolf.} The citation acknowledges that Definition~1 is Wolf's view. Wolf's lectures argue for three claims that the definition presupposes: that meaning is a distinct value, not reducible to happiness or to morality; that it has a subjective and an objective component, neither of which is sufficient; and that the objective component is worth that does not derive from the subject's attitudes. Her slogan, quoted in EXI-20-a, is that meaning arises when subjective attraction meets objective attractiveness. Definition~1 follows closely Wolf's own formula, in her earlier paper, that meaningful lives are lives of active engagement in projects of worth \citep{Wolf1997}.

\paragraph{What is added.} Definition~1 adds three things to Wolf's formulation. It makes the degree explicit, as a variable, where Wolf speaks of lives as meaningful \enquote{insofar as} they are engaged. It replaces Wolf's emphasis on fulfilment, a word that suggests a feeling, by \enquote{active engagement}, which the comment on \Lm{2.5} has read conatively; Wolf's own descriptions of engagement include affective terms, and the conative reading is a choice the commentary has defended, not a report of Wolf's view. And it is combined, in Assumption~2, with a definition of \enquote{objective} that is more restricted than the one Wolf's lectures suggest (\Lm{3.4}).

\paragraph{What is not taken.} Wolf's lectures also argue that meaning gives us reasons for action that are neither self-interested nor moral, so that a picture of practical reason divided between egoism and impartial morality leaves out a large part of what moves us. Definition~1 is neutral on this claim. It says what meaning is, not what reasons it gives. Wolf also declines to give a theory of objective worth, holding that our judgments of worth are fallible and revisable and that no theory of them is needed for the account of meaning to be correct. EXI-01-a follows her in giving no theory, but it is not neutral on the status of worth, since Assumption~2 commits it to the existence of objective worth in a specified sense. The collection supplies the theory Wolf declined to give, in ETH-01 and MTE-01, and Section~\ref{s:defence} examines how much of that theory the text needs.

\paragraph{Questions pressed against Wolf.} Three questions that have been pressed against Wolf's view, in the exchange published with her lectures and in the later literature, bear on Definition~1. They are whether the objective component can be specified without either elitism or emptiness, whether the subjective component asks too much of people whose circumstances do not allow absorption in their work, and whether the account matches what empirical psychology finds about the sources of a sense of meaning. The first question is the elitism objection, which EXI-01-d answers; the second is answered by the conative reading of \Lm{2.5} and by EXI-01-c; the third concerns the sense of meaning, which Definition~1 does not define, rather than meaning itself, and is considered at \Lm{2.4}, where we noted that a subject can be mistaken about the meaning of her life.

\subsection{Definition 1 and its rivals: the cases}\label{s:rivals}

The comments on the phrases of Definition~1 have tested it against cases one at a time. We now bring the cases together and compare the verdicts of Definition~1 with those of its principal rivals. The comparison is a commentary device of a familiar kind: it shows what each phrase of the definition does by showing what goes wrong when it is removed or replaced.

\paragraph{The rival schemata.} We state each rival in a simplified schematic form, sufficient to derive its verdicts on the cases. The schemata are ours; they capture the central condition of each family of views and are not offered as complete statements of any author's theory.

\begin{enumerate}[label=(R\arabic*)]
\item \emph{Subjectivism.} A life is meaningful to the extent that its subject does what she most wants to do, or is engaged with what she cares about, whatever its object. Taylor's conclusion about his Sisyphus is the classical statement \citep{Taylor1970}.
\item \emph{Pure objectivism.} A life is meaningful to the extent that its subject's activity is directed at, or realizes, objective value, whatever her attitude toward it.
\item \emph{The good cause schema.} A life is meaningful to the extent that it causes good outcomes, whether or not they are intended or known; this is the central condition of the account of \citet{Smuts2013}.
\item \emph{The transcendence schema.} A life is meaningful to the extent that it makes a difference to, and leaves traces in, something of value beyond itself; this is one central condition of the view of \citet[ch.~6]{Nozick1981}.
\item \emph{The purpose schema.} A life is meaningful only if it fulfils a purpose assigned to it by a creator; this is the purpose theory of EXI-11-a (its Definition~2.3).
\item \emph{Definition~1.} A life is meaningful to the extent that it consists in active engagement, read conatively, with projects of objective worth, read in the broad sense.
\end{enumerate}

\paragraph{The cases.} The cases are those that the comments have used, together with two from the family. In each, the question is whether the life is meaningful in respect of the activity described.

\begin{enumerate}[label=(c\arabic*)]
\item Taylor's Sisyphus, given a compulsive desire to roll stones (\Lm{2.7}).
\item Rawls's grass counter (\Lm{2.7}).
\item The mathematician who never solves her problem (EXI-17-a, Lemma~3.1).
\item The nurse on the ward where all patients die (EXI-17-a, Lemma~3.1).
\item The listener who studies music and produces nothing (EXI-17-a, Lemma~3.2).
\item The idle inventor, whose joke device later saves many lives (\Lm{2.4}).
\item The depressed nurse, who works with care but without feeling (\Lm{2.5}).
\item Ivan Ilyich, absorbed in his own advancement (Section~\ref{s:defence}).
\item The prisoner who cares for a fellow prisoner (EXI-01-c).
\item The life that fulfils a degrading purpose assigned by its creator (EXI-11-a, Lemma~3.1).
\end{enumerate}

\begin{center}\small
\begin{tabular}{@{}l c c c c c c c c c c@{}}
\toprule
 & c1 & c2 & c3 & c4 & c5 & c6 & c7 & c8 & c9 & c10 \\
\midrule
Reflective judgment & -- & -- & + & + & + & -- & + & -- & + & -- \\
\midrule
(R1) Subjectivism & + & + & + & + & + & -- & -- & + & + & ? \\
(R2) Pure objectivism & -- & -- & + & + & + & + & + & -- & + & -- \\
(R3) Good cause & -- & -- & -- & + & -- & + & + & -- & + & ? \\
(R4) Transcendence & -- & -- & -- & + & -- & + & + & -- & + & ? \\
(R5) Purpose & ? & ? & ? & ? & ? & ? & ? & ? & ? & + \\
(R6) Definition~1 & -- & -- & + & + & + & -- & + & -- & + & -- \\
\bottomrule
\end{tabular}

\smallskip
{\footnotesize + meaningful in respect of the activity; -- not meaningful in that respect; ? the schema gives no verdict without further facts.}
\end{center}

\paragraph{How the verdicts are derived.} The row of reflective judgments records the verdicts that the discipline accepts and that the Commentary has defended at the lemmata cited: EXI-17-a for c3--c5, the comments on \Lm{2.4}, \Lm{2.5} and \Lm{2.7} for c1, c2, c6 and c7, Section~\ref{s:defence} for c8, EXI-01-c for c9 and EXI-11-a for c10. The other rows are derived from the schemata. Under (R1), the depressed nurse is not doing what gives her satisfaction and is not, on the affective reading the schema requires, engaged; the inventor does not care about his device; Ivan Ilyich is doing what he wants. Under (R2), activity directed at or realizing value suffices, so the inventor and the depressed nurse qualify, while Sisyphus, the grass counter and Ivan Ilyich's advancement do not, since their objects have no value. Under (R3), what matters is good caused: the unsuccessful mathematician and the listener cause little or none, while the inventor causes much and the nurse on the second ward causes the relief of her patients' suffering. Under (R4), a difference made beyond oneself is required: the nurse on the second ward makes a difference to how her patients die, the listener and the unsuccessful mathematician leave no traces in the relevant sense. (R5) gives verdicts only given facts about a creator's purposes, except in c10, which is constructed so that the purpose is fulfilled. (R6) gives the verdicts derived in the comments.

\begin{proposition}[Fit with the cases]\label{p:fit}
Of the six schemata, only Definition~1 agrees with reflective judgment on all ten cases.
\end{proposition}

\begin{proof}
By inspection of the table. Subjectivism disagrees on c1, c2, c7 and c8; pure objectivism on c6; the good cause schema on c3, c5 and c6; the transcendence schema on c3, c5 and c6; the purpose schema on c10, and gives no verdict on the rest. Definition~1 agrees on every case.
\end{proof}

The proposition is not a proof that Definition~1 is true. The cases were chosen by the discipline, and some were chosen to support Definition~1; a defender of a rival can contest the reflective verdicts, as Taylor contested the verdict on c1 and as a defender of the good cause account may contest the verdict on c6. What the table shows is the structure of the dispute. Each rival departs from reflective judgment in a characteristic place: subjectivism where engagement meets worthlessness, objectivism where worth meets indifference, the consequence-based schemata where engagement meets failure. Definition~1 is the conjunction of the two conditions whose separate failures the table records, and the cases on which the rivals fail are, one by one, the cases that show why each conjunct is needed.

Two further observations. First, the only case on which pure objectivism fails is c6, the idle inventor, and it is the case on which reflective judgment is least secure; a reader who doubts the verdict there will find pure objectivism nearly as well supported by the table as Definition~1. The comment on \Lm{2.5} gave the reason for the verdict, and the depressed nurse (c7) shows that Definition~1 on the conative reading is not far from objectivism: it requires care, not feeling. Second, the shape-based views of \citet{Kauppinen2012} and of EXI-19 do not appear in the table, because they do not give verdicts on the meaning of a life in respect of an activity; they concern the meaning of a life as a whole. Objection~\ref{o:shape} in Section~\ref{s:whole} considers them.

\section{Assumption 2}\label{s:ass2}

\lem{3}{There exist projects of objective worth: projects whose value does not depend on the attitudes of those engaged in them \citep{Metz2013}.}

\paragraph{The sentence as a whole.} Assumption~2 has two parts separated by a colon. The first part is an existential claim: there are projects of objective worth. The second part says what objective worth is: a project has it if its value does not depend on the attitudes of those engaged in it. The second part is grammatically an apposition, but it functions as a definition, and it is the only place in the text where \enquote{objective} is explained. Definition~1, Assumption~3 and the proof of Theorem~4 all use the phrase \enquote{of objective worth}, and all of them inherit its sense from the colon of Assumption~2.

\paragraph{The role of the assumption.} One would expect an assumption in an axiomatic text to be used in the proof of its theorem. Assumption~2 is not.

\begin{proposition}[Redundancy in the derivation]\label{p:redundant}
Theorem~4 follows from Definition~1, Assumption~3 and the step from ability to an available life (Proposition~\ref{p:bridge}), without Assumption~2. Moreover, Assumption~3 entails the existential part of Assumption~2 provided that there is at least one agent capable of intentional action.
\end{proposition}

\begin{proof}
The proof of Theorem~4 (\Lm{6}--\Lm{8}) cites Assumption~3 and Definition~1 and nothing else, and each of its steps is licensed by one of these or by the step from ability to an available life. For the second claim: if there is an agent capable of intentional action, Assumption~3 gives a project of objective worth in which she is able to engage, and so there is a project of objective worth.
\end{proof}

The Lean module confirms the first claim mechanically: Assumption~2 is not among the axioms on which \texttt{meaning\_available} depends (Section~\ref{s:lean}). The second claim shows that the existential part of Assumption~2 adds nothing to Assumption~3 in any world that contains an agent.

Why, then, is the assumption stated? Proposition~\ref{p:redundant} answers the question by elimination. The assumption does three things, none of them in the derivation.

First, it \emph{interprets}. Its second part fixes the sense of \enquote{objective worth}, and without it Definition~1 would be indeterminate between the hybrid view and the subjective view: \enquote{objective} might mean no more than \enquote{real} or \enquote{genuine}, and a subjectivist could accept every sentence of the text. The formal derivation goes through under any interpretation of the predicate (Proposition~\ref{p:neutralfml}); the philosophical content of Theorem~4 depends on its interpretation, and Assumption~2 supplies it.

Second, it \emph{separates}. The existential part isolates, as a claim on its own, what the hybrid view needs from the theory of value: that there is something with worth of the required kind. A reader who doubts this has a reason to doubt the text that is independent of any doubt about Assumption~3, and the text lets her locate her doubt. The separation also lets the reader see that Assumption~3 is stronger than Assumption~2, since it entails it given an agent and adds a claim about every agent's abilities.

Third, it \emph{answers}. Remark~5 says that Assumption~2 denies the premise of the sceptical challenge that, \emph{sub specie aeternitatis}, nothing matters (\Lm{11.4}). The assumption is stated so that the text can say exactly where it stands on that challenge: it does not refute it, but it is committed to denying one of its premises, and the commitment is made in a sentence of its own.

\lem{3.1}{There exist}

The existential quantifier ranges over projects. Two questions about its domain should be settled.

The first is whether projects are types or tokens. A token project is a particular person's undertaking: this mathematician's work on this problem, this nurse's care of these patients. A project type is something that many people can engage in: mathematical research, the care of the sick. Assumption~2 is best read as a claim about types, or about possible tokens, since it must support Assumption~3, which says that every agent has a project of objective worth in which she is \emph{able} to engage. A project in which no one has yet engaged can exist in the relevant sense. When a project type has objective worth, the token undertakings that instantiate it inherit its worth, though they may add or subtract worth according to how they are carried out.

The second question is whether the existence asserted is modest or robust. The modest reading says only that some projects meet the condition stated after the colon. The robust reading adds that objective worth is a property with a place in the ultimate inventory of the world, perhaps a non-natural one. The text is committed only to the modest reading. Whether worth is natural or non-natural, reducible or primitive, is a question for the theory of value; Assumption~2 says that, however that question is answered, some projects have value that does not depend on the attitudes of those engaged in them. ETH-01 and MTE-01 give a naturalist answer, which Section~\ref{s:defence} shows to be sufficient.

\lem{3.2}{projects of objective worth: projects whose value}

\paragraph{The colon.} The phrase repeats \enquote{projects}, which signals that what follows the colon specifies the class introduced before it. The definition is of the form: a project is of objective worth if and only if its value does not depend on the attitudes of those engaged in it. The definition presupposes that the project has a value, and makes objectivity a property of that value: objective worth is objective value.

\paragraph{Whose value?} The value in question is the value of the project, but a project's value usually derives from the value of its object. The value of caring for a sick person derives from the value of her well-being; the value of a mathematical inquiry derives from the value of understanding. EXI-01-c makes the same point in the proof of its Lemma~2.3: the value of caring for a person or of coming to understand something derives from the value of the person or of the understanding, not from the agent's attitude. This matters for how Assumption~2 is defended. It need not be shown that projects, as such, have intrinsic value. It need only be shown that their objects have value, or that the projects serve value, independently of the attitudes of those engaged. The value of a project of objective worth may be wholly derivative from the value of what it serves, and its worth is no less objective for that.

\lem{3.3}{does not depend on}

\enquote{Depend} admits two readings, and the choice between them determines what Assumption~2 claims.

On the \emph{counterfactual} reading, the value of a project does not depend on certain attitudes if it would have the same value whatever those attitudes were. On the \emph{grounding} reading, the value of a project does not depend on certain attitudes if those attitudes are not among the facts in virtue of which the project has its value.

The counterfactual reading is too strong. Consider a friendship. Its value depends, counterfactually, on the attitudes of the friends: if they did not care about one another, there would be no friendship and so no value of friendship. On the counterfactual reading, friendship would therefore lack objective worth, and Definition~1 would imply that friendship confers no meaning. This is absurd, and it cannot be what the text means, since the family counts relationships of care among the paradigms of meaningful engagement (EXI-01-c, its Assumption~2.2; EXI-17-a, its Definition~2.2).

The grounding reading avoids this result if \enquote{attitudes} is read with care. The friends' care for one another is a constituent of the friendship, not an evaluation of it. What Assumption~2 excludes is that the value of the friendship should be grounded in the friends' \emph{evaluative} attitudes toward it: in their regarding it as valuable, or their being attached to it, or their caring about it as a project. A friendship is not valuable because the friends value it; they value it, if they see correctly, because it is valuable, and it is valuable in part because it consists in mutual care.

\begin{proposition}[Restricted objectivity]\label{p:restricted}
On the reading required by the text, a project is of objective worth if and only if it has value, and that value is not grounded in the evaluative attitudes, toward the project or its object, of those engaged in it.
\end{proposition}

\begin{proof}[Argument]
The counterfactual reading of \enquote{depend} implies that relationships of mutual care have no objective worth, which contradicts the family's treatment of such relationships as paradigms of meaningful engagement; so the text requires the grounding reading. On the grounding reading, if \enquote{attitudes} included every attitude, the same contradiction would follow, since mutual care consists of attitudes; so the attitudes excluded are evaluative attitudes, those by which the engaged agents regard, prize or care about the project or its object as worth their engagement. The resulting condition is what Definition~1 needs: it excludes the case in which engagement confers on its object the worth that makes the engagement meaningful.
\end{proof}

Proposition~\ref{p:restricted} identifies the target of Assumption~2 precisely. The target is the view that engagement \emph{bootstraps}: that by caring about something, one makes it worth caring about, so that every caring engagement is engagement with something of worth. That view is Taylor's, and in a qualified form Frankfurt's. Assumption~2 denies that bootstrapping is the only source of worth: it says that there are projects whose worth is not bootstrapped.

\lem{3.4}{the attitudes of those engaged in them}

\paragraph{The restriction.} Assumption~2 does not say that the value of a project of objective worth is independent of all attitudes. It says that it is independent of the attitudes of \emph{those engaged in it}. This is a restriction, and it is easy to miss.

The restricted notion is compatible with three kinds of attitude-dependence that a stronger notion would exclude. It is compatible, first, with the dependence of a project's value on the attitudes of those it serves: the value of giving a gift depends in part on the pleasure of the recipient, and the value of teaching depends in part on the pupils' coming to understand, which is an attitude. It is compatible, second, with the dependence of value on the attitudes of a community, as the value of preserving a language may depend on the attachment of its speakers to it. It is compatible, third, with theories on which value in general is constituted by idealized attitudes: theories on which something is good if and only if idealized agents would desire or endorse it. On such theories, the value of a project depends on attitudes, but not on the actual attitudes of the people engaged in it, which may differ from the idealized ones.

\paragraph{Rival readings in the collection.} The manuscripts that restate Assumption~2 do not all preserve the restriction. EXI-17-a defines a worthwhile object as one that has or serves value that does not depend on \enquote{the agent's attitudes toward it} (its Definition~2.2), which is the restricted notion, made more precise by the phrase \enquote{toward it}, which agrees with Proposition~\ref{p:restricted}. EXI-01-c adopts the same definition (its Definition~2.1). EXI-06-b, by contrast, describes meaning on the account of EXI-01 as engagement with what has worth \enquote{independently of anyone's attitudes} (see the testimonia, Appendix~\ref{app:test}). That is the unrestricted notion. EXI-06-b's argument does not need it: the argument is that the universe's lack of attitudes toward us takes nothing from the worth of what we engage with, and on the restricted notion the universe is not among those engaged in our projects, so its attitudes are irrelevant in any case. EXI-06-b's gloss is therefore a strengthening of the text that its argument does not require, and the Commentary reads Assumption~2 as restricted.

\begin{proposition}[Sufficiency of the restricted notion]\label{p:suff}
Every use that the collection makes of Assumption~2 is supported by the restricted notion of objective worth.
\end{proposition}

\begin{proof}[Argument]
The uses are four. (i) Definition~1 needs objective worth to exclude engagement with worthless objects such as Taylor's stones and Rawls's grass; the restricted notion does this, since the only source of worth for those objects would be the agent's own attitudes. (ii) EXI-17-a's argument that meaning is not hostage to luck in outcomes (its Corollary~3.6) needs only that the worth of the object belong to the object rather than to the agent's success, which the restricted notion secures. (iii) EXI-06-b needs only that worth not depend on the universe's attitudes, and the universe is not engaged in our projects. (iv) EXI-04-a needs, in its reply to its Objection~2, only that worth not be conferred by the seriousness with which we pursue our ends, and seriousness is an attitude of the engaged agent. The theistic case is the one that might seem to need more: EXI-11-a holds that worth is not conferred by assignment (its Principle~P1), so that a creator's purposes for us do not make our activities worthwhile. But Principle~P1 of EXI-11-a is a separate principle, argued for separately, and EXI-01-a does not need it to state its account.
\end{proof}

\paragraph{\enquote{those engaged}.} The plural noun phrase covers projects in which many people are engaged together. A scientific research program, an orchestra, a political movement are projects with many participants, and Assumption~2 requires that their value not be grounded in the evaluative attitudes of any of them. This excludes a form of collective bootstrapping, on which a shared enthusiasm makes its object worthwhile for those who share it. A community of enthusiasts for a worthless pursuit is not given meaning by its enthusiasm, however widely shared, any more than a single enthusiast is.

\lem{3.5}{\citep{Metz2013}}

The citation of Metz rather than Wolf at this point is significant. Wolf supplies the account of meaning in Definition~1, but she declines to give a theory of objective worth. Metz gives the most sustained recent argument that meaning depends on conditions that are not constituted by the subject's attitudes. His survey of theories argues, against subjectivism, that it cannot account for the judgment that some engaged lives lack meaning, and, against supernaturalism, that meaning does not require a relation to a spiritual being or realm; his own fundamentality theory holds that a life is more meaningful the more its subject employs her reason in positive ways toward conditions of human existence that are responsible for much else about it \citep{Metz2013}.

What EXI-01-a takes from Metz is the existential claim and the argument against subjectivism, not the fundamentality theory. The fundamentality theory is a rival to Definition~1 as an account of what meaning consists in: it makes the degree of meaning depend on the fundamentality of the conditions to which one's reason is directed, a notion that Definition~1 does not use. The citation is therefore a citation of an ally on one point, the existence of attitude-independent worth, and not of a source of the definition. A reader who checks the citation will find in Metz a theory that disagrees with EXI-01-a about the measure of meaning while agreeing with it about the existence of what meaning is a measure of.

\subsection{A defence of Assumption 2}\label{s:defence}

Remark~5 says that a defence of Assumption~2 is given in the Commentary (\Lm{12}). We give it here. It establishes the restricted claim of Proposition~\ref{p:restricted}: that there are projects whose value is not grounded in the evaluative attitudes of those engaged in them. It does not establish any stronger claim about the metaphysics of value, and it does not answer the demand, considered in Section~\ref{s:absurd}, that worth be vindicated from a standpoint that presupposes nothing. The defence has four parts.

\paragraph{(D1) The argument from cases.} Taylor's Sisyphus, given the desire to roll stones, is actively engaged in rolling them. Rawls's grass counter is actively engaged in counting. If the value of a project were grounded in the evaluative attitudes of those engaged in it, these projects would have as much value as any, since their subjects prize them as much as anyone prizes anything, and engagement with them would confer as much meaning as engagement with anything. Reflection judges that it does not: the lives of the stone-roller and the grass counter are not meaningful in respect of their rolling and counting, however engaged they are, and they are less meaningful than the lives of people equally engaged in the care of others or in the pursuit of understanding. The best explanation of this judgment is that the value of the care of others and of understanding is not grounded in the attitudes of those engaged in them, and the value of rolling stones and counting grass, if any, is. So there are projects of objective worth in the restricted sense.

The subjectivist can reject the judgment. Taylor did: he held that his Sisyphus's life was meaningful. But the cost of rejecting it is high, since it requires holding that there is no difference in meaning between the life of the grass counter and the life of an equally engaged physician, and few who consider the case are prepared to say this. Judgments about cases are evidence in this domain, as EXI-17-a assumes (its Assumption~2.4) and as ETH-01 holds for judgments about the good. The subjectivist's position is consistent, but it is not supported by the evidence that bears on the question.

\paragraph{(D2) The argument from the content of engagement.} \citet{Wiggins1976} observed that a purely subjective account of meaning cannot be held by someone who finds her life meaningful, because engagement presents its object as worth engaging with in a way that does not depend on the engagement. The mathematician does not care about her problem and conclude that, since she cares about it, it is important. She cares about it because she takes it to be important, and if she came to believe that its importance was nothing but a projection of her caring, her caring would be undermined. The same is true of the nurse's care for her patients and of the parent's love for a child: in each case the engaged agent takes the worth of the object to be the reason for her engagement and not its product.

If subjectivism were true, engagement would therefore be systematically mistaken about its own content: every engaged agent would take her object to have a worth it has only in virtue of her taking it to have it. This would make every meaningful life rest on an illusion. EXI-18 holds that some illusions are psychologically necessary, but that their necessity does not make them true, and EXI-18-a's reply to the objection that meaning needs illusion is that the serious agent need not believe that her ends are certified from the standpoint of the universe, only that they are worthwhile, which may be true (its reply to Objection~3). Subjectivism denies that it may be true in the sense in which the agent believes it. The hybrid view, and Assumption~2 with it, allows the engaged agent's view of her object to be correct.

\paragraph{(D3) The argument from error and discovery.} People discover that their lives have been wasted. Tolstoy's Ivan Ilyich, in \emph{The Death of Ivan Ilyich}, comes to see on his deathbed that the life he had lived, correct, successful and absorbed in its own advancement, had not been what it should have been. Mill records in his \emph{Autobiography} (ch.~V) how, at twenty, he asked himself whether the realization of all the reforms to which he had devoted himself would make him happy, answered that it would not, and found that the ends that had given his life its direction had lost their hold on him. Tolstoy's \emph{Confession} describes a crisis of the same structure \citep{Tolstoy1987}. In each case the subject had been engaged; in Ivan Ilyich's case the discovery is that the objects of his engagement had not deserved it.

On a subjectivist account, such discoveries are unintelligible as discoveries. If the worth of Ivan Ilyich's pursuits was grounded in his attitudes toward them, then while he held those attitudes his pursuits were worthwhile, and his deathbed judgment can only record a change of attitude, not a recognition that he had been wrong. But the judgment presents itself as a recognition, and it is so understood by those who read the story. The possibility of being wrong about whether one's life is meaningful, and of discovering it, requires that meaning depend on something that one's attitudes do not settle. Assumption~2 says what that something is.

Mill's case shows a different structure, and it is important not to assimilate it. Mill did not discover that the reforms to which he had devoted himself were worthless; he never doubted their worth. He discovered that his engagement had failed, that the ends no longer moved him. On Definition~1, his life during the crisis lacked meaning not because its objects lacked worth but because it lacked engagement, and his recovery, which he attributed to the cultivation of feeling, restored the subjective component. The pair of cases illustrates the two components of Definition~1: Ivan Ilyich had engagement without worth, Mill had worth without engagement.

\paragraph{(D4) The argument from the theory of value.} The collection contains a theory of value that entails Assumption~2. ETH-01 holds that the good is what idealized reflective endorsement converges on: one outcome is better than another if and only if every agent's idealized reflective endorsement would prefer the first (Definition~3.2 of ETH-01-a), where idealized reflective endorsement is the attitude an agent would take if she were fully informed, vividly imagined the alternatives, reasoned without error and gave equal weight to everyone affected (Definition~3.1 of ETH-01-a). MTE-01 holds that moral facts are natural facts about idealized reactions.

\begin{proposition}[ETH-01 entails Assumption~2]\label{p:eth01}
If the good is what idealized reflective endorsement converges on, and the good has the components listed in ETH-01, then there are projects of objective worth in the restricted sense.
\end{proposition}

\begin{proof}
Take a project that serves one of the components, for instance the care of the sick, which serves their well-being. On ETH-01, the outcomes this project brings about are better than their alternatives because the idealized reflective endorsement of every agent would prefer them. That fact does not depend on the actual evaluative attitudes of those engaged in the project: had the nurses who care for the sick regarded the care as worthless, or been indifferent to it, the idealized endorsement of every agent, theirs included, would still prefer the outcomes in which the sick are cared for, since idealization corrects for ignorance, failures of imagination and partiality, and nothing in the nurses' actual attitudes enters the definition except as material for correction. So the value of the project is not grounded in the evaluative attitudes of those engaged in it, and it is of objective worth in the sense of Proposition~\ref{p:restricted}.
\end{proof}

Proposition~\ref{p:eth01} shows that Assumption~2 is no more doubtful than the theory of the good of ETH-01, and it shows why the restriction in \Lm{3.4} matters. On ETH-01, value is attitude-dependent in a sense: it is constituted by idealized attitudes. An unrestricted notion of objective worth, on which worth is independent of all attitudes, would be inconsistent with ETH-01. The restricted notion is not. The text's restriction is therefore what allows EXI-01-a to rest on the theory of value of the collection rather than on a more contested non-naturalism.

The standard objections to that theory have answers elsewhere in the collection, which we record by their conclusions. The argument from queerness, that objective values would be entities of a strange sort known by a strange faculty \citep{Mackie1977}, proves too much: it would make mathematics false as well (MTE-05). The argument from evolutionary debunking, that our evaluative dispositions were selected for reasons unrelated to their truth \citep{Street2006}, fails because evolution selected for tracking reasons, and those reasons include evaluative ones (MTE-09, stated for moral reasons). The argument from disagreement fails because disagreement about value is no deeper than disagreement in other theoretical domains that are not thought to lack facts (MTE-12). Each of these is stated in the collection for moral facts, and each applies to the restricted worth of Assumption~2 at least as well, since the restricted claim is weaker than moral realism.

\paragraph{Objections to the defence.}

\begin{objection}[Frankfurt]\label{o:frankfurt}
\citet{Frankfurt1982,Frankfurt2004} argues that what is important to a person is important to her largely because she cares about it, and that love, in particular, does not respond to the value of its object but confers value on it for the lover. The parent does not love her child because the child is valuable; the child is valuable to her because she loves it. So the paradigm of meaningful engagement is a case in which worth is conferred by the attitude of the one engaged.
\end{objection}

\begin{reply}
Frankfurt's claim concerns importance \emph{to} the lover, and on that it may be right: the child matters to the parent in a way it matters to no one else, and the difference is made by her love. Assumption~2 concerns the worth of the project, here the care of the child, and that worth is not conferred by the parent's love. The child's well-being has value whether or not anyone loves her, as the case of an unloved child shows: it is not a matter of indifference whether an unloved child flourishes. A parent who cares for her child is engaged with something whose worth she did not create, and her love is fitting because of that worth. What her love adds is a particular relation to that worth, the relation that EXI-16 calls a reason of love not transferable to qualitatively similar persons; it does not add the worth. Frankfurt's view of importance can be accepted as an account of the subjective component of Definition~1, which is what the comment on \Lm{2.5} took from him, and rejected as an account of the objective component.
\end{reply}

\begin{objection}[Elitism]\label{o:elitism}
Every attempt to say which projects have objective worth privileges the projects of the educated and the privileged: art, science, causes. Assumption~2 therefore smuggles a class judgment into the theory of meaning.
\end{objection}

\begin{reply}
Assumption~2 is existential. It says that there are projects of objective worth and what their worth consists in; it does not say which projects have it. The list of goods that the collection uses, from ETH-01, begins with well-being and relationships, which are within the reach of every life, and EXI-01-d answers the objection directly: caring for others, honest work that serves others' needs and the pursuit of understanding at any level meet the condition of Assumption~2, and are open to everyone. The objection would hold against a theory that measured worth by prestige. It does not hold against a theory that measures it by idealized endorsement, which gives equal weight to everyone affected.
\end{reply}

\begin{objection}[Alienation]\label{o:alien}
If the worth of my projects is independent of my attitudes, it is not clear why engaging with them makes \emph{my life} meaningful rather than merely making the world better. Objective worth, so understood, is alien to me.
\end{objection}

\begin{reply}
The objection would hold against a purely objective theory, on which meaning consists in serving objective value whether or not one cares about it. On such a theory, the meaning of my life would be a fact about the world's ledger. Definition~1 is not such a theory. It makes meaning consist in a relation between the subject and objective worth, the relation of active engagement, and the relation is as much the subject's as anything in her life. Objective worth is what makes the relation fitting; engagement is what makes it hers. This is the reason Wolf gave for the hybrid view, and it is the reason Definition~1 has two components.
\end{reply}

\paragraph{What the defence establishes.} The defence establishes that Assumption~2 is supported by judgments about cases (D1), by the content of engagement itself (D2), by the intelligibility of discovering that one's life has been wasted (D3), and by the theory of value of the collection (D4), and that the objections from Frankfurt, from elitism and from alienation do not undermine it. It does not establish that the theory of value is true; it shows that Assumption~2 is supported to the extent that that theory is, and that the restricted form of the assumption is weaker than the theory. Nor does it answer the sceptic who grants all this and asks what vindicates the judgments about cases, the content of engagement and the theory of value from a standpoint that does not presuppose them. That sceptic is Nagel's, and Remark~5 says that Assumption~2 denies the premise of his challenge rather than refuting it. Section~\ref{s:absurd} examines what this means. Volume~II, Part~A, develops (D4) in detail, and Part~B examines the subjectivist and hybrid theories at length.

\subsection{Assumption 2 and the metaethical options}\label{s:metaeth}

The defence has shown that Assumption~2 follows from the theory of value of ETH-01. A reader who does not accept that theory will want to know whether Assumption~2 survives under the others. Because the assumption is restricted (Proposition~\ref{p:restricted}), it is compatible with more metaethical views than an unrestricted assumption of objective worth would be. We take the main options in turn, describing each by its central thesis and, where the collection has a family on it, by that family's conclusion.

\paragraph{Non-naturalist realism.} On the view that evaluative properties are sui generis and not natural, worth is independent of every attitude, and Assumption~2 holds in its unrestricted form and so in its restricted form. The restriction costs the non-naturalist nothing.

\paragraph{Naturalist realism.} On the view of MTE-01, moral facts are natural facts about idealized reactions; on the view of ETH-01, the good is what idealized reflective endorsement converges on. Proposition~\ref{p:eth01} has shown that Assumption~2 holds on these views because the idealized attitudes that constitute value are not the actual attitudes of those engaged. The restriction is what makes Assumption~2 consistent with them: an unrestricted assumption would deny that value is constituted by attitudes of any kind.

\paragraph{Subjective response-dependence.} On the view that something is valuable for a person if and only if she would desire it under some modest idealization of her own attitudes, value is relative to the person. Assumption~2 is false on this view if the idealization is so modest that the grass counter, idealized, would still desire to count grass; it is true if the idealization corrects for the narrowness that makes the grass counter's desire compulsive. The defence of (D1) argues that the case requires the second. A response-dependent view that cannot distinguish the grass counter from the physician is not a rival to Assumption~2 but a version of subjectivism, and is answered by (D1) to (D3).

\paragraph{Expressivism and quasi-realism.} On expressivist views, evaluative judgments express non-cognitive states; MTE-03 holds that moral sentences express states of being for, which compose according to the logic of plans. Such views have a characteristic way with claims of attitude-independence: the claim that the value of caring for the sick does not depend on the attitudes of the carers is itself an evaluative claim, expressing a plan to value the care whatever the carers' attitudes, and the expressivist can make it. Assumption~2 can therefore be asserted by an expressivist, as a first-order claim. MTE-17 adds that quasi-realism, once its explanatory commitments are made explicit, is indistinguishable from realism, and on that conclusion the expressivist who asserts Assumption~2 asserts what the realist asserts. Either way, the assumption is not hostage to the outcome of the dispute between realism and expressivism.

\paragraph{Constructivism.} On constructivist views, evaluative truths are the outcomes of an idealized procedure of practical reasoning; MTE-18 holds that constructivism is realist about procedures and not about their outcomes, and that the procedure itself is mandatory. On such views the worth of a project is the verdict the procedure delivers on it. The procedure is not the actual attitudes of those engaged in the project, and Assumption~2 in its restricted form holds whenever the procedure delivers a verdict of worth on some project independently of what its participants happen to think, which any procedure worth the name does.

\paragraph{Error theory.} On the view that all evaluative judgments are false, because they presuppose objective values that do not exist \citep{Mackie1977}, Assumption~2 is false. MTE-05 answers this view: the argument from queerness proves too much, since it would make mathematics false as well. The error theorist is the one metaethical opponent of Assumption~2, and the dispute with her is not special to the theory of meaning.

\begin{proposition}[Metaethical robustness]\label{p:robust}
Assumption~2 in its restricted form holds on every metaethical view considered above except error theory and those forms of subjective response-dependence that cannot distinguish idealized from actual desire.
\end{proposition}

\begin{proof}[Argument]
By the discussion of each view: non-naturalist realism entails the unrestricted form; naturalist realism of the kind of ETH-01 and MTE-01 entails the restricted form (Proposition~\ref{p:eth01}); expressivism permits it as a first-order claim and, on MTE-17, does not differ from realism in asserting it; constructivism entails it whenever its procedure is not identical with the actual attitudes of participants; error theory denies it; and response-dependence entails it just when its idealization corrects for actual desire.
\end{proof}

Proposition~\ref{p:robust} explains why EXI-01-a can state Assumption~2 without defending a metaethics. The assumption is a point on which almost every metaethical view that allows evaluative truth agrees, once the restriction of \Lm{3.4} is observed. Its opponents are the error theorist, who denies that there is evaluative truth, and the subjectivist about meaning, who holds that the only evaluative truth relevant to meaning is truth about what the engaged agent herself wants. The first is answered by MTE-05 and the second by the defence.

\section{Assumption 3}\label{s:ass3}

\lem{4}{For every agent capable of intentional action, there is at least one project of objective worth in which she is able to engage.}

\paragraph{The sentence as a whole.} Assumption~3 carries the weight of the availability claim. Its logical form is $\forall a\,\exists p\,(W(p) \wedge E(a,p))$: for every agent there is a project that is of objective worth and in which she is able to engage. Each conjunct of the matrix makes its own demand. The first requires that the project have worth in the restricted sense of Assumption~2. The second requires that the agent be able to engage in it. Neither alone would support Theorem~4: that there are projects of objective worth says nothing about any particular agent, and that every agent can engage in some project says nothing about worth, since the project might be the counting of grass.

EXI-01-c calls this the premise a critic will press, and it is. The comment examines its quantifiers, its domain, the modality of \enquote{able}, and the support it receives in the family.

\lem{4.1}{For every agent}

\paragraph{The universal quantifier.} The quantifier is unrestricted within its domain: it covers every agent capable of intentional action, without exception for circumstance, talent, wealth, health or location. This is what makes the conclusion of the family universal, and it is what makes Assumption~3 vulnerable, since one counterexample suffices to refute it.

\paragraph{The order of the quantifiers.} The universal quantifier precedes the existential one.

\begin{proposition}[Order of quantifiers]\label{p:order}
Assumption~3 does not imply that there is a single project of objective worth in which every agent is able to engage, and Theorem~4 does not require it.
\end{proposition}

\begin{proof}
The form $\forall a\,\exists p\,\phi(a,p)$ does not imply $\exists p\,\forall a\,\phi(a,p)$: each agent may have her own project. The proof of Theorem~4 instantiates the existential quantifier for each agent separately (\Lm{6}), and so uses only the weaker form.
\end{proof}

The proposition marks the difference between EXI-01-a and views on which meaning is available to all because there is one project open to all. On some religious views, the love of God is such a project; on some moral views, the moral life is. EXI-01-a does not need any such project. The prisoner's project and the scholar's may have nothing in common except that each has objective worth and each is within the reach of the one who engages in it. This is one reason why the availability claim is more plausible than it might seem: it does not require that the same opportunity be open to everyone, only that each person have some opportunity.

The stronger form is nonetheless arguably true, and EXI-01-c comes close to asserting it: its Assumption~2.2 says that every agent can direct caring activity at another person or at some part of the world she can come to understand. That is a disjunction of two kinds of project rather than a single project, but the kinds are so general that nearly every agent has both. The text does not need this, and the Commentary does not assert it.

\lem{4.2}{capable of intentional action}

\paragraph{The domain.} The domain of the quantifier is agents capable of intentional action. The phrase has a standard sense in the philosophy of action: an intentional action is one done for a reason, one to which the question \enquote{why?}, in the sense that asks for the agent's reason, applies \citep{Anscombe1957}. An agent capable of intentional action is one who can act for reasons. The restriction is natural, since engagement, on the conative reading of \Lm{2.5}, involves acting in ways directed at an object out of care for it, which is acting for a reason.

The same phrase defines the domain of Lemma~3.1 of EXI-04-a, which holds that any agent capable of intentional action takes some ends seriously, and of the theorems of EXI-01-c and EXI-01-d. The agreement is not accidental. The discipline treats the capacity for intentional action as the threshold at which questions of meaning and absurdity arise: below it there is no seriousness to collide with detachment (EXI-04-a) and no engagement to constitute meaning (EXI-01-a).

\paragraph{Capacity and exercise.} The phrase speaks of capacity, not exercise. An agent who is asleep, or unconscious under anesthesia, remains capable of intentional action and remains in the domain. An agent in a persistent vegetative state, if she has permanently lost the capacity, is outside it. The boundary is not sharp, and the text does not claim that it is: capacities for intentional action come in degrees and develop over time, as in children, and decline, as in dementia. The availability claim holds for each agent while she has the capacity, and EXI-01-c notes, in reply to its third objection, that an agent for whom nothing worthwhile is within reach of attention is, if possible at all, a case of the absence of agency rather than of its exercise.

\paragraph{\enquote{Everyone}.} The family conclusion says that meaning is available to \enquote{everyone}, and Remark~5 says that the theorem answers the question \enquote{whether meaning is available to everyone} (\Lm{9.2}). The theorem itself says \enquote{every such agent}. The two are reconciled by reading \enquote{everyone} as \enquote{every agent capable of intentional action}, and the reading is the natural one: the question whether meaning is available to someone is a question about someone who could engage with something. Human beings who are not, or not yet, or no longer, agents in this sense raise a different question, whether their lives can have meaning in a sense other than the one Definition~1 defines, for instance through the engagement of others with them. Volume~III, Part~A, takes up that question; the family conclusion does not answer it and does not need to.

\lem{4.3}{there is at least one project of objective worth}

\paragraph{The minimal claim.} \enquote{At least one} states the minimum that Theorem~4 needs. A single project of objective worth in which an agent can engage suffices for a life of positive meaning to be available to her. The assumption says nothing about how many such projects there are, how great their worth is or how wide the choice among them. These bear on the degree of meaning available, which EXI-01-d says varies with circumstance (its Proposition~2.4), and not on its availability.

The minimal form makes the assumption easier to defend than its stronger relatives. It is a much stronger claim that every agent has a \emph{wide} range of worthwhile projects available, or that every agent can attain a \emph{high} degree of meaning, and those claims are false: deprivation narrows the range, and injustice lowers the attainable degree. Assumption~3 does not deny either fact. The comment on \Lm{5.2} returns to the point.

\paragraph{\enquote{Project}, again.} If \enquote{project} were read in the narrow sense of an undertaking with a completion condition, Assumption~3 would be doubtful, as EXI-01-c observes: prisoners, those who labor without respite, the very ill and the very poor may have no such undertaking open to them. Proposition~\ref{p:broad} has argued that the broad sense is the correct reading of Definition~1, and Assumption~3 uses the word in the same sense. In the broad sense, the care of a fellow prisoner, the attention a laborer gives her children and the bedridden patient's attention to those who visit her are projects, provided they are sustained and caring, and they have objective worth because their objects do. EXI-01-c derives Assumption~3, so read, from its Assumption~2.2 and its Lemma~2.3, and so makes it a theorem rather than a premise.

\lem{4.4}{in which she is able to engage}

\paragraph{The modality.} \enquote{Able} is a modal word, and several modalities are candidates. A project is \emph{physically} open to an agent if nothing in the world prevents her from undertaking it. It is \emph{epistemically} open if she knows of it or could come to know of it. It is \emph{motivationally} open if she could come to care about it. It is \emph{socially} open if she is permitted and enabled by those around her to undertake it. The text does not say which is meant.

The best reading is the conjunction, restricted to what the agent can do from where she is. To be able to engage in a project, in the sense that matters for availability, is to be able, from one's actual situation, to come to act with care in ways directed at it. That requires that the project be physically within reach, that the agent be able to come to know of it, and that she be able to come to care about it, though not that she already knows of it or cares about it. The last condition is the one most often overlooked. A person may be surrounded by worthwhile objects and unable to care about any of them, as in severe depression; and Assumption~3 then requires that the inability be temporary or partial. The comment on \Lm{2.5} argued that the depressed agent who acts with care but without feeling is engaged; the agent who cannot act with care at all is a harder case, and Volume~III, Part~A, discusses it. Here we note only that \enquote{able} is a dispositional notion, and that a disposition can be masked without being lost.

\paragraph{\enquote{she}.} The pronoun is the generic feminine used throughout the discipline. It has no philosophical significance, but it does have a grammatical one: it is bound by \enquote{every agent} and shows that the ability is the agent's own. The availability claim is not that projects of objective worth are available to the agent in the sense that someone else could engage in them on her behalf, but that she can.

\paragraph{\enquote{engage}.} The verb is \enquote{engage}, without \enquote{actively}. The text is not careless here: the ability in question is the ability to engage, and active engagement in the sense of Definition~1 is its exercise. The proof of Theorem~4 must move from the ability to engage to a life that consists in active engagement, and the comment on \Lm{7} examines how it does so.

\paragraph{Frankl's third way.} \citet{Frankl1962} held, from his experience of the camps, that meaning can be found in three ways: by creating a work or doing a deed, by experiencing something or encountering someone, and by the attitude one takes toward unavoidable suffering. The first two correspond to engagement with projects in the broad sense. The third raises a question for Assumption~3. If an agent is so placed that no project, not even the care of another, is open to her, can she still find meaning in the way she bears her condition? On Definition~1, she can only if bearing one's condition with courage or dignity is a project of objective worth in which she is able to engage. We think it is: the exercise of courage is directed activity, it can be sustained and caring, and courage is a good that idealized reflection endorses. If so, Frankl's third way is a limiting case of Definition~1 rather than an exception to it, and it is the case that makes Assumption~3 hardest to refute, since an agent capable of intentional action is always able to take some attitude toward her situation. We record this reading without relying on it: the argument of EXI-01-c establishes Assumption~3 without appeal to it.

\paragraph{Support for the assumption.} The text does not argue for Assumption~3. The family supplies three supports.

The first is EXI-01-c, which derives Assumption~3 from the premise that every agent capable of intentional action can direct some intentional activity, with care, at a worthwhile object of engagement, another person or some part of the world she can come to understand (its Assumption~2.2), together with the lemma that sustained caring activity directed at such an object is a project of objective worth (its Lemma~2.3). The premise is supported by the ubiquity of persons and of the world as an object of understanding. It is weaker than Assumption~3 in what it asserts about projects and stronger in what it asserts about objects, and EXI-01-c argues that it is better supported.

The second is EXI-17-a, which shows that meaning does not require achievement (its Theorem~3.3) and that it is therefore not hostage to luck in outcomes (its Corollary~3.6). This removes one of the main threats to Assumption~3. If meaning required success, the projects in which an agent could engage would confer meaning only if they could succeed, and many agents have no project that can succeed. Since meaning requires only engagement, a project can confer meaning though it fails, and the projects open to an agent are not reduced by the world's refusal to cooperate.

The third is EXI-04-a's Lemma~3.1, that every agent capable of intentional action takes some ends seriously. This does not entail Assumption~3, since the ends one takes seriously may be worthless. But it shows that the subjective capacity on which engagement depends is present in every agent in the domain, so that a failure of Assumption~3 would have to be a failure of worth or of reach, not a failure of the capacity for engagement as such.

\paragraph{Objections.}

\begin{objection}[Radical deprivation]\label{o:depriv}
There are agents capable of intentional action who are so deprived, by imprisonment, poverty, illness or oppression, that nothing worthwhile is within their reach. Assumption~3 is false for them.
\end{objection}

\begin{reply}
The objection is the one EXI-01-c was written to answer, and it is answered by the broad reading of \enquote{project}. The prisoner can care for a fellow prisoner and can come to understand something; the exhausted laborer can attend to her children; the bedridden patient can attend to her visitors and to what she hears. These are projects of objective worth in the sense of Definition~1, and they are within reach. What deprivation takes away is the range of projects and the degree of meaning attainable, which EXI-01-d acknowledges (its Proposition~2.4). It does not take away the possibility of meaning, because the objects of worth that remain in every human situation, persons and the world as an object of understanding, are not things that deprivation can remove.
\end{reply}

\begin{objection}[Triviality]\label{o:trivialA3}
On the broad reading, Assumption~3 is trivially true, and Theorem~4 inherits its triviality. The text purchases universality by emptying \enquote{project} of content.
\end{objection}

\begin{reply}
The assumption is not trivial on any reading. It requires, for each agent, an object of objective worth, which a subjectivist denies exists; it requires that the agent be able to engage with it, which requires that it be within her physical, epistemic and motivational reach; and it requires, through the word \enquote{project}, that the engagement be sustained and caring, not momentary. What the broad reading removes is the requirement that a project have a completion condition, a public product or an institutional form, and these are requirements that, as EXI-17-a shows for achievement, have nothing to do with meaning. A claim that is easy to satisfy because its conditions are the right ones is not thereby trivial.
\end{reply}

\subsection{Hard cases}\label{s:hard}

The objection from radical deprivation is usually pressed with examples, and an assumption with a universal quantifier should be tested on the hardest of them. We take seven. In each we ask whether the agent, while she remains capable of intentional action, has a project of objective worth within her reach, and we read Assumption~3 as indexed to times, as the comment on \Lm{5.3} argues. A fuller treatment, with the empirical literature on each condition, belongs to Volume~III, Part~A; the purpose here is to show how the lemmata apply.

\paragraph{Solitary confinement.} The prisoner in solitary confinement lacks the first of EXI-01-c's two routes, care for another person present to her, for long stretches. Engagement with absent persons through memory and imagination is real attention, but Principle~P2 of EXI-17-a requires that engagement bear on its object, and it is doubtful that thinking of an absent friend bears on the friend. The second route remains: some part of the world can be attended to and understood. Prisoners in solitary confinement have worked through mathematics in their heads, reconstructed and studied what they remembered, learned the habits of insects in their cells; what is attended to is real, and understanding of it is a good that does not depend on the prisoner's attitudes. The project is meager, and the degree of meaning it can sustain is low. But it is a project of objective worth in the broad sense of Proposition~\ref{p:broad}, and the prisoner is able to engage in it. The case confirms the importance of the second route, which EXI-01-c includes precisely because the first can be taken away.

\paragraph{Severe chronic pain.} Pain consumes attention, and engagement requires attention. A person in severe chronic pain may be unable, for long periods, to sustain any directed activity. Two features of the lemmata bear on the case. First, \enquote{able} in \Lm{4.4} is dispositional, and the disposition to engage is masked, not lost, by pain: in the intervals in which the pain recedes, the person can engage, and the projects she takes up then are hers. Second, Proposition~\ref{p:neutral} requires only that engagement be sustained over a non-null part of the life, not over every part. The availability claim holds; the degree attainable is reduced, which is EXI-01-d's Proposition~2.4.

\paragraph{Advanced dementia.} Dementia erodes the capacity for intentional action gradually. In its earlier and middle stages, the person retains the capacity to act for reasons and to care about others, and relationships of affection with those who care for her remain within her reach and have worth. Assumption~3 holds. In its final stages, the capacity for intentional action may be lost, and the person then falls outside the domain of \Lm{4.2}. The availability claim does not say that meaning is available to her then, and it does not say that it is not; the comment on \Lm{4.2} has noted that the question whether lives outside the domain can have meaning in another sense is left open. What the indexed reading secures is that the boundary of availability is the boundary of agency, not a boundary set by the illness before agency is lost.

\paragraph{Early childhood.} A young child acquires the capacity for intentional action gradually, and with it the capacity for engagement. Play and the exploration of the world are the child's first engagements, and much of them is engagement with understanding and with other persons, which have worth. The domain of Assumption~3 grows with the child, and within it the assumption holds. Whether play as such is a project of objective worth is a further question, which the case does not require us to settle, since the child's play typically includes engagement with persons and with understanding.

\paragraph{Forced labor and enslavement.} The person whose labor is owned by another and directed to another's ends has the narrowest range of projects that human societies have imposed. The ends to which her labor is directed are not hers, and whether they have worth is irrelevant to the meaning of her life, since she is not engaged with them; EXI-11-a's principle that worth is not conferred by assignment (its Principle~P1) has an analogue here, that meaning is not conferred by assignment either. But the persons among whom she lives remain, and with them the care of children and kin, friendship, the transmission of what one knows, and the pursuit of fairness in the form of resistance. These are projects of objective worth, and they have been pursued under the worst conditions that history records. The case is the strongest illustration of EXI-01-d's insistence that availability does not imply indifference: that meaning remains available under oppression is no reason to tolerate oppression, which reduces the degree of meaning as it reduces every other good.

\paragraph{The last days.} A person who is dying, with days to live and little strength, has a short continuation before her. The indexed reading of Theorem~4 requires that some continuation available to her contain engagement with a project of objective worth to a positive extent. The last conversations of the dying with those they love are such engagement: directed, caring and bearing on real persons. On the proportional reading of the measure, they make up a positive share of the remaining life; on the absolute reading, they add a positive amount. The case shows why the indexed reading matters: it is not too late for the dying to engage, and Definition~1 counts what they do.

\paragraph{Severe depression.} The hardest case is the one in which the capacity to care is itself impaired. The depressed nurse of \Lm{2.5} acts with care though without feeling; a person in a severe depressive episode may be unable to act with care at all, not because nothing worthwhile is within reach but because nothing moves her. Assumption~3 then holds only if \enquote{able to engage} is read dispositionally, as the comment on \Lm{4.4} argued: the capacity to care is masked by the illness, and typically returns with treatment or with time. If there were a person in whom the capacity to care was permanently and wholly lost while the capacity for intentional action remained, she would be a counterexample to Assumption~3 as read here. We doubt that such a case is possible, since acting for reasons involves caring about something, but we record it as the point at which the assumption is most exposed, and Volume~III, Part~A, examines it.

\begin{proposition}[Reach in the hard cases]\label{p:reach}
In each of the seven cases, while the agent remains capable of intentional action and her capacity to care is not permanently lost, some project of objective worth is within her reach, and Assumption~3 holds of her at that time.
\end{proposition}

\begin{proof}[Argument]
By the discussion of each case: the second route of EXI-01-c (understanding some part of the world) remains in solitary confinement; the dispositional reading of \enquote{able} covers chronic pain and depression; the first route (care for persons present) remains in dementia short of the loss of agency, in childhood once agency begins, under forced labor and in the last days.
\end{proof}

The proposition does not say that the meaning available in these cases is great. In most of them it is small, and the circumstances that make it small are evils. It says that it is not zero, and that is what Theorem~4 claims.

\section{Theorem 4}\label{s:thm}

\lem{5}{For every such agent, a life of positive meaning is available.}

\paragraph{The sentence as a whole.} Theorem~4 is the only result of EXI-01-a, and it states the second half of the family conclusion: meaning in the sense of Definition~1 is available to everyone. The first half, that life is meaningful to the extent that one engages with projects of objective worth, is stated by Definition~1 and is not a theorem. The text thus divides the family conclusion between a definition and a theorem, and Remark~5 treats the theorem, with its definition, as answering both of the questions the title raises (\Lm{9}).

The theorem has the logical form $\forall a\,\exists l\,(l \text{ is a life of } a \wedge M(l) > 0)$: for every agent in the domain, there is a life of hers that has positive meaning. It is weaker than it sounds in three ways and stronger in one. It is weaker because \enquote{positive} means no more than \enquote{greater than zero}, because \enquote{available} means no more than \enquote{possible for her from her situation}, and because it says nothing about actual lives. It is stronger because its quantifier is universal and admits no exception.

\lem{5.1}{For every such agent}

\enquote{Such} refers back to the domain of Assumption~3, agents capable of intentional action. The theorem therefore has exactly the scope of the assumption from which it is derived, and no wider. In particular, it has nothing to say about beings that are not agents in this sense; the comment on \Lm{4.2} has discussed what follows for the reading of \enquote{everyone}. The anaphor also shows that Theorem~4 is meant to be read together with Assumption~3, and the proof confirms this: the theorem is, in effect, Assumption~3 transposed from projects to lives.

\lem{5.2}{a life of positive meaning}

\paragraph{\enquote{a life}.} The indefinite article is important. The theorem does not say that the agent's life has positive meaning, or that it will; it says that there is a life, among those available to her, that does. Her actual life may be one of the others.

\paragraph{\enquote{positive meaning}.} The phrase is defined through Definition~1: a life has positive meaning if it is meaningful to some degree greater than zero. The theorem claims nothing about how much. This is the source of the objection, which EXI-01-d considers, that availability so secured is trivial: if a life of the barest positive meaning counts, it is no surprise that one is available to everyone.

The reply has three parts. First, the claim is not trivial, since it is denied by every view on which meaning requires success, talent, a cosmic purpose or an exceptional life, and these views have been widely held. EXI-17-a argues against the first, EXI-01-d against the second, and EXI-11-a against the third (its Theorem~3.3). Second, a minimal claim was what the question required. The question whether meaning is available to everyone is a question whether anyone is shut out of it, and to show that no one is, one needs to show only that a positive degree is available to each. Whether a high degree is available to each is a different question, to which the answer is no. Third, \enquote{positive} meaning in the sense of Definition~1 is not negligible meaning. A life that contains sustained, caring engagement with the well-being of one other person, or with the understanding of one part of the world, has positive meaning, and reflection does not judge such a life to be meaningful only in a degenerate sense. The minimal case is not a borderline case.

\paragraph{Positive meaning and equality.} The theorem implies nothing about equality of meaning. EXI-01-d draws the distinction: availability is universal, degree is not, and it would be an invalid inference from the first to the claim that all lives are equally meaningful, or that deprivation is harmless to meaning (its section on what availability does not imply). The comment on \Lm{4.3} has noted that the minimal form of Assumption~3 is compatible with great inequality in the range of projects open to different agents. Theorem~4 inherits that compatibility. It is a floor, not a level.

\lem{5.3}{is available}

\paragraph{Availability and actuality.} \enquote{Available} is the theorem's modal operator. A life of positive meaning is available to an agent if it is among the lives she can lead. The text does not define availability, and the proof relies on an unstated principle connecting it with Assumption~3's ability to engage (Proposition~\ref{p:bridge}). What the word excludes is clear: it does not say that the agent leads, will lead or is likely to lead such a life. Remark~5 makes the point explicit (\Lm{10}).

\paragraph{Availability at a time.} A question the text does not address is whether availability is indexed to times. Read without an index, the theorem says that among the lives an agent could have led, from the beginning, some have positive meaning. That is weaker than what the family conclusion seems to intend, which is that meaning is available to her now, whatever her situation. The indexed reading says: for every agent and every time at which she is an agent, there is a continuation of her life from that time in which it has positive meaning. On the indexed reading, the theorem says that it is never too late; on the unindexed reading, it says only that it was not always too late.

The proof supports the indexed reading if Assumption~3 is read as indexed: if at every time at which she is an agent there is a project of objective worth in which she is then able to engage. EXI-01-c's argument supports the indexed assumption, since persons and the world as an object of understanding are present at every time in a human life. And Definition~1 permits the inference: a life that, from some time on, includes sustained engagement with a worthwhile project consists to a positive extent in such engagement, on any measure satisfying Proposition~\ref{p:neutral}. We therefore read the theorem as indexed. It says that a life of positive meaning remains available to an agent for as long as she remains one.

\paragraph{Availability and luck.} The indexed reading connects Theorem~4 with EXI-17-a's Corollary~3.6, that meaning is not hostage to luck in outcomes, which EXI-17-a says supports, without by itself establishing, the availability claim of EXI-01. Luck enters a life in two ways: in the circumstances from which one acts, and in the outcomes of one's actions. EXI-17-a addresses the second: since meaning does not require success, bad luck in outcomes cannot remove it. Theorem~4, read with EXI-01-c, addresses the first: since every situation contains objects of worth within reach, bad luck in circumstances cannot remove the possibility of meaning either, though it can lower its attainable degree. Together they give the sense in which meaning, on the account of EXI-01, is immune to fortune: its possibility is, and its degree is not. \citet{Williams1981b} argued that the justification of some choices depends on luck in outcomes, with the example of a painter who abandons his family for his art and whose choice is vindicated only if he succeeds. If meaning depended on such vindication, the painter's life would be meaningful only if he were lucky. On Definition~1 it is meaningful to the extent of his engagement with painting, whether he succeeds or not; what depends on luck is the justification of his choice and the value his work adds to the world, not the meaning of his life. Section~\ref{s:whole} returns to the case.

\subsection{Availability on the rival views}\label{s:availrivals}

Theorem~4 is distinctive of EXI-01-a only if the rival accounts of meaning do not deliver it, or deliver it at a cost Definition~1 avoids. We consider the rival schemata of Section~\ref{s:rivals} in turn.

On subjectivism (R1), meaning is available to every agent trivially, since every agent capable of intentional action wants something and can do something she wants. Availability is secured, but at the cost of the cases: the grass counter and Taylor's Sisyphus have as much meaning as anyone. Subjectivism makes meaning universally available by making it indiscriminate.

On pure objectivism (R2), meaning is available to every agent whose activity can be directed at objective value, which, given the reach of persons and of the world as an object of understanding, is every agent. Availability is secured at the cost of the alienation objection (Objection~\ref{o:alien}): the meaning of a life could be present without its subject's caring about anything, and the subject who discovers that her indifferent activity was meaningful has discovered something about the world rather than about her life.

On the good cause schema (R3), meaning is available only to those whose activity can cause good outcomes, and whether it does depends on luck in outcomes. EXI-17-a's argument applies directly: the unsuccessful mathematician has no meaning in respect of her work on this schema, and whether a given agent's efforts will produce good is not in her power. The schema does not secure availability, and the failure is principled: on it, meaning is hostage to luck in outcomes.

On the transcendence schema (R4), meaning is available to those who can make a difference to something of value beyond themselves. Whether one can make a difference is less dependent on luck than whether one can cause good outcomes, since a difference can be small, but the prisoner in solitary confinement and the listener who produces nothing (Section~\ref{s:hard}; EXI-17-a, Lemma~3.2) make no difference that leaves traces, and the schema does not secure availability for them.

On the purpose schema (R5), meaning is available to every agent if and only if a creator has assigned to every human life a purpose that it can fulfil. EXI-11 holds that meaning does not require a cosmic purpose; on the purpose schema, availability is secured only if theism is true and the purposes assigned are within every agent's reach, and EXI-11-a's Lemma~3.1 shows that even then the assigned purpose confers meaning only if it is good.

\begin{proposition}[Availability without its costs]\label{p:availcost}
Of the six schemata, Definition~1 is the only one that secures the availability of meaning to every agent capable of intentional action without making meaning indiscriminate, without making it independent of the subject's care, without making it depend on luck in outcomes, and without making it depend on theism.
\end{proposition}

\begin{proof}[Argument]
By the discussion above: subjectivism secures availability and makes meaning indiscriminate; pure objectivism secures it and makes meaning independent of care; the good cause and transcendence schemata do not secure it, because they make meaning depend on outcomes or traces; the purpose schema secures it only if theism is true. Definition~1 secures it by Theorem~4, with Assumption~3 supported by EXI-01-c, and it avoids the four costs by its two components: worth excludes indiscriminateness, engagement excludes independence of care, the constitutive reading of \enquote{consists in} excludes dependence on outcomes, and Assumption~2 makes no reference to a creator.
\end{proof}

The proposition gives the strongest reason for the form of Definition~1. Each of its components is required to avoid one of the costs, and the two together secure availability without any of them. The family conclusion's two halves, which the text divides between a definition and a theorem, are in this sense connected: the account of what meaning is was not chosen to make meaning available, but it is the account on which meaning is available for the right reasons.

\section{The proof of Theorem 4}\label{s:proof}

The proof has three sentences. The first instantiates Assumption~3, the second infers that a life containing the resulting engagement consists in engagement to a positive extent, and the third applies Definition~1. The proof is valid given one principle that it uses without stating, and the comment on \Lm{7} identifies it.

\lem{6}{By Assumption 3, there is a project $p$ of objective worth in which the agent can engage.}

\paragraph{Instantiation.} The sentence performs two instantiations. It instantiates the universal quantifier of Assumption~3 to \enquote{the agent}, the arbitrary agent of Theorem~4's \enquote{every such agent}, and it introduces a name, $p$, for a project whose existence Assumption~3 then guarantees. The second is existential instantiation, and the proof treats $p$ as arbitrary among the projects that meet the condition: nothing is assumed about $p$ except that it is of objective worth and that the agent can engage in it. The Lean module performs the same step with \texttt{obtain} (Section~\ref{s:lean}).

\paragraph{\enquote{can} for \enquote{is able to}.} Assumption~3 says that the agent \enquote{is able to engage}; the proof says that she \enquote{can engage}. The substitution is harmless, but it marks that the modality of Assumption~3 is the modality of the ordinary \enquote{can} of ability, which the comment on \Lm{4.4} has analyzed.

\lem{7}{A life that includes engagement with $p$ consists, to some positive extent, in active engagement with a project of objective worth.}

This is the central step of the proof, and it contains more than its surface suggests. It is a universally quantified conditional: any life that includes engagement with $p$ consists to a positive extent in active engagement with a project of objective worth. For the proof to reach its conclusion, it must also be true that there is such a life among those available to the agent. The sentence does not say this; the proof assumes it.

\lem{7.1}{A life that includes engagement with $p$}

\paragraph{The hidden premise.} The phrase \enquote{a life that includes engagement with $p$} is used as if there were such a life available to the agent. Nothing in Definition~1 or Assumption~3 says that there is. Assumption~3 says that the agent is able to engage in $p$, which is a fact about her abilities; Theorem~4 says that a life of a certain kind is available to her, which is a fact about the lives she can lead. The step between them is a principle connecting abilities with available lives.

\begin{proposition}[The bridge]\label{p:bridge}
The proof of Theorem~4 requires the following principle, which is not stated in EXI-01-a: if an agent is able to engage in a project, then some life available to her includes engagement with that project.
\end{proposition}

\begin{proof}
Without the principle, Assumption~3 and Definition~1 are consistent with the falsity of Theorem~4. Suppose that an agent is able to engage in a worthwhile project, but that no life available to her includes engagement with it. Then no life available to her consists to any extent in engagement with that project, and if it is her only such project, Definition~1 gives no life available to her a positive degree of meaning. So the principle is required. It suffices with the other premises, since it supplies the life whose existence the second sentence of the proof assumes.
\end{proof}

The Lean module makes the principle an axiom, \texttt{engagement\_realizable}, and \texttt{\#print axioms} lists it among the axioms of the theorem (Section~\ref{s:lean}). The formalization thus exposes a premise that the prose leaves implicit. That is one of the services formalization performs, and in this case it shows where the modal content of \enquote{available} comes from.

\paragraph{Is the bridge true?} The principle connects two modalities: what an agent is able to do and what lives she can lead. It is true on any reading on which they are the same modality, since to be able to engage in a project is then to be able to lead a life that includes the engagement. It could fail only if the agent could engage in a project but could not lead a life that included it, which would require that abilities be exercisable only in lives the agent cannot lead. On the reading of \enquote{able} given at \Lm{4.4}, ability is what the agent can do from her actual situation, and what she can do from her situation is what some continuation of her life available to her includes. The principle is then analytic of the notions involved, and we think it is properly regarded as a clarification of \enquote{available} rather than as a substantive premise. Its formalization as an axiom is nonetheless correct, since the module does not define availability in terms of ability, and an undefined connection must be stated.

\paragraph{\enquote{engagement} and \enquote{active engagement}.} The phrase says \enquote{engagement with $p$}; the rest of the sentence says \enquote{active engagement}. The proof moves from a life that includes engagement to a life that consists in active engagement. The move is valid only if the engagement included is active. The text supplies no argument for this, and none is needed if the bridge principle is read as concerning active engagement, since Assumption~3's \enquote{engage}, as the comment on \Lm{4.4} noted, refers to the ability whose exercise is active engagement in the sense of Definition~1. A life that included only half-hearted or perfunctory engagement with $p$ would not be the life the proof needs, and the bridge principle should be read as supplying a life of active engagement. The Lean module's predicate \texttt{Includes} carries this reading in its docstring: a life includes \enquote{active engagement} with a project.

\lem{7.2}{consists, to some positive extent, in}

\paragraph{From inclusion to extent.} The inference from \enquote{includes} to \enquote{consists to some positive extent in} is the second inference that the sentence contains. It depends on the measure of Definition~1. On a reading on which momentary engagement could be included in a life without making up any positive extent of it, the inference would fail: a single moment of engagement might be a part of measure zero. Proposition~\ref{p:neutral} states the condition under which the inference holds: the engagement included must be sustained over a non-null part of the life. The word \enquote{project}, read in the broad sense, supplies this condition, since a project is sustained engagement (Proposition~\ref{p:broad}), and so does the reading of \enquote{active engagement} as caring activity with a temporal structure.

\paragraph{\enquote{to some positive extent}.} The phrase states only positivity, which is all that Theorem~4 needs. It is the textual counterpart of the minimal claim of \Lm{5.2}. It also shows that the proof is insensitive to the worth of $p$, so long as $p$ has some: a project of slight worth and a project of great worth both yield a life of positive meaning, though not of the same degree.

\lem{8}{By Definition 1, it is meaningful to a positive degree.}

The last sentence applies Definition~1 from right to left: since the life consists to a positive extent in active engagement with a project of objective worth, the degree $d$ that Definition~1 assigns it is positive. The pronoun \enquote{it} refers to the life of \Lm{7}. Since that life is available to the agent by the bridge principle, a life of positive meaning is available, which is Theorem~4. The square that closes the proof marks the end of the derivation and not of the argument: the derivation is complete, but its premises, as Remark~5 goes on to say, are where the philosophical questions lie.

\paragraph{The proof as a whole.} The proof is valid, given the bridge principle, and short, because the theorem says little more than its premises. That is not a defect. The text was written to make the structure of the availability claim visible, and the proof shows that structure: the claim holds if and only if, for every agent, some project of objective worth is within her reach (given the bridge principle and Definition~1). Every objection to the availability of meaning must therefore be an objection to Definition~1, to Assumption~3 or to the bridge, and the Commentary has examined each. The proof also shows that Assumption~2 plays no role in the derivation (Proposition~\ref{p:redundant}), so that an objection to the availability claim that attacks the existence of objective worth is, strictly, an objection to Assumption~3, which entails it.

\subsection{Other routes to Theorem 4}\label{s:routes}

The proof of EXI-01-a is one route to Theorem~4. The collection contains the materials for others, and comparing them shows which premises are essential.

\paragraph{The route of EXI-01-c.} The alternative proof replaces Assumption~3 by Reach and Lemma~2.3 of EXI-01-c, and derives Assumption~3 as a corollary. It keeps Definition~1, and it needs the bridge principle at the same point as the original proof. Its premises are not fewer than those of EXI-01-a, but its premise about agents is a claim about objects within reach rather than about projects, and the cases discussed in Section~\ref{s:hard} show that it is the better supported of the two.

\paragraph{The route through achievement.} One might try to derive Theorem~4 from EXI-17-a's Corollary~3.6, that meaning is not hostage to luck in outcomes, together with EXI-04-a's Lemma~3.1, that every agent capable of intentional action takes some ends seriously. The attempt fails, and the way it fails is instructive. Corollary~3.6 removes one obstacle to availability, the dependence of meaning on success, but it does not supply a worthwhile project for every agent. Lemma~3.1 of EXI-04-a supplies, for every agent, ends taken seriously, but not ends of worth. To reach Theorem~4, one needs in addition that every agent can take seriously some end of worth, which is Assumption~3 again, or Reach. The route through achievement therefore shows that no combination of results about engagement and success suffices without a premise about the reach of worth. EXI-17-a says as much when it states that its corollary supports the availability claim without by itself establishing it.

\paragraph{The route through the attitude toward one's condition.} The comment on \Lm{4.4} considered Frankl's third way: the attitude one takes toward unavoidable suffering. If bearing one's condition with courage is a project of objective worth, and every agent capable of intentional action is able to take some attitude toward her condition, then every agent is able to engage in a project of objective worth, and Theorem~4 follows. This route has the weakest premise about reach, since an attitude toward one's condition is available whenever agency is, and the strongest premise about worth, since it requires that the manner in which one bears one's condition be itself a project of objective worth. We think the premise about worth is true, since courage is a good, but the route is less secure than EXI-01-c's because that premise is more contested than the worth of persons and of understanding. It has the merit of covering the limiting cases, such as prolonged solitary confinement, in which EXI-01-c's first route is closed and its second is narrow.

\paragraph{What every route needs.} Every route to Theorem~4 needs three things: an account of meaning on which engagement with worth suffices for positive meaning (Definition~1, read right to left), a premise that every agent has something of worth within her reach (Assumption~3, Reach, or the premise of the third route), and the bridge from what an agent can do to the lives available to her. None needs Assumption~2 in its derivation. This is a structural fact about the availability claim, and the Lean module displays it for the route of EXI-01-a.

\section{Remark 5}\label{s:remark}

Remark~5 has four sentences. The first says what the theorem answers, the second and third what it does not, and the fourth where the missing parts are to be found. The remark is the only place in EXI-01-a where the text speaks about itself, and it does so with precision: each of its claims about what the theorem does and does not answer can be checked against the definition, the assumptions and the proof, and each is correct.

\lem{9}{The theorem answers the question of what the meaning of life is, and of whether meaning is available to everyone.}

\lem{9.1}{The theorem answers the question of what the meaning of life is}

\paragraph{The theorem or the definition?} Strictly, the question what the meaning of life is is answered by Definition~1, not by Theorem~4. The theorem says that a life of positive meaning is available to every agent; it does not say what meaning is. The sentence is best read with \enquote{the theorem} understood as the result together with the definition in whose terms it is stated, as one speaks of a theorem of arithmetic as answering a question that its definitions make precise. The alternative reading, on which the remark claims that the theorem alone answers the first question, would make the remark false, and the principle of charity forbids it where another reading is available.

\paragraph{What kind of answer.} The answer to \enquote{what is the meaning of life?} that Definition~1 gives is a real definition of the property in virtue of which lives are meaningful. The comment on \Lm{1} has argued that this is the second of the three questions the title can raise. The remark confirms the reading: it treats the question of what the meaning of life is as answered, and the only answer in the text is Definition~1. A reader who asked the first question, what human life as such is for, will find the remark's claim surprising, and the text owes her the explanation, given in the comment on \Lm{1}, that the second question does not depend on the first. EXI-11 and EXI-06 support the explanation from two sides: meaning does not require a cosmic purpose (Theorem~3.3 of EXI-11-a), and the universe's indifference takes nothing from it (EXI-06-b, its Proposition~2.1).

\lem{9.2}{and of whether meaning is available to everyone}

The second question is answered by Theorem~4, with \enquote{everyone} read as every agent capable of intentional action (\Lm{4.2}). The answer is yes. The phrase \enquote{available to everyone} is the phrase of the family conclusion, and the remark thus identifies Theorem~4 as the text's statement of the second half of that conclusion. The comment on \Lm{5.3} has argued that the availability is best read as indexed to times, so that meaning remains available to an agent for as long as she is one.

\lem{10}{It does not answer the question of whether any particular life is meaningful, which depends on what is done with it.}

\paragraph{Availability is not possession.} The sentence draws the distinction on which EXI-01-d builds its account of what the availability claim does and does not imply: from availability one cannot infer actuality. Whether a particular life is meaningful depends on whether it in fact consists in active engagement with projects of objective worth, and Theorem~4 says only that some life available to its subject does.

\paragraph{\enquote{what is done with it}.} The relative clause makes the meaning of a particular life depend on what is done with it. The passive voice leaves open who does it, and the omission is apt. Much of what is done with a life is done by its subject: whether she engages with what is worthwhile is largely up to her. But some is done by others and by circumstance: the range of projects within her reach, the time and strength she has for them, the education that lets her see their worth. A voluntarist reading, on which the meaning of a life is wholly up to its subject, would overstate the case; a fatalist reading, on which it is wholly up to circumstance, would contradict Theorem~4. The passive voice accommodates both contributions, and EXI-01-d states how they are divided: the possibility of meaning is secured whatever circumstances do, and its degree depends on what the subject does and on the conditions under which she does it (its Proposition~2.4).

\paragraph{The left-to-right direction.} The sentence uses Definition~1 in the direction the proof does not use: a life in which nothing worthwhile is engaged with is not meaningful. The comment on \Lm{2.3} noted that the definition is a biconditional and that the remark depends on this. The Lean module, which formalizes only the positive case of Definition~1, cannot represent this sentence (Section~\ref{s:lean}).

\paragraph{The first-person question.} The question that Remark~5 sets aside, whether a particular life is meaningful, is the question a person most often asks about her own life, and Definition~1 says how it is to be answered. It divides into two questions that can be asked separately. The first is whether, and how far, her life consists in active engagement: whether there are things she cares about in a way that issues in sustained, directed activity. This is a question about her, and she is in a good position to answer it, though not an infallible one, since a person can mistake habit for care and care for habit. The second is whether the objects of her engagement have worth that does not depend on her caring about them. This is a question about the objects, and she is in no better position to answer it than anyone else who knows the facts and reflects well. The two questions have different epistemologies, and the division explains a familiar feature of crises of meaning: they take two forms. In one, the person finds that she no longer cares, as Mill did; the remedy, if there is one, lies in her. In the other, she finds that what she cared about did not deserve it, as Ivan Ilyich did; the remedy lies in finding something that does. Definition~1 gives each form its own diagnosis.

\paragraph{The comparative question.} Remark~5 says that whether a particular life is meaningful depends on what is done with it. The same holds for how meaningful it is, and for whether one life is more meaningful than another. The comment on \Lm{2.3} has argued that the degree depends on at least three variables, the share or amount of the life given to engagement, the depth of the engagement and the worth of its objects, and that some comparisons may be indeterminate. A particular life can therefore be made more meaningful in three ways: by giving more of it to engagement, by engaging more deeply with what one already engages with, or by turning to objects of greater worth. EXI-11-a's Principle~P2 states the third; EXI-17-a's remark that achievement may add to the degree of meaning (its Remark~3.5) bears on the second, since success can deepen engagement and enlarge the worth realized through it. The first is the one most within a person's control and the one most limited by circumstance, and it is the one on which EXI-01-d's Proposition~2.4 bears most directly.

\lem{11}{Nor does it answer the sceptical challenge that, \emph{sub specie aeternitatis}, nothing matters \citep{Nagel1971}: Assumption 2 denies the premise of that challenge rather than refuting it.}

This is the most important sentence of the remark and the one whose interpretation is most disputed in the discipline. It makes two claims. The first is that the theorem does not answer the challenge. The second, after the colon, says why: Assumption~2 denies the premise of the challenge but does not refute it. We take the phrases in turn and then give the analysis that the sentence requires.

\lem{11.1}{the sceptical challenge}

\paragraph{\enquote{sceptical}.} The adjective places the challenge in the family of sceptical arguments, and it does so on Nagel's own authority. \citet{Nagel1971} compared the sense of the absurd to epistemological scepticism: as the sceptic about knowledge steps back from our beliefs and finds that every justification we can offer depends on assumptions we cannot justify without circularity, so the reflective subject steps back from her concerns and finds that every justification of them depends on further concerns open to the same question. In both cases, Nagel argued, we cannot refute the sceptic and we cannot stop believing, or caring; and in both cases the right response is neither to pretend to a refutation nor to abandon what cannot be justified, but to live with the discrepancy.

\paragraph{\enquote{challenge}.} The noun is chosen with care. A challenge is not a thesis. The sceptic about knowledge need not assert that we know nothing; she asks how we know, and finds every answer wanting. The sceptic about meaning need not assert that nothing matters; she asks what justifies our taking anything to matter, and finds every answer open to the same question. The distinction matters for what Assumption~2 can do against the challenge. Against a thesis that nothing matters, Assumption~2 is a contradictory thesis. Against a challenge, it is an answer that the challenger will not accept.

\lem{11.2}{\emph{sub specie aeternitatis}}

The phrase is Spinoza's: the mind conceives things under a species of eternity when it conceives them as following from the necessity of the divine nature, and not in relation to a time and place (\emph{Ethics} V, propositions 29--31). Nagel adopts it for the standpoint from which the reflective subject regards her life as one among others, without presupposing the importance of anything in it, and he observes that it is a standpoint we can occupy without leaving our lives, since it is the standpoint of reflection itself \citep{Nagel1971}. In \emph{The View from Nowhere} he gave the standpoint a general theory, as the objective standpoint that we reach by abstracting from the particularities of our own point of view, and returned in its last chapter to the meaning of life \citep{Nagel1986}. EXI-04-a defines the corresponding notion as the detached standpoint: the standpoint from which an agent considers her ends as the ends of one creature among others, without presupposing their importance, and asks what justifies her concern with them (its Definition~2.2).

The Spinozan phrase carries an implication that Nagel's use drops. For Spinoza, to conceive things under the species of eternity is to conceive them truly, and the highest virtue of the mind lies in knowledge of this kind (\emph{Ethics} V, proposition 25). For Nagel, the standpoint discloses not the truth about our concerns but the absence of a vindication of them. EXI-01-a uses the phrase in Nagel's sense, and its citation of Nagel confirms this.

\lem{11.3}{nothing matters}

\paragraph{The formula.} \enquote{Nothing matters} is the slogan of the challenge in its strongest form. It has a history in the analytic literature before Nagel. \citet{Hare1972} described a young visitor who, after reading Camus's \emph{L'\'Etranger}, came to the conclusion that nothing matters and was distressed by it, and argued that the conclusion rests on a confusion about the word \enquote{matters}: to say that something matters is to express concern about it, and since the visitor was concerned about many things, the conclusion that nothing matters to him was false, and the conclusion that nothing matters \emph{simpliciter} was not something he could make sense of. Hare's answer is expressivist, and it is not the answer EXI-01-a gives, since on Assumption~2 whether something has worth is not settled by whether anyone is concerned about it. But Hare's diagnosis, that the challenge trades on an ambiguity between \enquote{matters to someone} and \enquote{matters}, survives the change of metaethics, and the analysis below uses it.

\paragraph{\enquote{matters} and \enquote{worth}.} The challenge says that nothing \emph{matters}; Assumption~2 says that some projects have \emph{worth}. The words are not synonyms, but in the context of the challenge they are equivalent: to matter, \emph{sub specie aeternitatis}, is to have importance that does not depend on its mattering to anyone in particular, and that is objective worth in the sense of Assumption~2 or a stronger one. The restriction of Proposition~\ref{p:restricted} does not affect the equivalence, since the detached standpoint abstracts from the attitudes of the subject who occupies it, and worth that does not depend on the attitudes of those engaged is worth that the detached subject has not presupposed by abstracting from her own engagement.

\lem{11.4}{Assumption 2 denies the premise of that challenge rather than refuting it}

\paragraph{Which premise?} The sentence speaks of \enquote{the premise} of the challenge, in the singular, and does not say which premise it is. A challenge, as the comment on \Lm{11.1} noted, need not have premises in the way an argument does. But a challenge that is to have the force of an objection to Assumption~2 must be reconstructible as an argument for the negation of something Assumption~2 asserts. The question is which of its premises Assumption~2 denies.

We reconstruct the challenge as follows.

\begin{enumerate}[label=(C\arabic*)]
\item If anything matters, its mattering can be vindicated from the detached standpoint: there is a justification of it that the detached standpoint is compelled to accept.
\item Nothing's mattering can be vindicated from the detached standpoint: every justification of concern is open to the question what justifies it, and no answer is forced on that standpoint.
\item Therefore nothing matters.
\end{enumerate}

(C2) is Nagel's central claim and the content of Lemma~3.2 of EXI-04-a, which proves it for any reflective agent: if the agent answers that a consideration has objective worth, the detached standpoint can register the claim but cannot be compelled to endorse it, since what she would exhibit is her own conviction, which the step back has placed in question. (C1) is the requirement that mattering be vindicable from the detached standpoint. The argument is valid.

\begin{proposition}[The premise denied]\label{p:vind}
Given (C2), Assumption~2 entails the negation of (C1). The premise of the sceptical challenge that Assumption~2 denies is (C1), the requirement that whatever matters be vindicable from the detached standpoint.
\end{proposition}

\begin{proof}
Assumption~2 entails that something has objective worth, and so, by the equivalence noted at \Lm{11.3}, that something matters; this contradicts (C3). Since the argument is valid, Assumption~2 is inconsistent with the conjunction of (C1) and (C2). The text accepts (C2): it is not mentioned in EXI-01-a, but it is proved in EXI-04-a, whose conclusion the family accepts, and Remark~5 does not dispute it. So the premise Assumption~2 denies is (C1).
\end{proof}

\paragraph{Denying and refuting.} A premise is refuted when its falsity is derived from premises that the one who holds it accepts, or that she is rationally required to accept. It is denied when its negation is asserted, or entailed by what is asserted, without such a derivation. Assumption~2 denies (C1): it entails, together with (C2), that (C1) is false, but it does not derive the falsity of (C1) from anything the sceptic accepts, since the sceptic does not accept Assumption~2. Remark~5 is therefore exact. The theorem does not answer the challenge, because the challenge calls into question a premise of the theorem's interpretation; and Assumption~2 denies a premise of the challenge without refuting it, because the denial rests on Assumption~2 itself.

\paragraph{Is the denial question-begging?} A denial that rests on the very claim the challenge calls into question may seem to beg the question. It does not, for two reasons. First, the denial is not offered as an answer to the sceptic; Remark~5 says expressly that the theorem does not answer her. It is offered as a statement of where the text stands, so that the reader can see that the text's conclusion is conditional on rejecting (C1), and can assess (C1) for herself. Second, (C1) has a cost that can be assessed independently of Assumption~2. It is a vindication requirement: it says that a claim of worth is true only if it can be vindicated from a standpoint that presupposes no worth. The analogous requirement in epistemology, that a belief is knowledge only if it can be vindicated from a standpoint that presupposes no belief, leads to scepticism about everything, including the beliefs on which the sceptical argument itself relies, and Nagel's analogy is an analogy in exactly this respect. A requirement that, applied in the one domain, would eliminate all knowledge, and that has no better support in the other, may be rejected without a refutation. MTE-05 makes a parallel point about the argument from queerness: an argument against objective value that would equally show mathematics false proves too much.

\paragraph{The testimony of EXI-04-a.} EXI-04-a gives the same reading of Remark~5. In its reply to the objection that objective worth dissolves the absurd, it says that seriousness is justified if the projects are worthwhile, that the absurd survives because that justification is not one the detached standpoint must accept, and that its definition of the absurd concerns the availability of a forced justification, not the existence of a sound one; and it concludes that this is the sense in which, as the remark at the end of EXI-01-a puts it, the assumption of objective worth denies the premise of the sceptical challenge without refuting it. In the terms of Proposition~\ref{p:vind}: the existence of a sound justification is what Assumption~2 asserts; the availability of a forced one is what (C1) requires and (C2) denies. EXI-04-a and EXI-01-a agree that (C2) is true and (C1) is false, and they agree that the falsity of (C1) is asserted rather than proved.

\subsection{The absurd}\label{s:absurd}

Remark~5 says that a discussion of the absurd is given in the Commentary (\Lm{12}). The analysis of \Lm{11.4} has shown where EXI-01-a stands on Nagel's challenge. We now relate that position to the principal treatments of the absurd and state what follows for meaning.

\paragraph{Nagel.} \citet{Nagel1971} rejected the standard arguments for the absurdity of life: that our lives are short, that we are small in the universe, that what we do will not matter in a million years, that every justification ends somewhere. Each, he argued, either fails to show what it claims or presupposes the sense of absurdity it purports to explain: if our lives are absurd at seventy years, they would be no less absurd if they lasted through eternity; if nothing we do now will matter in a million years, then nothing that will happen in a million years matters now. The real source of the absurd is the collision between the seriousness with which we live and the detachment of which we are capable. His conclusion was that if, \emph{sub specie aeternitatis}, there is no reason to believe that anything matters, then that does not matter either, and we can approach our lives with irony rather than with heroism or despair. EXI-04 takes over the conclusion and gives it a proof: the absurd is ineliminable for any reflective agent (Theorem~3.3 of EXI-04-a), despair and revolt do not fit it (Propositions~3.4 and~3.5), and irony does (Theorem~3.6).

EXI-01-a is consistent with this conclusion, and Proposition~\ref{p:vind} shows exactly how. Nagel's argument establishes (C2), which EXI-01-a accepts. It does not establish (C1), and Nagel's own conclusion does not need it: Nagel says that the detached standpoint gives no reason to believe that anything matters, not that it gives a reason to believe that nothing does. The sceptic who infers (C3) has added (C1), and it is (C1) that Assumption~2 denies. The absurd, in Nagel's sense, and meaning, in the sense of Definition~1, are therefore compatible.

\begin{proposition}[Meaning and absurdity]\label{p:compat}
Every life of positive meaning lived by a reflective agent is absurd in the sense of EXI-04-a, and its absurdity does not reduce its meaning.
\end{proposition}

\begin{proof}
By Theorem~3.3 of EXI-04-a, the life of every reflective agent capable of intentional action is absurd, whatever its content; so a meaningful life of such an agent is absurd. Absurdity consists in the absence of a forced justification of seriousness from the detached standpoint (Definition~2.3 of EXI-04-a). The degree of meaning of a life, by Definition~1, depends on the extent to which it consists in active engagement with projects of objective worth, and neither the extent of engagement nor the worth of its objects depends on whether that worth can be vindicated from the detached standpoint, since Assumption~2, which fixes what worth is, makes no reference to vindication. So absurdity does not reduce meaning.
\end{proof}

\paragraph{Camus.} \citet{Camus1955} located the absurd in the confrontation between the human longing for clarity and meaning and the unreasonable silence of the world, and he took three responses into view: physical suicide, philosophical suicide, the leap of faith by which one escapes the absurd into a meaning beyond it, and revolt, the lucid refusal either to escape or to submit. His Sisyphus, conscious of the futility of his labor and scornful of the gods who imposed it, is the figure of revolt, and Camus concludes that one must imagine him happy. EXI-04-a argues that revolt does not fit the absurd, because it presents the world as an adversary, and an indifferent world is not one (its Proposition~3.5 and Remark~3.7; on indifference, EXI-06-a). EXI-04-b distinguishes the absurdity of a life from the futility of a task: Sisyphus's task is futile, both because it cannot bring about its outcome and because the outcome has no value (its Definition~1.1), while the absurdity of a life is compatible with its projects' succeeding and having worth (its Proposition~1.2).

The distinction is essential to EXI-01-a. Camus's Sisyphus is a figure of the absurd and of futility together, and the image draws its power from the futility. On Definition~1, his life is not meaningful, however lucid and however happy, because the object of his labor has no worth; EXI-04-b makes the same point and notes that Taylor's Sisyphus is in the same position. The human situation, on the analysis of EXI-04, is absurd without being futile: our projects can succeed, and many of their objects have worth. The meaning that Definition~1 describes is meaning that Camus's Sisyphus lacks and that the absurd does not take from us. What remains of Camus's conclusion, EXI-04-b argues, is the claim that full engagement with one's ends is compatible with lucidity about the absence of their vindication (its Proposition~2.1), and that is the ironic response that Proposition~\ref{p:compat} shows to be compatible with meaning.

\paragraph{Tolstoy.} Tolstoy's crisis \citep{Tolstoy1987} has a different structure from Nagel's and Camus's. It is driven not by detachment but by death: the thought that everything he did would come to nothing, since he would die and his works would be forgotten, made his goods taste of nothing, as the honey in the fable no longer tasted sweet to the traveller who saw the dragon and the mice. His resolution was religious: he found in the faith of ordinary people a meaning that death does not destroy, because it relates the finite life to the infinite.

Tolstoy's argument is that what comes to nothing was nothing. EXI-11-a states it as an objection and replies that the inference from transience to worthlessness is invalid: a friendship that ends with death was a friendship, and the relief of a patient's pain was relief while it lasted (its reply to Objection~1). EXI-19-b adds that the limit death sets is not a source of meaning and does not make anything engaged with more or less worthwhile (its Proposition~1.1). On Definition~1, the meaning of a life consists in what the life consists in, and the comment on \Lm{2.4} has drawn the consequence that it is fixed by the life's course and is not undone by what happens after it. Tolstoy's premise, that worth realized at a time is cancelled by later loss, is therefore denied by the constitutive reading of Definition~1, and the denial is supported by the theory of value: on ETH-01, the value of an outcome is fixed by idealized endorsement of it, and idealized endorsement of the relief of pain does not wait on whether the relief is remembered.

There remains something in Tolstoy's case that this reply does not touch. His goods lost their taste: the engagement that had filled his life failed. On Definition~1, his life during the crisis lost meaning, not because its objects lost their worth, but because he could no longer engage with them. This is the structure of Mill's crisis (Section~\ref{s:defence}, D3), and it shows that the thought of death can threaten meaning through the subjective component, by undermining engagement, even if it cannot threaten it through the objective component. EXI-05 holds that fear of death is not rational, though regret at its prospect may be. If the fear is not fitting, its power to undermine engagement is the power of an unfitting attitude, and the remedy is the clarity about what death does and does not take that EXI-02, EXI-05 and EXI-19 together supply. Volume~III, Part~C, develops the point.

\paragraph{Hume's remedy.} Nagel's analogy between the absurd and scepticism has a precedent that EXI-04-a considers as an objection. Hume, at the end of the first book of the \emph{Treatise} (1.4.7), describes how the sceptical reflections that had reduced him to melancholy were dispelled when he dined, played backgammon and conversed with his friends, and how, returning to them afterward, he found them cold and strained. The objection is that forgetting, not irony, is the appropriate response to the absurd. EXI-04-a replies that forgetting is not a response to the absurd but an interruption of the reflection that discloses it, and that the question is what attitude fits the absurd when one attends to it (its reply to Objection~4). For EXI-01-a, Hume's case has a further interest. Hume's return to ordinary life is a return to engagement, with friends, with games and with conversation, and his melancholy was a failure of engagement induced by reflection. The case therefore illustrates, as Tolstoy's and Mill's do, the way in which reflection can threaten meaning through its subjective component without touching its objective one. What reflection showed Hume was that his beliefs could not be vindicated from the standpoint of the sceptic; it did not show that they were false, and his friends, his backgammon and his conversation were no less worth engaging with for his having, for a time, been unable to engage with them.

\paragraph{The point of view of the universe.} Sidgwick held that the good of any one individual is of no more importance, from the point of view of the universe, than the good of any other \citep[III.xiii]{Sidgwick1907}. \citet{Williams1982} criticized the phrase, arguing that there is no such point of view and that ethics cannot take its standards from it. The criticism bears on (C1). If there is no point of view of the universe from which our concerns could be vindicated, the demand that they be so vindicated is not a demand that could be met by anything, and its failure to be met shows nothing about whether anything matters. The detached standpoint of EXI-04-a is not the point of view of the universe; it is a standpoint the reflective subject can occupy. But (C1) requires vindication from that standpoint in its most detached form, and Williams's criticism suggests that the most detached form is not a standpoint from which anything could be vindicated or condemned.

\paragraph{Cosmic insignificance.} \citet{Kahane2014} distinguished the claim that we are insignificant from the cosmic standpoint from the claim that our lives lack meaning, and argued that the first, if true, depends on the claim that we are surrounded by much greater value elsewhere in the universe, not on our smallness. EXI-06 holds that the universe does not care about us, and EXI-06-b argues that its indifference takes nothing from meaning. Both are consistent with Definition~1, which makes meaning depend on the worth of what we engage with, not on how that worth compares with the worth of things elsewhere. A small contribution to a great good is engagement with something of objective worth, and its meaning is not reduced by the existence of greater goods, though its significance, in Kahane's sense, may be.

\paragraph{Summary.} The absurd, in Nagel's sense, is real and ineliminable, and EXI-01-a accepts it. It does not show that nothing matters, because that conclusion requires (C1), which Assumption~2 denies. It does not reduce the meaning of a life, because meaning depends on engagement and worth, not on vindication (Proposition~\ref{p:compat}). It is distinct from futility, which does remove meaning, and Camus's image confuses the two. It is distinct from the threat of death, which can undermine engagement but not worth, and Tolstoy's argument from transience to worthlessness is invalid. The fitting response to it is irony, as EXI-04 holds, and irony is compatible with the full engagement on which meaning depends. Volume~II, Part~C, gives the history of the absurd at length and develops the analysis of (C1).

\lem{12}{A defence of Assumption 2, and a discussion of the absurd, are given in the commentary (EXI-01-b).}

The last sentence refers to the present work. It makes two commitments, and both are met in this volume: the defence of Assumption~2 is Section~\ref{s:defence}, and the discussion of the absurd is the analysis of \Lm{11.4} together with Section~\ref{s:absurd}. Both are extended in Volume~II.

The sentence also tells the reader something about the text. It is a text that knows it is incomplete and says where its completion is to be found. The pattern is that of axiomatic texts with commentaries, in which the axioms are stated without defence and the defence is given separately. It shows that the brevity of EXI-01-a is a choice of form, not a judgment that its assumptions need no defence: the text separates the derivation, which fits on one page, from the defence of its premises, which does not.

\section{The bibliography of EXI-01-a}\label{s:bib}

The bibliography of EXI-01-a has three items, one for each of the three positions that the text must locate itself among: \citet{Wolf2010} for the hybrid view it adopts, \citet{Metz2013} for the objectivism about worth on which it relies, and \citet{Nagel1971} for the challenge it does not answer. The selection is a summary of the text's dialectical situation. Each item is cited once, at the point where the position it represents enters the argument: Wolf at the definition, Metz at the assumption of objective worth, Nagel at the remark on what the theorem leaves open.

\paragraph{Wolf.} The citation is of the Tanner Lectures, not of the earlier paper \citep{Wolf1997} in which the fitting-fulfilment view was first stated. The lectures are the fuller statement, they include the defence against subjectivism on which Definition~1 relies, and they are the version most often cited in the discipline: EXI-17-a, EXI-20-a, EXI-01-c and EXI-01-d all cite them. The Commentary has cited the earlier paper where its formulations differ (\Lm{2.8}).

\paragraph{Metz.} The citation is of Metz's book, which is both the standard survey of theories of meaning and the statement of his own theory. The comment on \Lm{3.5} has noted that Metz is an ally of EXI-01-a on the existence of objective worth and a rival on the measure of meaning. A reader who wants Metz's survey of the debate will also find his classification of the views from which EXI-01-a distinguishes itself, and the Commentary's introduction follows that classification.

\paragraph{Nagel.} The citation is of the 1971 paper, not of \emph{The View from Nowhere}. The paper is where the sense of the absurd is analyzed as a collision between seriousness and detachment and where the ironic response is proposed, and it is the paper that EXI-04-a reconstructs. The book generalizes the analysis and adds a treatment of death; the Commentary cites it at \Lm{11.2} and Volume~II examines the relation between the two.

\paragraph{What is not cited.} Four works that the Commentary has found indispensable for understanding EXI-01-a are absent from its bibliography. Taylor's chapter on the meaning of life \citep[ch.~18]{Taylor1970} is the source of the subjectivist view that the qualifier \enquote{of objective worth} excludes, and of the figure of Sisyphus that the discipline uses throughout; EXI-17-a cites it. Nozick's chapter \citep[ch.~6]{Nozick1981} is the principal statement of the view that meaning requires connection with value beyond oneself; Definition~1 secures connection through objective worth, and EXI-17-a's reply to its Objection~2 explains why connection does not require success. Camus's essay \citep{Camus1955} and Tolstoy's \emph{Confession} \citep{Tolstoy1987} are the classical texts of the absurd and of the crisis of meaning. Their absence from a one-page text is not a judgment of their importance, and the Commentary has supplied them where they bear on the lemmata.

\section{The Lean module}\label{s:lean}

Theorem~4 is formalized in the Lean~4 module \texttt{lean/Philosophy/EXI/EXI\_01.lean} \citep{deMouraUllrich2021}. The module is short, and we comment on every declaration in order, quoting each exactly. The quotations are checked against the module by the program of Appendix~\ref{app:conc}. The module uses only the core of Lean, with no external library, and contains no unproved step.

\subsection{The header}

\begin{lstlisting}
/-
EXI-01 · On the Meaning of Life
Formalization of Theorem 4.
\end{lstlisting}

The header names the family and the text and says what is formalized: Theorem~4, not the whole of EXI-01-a. Definition~1 is formalized only in its positive case, and Assumption~2 and Remark~5 are not formalized at all.

\begin{lstlisting}
What this file establishes. Given Definition 1 and Assumption 3, a life of
positive meaning is available to every agent. The proof is immediate; the
substance of the manuscript lies in its assumptions, which are discussed in
the commentary (EXI-01-b).
-/
\end{lstlisting}

The header makes three claims. The first, that the module establishes Theorem~4 given Definition~1 and Assumption~3, omits the bridge principle, which the module states as an axiom; the omission is in the header and not in the module, and the \texttt{\#print axioms} output corrects it. The second, that the proof is immediate, is accurate: it has three steps, one for each sentence of the prose proof. The third, that the substance lies in the assumptions, is the formal counterpart of Remark~5, and it is the reason for the present work. The header refers to the Commentary, as Remark~5 does.

\subsection{The primitives}

\begin{lstlisting}
namespace Philosophy.EXI_01
\end{lstlisting}

The namespace follows the convention of the collection's modules: one namespace per family.

\begin{lstlisting}
/-- Agents capable of intentional action. -/
axiom Agent : Type
\end{lstlisting}

The domain of Assumption~3 and Theorem~4 is a type. The docstring carries the restriction of \Lm{4.2}: the type is not of all beings, or of all persons, but of agents capable of intentional action. Nothing in the formal system enforces this; the restriction lives in the interpretation. The type may be empty, and Proposition~\ref{p:indep} uses that possibility.

\begin{lstlisting}
/-- Projects. -/
axiom Project : Type
\end{lstlisting}

Projects form a type without internal structure. The module does not distinguish project types from token undertakings, or projects in the narrow sense from projects in the broad sense; the comments on \Lm{2.6} and \Lm{3.1} have argued for the broad reading and the type reading, and both are consistent with the module.

\begin{lstlisting}
/-- Lives available to an agent. -/
axiom Life : Agent → Type
\end{lstlisting}

This is the most interesting of the primitives. \texttt{Life} is not a type but a family of types indexed by agents: for each agent \texttt{a}, \texttt{Life a} is the type of lives available to \texttt{a}. Two features of the conception of a life discussed at \Lm{2.1} are built in. A life is someone's life, since there are no lives except as elements of \texttt{Life a} for some \texttt{a}. And the lives of an agent are those \emph{available} to her: the docstring identifies the elements of \texttt{Life a} with the possibilities open to the agent, so that existential quantification over \texttt{Life a} expresses availability. The modal content of \enquote{is available} in \Lm{5.3} is thereby represented without a modal operator, by quantification over a type of possibilities. The indexing of availability to times, which the comment on \Lm{5.3} argued for, is not represented; a time-indexed formalization would replace \texttt{Life a} by a type of continuations of the agent's life from a given time.

\begin{lstlisting}
/-- `ObjectivelyWorthy p`: the value of `p` does not depend on the attitudes
of those engaged in it (Assumption 2 asserts that such projects exist). -/
axiom ObjectivelyWorthy : Project → Prop
\end{lstlisting}

The predicate of objective worth is primitive, and its docstring gives the colon clause of Assumption~2 as its intended interpretation. The docstring is the only place in the module where Assumption~2 appears, and it appears as a gloss on a predicate, not as an axiom. This is the formal counterpart of the first finding of Section~\ref{s:ass2}: Assumption~2's role is to interpret, not to derive. The formal system does not know what the predicate means, and the theorem holds under every interpretation of it that satisfies the axioms (Proposition~\ref{p:neutralfml}).

\begin{lstlisting}
/-- `CanEngage a p`: agent `a` is able to engage actively in project `p`. -/
axiom CanEngage : (a : Agent) → Project → Prop
\end{lstlisting}

The ability of \Lm{4.4} is a primitive relation between agents and projects. The docstring says \enquote{engage actively}, which settles in the module the question left open in the text, whether the ability of Assumption~3 is an ability to engage actively (\Lm{7.1}). The binder \texttt{(a : Agent)} is written in dependent form, though the result type does not depend on \texttt{a}; it is equivalent to \texttt{Agent} $\to$ \texttt{Project} $\to$ \texttt{Prop}.

\begin{lstlisting}
/-- `Includes l p`: life `l` includes active engagement with `p`. -/
axiom Includes : {a : Agent} → Life a → Project → Prop
\end{lstlisting}

The relation between a life and a project that the prose expresses by \enquote{a life that includes engagement with $p$} (\Lm{7.1}) is primitive, and its docstring again says \enquote{active engagement}. The agent is an implicit argument, inferred from the life. The relation is binary in the life and the project and carries no measure: it records that a life includes active engagement with a project, not how much of the life it occupies.

\subsection{The premises}

\begin{lstlisting}
/-- An agent who can engage in a project has a life that includes it. -/
axiom engagement_realizable :
  ∀ (a : Agent) (p : Project), CanEngage a p → ∃ l : Life a, Includes l p
\end{lstlisting}

This is the bridge principle of Proposition~\ref{p:bridge}, stated as an axiom. It does not correspond to any numbered statement of EXI-01-a, and its docstring cites none. The comment on \Lm{7.1} argued that it is properly regarded as a clarification of \enquote{available} rather than as a substantive premise, since on the reading of \enquote{able} adopted there, what an agent is able to do from her situation is what some life available to her includes. In the module, where ability and availability are independent primitives, the connection must be an axiom, and its presence in the axiom list is the clearest case in the collection of a formalization exposing a step that the prose takes for granted.

\begin{lstlisting}
/-- Definition 1, positive case. A life has positive meaning if it includes
active engagement with some objectively worthy project. -/
def PositivelyMeaningful {a : Agent} (l : Life a) : Prop :=
  ∃ p, ObjectivelyWorthy p ∧ Includes l p
\end{lstlisting}

Definition~1 is formalized as a definition, not an axiom, and only in its positive case. Three features deserve comment.

First, the degree variable is gone. The formal definition is of a one-place predicate, positive meaning, not of the two-place predicate of \Lm{2.2}. This is sufficient for Theorem~4, which concerns only positive meaning, and it is the reason the header says that Definition~1 is formalized in its positive case.

Second, the step from inclusion to positive extent (\Lm{7.2}) is absorbed into the definition. In the prose, a life that includes engagement with a worthy project is inferred to consist in such engagement to a positive extent, and then, by Definition~1, to be meaningful to a positive degree. In the module, a life that includes active engagement with a worthy project is positively meaningful by definition. The condition of Proposition~\ref{p:neutral}, that the engagement be sustained over a non-null part of the life, is thereby assumed in the interpretation of \texttt{Includes}: a life \enquote{includes} active engagement with a project, in the sense of the module, only if the engagement occupies a non-null part of it.

Third, only the right-to-left direction of Definition~1 is captured, though in a definition that makes it a biconditional for positive meaning. The module cannot express the comparative claims of Definition~1, or the claim of \Lm{10} that a life with no engagement with worthy projects has no meaning, except in the trivial form that such a life is not \texttt{PositivelyMeaningful}.

\begin{lstlisting}
/-- Assumption 3. Every agent can engage in some objectively worthy project. -/
axiom assumption_3 : ∀ a : Agent, ∃ p, ObjectivelyWorthy p ∧ CanEngage a p
\end{lstlisting}

Assumption~3 is the one axiom of the module that corresponds to a numbered statement of EXI-01-a, and its form is exactly the form of \Lm{4}: universal over agents, existential over projects, as Proposition~\ref{p:order} requires.

\subsection{The theorem}

\begin{lstlisting}
/-- Theorem 4. For every agent, a life of positive meaning is available. -/
theorem meaning_available (a : Agent) : ∃ l : Life a, PositivelyMeaningful l := by
  obtain ⟨p, hw, hc⟩ := assumption_3 a
  obtain ⟨l, hl⟩ := engagement_realizable a p hc
  exact ⟨l, p, hw, hl⟩
\end{lstlisting}

The statement is Theorem~4: for an arbitrary agent \texttt{a}, there is a life in \texttt{Life a} that is positively meaningful. The proof has three lines, which correspond to the three sentences of the prose proof with the bridge made explicit. The first line instantiates Assumption~3 at \texttt{a} and names the project \texttt{p}, its worth \texttt{hw} and the agent's ability \texttt{hc}; this is \Lm{6}. The second line applies the bridge to obtain a life \texttt{l} and the fact \texttt{hl} that it includes engagement with \texttt{p}; this is the existence that \Lm{7.1} assumes. The third line assembles the witness: \texttt{l} is a life of \texttt{a}, and \texttt{p}, \texttt{hw} and \texttt{hl} show that it is positively meaningful; this is \Lm{7.2} and \Lm{8}, combined by the formal definition.

\begin{lstlisting}
end Philosophy.EXI_01

#print axioms Philosophy.EXI_01.meaning_available
\end{lstlisting}

The checker's output for the last command is:

\begin{lstlisting}
'Philosophy.EXI_01.meaning_available' depends on axioms:
  [Philosophy.EXI_01.Agent, Philosophy.EXI_01.CanEngage,
   Philosophy.EXI_01.Includes, Philosophy.EXI_01.Life,
   Philosophy.EXI_01.ObjectivelyWorthy, Philosophy.EXI_01.Project,
   Philosophy.EXI_01.assumption_3,
   Philosophy.EXI_01.engagement_realizable]
\end{lstlisting}

Six of the eight axioms are primitives: types and relations without content. Two are premises: \texttt{assumption\_3}, which is Assumption~3, and \texttt{engagement\_realizable}, which is the bridge. None of Lean's own logical axioms (propositional extensionality, quotient soundness, choice) appears. The proof is constructive: it extracts the witness from the premises without appeal to excluded middle, which is fitting for an existence claim whose philosophical point is that the witness can be found, since a life of positive meaning is one the agent can lead by engaging with a project she can identify.

\subsection{What the formalization shows}

\begin{proposition}[Neutrality of the formal skeleton]\label{p:neutralfml}
The formal theorem \texttt{meaning\_available} holds under every interpretation of the primitives that satisfies the two premises of the module, \texttt{assumption\_3} and the bridge axiom. In particular, it holds under interpretations on which \texttt{ObjectivelyWorthy} is true of every project.
\end{proposition}

\begin{proof}
The theorem is derived from the two premises by the rules of Lean's logic, and the derivation is valid under every interpretation of the primitives. If \texttt{ObjectivelyWorthy} is true of every project, the theorem holds under any interpretation that also satisfies the premises, as every derived theorem does.
\end{proof}

The proposition shows what the formalization does not do. It does not distinguish the hybrid view from the subjective view: a subjectivist who read \texttt{ObjectivelyWorthy} as \enquote{is a project at all} would find the module's theorem true and the module's axioms acceptable. The distinction between EXI-01-a and Taylor lies entirely in the docstring of \texttt{ObjectivelyWorthy}, that is, in Assumption~2's gloss. This is not a defect of the formalization, which formalizes the logical skeleton of the argument, as it should; it is a precise statement of where the philosophy is.

\begin{proposition}[Independence of Assumption~2]\label{p:indep}
The existential part of Assumption~2, $\exists p$, \texttt{ObjectivelyWorthy p}, is not derivable from the axioms of the module.
\end{proposition}

\begin{proof}
Interpret \texttt{Agent} as an empty type, \texttt{Project} as a type with one element, \texttt{Life} as the function assigning the empty type to every agent (vacuously, since there are none), and \texttt{ObjectivelyWorthy}, \texttt{CanEngage} and \texttt{Includes} as everywhere false. \texttt{assumption\_3} and \texttt{engagement\_realizable} quantify over agents and are vacuously true. $\exists p$, \texttt{ObjectivelyWorthy p} is false. So the axioms do not entail it.
\end{proof}

Proposition~\ref{p:indep} is the formal counterpart of the second part of Proposition~\ref{p:redundant}: Assumption~3 entails Assumption~2 only if there is an agent. In a world without agents, Assumption~3 is true and Assumption~2 may be false. The text's decision to state Assumption~2 separately is therefore not formally idle: it is the one premise of the text that says something about worth whether or not anyone exists to engage with it.

\begin{proposition}[Consistency]\label{p:consistent}
The axioms of the module are jointly satisfiable, and in a model in which there is at least one agent.
\end{proposition}

\begin{proof}
Interpret \texttt{Agent}, \texttt{Project} and every \texttt{Life a} as one-element types, and \texttt{ObjectivelyWorthy}, \texttt{CanEngage} and \texttt{Includes} as everywhere true. Both premises then hold, and so does \texttt{meaning\_available}.
\end{proof}

The consistency of the axioms is cheap, as it should be for a module whose axioms are few and weak. What is not cheap is the truth of the axioms under their intended interpretation, and that, as the header says, is what the Commentary has discussed.

\paragraph{What a fuller formalization would add.} A formalization of the whole of EXI-01-a would need a type of degrees with a distinguished zero and an order, a measure assigning degrees to lives, and an axiom connecting inclusion of sustained engagement with positive measure; it could then state Definition~1 in full, prove Proposition~\ref{p:noneg} and Proposition~\ref{p:neutral}, and express the comparative claims of EXI-01-d. It would need a predicate for vindication from the detached standpoint to state the reconstruction of the sceptical challenge and prove Proposition~\ref{p:vind}. Each of these would add premises, not remove them, and each would make the module's dependence on its interpretation more visible, not less. Volume~III, Part~D, states the measure-theoretic version of Definition~1 and proves the propositions of Volume~I that concern degree under it.

\section{The argument as a whole}\label{s:whole}

The commentary on the lemmata has examined the text sentence by sentence. This section considers it as a whole: the structure of the argument, the readings of the whole that the text admits, the objections that are not objections to any one lemma, and what the text, with the Commentary, does not establish.

\subsection{Structure}

The argument of EXI-01-a, as the Commentary has reconstructed it, has the following structure.

\begin{enumerate}[label=(\roman*)]
\item Definition~1 says what meaning is: the extent to which a life consists in active engagement with projects of objective worth. It is a real definition, defended by Wolf, read conatively in its subjective component (\Lm{2.5}), broadly in its notion of a project (Proposition~\ref{p:broad}), and with a measure left unspecified but constrained (Propositions~\ref{p:noneg} and~\ref{p:neutral}).
\item Assumption~2 says what objective worth is, in a restricted sense (Proposition~\ref{p:restricted}), and that there is some. It plays no role in the derivation (Proposition~\ref{p:redundant}) and is independent of the other premises (Proposition~\ref{p:indep}); its role is to interpret the theorem and to fix the text's position on the sceptical challenge. It is defended in Section~\ref{s:defence} and entailed by the theory of value of ETH-01 (Proposition~\ref{p:eth01}).
\item Assumption~3 says that every agent capable of intentional action can engage in some project of objective worth. It carries the weight of the availability claim, it is derived in EXI-01-c from a premise about the reach of worthwhile objects, and it is supported by EXI-17-a's demonstration that meaning does not require success.
\item The bridge principle says that what an agent can engage in, some life available to her includes. It is unstated in the text, explicit in the Lean module, and analytic on the reading of ability adopted here (Proposition~\ref{p:bridge}).
\item Theorem~4 follows from (i), (iii) and (iv): a life of positive meaning is available to every agent, at every time at which she is one (\Lm{5.3}).
\item Remark~5 says what the theorem answers and what it does not. The sceptical challenge is not answered; its vindication premise is denied by Assumption~2 (Proposition~\ref{p:vind}). The absurd in Nagel's sense survives and is compatible with meaning (Proposition~\ref{p:compat}).
\end{enumerate}

\subsection{Readings of the whole}

Three readings of the text as a whole should be distinguished.

\paragraph{The deflationary reading.} On this reading, EXI-01-a is a trivial text. Its theorem follows in three lines from assumptions that say almost the same thing, and the module's header concedes that the proof is immediate. The text adds nothing to Wolf except the claim of availability, and that claim is purchased by an assumption that simply asserts it.

The reading is mistaken, though it starts from a truth. It is true that the derivation is short. But the shortness of a derivation is not the triviality of a text, and the deflationary reading overlooks three things the text does. It separates the availability claim from the definition of meaning, which Wolf's lectures do not; it isolates the assumption on which availability depends, which makes the claim testable and shows exactly what a critic must deny; and it locates its own position on the sceptical challenge with a precision that the literature on Nagel rarely achieves. The comments on the lemmata have shown that each of these is substantive.

\paragraph{The programmatic reading.} On this reading, EXI-01-a is the statement of a position whose defence is distributed: the definition is defended by Wolf, the existence of objective worth by Metz and by ETH-01, the availability claim by EXI-01-c and EXI-17-a, its consequences by EXI-01-d, its relation to the absurd by EXI-04 and by this Commentary, its relation to death by EXI-02, EXI-03 and EXI-19, its relation to theism by EXI-11. The page states the position in its barest form, and the family and the discipline supply the rest.

We think this is the correct reading. It explains the text's form: an axiomatic statement is the natural form for a position whose parts are defended elsewhere, since it shows how the parts fit together. It explains the citations, which are not of authorities for the claims but of the places where the claims are defended. And it explains Remark~5, which reads as a guide to where the defence of the text is to be found.

\paragraph{The dialectical reading.} On this reading, the purpose of EXI-01-a is to fix the terms of debate: to say what a theory of meaning must contain (a definition, an assumption about value, an assumption about reach) and to show where each rival theory disagrees. The subjectivist denies Assumption~2's gloss; the supernaturalist denies that Assumption~2 can be true without God; the sceptic of the absurd denies Assumption~2 itself, or demands its vindication; the pessimist denies Assumption~3 for some agents. The reading is compatible with the programmatic one, and the Commentary has used it in identifying, at each lemma, the rival positions that deny it.

\subsection{The family read with the lemmata}\label{s:family}

The two other manuscripts of the family were written after EXI-01-a and with it in view. Read with the lemmata, they show which parts of the text the family regarded as needing support and how it supplied it.

\paragraph{EXI-01-c.} The alternative proof keeps Definition~1 and Assumption~2 and replaces Assumption~3 (\Lm{4}). Its Definition~2.1 introduces objects of engagement, which include persons, living creatures, practices, inquiries and parts of the world as objects of understanding, and calls them worthwhile if they have or serve attitude-independent value. Its Assumption~2.2 (Reach) says that every agent capable of intentional action can direct caring activity at such an object. Its Lemma~2.3 says that sustained caring activity directed at a worthwhile object is a project of objective worth in the sense of Definition~1. Its Theorem~2.4 is Theorem~4.

Three lemmata of EXI-01-a are affected. At \Lm{2.6}, EXI-01-c's Lemma~2.3 is the argument for the broad reading of \enquote{projects} (Proposition~\ref{p:broad}). At \Lm{4.3}, the quantifier over projects is replaced by a quantifier over objects, and Assumption~3 becomes a corollary. At \Lm{4.4}, \enquote{able to engage} is analyzed as \enquote{able to direct caring activity at}, which is the conative reading of engagement defended at \Lm{2.5}. The alternative proof does not affect the bridge principle (Proposition~\ref{p:bridge}), which it also uses without stating, or the role of Assumption~2, which it also does not use in its derivation. The Lean module could be adapted to it by replacing the axiom \texttt{assumption\_3} with two axioms, one for Reach and one for the lemma, and deriving \texttt{assumption\_3} as a theorem; the axiom list of \texttt{meaning\_available} would then contain Reach, the lemma and the bridge.

\paragraph{EXI-01-d.} The consequence manuscript draws four propositions from the availability claim. Its Proposition~2.1 (independence from outcome) is EXI-17-a's Corollary~3.6, and bears on \Lm{5.3}. Its Proposition~2.2 (independence from talent and scale) bears on \Lm{2.7}: it uses the degree reading of objective worth, on which scale affects degree but not presence. Its Proposition~2.3 (no condemnation by circumstance) is Theorem~4 restated as a denial, with the strengthening of EXI-01-c. Its Proposition~2.4 (degree, not possibility) bears on \Lm{2.3} and \Lm{10}: circumstances affect the degree of meaning attainable but not its availability. Its section on what availability does not imply corresponds to Remark~5's second sentence (\Lm{10}). Its reply on the absurd repeats Remark~5's third sentence and cites EXI-04.

\paragraph{The division of labor.} Read together, the three manuscripts divide the family conclusion as follows. Its first half, that life is meaningful to the extent that one engages with projects of objective worth, is Definition~1, with its objective component interpreted by Assumption~2 and its project component interpreted by EXI-01-c's Lemma~2.3. Its second half, that meaning so understood is available to everyone, is Theorem~4, derived in EXI-01-a from Assumption~3 and in EXI-01-c from the weaker-seeming Reach, and its consequences and limits are drawn in EXI-01-d. The present work supplies what Remark~5 says is missing: the defence of Assumption~2 and the discussion of the absurd.

\subsection{Objections to the whole}

\begin{objection}[Circularity]\label{o:circ}
The availability of meaning is secured by stipulation. Definition~1 is tailored so that the conditions of meaning are within everyone's reach, and Assumption~3 then says that they are. A theory of meaning should be tested against our judgments of which lives are meaningful, not designed to deliver a desired distribution.
\end{objection}

\begin{reply}
Definition~1 was not designed to deliver availability. It is Wolf's definition, defended by her on the basis of judgments about cases, and Wolf does not argue for availability. The availability claim is a consequence, derived in the family by a separate argument (EXI-01-c) from a separate premise. That the definition, so supported, has this consequence is a discovery about it, and a welcome one; it is not a stipulation. The test of the definition against cases is the one the comments on \Lm{2.5} and \Lm{2.7} have applied, and it is passed independently of availability.
\end{reply}

\begin{objection}[Meaning and morality]\label{o:gauguin}
Williams's painter abandons his family to pursue his art \citep{Williams1981b}. On Definition~1, his life gains meaning from his engagement with painting, which has objective worth, while the wrong he does his family is irrelevant to meaning. A theory on which a life can be made more meaningful by a moral wrong is defective.
\end{objection}

\begin{reply}
The objection correctly reports what Definition~1 implies, and the implication is correct. Wolf argued that meaning is a value distinct from morality, and the painter's case is a case in which the two come apart: his life gains meaning from his engagement with painting and is morally worse for the abandonment. To say that his life is more meaningful is not to say that it is better all things considered, or that the abandonment was justified, or that anyone should admire it. A theory that refused to recognize the meaning of his engagement because of the wrong would be confusing meaning with morality. The case also illustrates EXI-17-a's thesis. Williams argued that whether the painter's choice was justified depends on whether he succeeds; on Definition~1, the meaning of his life does not depend on it. What his success affects is the justification of his choice and the value his work adds to the world, not the extent to which his life consisted in engagement with a worthwhile project.
\end{reply}

\begin{objection}[Shape]\label{o:shape}
The meaning of a life depends on its shape: on whether it progresses, whether its later parts redeem its earlier failures, whether it ends well \citep{Velleman1991,Kauppinen2012}. Definition~1 measures meaning by extent, which is indifferent to order, and so cannot be the right account.
\end{objection}

\begin{reply}
Definition~1 does not say that extent is indifferent to order; the comment on \Lm{2.3} has argued that the measure is left open, and Proposition~\ref{p:neutral} that the availability claim is independent of it. EXI-19 holds that the value of a life as a whole depends in part on its shape (Principle~P1 of EXI-19-a), and EXI-19-b that finitude is a condition of the meaning of a life as a whole, which depends on its shape (its Proposition~1.2). The view that shape matters to meaning is therefore part of the discipline, and it is compatible with Definition~1 on any measure on which the extent to which a life consists in engagement is assessed over the whole and not summed over moments. What would be incompatible with Definition~1 is the view that shape could confer meaning on a life that contains no engagement with anything of worth, and that view has no defenders: a life that progresses steadily toward the counting of more blades of grass has a shape but no meaning.
\end{reply}

\begin{objection}[Transcendence]\label{o:transc}
\citet{Nozick1981} argued that a life is meaningful only through connection with something beyond itself, and that the question of meaning recurs for whatever it connects with until it reaches something unlimited. Definition~1 stops the regress arbitrarily at projects of objective worth.
\end{objection}

\begin{reply}
The regress arises only if the worth of what one connects with must itself be explained by connection with something further. On Assumption~2, the worth of a project is not conferred by its connection with anything; it belongs to the project, or to the object it serves, in virtue of what it is. The regress stops where worth is found, and it is found in the care of persons, in understanding and in the other components of the good. EXI-17-a accepts that meaning requires connection with value beyond oneself (its reply to Objection~2), and Definition~1 secures it through the objective component. What Nozick's regress adds is the demand that the endpoint be unlimited, and that is a demand for the kind of meaning that EXI-11-a says a theistic framework would add (its Proposition~4.1), not a condition on meaning as such.
\end{reply}

\begin{objection}[Theism]\label{o:theism}
If REL-01 is right that a maximally great being exists, the meaning of a life consists in its relation to that being, and Definition~1 describes at best a secondary meaning.
\end{objection}

\begin{reply}
EXI-11 holds that meaning does not require God, though a theistic framework would add to it. The addition is accommodated by Definition~1's measure: engagement with a being of maximal worth would be engagement with the greatest object of worth there is, and by the principle that engagement with what is of greater worth makes a life more meaningful (Principle~P2 of EXI-11-a), it would add to the degree of meaning. It would not change what meaning is. If REL-01 is right, the additional meaning is available; if it is wrong, the meaning Definition~1 describes is all there is, and it is available to everyone. Either way, the conclusion of EXI-01 stands. Volume~II, Part~D, develops the point.
\end{reply}

\subsection{What the text does not establish}

EXI-01-a, read with the Commentary, does not establish five things that might be expected of a text with its title.

It does not say what human life as such is for. The comment on \Lm{1} has explained why the second question, what makes a life meaningful, does not depend on the answer to the first.

It does not say which projects have objective worth. It relies on the theory of the good of ETH-01 for that, and its availability claim depends only on there being, within every agent's reach, at least one.

It does not say how meaningful any particular life is, or how meaningful a life any particular agent can attain. These depend on what is done with the life and on the circumstances in which it is done, and Remark~5 sets them aside.

It does not refute the sceptic who demands that worth be vindicated from the detached standpoint. It denies her vindication premise and accepts the absurdity that her challenge reveals, and it holds that the absurdity is compatible with meaning.

It does not settle whether beings who are not agents capable of intentional action can have meaningful lives in some other sense. Its availability claim is a claim about agents.

\subsection{Conclusion}

EXI-01-a states, on one page, that life is meaningful to the extent that one engages with projects of objective worth, and that meaning in this sense is available to everyone. The Commentary has shown that the page is exact. Its definition is Wolf's, read so as to include love, receptive engagement and the work of those whose engagement is not felt as absorption. Its notion of objective worth is restricted, requiring independence only from the evaluative attitudes of those engaged, and it is entailed by the theory of the good on which the collection rests. Its theorem follows from its third assumption and a bridge principle, not from its second, whose role is to fix what the theorem means and where the text stands on the sceptical challenge. Its remark identifies the premise of that challenge that the text denies and correctly declines to claim a refutation. And its conclusion is compatible with the absurdity of life in Nagel's sense, which is real but does not touch meaning. Meaning, on this account, is the meeting of engagement and worth, and since engagement is within every agent's power and worth is present in every human situation, it is available to every agent who can engage, for as long as she can.


\appendix
\section{Concordance}\label{app:conc}

\subsection*{Preface}

\paragraph{What is enumerated.} This appendix is a concordance of EXI-01-a in the corpus of its discipline. Its object is every word type that occurs in the title and the eleven sentences of EXI-01-a, \cnTypes{} types in all, of which \cnContent{} are content words and the rest function words. Its corpus is the source text of the \cnOtherMss{} other manuscripts of the discipline (EXI-01-c to EXI-20-b; the present work is excluded), \cnCorpusWords{} words after the removal of markup. It has three parts. Table~A.1 lists every word type of EXI-01-a with the lemmata in which it occurs, its number of occurrences in the corpus, the number of manuscripts in which it occurs, and, for content words, its number of occurrences in each manuscript. Table~A.2 counts fifteen term families, defined by the patterns listed below, in every manuscript including EXI-01-a. Section~A.3 is a keyword-in-context concordance: every occurrence in the corpus of every content word of EXI-01-a, \cnKwic{} lines in all, with forty characters of context on each side, the manuscript and the section in which it occurs.

\paragraph{The program and the encoding.} The program \texttt{gen/generate.py} (Python~3, standard library only, deterministic) reads each manuscript's \texttt{main.tex} and \texttt{main.aux}. It keeps the text between the beginning and end of the document, removes the bibliography, code listings, comments and mathematics, resolves cross-references to their numbers using the manuscript's own \texttt{.aux} file and citations to author and year using its own bibliography, keeps the arguments of text commands such as \texttt{\textbackslash emph}, and removes all other commands. It divides the text into sections, numbered as in the manuscript (\enquote{Abs} for the abstract, letters for appendices, 0 for material before the first section), and the sections into sentences. A word is a maximal string of letters, possibly with internal apostrophes or hyphens; identifiers of manuscripts count as single words. Words are compared after conversion to lower case. For EXI-01-a, the program uses the lemmata of Section~\ref{s:text}, with citations removed, so that the occurrences in EXI-01-a are recorded by lemma number. The function words, listed in the program, are given totals only; five words that name statements of EXI-01-a or the present work (\emph{definition}, \emph{assumption}, \emph{theorem}, \emph{commentary} and the identifier of this work) are indexed in Table~A.1 but have no lines in A.3, since their occurrences in the corpus are references to statements rather than uses of terms. The families of Table~A.2 are defined by these patterns, matched against whole words: \textit{mean-} (\emph{mean}, \emph{means}, \emph{meaning}, \emph{meanings}, \emph{meaningful}, \emph{meaningfully}, \emph{meaningfulness}, \emph{meaningless}, \emph{meaningly}), \textit{life-} (\emph{life}, \emph{lives}, \emph{lived}, \emph{living}), \textit{engag-} (every word beginning \emph{engag}), \textit{activ-} (\emph{active}, \emph{actively}, \emph{activity}, \emph{activities}), \textit{project-} (\emph{project}, \emph{projects}), \textit{worth-} (\emph{worth}, \emph{worthwhile}, \emph{worthy}, \emph{worthiness}), \textit{objectiv-} (\emph{objective}, \emph{objectively}, \emph{objectivity}), \textit{valu-} (every word beginning \emph{valu}), \textit{attitud-} (\emph{attitude}, \emph{attitudes}), \textit{agent-} (\emph{agent}, \emph{agents}, \emph{agency}), \textit{availab-} (every word beginning \emph{availab}), \textit{degree-} (\emph{degree}, \emph{degrees}, \emph{extent}), \textit{matter-} (\emph{matter}, \emph{matters}, \emph{mattered}, \emph{mattering}), \textit{absurd-} (every word beginning \emph{absurd}) and \textit{sceptic-} (every word beginning \emph{sceptic} or \emph{skeptic}).

\paragraph{Checks.} Before writing the tables, the program checks that every lemma in its list occurs verbatim in EXI-01-a (the title against the title), that every phrase lemma occurs in its sentence lemma, that the sentence lemmata occur in the order of the text, that every lemma quoted in the chapters of this volume is quoted exactly and under its correct number, that every lemma in the list is quoted at least once, and that every listing of Lean code in Section~\ref{s:lean} occurs in the module. It halts if any check fails. The counts are exhaustive for the corpus as defined: every word of every sentence of every manuscript of the discipline other than this one is examined. Two limits should be stated. Text in generated tables included by a manuscript from separate files, which occurs only in EXI-10-b, is not part of the corpus, since it summarizes the work of other authors rather than the claims of the discipline. And the segmentation into sentences is mechanical, so that a line of A.3 may occasionally begin inside a sentence that contains an abbreviation.

\paragraph{Summary.} Of the \cnContent{} content words of EXI-01-a, \cnUnshared{} do not occur in the corpus: \cnUnsharedList. Three of these are the phrase \emph{sub specie aeternitatis}, which no other manuscript of the discipline uses, though EXI-04-a analyzes the standpoint it names under the name \enquote{detached standpoint}. The content words with the most concordance lines are \cnTopKwic. In the corpus, the family \textit{mean-} occurs \cnFammean{} times, \textit{engag-} \cnFamengag{} times, \textit{worth-} \cnFamworth{} times and \textit{objectiv-} only \cnFamobjectiv{} times.

\paragraph{Entries discussed.} Five entries bear on the commentary. (i) The entry \emph{objective} in Table~A.1 records \cnFamobjectiv{} occurrences of the family in the corpus against \cnFamworth{} of \textit{worth-}: the discipline speaks of the objective component of Definition~1 mainly through \emph{worth} and \emph{worthwhile}, the vocabulary of EXI-17-a's Definition~2.2, and rarely through the word \emph{objective}. The lines for \emph{objective} in A.3 show that most of its occurrences are in EXI-01-c, EXI-01-d and EXI-04-a, the manuscripts that discuss Assumption~2 directly. (ii) The entry \emph{project} records its occurrences as concentrated in EXI-01-c, EXI-03-a and EXI-03-b and in EXI-09-a and EXI-09-b. The lines in A.3 for EXI-09-a and EXI-09-b show the word in the sense of the defining project, which the comment on \Lm{2.6} distinguished from the projects of Definition~1; the lines for EXI-03-a show it in the sense of the projects for which categorical desires are desires, which is the sense of Definition~1. (iii) The entry \emph{matters} has fourteen of its occurrences in EXI-03-a, where the lines in A.3 show the phrase \enquote{what matters in survival}, a sense unrelated to the \enquote{nothing matters} of \Lm{11.3}; the concordance thus separates two senses that a count of the word would conflate. (iv) The entry \emph{sceptical} records two occurrences: one in EXI-04-a, in the sentence that quotes Remark~5 (see Appendix~\ref{app:test}), and one in EXI-08-a. The second concerns arguments about authenticity and is unrelated to the challenge of \Lm{11}. (v) The entry \emph{everyone} has five of its occurrences in EXI-01-d and the rest in EXI-02-a, EXI-07-a, EXI-11-a, EXI-17-a and EXI-20-a; the lines in A.3 for the last three are statements of the availability claim of EXI-01, and they show that the discipline uses the family conclusion's word, not Theorem~4's \enquote{every such agent}, when it reports the claim.

\subsection*{A.1 Index of the words of EXI-01-a}

{\footnotesize Columns: the word; the sentence lemmata of EXI-01-a in which it occurs (with multiplicity); its total occurrences in the other \cnOtherMss{} manuscripts; the number of those manuscripts in which it occurs; its occurrences by manuscript, with the prefix \enquote{EXI-} omitted.\par}

{\scriptsize% Generated by gen/generate.py; do not edit.
\begin{longtable}{@{}>{\raggedright}p{0.13\textwidth}>{\raggedright}p{0.1\textwidth}r r>{\raggedright\arraybackslash}p{0.55\textwidth}@{}}
\toprule Word & Lemmata & Occ. & Mss & Occurrences in the other EXI manuscripts \\ \midrule \endhead
\bottomrule \endfoot
\textbf{EXI-01-b} & 12 & 0 & 0 & --- \\
\textbf{a} & 2, 5, 6, 7\,(\(\times\)2), 8, 12\,(\(\times\)2) & 1692 & 36 & \emph{function word: totals only} \\
\textbf{able} & 4 & 2 & 1 & 01-c\,(2) \\
\textbf{absurd} & 12 & 46 & 5 & 01-d\,(1); 04-a\,(31); 04-b\,(12); 06-a\,(1); 06-b\,(1) \\
\textbf{action} & 4 & 19 & 9 & 01-c\,(6); 01-d\,(2); 04-a\,(2); 07-a\,(1); 08-a\,(2); 10-a\,(2); 14-a\,(2); 15-b\,(1); 16-b\,(1) \\
\textbf{active} & 2, 7 & 21 & 11 & 01-c\,(2); 01-d\,(1); 04-a\,(1); 07-a\,(1); 10-a\,(1); 10-b\,(2); 11-a\,(2); 11-b\,(1); 17-a\,(7); 19-b\,(1); 20-a\,(2) \\
\textbf{aeternitatis} & 11 & 0 & 0 & --- \\
\textbf{agent} & 4, 5, 6 & 104 & 19 & 01-c\,(15); 01-d\,(4); 04-a\,(15); 04-b\,(3); 06-a\,(2); 06-b\,(1); 07-a\,(15); 07-b\,(5); 08-a\,(13); 09-a\,(1); 10-a\,(5); 10-b\,(3); 11-a\,(1); 15-a\,(4); 15-b\,(7); 17-a\,(2); 18-a\,(6); 18-b\,(1); 19-a\,(1) \\
\textbf{and} & 9, 12 & 1444 & 36 & \emph{function word: totals only} \\
\textbf{answer} & 10, 11 & 34 & 20 & 01-d\,(1); 02-a\,(2); 02-b\,(1); 03-a\,(1); 04-a\,(2); 07-a\,(3); 08-a\,(1); 10-a\,(1); 10-b\,(2); 11-a\,(6); 11-b\,(1); 12-a\,(2); 14-a\,(1); 15-a\,(1); 15-b\,(1); 16-a\,(2); 16-b\,(2); 18-b\,(1); 19-a\,(1); 20-a\,(2) \\
\textbf{answers} & 9 & 32 & 14 & 02-a\,(1); 02-b\,(2); 03-a\,(1); 04-a\,(2); 07-a\,(1); 10-b\,(1); 11-a\,(1); 12-a\,(1); 13-a\,(2); 13-b\,(2); 16-a\,(7); 16-b\,(8); 19-a\,(2); 20-a\,(1) \\
\textbf{any} & 10 & 70 & 27 & \emph{function word: totals only} \\
\textbf{are} & 12 & 346 & 36 & \emph{function word: totals only} \\
\textbf{assumption} & 6, 11, 12 & 67 & 16 & 01-c\,(22); 01-d\,(2); 02-a\,(2); 02-b\,(1); 03-a\,(11); 03-b\,(4); 04-a\,(1); 06-a\,(3); 06-b\,(1); 10-a\,(2); 10-b\,(4); 13-a\,(5); 13-b\,(4); 17-a\,(1); 19-a\,(3); 19-b\,(1) \\
\textbf{at} & 4 & 197 & 31 & \emph{function word: totals only} \\
\textbf{attitudes} & 3 & 31 & 8 & 01-c\,(3); 01-d\,(1); 04-a\,(4); 06-a\,(4); 06-b\,(9); 07-a\,(5); 17-a\,(1); 18-a\,(4) \\
\textbf{available} & 5, 9 & 61 & 18 & 01-c\,(2); 01-d\,(8); 02-a\,(1); 03-a\,(1); 04-a\,(5); 04-b\,(2); 05-a\,(1); 07-a\,(4); 08-a\,(1); 10-a\,(4); 10-b\,(2); 11-a\,(8); 11-b\,(8); 14-a\,(9); 17-a\,(1); 18-a\,(1); 19-a\,(1); 20-a\,(2) \\
\textbf{by} & 6, 8 & 460 & 36 & \emph{function word: totals only} \\
\textbf{can} & 6 & 188 & 33 & \emph{function word: totals only} \\
\textbf{capable} & 4 & 15 & 4 & 01-c\,(6); 01-d\,(2); 04-a\,(2); 09-a\,(5) \\
\textbf{challenge} & 11\,(\(\times\)2) & 10 & 4 & 01-d\,(1); 04-a\,(1); 10-a\,(1); 10-b\,(7) \\
\textbf{commentary} & 12 & 0 & 0 & --- \\
\textbf{consists} & 2, 7 & 18 & 10 & 01-c\,(2); 04-a\,(2); 06-b\,(2); 07-a\,(4); 11-a\,(1); 15-a\,(2); 15-b\,(1); 16-a\,(1); 16-b\,(2); 18-a\,(1) \\
\textbf{defence} & 12 & 1 & 1 & 08-a\,(1) \\
\textbf{definition} & 8 & 153 & 31 & 01-c\,(9); 01-d\,(3); 02-a\,(6); 03-a\,(3); 03-b\,(1); 04-a\,(9); 04-b\,(7); 05-a\,(1); 06-a\,(4); 06-b\,(3); 07-a\,(6); 07-b\,(4); 08-a\,(5); 09-a\,(15); 09-b\,(13); 10-a\,(6); 10-b\,(10); 11-a\,(9); 11-b\,(3); 12-a\,(2); 13-a\,(2); 14-a\,(3); 15-a\,(5); 15-b\,(8); 16-a\,(5); 16-b\,(1); 17-a\,(5); 18-a\,(1); 18-b\,(1); 19-a\,(2); 19-b\,(1) \\
\textbf{degree} & 2, 8 & 31 & 10 & 01-c\,(2); 01-d\,(8); 02-a\,(1); 03-a\,(1); 05-a\,(2); 11-b\,(5); 15-a\,(5); 16-a\,(4); 16-b\,(2); 17-a\,(1) \\
\textbf{denies} & 11 & 14 & 9 & 01-d\,(1); 04-a\,(4); 06-a\,(1); 06-b\,(2); 09-a\,(2); 11-a\,(1); 12-a\,(1); 12-b\,(1); 16-a\,(1) \\
\textbf{depend} & 3 & 36 & 18 & 01-c\,(4); 01-d\,(3); 05-a\,(2); 06-b\,(4); 07-a\,(3); 08-a\,(1); 09-a\,(2); 09-b\,(1); 10-a\,(1); 11-a\,(3); 11-b\,(3); 15-b\,(1); 17-a\,(1); 18-a\,(2); 19-a\,(2); 19-b\,(1); 20-a\,(1); 20-b\,(1) \\
\textbf{depends} & 10 & 53 & 19 & 01-c\,(1); 01-d\,(3); 02-a\,(3); 03-a\,(2); 06-a\,(3); 07-a\,(7); 07-b\,(3); 08-a\,(3); 09-a\,(2); 09-b\,(1); 11-b\,(1); 15-a\,(2); 16-b\,(1); 17-a\,(3); 18-a\,(4); 19-a\,(4); 19-b\,(4); 20-a\,(4); 20-b\,(2) \\
\textbf{discussion} & 12 & 2 & 2 & 02-a\,(1); 04-a\,(1) \\
\textbf{does} & 3, 10, 11 & 306 & 35 & \emph{function word: totals only} \\
\textbf{done} & 10 & 11 & 7 & 01-d\,(3); 03-a\,(2); 03-b\,(1); 04-b\,(1); 13-a\,(1); 14-a\,(1); 17-a\,(2) \\
\textbf{engage} & 4, 6 & 27 & 9 & 01-c\,(2); 03-a\,(5); 06-b\,(2); 10-a\,(7); 10-b\,(4); 11-a\,(1); 11-b\,(1); 19-b\,(4); 20-a\,(1) \\
\textbf{engaged} & 3 & 34 & 14 & 01-c\,(1); 01-d\,(1); 03-a\,(1); 04-a\,(3); 06-b\,(2); 07-a\,(1); 07-b\,(1); 10-a\,(3); 10-b\,(1); 11-a\,(4); 17-a\,(7); 19-b\,(1); 20-a\,(5); 20-b\,(3) \\
\textbf{engagement} & 2, 7\,(\(\times\)2) & 148 & 17 & 01-c\,(14); 01-d\,(10); 03-a\,(3); 04-a\,(3); 04-b\,(8); 06-b\,(4); 07-a\,(14); 07-b\,(4); 10-a\,(5); 10-b\,(11); 11-a\,(13); 11-b\,(9); 14-a\,(1); 17-a\,(24); 19-b\,(11); 20-a\,(9); 20-b\,(5) \\
\textbf{every} & 4, 5 & 71 & 26 & \emph{function word: totals only} \\
\textbf{everyone} & 9 & 11 & 6 & 01-d\,(5); 02-a\,(2); 07-a\,(1); 11-a\,(1); 17-a\,(1); 20-a\,(1) \\
\textbf{exist} & 3 & 29 & 10 & 02-a\,(17); 02-b\,(1); 04-a\,(1); 09-a\,(1); 11-b\,(1); 12-a\,(1); 12-b\,(1); 13-a\,(3); 14-a\,(2); 19-a\,(1) \\
\textbf{extent} & 2, 7 & 29 & 17 & 01-c\,(1); 01-d\,(2); 02-a\,(1); 03-a\,(1); 04-a\,(2); 06-b\,(1); 07-a\,(4); 09-a\,(1); 11-a\,(4); 11-b\,(1); 16-b\,(1); 17-a\,(3); 18-a\,(1); 18-b\,(1); 19-b\,(1); 20-a\,(3); 20-b\,(1) \\
\textbf{for} & 4, 5 & 499 & 36 & \emph{function word: totals only} \\
\textbf{given} & 12 & 25 & 17 & 01-c\,(1); 01-d\,(2); 03-a\,(2); 03-b\,(1); 05-a\,(1); 06-a\,(1); 07-b\,(1); 09-a\,(1); 10-b\,(2); 11-a\,(1); 11-b\,(1); 12-a\,(2); 15-b\,(2); 17-a\,(1); 18-b\,(2); 19-a\,(2); 19-b\,(2) \\
\textbf{if} & 2 & 226 & 34 & \emph{function word: totals only} \\
\textbf{in} & 2, 3, 4, 6, 7, 12 & 752 & 36 & \emph{function word: totals only} \\
\textbf{includes} & 7 & 15 & 9 & 01-c\,(1); 04-a\,(1); 07-a\,(2); 09-b\,(2); 10-a\,(2); 10-b\,(2); 15-b\,(3); 17-a\,(1); 18-a\,(1) \\
\textbf{intentional} & 4 & 18 & 6 & 01-c\,(10); 01-d\,(2); 04-a\,(2); 09-a\,(2); 10-a\,(1); 10-b\,(1) \\
\textbf{is} & 2\,(\(\times\)2), 4\,(\(\times\)2), 5, 6, 8, 9\,(\(\times\)2), 10\,(\(\times\)2) & 1896 & 36 & \emph{function word: totals only} \\
\textbf{it} & 8, 10\,(\(\times\)2), 11\,(\(\times\)2) & 1065 & 36 & \emph{function word: totals only} \\
\textbf{least} & 4 & 14 & 8 & 01-c\,(2); 02-b\,(1); 09-a\,(1); 10-b\,(1); 15-a\,(1); 16-a\,(5); 16-b\,(2); 17-a\,(1) \\
\textbf{life} & 1, 2, 5, 7, 9, 10 & 450 & 33 & 01-c\,(7); 01-d\,(14); 02-a\,(40); 02-b\,(2); 03-a\,(62); 03-b\,(24); 04-a\,(19); 04-b\,(12); 05-a\,(4); 06-a\,(3); 06-b\,(3); 07-a\,(24); 07-b\,(2); 08-a\,(1); 09-a\,(13); 09-b\,(9); 10-a\,(2); 10-b\,(5); 11-a\,(43); 11-b\,(4); 12-a\,(25); 12-b\,(6); 13-a\,(5); 13-b\,(1); 14-a\,(1); 15-b\,(1); 17-a\,(35); 18-a\,(4); 18-b\,(1); 19-a\,(47); 19-b\,(24); 20-a\,(6); 20-b\,(1) \\
\textbf{matters} & 11 & 37 & 15 & 01-c\,(1); 01-d\,(2); 02-a\,(2); 02-b\,(1); 03-a\,(14); 03-b\,(2); 06-a\,(1); 06-b\,(1); 07-a\,(2); 10-a\,(1); 10-b\,(3); 12-a\,(1); 13-a\,(2); 17-a\,(3); 19-b\,(1) \\
\textbf{meaning} & 1, 5, 9\,(\(\times\)2) & 229 & 17 & 01-c\,(9); 01-d\,(28); 03-a\,(2); 04-a\,(5); 04-b\,(4); 06-b\,(16); 07-a\,(3); 10-a\,(8); 10-b\,(13); 11-a\,(49); 11-b\,(22); 14-a\,(1); 17-a\,(31); 18-a\,(1); 19-a\,(13); 19-b\,(21); 20-a\,(3) \\
\textbf{meaningful} & 2, 8, 10 & 92 & 16 & 01-c\,(3); 01-d\,(9); 03-a\,(1); 04-a\,(3); 04-b\,(2); 06-b\,(2); 09-a\,(1); 09-b\,(1); 10-b\,(2); 11-a\,(27); 11-b\,(3); 14-a\,(1); 17-a\,(23); 19-b\,(8); 20-a\,(5); 20-b\,(1) \\
\textbf{nor} & 11 & 23 & 14 & \emph{function word: totals only} \\
\textbf{not} & 3, 10 & 809 & 36 & \emph{function word: totals only} \\
\textbf{nothing} & 11 & 79 & 25 & 01-c\,(2); 01-d\,(3); 02-a\,(7); 02-b\,(4); 03-b\,(1); 04-a\,(5); 04-b\,(6); 05-a\,(2); 06-b\,(5); 07-a\,(4); 08-a\,(1); 09-a\,(2); 10-a\,(2); 10-b\,(1); 11-a\,(5); 12-a\,(1); 13-a\,(4); 13-b\,(2); 14-a\,(3); 15-a\,(1); 15-b\,(6); 16-a\,(1); 17-a\,(7); 19-a\,(1); 19-b\,(3) \\
\textbf{objective} & 2, 3, 4, 6, 7 & 25 & 9 & 01-c\,(7); 01-d\,(3); 03-a\,(1); 04-a\,(6); 06-b\,(2); 09-a\,(1); 11-a\,(2); 17-a\,(2); 20-a\,(1) \\
\textbf{of} & 1, 2, 3\,(\(\times\)2), 4\,(\(\times\)2), 5, 6, 7, 9\,(\(\times\)3), 10, 11, 12\,(\(\times\)2) & 1794 & 36 & \emph{function word: totals only} \\
\textbf{on} & 1, 3, 10 & 418 & 36 & \emph{function word: totals only} \\
\textbf{one} & 4 & 279 & 34 & \emph{function word: totals only} \\
\textbf{particular} & 10 & 27 & 15 & 01-d\,(2); 02-a\,(1); 03-b\,(1); 04-a\,(1); 05-a\,(3); 06-b\,(1); 07-a\,(3); 09-a\,(1); 11-a\,(1); 11-b\,(2); 13-a\,(2); 16-a\,(4); 19-a\,(3); 19-b\,(1); 20-b\,(1) \\
\textbf{positive} & 5, 7, 8 & 23 & 9 & 01-c\,(4); 01-d\,(3); 02-a\,(4); 07-a\,(1); 10-a\,(1); 14-a\,(2); 15-b\,(1); 18-a\,(3); 18-b\,(4) \\
\textbf{premise} & 11 & 18 & 10 & 01-c\,(3); 01-d\,(1); 02-a\,(5); 02-b\,(2); 03-a\,(1); 04-a\,(1); 05-a\,(2); 06-a\,(1); 10-a\,(1); 11-a\,(1) \\
\textbf{project} & 4, 6, 7 & 77 & 10 & 01-c\,(13); 01-d\,(2); 03-a\,(12); 03-b\,(8); 04-a\,(1); 06-b\,(1); 09-a\,(21); 09-b\,(17); 17-a\,(1); 18-a\,(1) \\
\textbf{projects} & 2, 3\,(\(\times\)2) & 61 & 15 & 01-c\,(10); 01-d\,(9); 03-a\,(8); 03-b\,(4); 04-a\,(5); 04-b\,(3); 06-b\,(1); 09-a\,(3); 11-a\,(2); 11-b\,(4); 13-a\,(4); 18-a\,(1); 19-a\,(3); 19-b\,(3); 20-a\,(1) \\
\textbf{question} & 9, 10 & 80 & 25 & 02-a\,(2); 03-a\,(2); 04-a\,(14); 04-b\,(2); 05-a\,(1); 06-a\,(3); 06-b\,(3); 07-a\,(1); 07-b\,(1); 08-a\,(10); 10-a\,(5); 10-b\,(2); 11-a\,(10); 12-a\,(2); 13-a\,(1); 13-b\,(2); 14-a\,(2); 15-a\,(1); 16-a\,(1); 17-a\,(2); 18-a\,(1); 18-b\,(1); 19-a\,(1); 20-a\,(7); 20-b\,(3) \\
\textbf{rather} & 11 & 39 & 20 & 01-c\,(2); 01-d\,(2); 04-a\,(1); 04-b\,(1); 06-b\,(1); 07-b\,(1); 08-a\,(6); 09-a\,(1); 10-a\,(2); 10-b\,(5); 11-a\,(1); 13-a\,(2); 14-a\,(1); 15-b\,(2); 16-a\,(2); 16-b\,(2); 17-a\,(2); 18-a\,(1); 19-a\,(2); 20-a\,(2) \\
\textbf{refuting} & 11 & 2 & 2 & 01-d\,(1); 04-a\,(1) \\
\textbf{sceptical} & 11 & 2 & 2 & 04-a\,(1); 08-a\,(1) \\
\textbf{she} & 4 & 193 & 27 & \emph{function word: totals only} \\
\textbf{some} & 7 & 59 & 20 & \emph{function word: totals only} \\
\textbf{specie} & 11 & 0 & 0 & --- \\
\textbf{sub} & 11 & 0 & 0 & --- \\
\textbf{such} & 5 & 92 & 26 & \emph{function word: totals only} \\
\textbf{than} & 11 & 139 & 30 & \emph{function word: totals only} \\
\textbf{that} & 7, 11\,(\(\times\)2) & 1327 & 36 & \emph{function word: totals only} \\
\textbf{the} & 1, 2, 3, 6, 9\,(\(\times\)3), 10, 11\,(\(\times\)2), 12\,(\(\times\)2) & 3437 & 36 & \emph{function word: totals only} \\
\textbf{them} & 3 & 122 & 31 & \emph{function word: totals only} \\
\textbf{theorem} & 9 & 81 & 33 & 01-c\,(2); 01-d\,(4); 02-a\,(2); 02-b\,(1); 03-a\,(2); 03-b\,(4); 04-a\,(4); 04-b\,(3); 05-a\,(2); 06-a\,(5); 06-b\,(3); 07-a\,(6); 07-b\,(3); 08-a\,(1); 09-a\,(1); 10-a\,(1); 10-b\,(4); 11-a\,(5); 11-b\,(2); 12-a\,(4); 12-b\,(2); 13-a\,(2); 13-b\,(1); 14-a\,(2); 15-a\,(1); 15-b\,(1); 16-a\,(3); 16-b\,(1); 17-a\,(2); 18-a\,(1); 18-b\,(1); 19-a\,(4); 19-b\,(1) \\
\textbf{there} & 3, 4, 6 & 72 & 24 & \emph{function word: totals only} \\
\textbf{those} & 3 & 55 & 24 & \emph{function word: totals only} \\
\textbf{to} & 2\,(\(\times\)2), 4, 7, 8, 9 & 1050 & 36 & \emph{function word: totals only} \\
\textbf{value} & 3 & 93 & 26 & 01-c\,(6); 01-d\,(2); 02-a\,(11); 03-a\,(2); 03-b\,(1); 04-a\,(3); 04-b\,(2); 05-a\,(2); 06-a\,(2); 06-b\,(2); 07-a\,(8); 09-b\,(2); 10-a\,(1); 10-b\,(3); 11-a\,(3); 12-b\,(1); 13-a\,(2); 14-a\,(1); 15-a\,(5); 15-b\,(4); 16-a\,(7); 16-b\,(3); 17-a\,(7); 18-a\,(1); 19-a\,(11); 19-b\,(1) \\
\textbf{what} & 9, 10 & 311 & 36 & \emph{function word: totals only} \\
\textbf{whether} & 9, 10 & 102 & 26 & \emph{function word: totals only} \\
\textbf{which} & 2, 4, 6, 10 & 345 & 35 & \emph{function word: totals only} \\
\textbf{whose} & 3 & 63 & 26 & \emph{function word: totals only} \\
\textbf{with} & 2, 7\,(\(\times\)2), 10 & 348 & 36 & \emph{function word: totals only} \\
\textbf{worth} & 2, 3, 4, 6, 7 & 82 & 17 & 01-c\,(7); 01-d\,(7); 02-a\,(2); 03-a\,(2); 03-b\,(2); 04-a\,(7); 04-b\,(3); 06-b\,(9); 09-a\,(2); 10-a\,(2); 11-a\,(16); 11-b\,(8); 15-a\,(1); 17-a\,(3); 19-b\,(7); 20-a\,(2); 20-b\,(2) \\
\end{longtable}
}

\subsection*{A.2 Term families by manuscript}

{\footnotesize Occurrences of the fifteen families defined in the preface, by manuscript; \enquote{Words} is the number of words of the manuscript's text after the removal of markup (for EXI-01-a, of its eleven sentences). A dot stands for zero. The last row totals the other manuscripts.\par}

{\footnotesize% Generated by gen/generate.py; do not edit.
\begin{longtable}{@{}l r r r r r r r r r@{}}
\toprule Ms & Words & \textit{mean-} & \textit{life-} & \textit{engag-} & \textit{activ-} & \textit{project-} & \textit{worth-} & \textit{objectiv-} & \textit{valu-} \\ \midrule \endhead
\bottomrule \endfoot
EXI-01-a & 200 & 6 & 5 & 6 & 2 & 6 & 5 & 5 & 1 \\
EXI-01-c & 1130 & 12 & 11 & 18 & 10 & 23 & 16 & 7 & 6 \\
EXI-01-d & 1031 & 38 & 18 & 15 & 1 & 11 & 13 & 3 & 3 \\
EXI-02-a & 2968 & . & 56 & . & . & . & 2 & . & 13 \\
EXI-02-b & 651 & . & 4 & . & . & . & . & . & . \\
EXI-03-a & 2359 & 4 & 73 & 11 & 2 & 20 & 9 & 1 & 4 \\
EXI-03-b & 772 & . & 24 & . & . & 12 & 4 & . & 1 \\
EXI-04-a & 2421 & 10 & 23 & 7 & 1 & 6 & 15 & 6 & 3 \\
EXI-04-b & 1182 & 6 & 15 & 8 & 2 & 3 & 8 & . & 2 \\
EXI-05-a & 1189 & 1 & 4 & . & . & . & . & . & 2 \\
EXI-06-a & 1197 & . & 3 & . & . & . & . & . & 2 \\
EXI-06-b & 845 & 18 & 10 & 9 & . & 2 & 10 & 2 & 3 \\
EXI-07-a & 2302 & 8 & 27 & 15 & 12 & . & 1 & . & 20 \\
EXI-07-b & 747 & 2 & 2 & 5 & 1 & . & . & . & 5 \\
EXI-08-a & 1505 & 5 & 2 & . & . & . & . & . & . \\
EXI-09-a & 1869 & 3 & 19 & . & 3 & 24 & 2 & 1 & 2 \\
EXI-09-b & 1013 & 1 & 11 & 1 & 1 & 17 & 1 & . & 2 \\
EXI-10-a & 1734 & 9 & 2 & 16 & 29 & . & 3 & . & 1 \\
EXI-10-b & 4261 & 15 & 5 & 17 & 19 & . & 1 & . & 3 \\
EXI-11-a & 2382 & 77 & 45 & 23 & 11 & 2 & 33 & 2 & 5 \\
EXI-11-b & 726 & 25 & 9 & 12 & 1 & 4 & 12 & . & . \\
EXI-12-a & 1312 & . & 49 & . & . & . & . & . & . \\
EXI-12-b & 758 & . & 8 & . & . & . & . & . & 1 \\
EXI-13-a & 2107 & . & 22 & . & . & 4 & . & . & 2 \\
EXI-13-b & 746 & . & 2 & . & . & . & . & . & . \\
EXI-14-a & 1210 & 2 & 1 & 1 & . & . & . & . & 1 \\
EXI-15-a & 1514 & . & . & . & . & . & 1 & . & 6 \\
EXI-15-b & 1158 & . & 1 & . & . & . & . & . & 4 \\
EXI-16-a & 2503 & . & . & . & . & . & . & . & 12 \\
EXI-16-b & 827 & . & . & . & . & . & . & . & 3 \\
EXI-17-a & 2018 & 57 & 46 & 32 & 28 & 1 & 20 & 2 & 8 \\
EXI-18-a & 1181 & 1 & 6 & . & . & 2 & 1 & . & 2 \\
EXI-18-b & 300 & . & 1 & . & . & . & . & . & . \\
EXI-19-a & 1470 & 15 & 48 & . & 1 & 3 & . & . & 15 \\
EXI-19-b & 861 & 29 & 24 & 19 & 2 & 3 & 10 & . & 4 \\
EXI-20-a & 735 & 9 & 6 & 15 & 11 & 1 & 11 & 1 & . \\
EXI-20-b & 412 & 1 & 1 & 8 & . & . & 4 & . & . \\
\midrule
Others & 51396 & 348 & 578 & 232 & 135 & 138 & 177 & 25 & 135 \\
\end{longtable}
}

{\footnotesize% Generated by gen/generate.py; do not edit.
\begin{longtable}{@{}l r r r r r r r r@{}}
\toprule Ms & Words & \textit{attitud-} & \textit{agent-} & \textit{availab-} & \textit{degree-} & \textit{matter-} & \textit{absurd-} & \textit{sceptic-} \\ \midrule \endhead
\bottomrule \endfoot
EXI-01-a & 200 & 1 & 3 & 2 & 4 & 1 & 1 & 1 \\
EXI-01-c & 1130 & 4 & 18 & 7 & 3 & 1 & . & . \\
EXI-01-d & 1031 & 1 & 5 & 22 & 11 & 3 & 1 & . \\
EXI-02-a & 2968 & 1 & . & 1 & 2 & 4 & . & . \\
EXI-02-b & 651 & 1 & . & . & . & 2 & . & . \\
EXI-03-a & 2359 & . & . & 1 & 2 & 14 & . & . \\
EXI-03-b & 772 & . & . & . & . & 2 & . & . \\
EXI-04-a & 2421 & 14 & 15 & 6 & 2 & 3 & 37 & 1 \\
EXI-04-b & 1182 & 3 & 3 & 2 & . & . & 15 & . \\
EXI-05-a & 1189 & 5 & . & 1 & 2 & . & . & . \\
EXI-06-a & 1197 & 8 & 3 & . & . & 2 & 1 & . \\
EXI-06-b & 845 & 9 & 1 & . & 1 & 2 & 1 & . \\
EXI-07-a & 2302 & 6 & 16 & 4 & 4 & 3 & . & . \\
EXI-07-b & 747 & . & 5 & . & . & 1 & . & . \\
EXI-08-a & 1505 & . & 15 & 1 & . & 1 & . & 3 \\
EXI-09-a & 1869 & . & 2 & . & 1 & . & . & . \\
EXI-09-b & 1013 & . & . & . & . & . & . & . \\
EXI-10-a & 1734 & . & 6 & 4 & . & 2 & . & . \\
EXI-10-b & 4261 & . & 6 & 2 & . & 6 & . & . \\
EXI-11-a & 2382 & . & 2 & 8 & 6 & 4 & . & . \\
EXI-11-b & 726 & . & . & 9 & 7 & . & . & . \\
EXI-12-a & 1312 & . & . & . & . & 1 & . & . \\
EXI-12-b & 758 & 1 & . & . & . & . & . & . \\
EXI-13-a & 2107 & . & . & . & . & 3 & . & . \\
EXI-13-b & 746 & . & . & . & . & . & . & . \\
EXI-14-a & 1210 & 1 & . & 9 & . & . & . & . \\
EXI-15-a & 1514 & . & 12 & . & 5 & . & . & . \\
EXI-15-b & 1158 & 1 & 7 & . & . & . & . & . \\
EXI-16-a & 2503 & 3 & 1 & . & 4 & 1 & . & . \\
EXI-16-b & 827 & . & . & . & 3 & . & . & . \\
EXI-17-a & 2018 & 1 & 2 & 2 & 4 & 5 & . & . \\
EXI-18-a & 1181 & 4 & 8 & 1 & 1 & 2 & . & . \\
EXI-18-b & 300 & . & 2 & . & 1 & . & . & . \\
EXI-19-a & 1470 & . & 1 & 1 & . & . & . & . \\
EXI-19-b & 861 & . & . & . & 1 & 4 & . & . \\
EXI-20-a & 735 & . & . & 2 & 3 & . & . & . \\
EXI-20-b & 412 & . & . & . & 1 & . & . & . \\
\midrule
Others & 51396 & 63 & 130 & 83 & 64 & 66 & 55 & 4 \\
\end{longtable}
}

\subsection*{A.3 Keyword-in-context concordance}

{\footnotesize Every occurrence, in the other manuscripts of the discipline, of every content word of EXI-01-a, in alphabetical order of the word and, within a word, in the order of the manuscripts and of their text. Each line gives the manuscript (without the prefix \enquote{EXI-}) and the section, then the context. The number after each heading is the number of lines.\par}

{\scriptsize% Generated by gen/generate.py; do not edit.
\kwichead{able}{2}
\begin{longtable}{@{}>{\raggedright}p{0.07\textwidth}>{\raggedleft}p{0.41\textwidth}@{\,}>{\raggedright\arraybackslash}p{0.47\textwidth}@{}}
01-c Abs & ...the assumption that every such agent is  & \textbf{able} to engage in at least one project of ob... \\
01-c 1 & ... agent capable of intentional action is  & \textbf{able} to engage in at least one of them (Assu... \\
\end{longtable}
\kwichead{absurd}{46}
\begin{longtable}{@{}>{\raggedright}p{0.07\textwidth}>{\raggedleft}p{0.41\textwidth}@{\,}>{\raggedright\arraybackslash}p{0.47\textwidth}@{}}
01-d 4 & EXI-04 grants that life is  & \textbf{absurd} in Nagel's sense and argues that the fi... \\
04-a Abs & Is life  & \textbf{absurd}, and if so, how should one respond? \\
04-a Abs & Nagel (1971) located the  & \textbf{absurd} not in the smallness or brevity of huma... \\
04-a Abs & We argue that life is  & \textbf{absurd} in this sense, because both of the coll... \\
04-a Abs & ...e key move is to treat responses to the  & \textbf{absurd} as attitudes assessable for fit. \\
04-a Abs & ...nd that discrepancy is exactly what the  & \textbf{absurd} consists in. \\
04-a 1 & ...veals gives no reason for anything, the  & \textbf{absurd} does not warrant despair; and the heroi... \\
04-a 1 & ...recommendation was that we approach our  & \textbf{absurd} lives with irony. \\
04-a 1 & One asks whether the  & \textbf{absurd} is a defect at all: if it follows from ... \\
04-a 1 & ...e value, if there is any, dissolves the  & \textbf{absurd}: on theories of meaning such as those o... \\
04-a 1 & We argue that life is  & \textbf{absurd} in Nagel's sense even if those theories... \\
04-a 2 & A life is  & \textbf{absurd} if its subject takes some ends seriousl... \\
04-a 2 & ...tion of the appropriate response to the  & \textbf{absurd} is then the question of which attitude ... \\
04-a 2 & ...question of which attitude presents the  & \textbf{absurd} situation as it is. \\
04-a 3 & ...luative convictions is another, and the  & \textbf{absurd} concerns the second. \\
04-a 3 & ... agent capable of intentional action is  & \textbf{absurd} in the sense of Definition 2.3. \\
04-a 3 & By Definition 2.3, her life is  & \textbf{absurd}. \\
04-a 3 & The  & \textbf{absurd} of Definition 2.3 is not a relation bet... \\
04-a 3 & ...ase and which Proposition 3.4 shows the  & \textbf{absurd} gives no reason to abandon. \\
04-a 4 & Irony is a way of not taking the  & \textbf{absurd} seriously. \\
04-a 4 & The demand that the  & \textbf{absurd} be met with constant tension presuppose... \\
04-a 4 & To take the  & \textbf{absurd} exactly as seriously as it deserves is ... \\
04-a 4 & ... life of someone engaged in them is not  & \textbf{absurd}. \\
04-a 4 & ...ed, if the projects are worthwhile; the  & \textbf{absurd} survives because that justification is ... \\
04-a 4 & ...ther from fitting it, but it is no less  & \textbf{absurd}, and irony fits it as well. \\
04-a 4 & The appropriate response to the  & \textbf{absurd} is not irony but forgetting. \\
04-a 4 & Forgetting is not a response to the  & \textbf{absurd} but an interruption of the reflection t... \\
04-a 4 & ...the question was what attitude fits the  & \textbf{absurd} when one does attend to it, and to that... \\
04-a 5 & ...n to despair, but the reason is not the  & \textbf{absurd}. \\
04-a 5 & ... reflects in the way that discloses the  & \textbf{absurd} has no occasion for irony. \\
04-a 6 & Life is  & \textbf{absurd} in Nagel's sense: a reflective agent ca... \\
04-a 6 & The  & \textbf{absurd} so understood is the absence of a vindi... \\
04-b Abs & ...s the condemned king as the hero of the  & \textbf{absurd}. \\
04-b Abs & EXI-04-a argues that life is  & \textbf{absurd} in Nagel's sense and that the fitting r... \\
04-b Abs & ...hat his situation is not a model of the  & \textbf{absurd} but of futility, which is a different c... \\
04-b 1 & He is the  & \textbf{absurd} hero because he is lucid about his fate... \\
04-b 1 & EXI-04-a defines the  & \textbf{absurd} as the conjunction of two features: the... \\
04-b 1 & A life can be  & \textbf{absurd} without being futile, and a task can be... \\
04-b 1 & ...ask can be futile in a life that is not  & \textbf{absurd}. \\
04-b 1 & ..., the life of every reflective agent is  & \textbf{absurd}, including the life of an agent whose p... \\
04-b 1 & ...ndpoint, and her life would then not be  & \textbf{absurd} in the sense of Definition 2.3 of EXI-0... \\
04-b 2 & ...revolt is not a fitting response to the  & \textbf{absurd}. \\
04-b 2 & ...-04-a it is the fitting response to the  & \textbf{absurd}. \\
04-b 5 & ...emblem of the human condition, which is  & \textbf{absurd} without being futile, faces no adversar... \\
06-a 5 & This bears on the response to the  & \textbf{absurd} (EXI-04): a world that is not an advers... \\
06-b 2 & ...revolt is not a fitting response to the  & \textbf{absurd}, since the world it scorns is indiffere... \\
\end{longtable}
\kwichead{action}{19}
\begin{longtable}{@{}>{\raggedright}p{0.07\textwidth}>{\raggedleft}p{0.41\textwidth}@{\,}>{\raggedright\arraybackslash}p{0.47\textwidth}@{}}
01-c Abs & ...e to every agent capable of intentional  & \textbf{action}. \\
01-c 1 & ...that every agent capable of intentional  & \textbf{action} is able to engage in at least one of th... \\
01-c 2 & Every agent capable of intentional  & \textbf{action} can direct some intentional activity, w... \\
01-c 2 & For every agent capable of intentional  & \textbf{action}, a life of positive meaning is availabl... \\
01-c 4 & Some agents capable of intentional  & \textbf{action} are wholly isolated and wholly unable t... \\
01-c 4 & An agent capable of intentional  & \textbf{action} acts on something, and anything she act... \\
01-d 1 & ... for every agent capable of intentional  & \textbf{action} a life of positive meaning is available... \\
01-d 2 & No agent capable of intentional  & \textbf{action} is condemned by her circumstances alone... \\
04-a 3 & Any agent capable of intentional  & \textbf{action} takes some ends seriously. \\
04-a 3 & ...reflective agent capable of intentional  & \textbf{action} is absurd in the sense of Definition 2.... \\
07-a 4 & ... claim that the good life is the end of  & \textbf{action} is compatible with our thesis, and inde... \\
08-a Abs & ... one's own is kept open and answered in  & \textbf{action}. \\
08-a 5 & ... ownership is kept open and answered in  & \textbf{action}, and this is what authenticity, read as... \\
10-a 3 & An aversive state motivates  & \textbf{action} that ends it, and the action that ends ... \\
10-a 3 & ... motivates action that ends it, and the  & \textbf{action} that ends boredom is a change of activi... \\
14-a 4 & An emotion that motivates no possible  & \textbf{action} is irrational. \\
14-a 4 & ...does not require that an emotion prompt  & \textbf{action}, as grief and admiration show. \\
15-b 3 & ... familiar in medicine and in collective  & \textbf{action}. \\
16-b 2 & ... 3.3 of EXI-16-a applied to reasons for  & \textbf{action}: a reason that consists in features any... \\
\end{longtable}
\kwichead{active}{21}
\begin{longtable}{@{}>{\raggedright}p{0.07\textwidth}>{\raggedleft}p{0.41\textwidth}@{\,}>{\raggedright\arraybackslash}p{0.47\textwidth}@{}}
01-c 1 & ...ngful to the degree that it consists in  & \textbf{active} engagement with projects of objective w... \\
01-c 2 & ...gement consists to a positive extent in  & \textbf{active} engagement with a project of objective ... \\
01-d 1 & ... by a different route: meaning requires  & \textbf{active} engagement with something worthwhile an... \\
04-a 4 & Meaning requires  & \textbf{active} engagement with something worthwhile (E... \\
07-a 3 & ...and EXI-17, on which meaning comes from  & \textbf{active} engagement with what is worthwhile. \\
10-a 4 & ...ng in EXI-17, on which meaning requires  & \textbf{active} engagement with something worthwhile, a... \\
10-b 5 & ... default mode network, the network more  & \textbf{active} when attention is not directed to an ex... \\
10-b 8 & ...7 holds that a meaningful life requires  & \textbf{active} engagement with something worthwhile; t... \\
11-a 1 & ...brid: meaning arises when the subject's  & \textbf{active} engagement meets the objective worth of... \\
11-a 7 & Meaning consists in  & \textbf{active} engagement with what is worthwhile, and... \\
11-b 1 & EXI-11-a defines meaning as  & \textbf{active} engagement with what is worthwhile (Def... \\
17-a Abs & A meaningful life requires  & \textbf{active} engagement with something worthwhile, i... \\
17-a 1 & ...ue that it should not: meaning requires  & \textbf{active} engagement with something worthwhile, a... \\
17-a 2 &  & \textbf{Active} engagement in the sense relevant to Pri... \\
17-a 3 & Her engagement is  & \textbf{active}, since study and attentive listening ar... \\
17-a 3 & A meaningful life requires  & \textbf{active} engagement with something worthwhile; i... \\
17-a 4 & ...ce of a difficult activity counts, then  & \textbf{active} engagement in a demanding worthwhile ac... \\
17-a 5 & A meaningful life requires  & \textbf{active} engagement with something worthwhile. \\
19-b 1 & ...e without categorical desires lacks the  & \textbf{active} engagement that meaning requires (EXI-1... \\
20-a Abs & ...eaningful activity within one, requires  & \textbf{active} engagement with something worthwhile. \\
20-a 1 & ... version of this view: meaning requires  & \textbf{active} engagement with something worthwhile, a... \\
\end{longtable}
\kwichead{aeternitatis}{0}
\kwicnone
\kwichead{agent}{104}
\begin{longtable}{@{}>{\raggedright}p{0.07\textwidth}>{\raggedleft}p{0.41\textwidth}@{\,}>{\raggedright\arraybackslash}p{0.47\textwidth}@{}}
01-c Abs & ... positive meaning is available to every  & \textbf{agent} capable of intentional action. \\
01-c Abs & ...rests on the assumption that every such  & \textbf{agent} is able to engage in at least one proje... \\
01-c Abs & ...h objects are within the reach of every  & \textbf{agent}, and show that sustained engagement wit... \\
01-c 1 & ... (Assumption 2), and assumes that every  & \textbf{agent} capable of intentional action is able t... \\
01-c 2 & ...t of engagement is anything at which an  & \textbf{agent} can direct intentional activity and abo... \\
01-c 2 & Every  & \textbf{agent} capable of intentional action can direc... \\
01-c 2 & If an  & \textbf{agent} directs intentional activity with care ... \\
01-c 2 & ...f a project is that it be something the  & \textbf{agent} engages in, actively, and that its valu... \\
01-c 2 & ...direction, extended over time, that the  & \textbf{agent} can pursue well or badly. \\
01-c 2 & For every  & \textbf{agent} capable of intentional action, a life o... \\
01-c 2 & By Assumption 2.2, the  & \textbf{agent} can direct intentional activity with ca... \\
01-c 3 & Almost every  & \textbf{agent} lives among other persons, whose good h... \\
01-c 3 & ...independent of her attitudes, and every  & \textbf{agent} lives in a world some part of which she... \\
01-c 4 & An  & \textbf{agent} capable of intentional action acts on s... \\
01-c 4 & The case of an  & \textbf{agent} for whom nothing worthwhile is within r... \\
01-d 1 & ...rem 4 of EXI-01-a states that for every  & \textbf{agent} capable of intentional action a life of... \\
01-d 2 & No  & \textbf{agent} capable of intentional action is condem... \\
01-d 3 & ...ility to actuality: a life in which the  & \textbf{agent} engages with nothing worthwhile, though... \\
01-d 5 & ...s reach: engagement, which is up to the  & \textbf{agent}, and worth, which belongs to what she e... \\
04-a Abs & ...res are ineliminable for any reflective  & \textbf{agent}, and that the appropriate response is i... \\
04-a 2 & An  & \textbf{agent} takes an end seriously if she devotes e... \\
04-a 3 & Any  & \textbf{agent} capable of intentional action takes som... \\
04-a 3 & An  & \textbf{agent} who took no end seriously in the sense ... \\
04-a 3 & For any reflective  & \textbf{agent} and any justification she can offer for... \\
04-a 3 & ...r an end cites a consideration that the  & \textbf{agent} treats as mattering: the end serves ano... \\
04-a 3 & From the detached standpoint the  & \textbf{agent} can ask, of that consideration, why it ... \\
04-a 3 & ...compelled to endorse it by anything the  & \textbf{agent} can exhibit, since what she would exhib... \\
04-a 3 & ...forbids the step back, and a reflective  & \textbf{agent} is one who can take it. \\
04-a 3 & The life of any reflective  & \textbf{agent} capable of intentional action is absurd... \\
04-a 3 & By Lemma 3.1 the  & \textbf{agent} takes some ends seriously. \\
04-a 3 & ...ition 2.3 is not a relation between the  & \textbf{agent} and the world but between two standpoin... \\
04-a 3 & ...orld but between two standpoints of the  & \textbf{agent}. \\
04-a 3 & ...nds, which Lemma 3.1 shows a reflective  & \textbf{agent} will do in any case and which Propositi... \\
04-a 6 & ...s absurd in Nagel's sense: a reflective  & \textbf{agent} cannot avoid taking her ends seriously,... \\
04-b 1 & ... EXI-04-a, the life of every reflective  & \textbf{agent} is absurd, including the life of an age... \\
04-b 1 & ...ent is absurd, including the life of an  & \textbf{agent} whose projects succeed and whose object... \\
04-b 1 & Conversely, an unreflective  & \textbf{agent} could be set to a futile task without e... \\
06-a 3 & ... a whole and its laws, which are not an  & \textbf{agent} and are not composed so as to aggregate... \\
06-a 4 & If the universe is not an  & \textbf{agent}, it is not indifferent either; indiffer... \\
06-b 3 & ...sed without reference to any particular  & \textbf{agent}, then the worth of objects of engagemen... \\
07-a Abs & An  & \textbf{agent} whose only final end is her own happine... \\
07-a 2 & An  & \textbf{agent} has as a final end if she desires for i... \\
07-a 2 & An  & \textbf{agent} aims directly at happiness if her own h... \\
07-a 2 & ...he pursuits contribute to only when the  & \textbf{agent} values their objects for their own sake... \\
07-a 2 & ...gement in Principle P1 require that the  & \textbf{agent} care about their objects for their own ... \\
07-a 2 & ...s that its being solved matters, and an  & \textbf{agent} to whom it matters only as a source of ... \\
07-a 3 & ... contributes to happiness only when the  & \textbf{agent} values its objects for their own sake. \\
07-a 3 & An  & \textbf{agent} who aims directly at happiness does not... \\
07-a 3 & ...ctly at happiness is self-defeating: an  & \textbf{agent} who does so forgoes the large part of h... \\
07-a 3 & By Lemma 3.2, the  & \textbf{agent} who aims directly at happiness does not... \\
07-a 3 & An  & \textbf{agent} who values other things for their own s... \\
07-a 3 & ...eaker thesis is true but contingent: an  & \textbf{agent} could perhaps aim at happiness while ke... \\
07-a 3 & An  & \textbf{agent} who never once thinks about her happine... \\
07-a 4 & Railton's  & \textbf{agent} can stably value other things for their... \\
07-a 6 & ...eople valued for their own sake, and an  & \textbf{agent} whose only final end is her own happine... \\
07-b 1 & ...e lies not in attention but in what the  & \textbf{agent} values: an agent whose sole final end i... \\
07-b 1 & ...ention but in what the agent values: an  & \textbf{agent} whose sole final end is her own happine... \\
07-b 1 & ...s indirectly self-defeating if, when an  & \textbf{agent} successfully follows it, its aims are w... \\
07-b 1 & The  & \textbf{agent} who aims directly at happiness does not... \\
07-b 1 & By Theorem 3.3 of EXI-07-a, an  & \textbf{agent} who does so attains less happiness than... \\
08-a 2 & An  & \textbf{agent} is authentic in the achievement sense a... \\
08-a 2 & An  & \textbf{agent} is authentic in the orientation sense o... \\
08-a 2 & ...uestion, and it is satisfied by how the  & \textbf{agent} stands to the question over time. \\
08-a 2 & Every interpretation an  & \textbf{agent} can give of her commitments, including ... \\
08-a 2 & No  & \textbf{agent} coincides with her commitments: whateve... \\
08-a 3 & Suppose an  & \textbf{agent} were authentic in the achievement sense... \\
08-a 3 & ...eted, cannot be held is not one that an  & \textbf{agent} can be in. \\
08-a 3 & ...f Definition 2.2 requires only that the  & \textbf{agent} treat a question as open, which Princip... \\
08-a 3 & ...is not trivial, because it can fail: an  & \textbf{agent} who meets every such occasion by retrea... \\
08-a 4 & ...ponding that must be maintained, and an  & \textbf{agent} who has them can lose them. \\
08-a 4 & An  & \textbf{agent} can believe she is keeping the question... \\
08-a 4 & ...matter of conduct, which others and the  & \textbf{agent} herself can observe over time. \\
08-a 5 & ...on draws on shared meanings and that no  & \textbf{agent} coincides with her commitments, show th... \\
09-a Abs & ...ect for a life can be fixed only by the  & \textbf{agent} who lives it. \\
10-a 3 & An activity fails to engage an  & \textbf{agent} who is not otherwise impaired just when... \\
10-a 3 & ...has the adaptive function of moving the  & \textbf{agent} toward activities that match her capaci... \\
10-a 3 & ...redom produces the effect of moving the  & \textbf{agent} toward matching activities. \\
10-a 3 & By Definition 2.3, moving the  & \textbf{agent} toward matching activities is its funct... \\
10-a 3 & ...ignal that cannot be acted on, when the  & \textbf{agent} is confined to mismatched activity. \\
10-b Abs & ...has the adaptive function of moving the  & \textbf{agent} toward activities that match them. \\
10-b 1 & ... selected-effect function of moving the  & \textbf{agent} toward matching activities (its Lemma 3... \\
10-b 6 & ...nction thesis is that boredom moves the  & \textbf{agent} toward activities that match her capaci... \\
11-a 2 & ...hat an activity has been assigned to an  & \textbf{agent} as a purpose by another, however powerf... \\
15-a 2 & An  & \textbf{agent} hopes that if (i) she desires that ; (i... \\
15-a 2 & ... the agent's evidence (EPI-17), and the  & \textbf{agent} does not control her credence directly ... \\
15-a 3 & ...tance is insensitive to probability, an  & \textbf{agent} invests as much attention, planning and... \\
15-a 4 & It exposes the  & \textbf{agent} to disappointment, but so does every at... \\
15-b Abs & ...what is hoped for, by changing what the  & \textbf{agent} does, from the illegitimate way in whic... \\
15-b 1 & On the account of EXI-15-a, an  & \textbf{agent} hopes that if she desires , has a crede... \\
15-b 1 & ...r the credence, for the value of to the  & \textbf{agent} and for her capacity to make more likel... \\
15-b 2 & ...ful stance passes into expectation: the  & \textbf{agent} plans on rather than for its possibilit... \\
15-b 2 & When is close to one, the rational  & \textbf{agent} plans on as she plans on what she expec... \\
15-b 3 & ...s probability only by changing what the  & \textbf{agent} does, hope is to be preferred, since it... \\
15-b 4 & ... and it does not exclude the despairing  & \textbf{agent} who makes the right choices but ceases ... \\
17-a 2 & An  & \textbf{agent} is actively engaged with an object if s... \\
17-a 2 & An  & \textbf{agent} achieves an outcome if she competently ... \\
18-a Abs & ...he sense that no one can function as an  & \textbf{agent} without some of them and their wholesal... \\
18-a 2 & ...ns, is psychologically necessary for an  & \textbf{agent} if she could not function as an agent, ... \\
18-a 2 & ...n agent if she could not function as an  & \textbf{agent}, in the sense of forming and pursuing p... \\
18-a 2 & ...r direct voluntary control (EPI-11), an  & \textbf{agent} cannot remove an illusion by deciding t... \\
18-a 3 & ...elief is necessary for functioning, the  & \textbf{agent} has a decisive practical reason to have... \\
18-a 3 & ...t involves no false belief: the serious  & \textbf{agent} need not believe that her ends are cert... \\
18-b 1 & ...s psychologically necessary for a given  & \textbf{agent} is an empirical question, and the answe... \\
19-a 3 & ...quires that it fulfil the purpose of an  & \textbf{agent} who imposes it or that it lead to a fur... \\
\end{longtable}
\kwichead{answer}{34}
\begin{longtable}{@{}>{\raggedright}p{0.07\textwidth}>{\raggedleft}p{0.41\textwidth}@{\,}>{\raggedright\arraybackslash}p{0.47\textwidth}@{}}
01-d Abs & We  & \textbf{answer} the objections that the view is elitist... \\
02-a Abs & We then  & \textbf{answer} the symmetry argument of Lucretius, the... \\
02-a 1 & Sections 4 and 5  & \textbf{answer} the symmetry argument and the main obje... \\
02-b Abs & ...in of two inferences and show where the  & \textbf{answer} bears on it. \\
03-a 2 & ... complexity, and the project types that  & \textbf{answer} to it are those that engage that struct... \\
04-a 3 & ...at justifies that justification, and no  & \textbf{answer} is forced on that standpoint. \\
04-a 4 & ...nd to that question forgetting gives no  & \textbf{answer}. \\
07-a Abs & We argue that the  & \textbf{answer} is no. \\
07-a Abs & ...itoring one's happiness is distracting,  & \textbf{answer} the reply that a sophisticated hedonist... \\
07-a 4 & Two points  & \textbf{answer} it. \\
08-a 3 & ...s hers or her parents', and acts on the  & \textbf{answer}, satisfies (ii) and (iii) on that occas... \\
10-a 5 & ...ome a correct signal that nothing could  & \textbf{answer}. \\
10-b Abs & Both claims  & \textbf{answer} to empirical work. \\
10-b 8 & ...ome a correct signal that nothing could  & \textbf{answer}, and the trait findings of Corollary 3.... \\
11-a Abs & ...on as one about dependence, so that the  & \textbf{answer} is independent of whether God exists. \\
11-a 1 & ... everything ended in death there was no  & \textbf{answer} (Tolstoy 1882). \\
11-a 1 & ...such a world is a counterpossible whose  & \textbf{answer} is unclear. \\
11-a 1 & So read, the question has a determinate  & \textbf{answer} whether or not REL-01 is correct, and t... \\
11-a 1 & ...ether or not REL-01 is correct, and the  & \textbf{answer} also settles a question of evident prac... \\
11-a 6 & ...nd our conclusion is independent of its  & \textbf{answer}. \\
11-b 1 & ...d if true, and so it holds whatever the  & \textbf{answer} to REL-01. \\
12-a Abs & This manuscript completes the  & \textbf{answer}. \\
12-a Abs & Second, we  & \textbf{answer} Benatar's independent argument that our... \\
14-a 4 & ...conditions of its rationality, but they  & \textbf{answer} the charge that it is idle. \\
15-a 1 & The  & \textbf{answer} depends on what hope is. \\
15-b Abs & And we  & \textbf{answer} the objection that the virtue so descri... \\
16-a 1 & The oldest  & \textbf{answer} to the question why we love is that we ... \\
16-a 3 & ...structed, loves something that does not  & \textbf{answer} to the person. \\
16-b Abs & We extend the  & \textbf{answer} to three cases of substitution: trading... \\
16-b 1 & ...red to the reasons of the old, it would  & \textbf{answer} to relational reasons constituted by th... \\
18-b 1 & ...agent is an empirical question, and the  & \textbf{answer} may differ across agents and across the... \\
19-a Abs & We argue for a third  & \textbf{answer}. \\
20-a Abs & ...ion of the collection and show that the  & \textbf{answer} turns on two questions: whether the pro... \\
20-a 4 & ... a position to know either fact, and an  & \textbf{answer} to the question from that position woul... \\
\end{longtable}
\kwichead{answers}{32}
\begin{longtable}{@{}>{\raggedright}p{0.07\textwidth}>{\raggedleft}p{0.41\textwidth}@{\,}>{\raggedright\arraybackslash}p{0.47\textwidth}@{}}
02-a 4 & It  & \textbf{answers} the argument without denying the symmet... \\
02-b Abs & EXI-02-a  & \textbf{answers} the first claim: death is bad for the o... \\
02-b 2 & EXI-02-a  & \textbf{answers} (E1). \\
03-a 2 & ...apacities and values to which that type  & \textbf{answers}. \\
04-a 3 & If she  & \textbf{answers} by citing a further consideration, the ... \\
04-a 3 & If she  & \textbf{answers} that the consideration has objective wo... \\
07-a Abs & A long tradition  & \textbf{answers} yes: whatever else we pursue, we pursue... \\
10-b 1 & Section 9  & \textbf{answers} objections to the use we make of the st... \\
11-a 1 & Section 5  & \textbf{answers} objections. \\
12-a 2 & Principle P2  & \textbf{answers} the familiar worry that coming into exi... \\
13-a 1 & ...s can affect whether she was happy, and  & \textbf{answers} that they can, though only slightly (Ar... \\
13-a 2 & Assumption 2.2  & \textbf{answers} that question for the case that matters... \\
13-b Abs & The distinction  & \textbf{answers} the question, left open by ETH-16-b, wh... \\
13-b 2 & This  & \textbf{answers} ETH-16-b in the sense it requires: the ... \\
16-a 2 & ... place of with no loss in what her love  & \textbf{answers} to. \\
16-a 3 & ...therefore lose a reason that 's love of  & \textbf{answers} to. \\
16-a 3 & Every established love  & \textbf{answers} to relational reasons. \\
16-a 3 & By Lemma 3.4, every established love  & \textbf{answers} to relational reasons, and by Lemma 3.3... \\
16-a 3 & ... reasons, and by Lemma 3.3, a love that  & \textbf{answers} to a relational reason is not transfera... \\
16-a 3 & ...rplus of qualities, since what the love  & \textbf{answers} to is not on offer elsewhere. \\
16-a 4 & ...oes not share the history that the love  & \textbf{answers} to. \\
16-b Abs & EXI-16-a  & \textbf{answers} that established love answers to relati... \\
16-b Abs & EXI-16-a answers that established love  & \textbf{answers} to relational reasons, which no qualita... \\
16-b Abs & ...l reason, and in the third the new love  & \textbf{answers} to a relational reason of its own rathe... \\
16-b 1 & Since every established love  & \textbf{answers} to relational reasons (Lemma 3.4), esta... \\
16-b 1 & ...3.4 of EXI-16-a, 's established love of  & \textbf{answers} to relational reasons. \\
16-b 1 & It  & \textbf{answers} to a relational reason of its own, cons... \\
16-b 1 & ...a, once it is established, the new love  & \textbf{answers} to a relational reason constituted by t... \\
16-b 3 & ... history, and a new love, if one comes,  & \textbf{answers} to a history of its own. \\
19-a Abs & Two  & \textbf{answers} are common. \\
19-a 4 & Corollary 3.3  & \textbf{answers} this. \\
20-a 3 & ...on, not philosophy, and a collection of  & \textbf{answers} is worthwhile to the extent that it con... \\
\end{longtable}
\kwichead{attitudes}{31}
\begin{longtable}{@{}>{\raggedright}p{0.07\textwidth}>{\raggedleft}p{0.41\textwidth}@{\,}>{\raggedright\arraybackslash}p{0.47\textwidth}@{}}
01-c 2 & ...lue that does not depend on the agent's  & \textbf{attitudes} toward it (Definition 2.2 of EXI-17-a). \\
01-c 2 & ...y, and that its value not depend on the  & \textbf{attitudes} of those engaged in it (Assumption 2 of... \\
01-c 3 & ...whose good has value independent of her  & \textbf{attitudes}, and every agent lives in a world some ... \\
01-d 4 & ...he value of a project not depend on the  & \textbf{attitudes} of those engaged in it. \\
04-a Abs & ... is to treat responses to the absurd as  & \textbf{attitudes} assessable for fit. \\
04-a 1 & ... may be worthwhile independently of our  & \textbf{attitudes} toward them, and then the step back see... \\
04-a 2 & ...o understood, is a relation between two  & \textbf{attitudes} of the same subject. \\
04-a 4 & ...cts are worthwhile independently of our  & \textbf{attitudes}, as EXI-01 holds, then the seriousness ... \\
06-a Abs & Care and hostility are  & \textbf{attitudes}, and both require a subject whose respo... \\
06-a 3 & Their  & \textbf{attitudes} are realized physically, and the theore... \\
06-a 3 & ...are not composed so as to aggregate the  & \textbf{attitudes} of the agents within them into an attit... \\
06-a 4 & ... the absence of the responsiveness that  & \textbf{attitudes} of care and hostility require. \\
06-b Abs & ...hat has worth independently of anyone's  & \textbf{attitudes}, and the universe's want of attitudes t... \\
06-b Abs & ...s attitudes, and the universe's want of  & \textbf{attitudes} toward us takes nothing from that worth... \\
06-b 2 & ...rth, whose value does not depend on the  & \textbf{attitudes} of those engaged in them (Definition 1 ... \\
06-b 2 & ...ts value does not depend on the agent's  & \textbf{attitudes} toward it (Definition 2.2 of EXI-17-a). \\
06-b 2 & ...rth of a project does not depend on the  & \textbf{attitudes} of those engaged in it, it does not dep... \\
06-b 2 & ...ngaged in it, it does not depend on the  & \textbf{attitudes} of the universe either, and a fortiori ... \\
06-b 2 & ... fortiori not on the universe's lack of  & \textbf{attitudes}. \\
06-b 2 & Both  & \textbf{attitudes} misrepresent their object. \\
06-b 3 & ...d cannot be met by any universe without  & \textbf{attitudes}. \\
07-a Abs & ...for its own sake, and so lacks the very  & \textbf{attitudes} on which most of her happiness depends. \\
07-a 2 & ...t is a claim about the structure of the  & \textbf{attitudes}, not about their phenomenology. \\
07-a 3 & ... else is a means to it, still lacks the  & \textbf{attitudes} that Principle P2 requires. \\
07-a 4 & ... other things, she has not acquired the  & \textbf{attitudes} on which engagement depends, and Theore... \\
07-a 6 & She therefore lacks the  & \textbf{attitudes} on which most of her happiness depends,... \\
17-a 2 & ...lue that does not depend on the agent's  & \textbf{attitudes} toward it. \\
18-a Abs & ...tment, which do not misrepresent, those  & \textbf{attitudes} are to be preferred. \\
18-a 2 & ... to positive illusions can be served by  & \textbf{attitudes} that do not misrepresent. \\
18-a 2 & To the extent that these  & \textbf{attitudes} serve the functions, the range of neces... \\
18-a 4 & ...cal commitment can do their work, those  & \textbf{attitudes}, which do not misrepresent, are better. \\
\end{longtable}
\kwichead{available}{61}
\begin{longtable}{@{}>{\raggedright}p{0.07\textwidth}>{\raggedleft}p{0.41\textwidth}@{\,}>{\raggedright\arraybackslash}p{0.47\textwidth}@{}}
01-c Abs & ...shes that a life of positive meaning is  & \textbf{available} to every agent capable of intentional a... \\
01-c 2 & ...l action, a life of positive meaning is  & \textbf{available}. \\
01-d Abs & ...orth, and that meaning in this sense is  & \textbf{available} to everyone. \\
01-d 1 & ...al action a life of positive meaning is  & \textbf{available}. \\
01-d 2 & ...f meaning required success, it would be  & \textbf{available} only to those whom the world allows to ... \\
01-d 2 & ...it, but not whether positive meaning is  & \textbf{available}. \\
01-d 4 & ...causes, and so make meaning in practice  & \textbf{available} mainly to the privileged. \\
01-d 4 & If meaning is  & \textbf{available} to everyone, the claim that a life is m... \\
01-d 4 & That positive meaning is  & \textbf{available} to everyone is compatible with its bein... \\
01-d 4 & ...andpoint nothing matters, so meaning is  & \textbf{available} to no one. \\
02-a 2 & ... theory of persons but that no evidence  & \textbf{available} before death could confirm it. \\
03-a 4 & Over an endless life the new theorems  & \textbf{available} to her become, in the respects that eng... \\
04-a Abs & ...ness of our commitments and the support  & \textbf{available} for them from the most detached standpo... \\
04-a 2 & ...no justification of that seriousness is  & \textbf{available} that is not itself open to the same que... \\
04-a 2 & ...em and the support for that seriousness  & \textbf{available} from the most detached standpoint, whil... \\
04-a 3 & ...no justification of that seriousness is  & \textbf{available} that is not open to the same question. \\
04-a 3 & ...holds one's ends and the support for it  & \textbf{available} from the most detached standpoint. \\
04-b 1 & ...no justification of that seriousness is  & \textbf{available} that is not itself open to the same que... \\
04-b 2 & ...ness of that engagement and the support  & \textbf{available} for it from the most detached standpoin... \\
05-a 3 & ...at end is coherent but that no evidence  & \textbf{available} before death could confirm it; the poss... \\
07-a 2 & ...t pleasure in the friend, and it is not  & \textbf{available} to someone who treats the friend as an ... \\
07-a 3 & ... that is realized through engagement is  & \textbf{available} only to agents who do. \\
07-a 3 & ...s yields less of its own target than an  & \textbf{available} alternative, which is what it is for a ... \\
07-a 4 & First, no such route is  & \textbf{available} to human beings, and the question wheth... \\
08-a 1 & ...ublic interpretation of the world makes  & \textbf{available} (Heidegger 1927). \\
10-a 2 & ...y when the activity does not occupy her  & \textbf{available} attention and abilities in a way she fi... \\
10-a 3 & ...n engagement is restored selects, among  & \textbf{available} alternatives, those that engage. \\
10-a 3 & ...correct but the motivated search has no  & \textbf{available} alternatives to find. \\
10-a 4 & ...state with the effects of the responses  & \textbf{available} to it. \\
10-b 2 & ... mismatch between cognitive demands and  & \textbf{available} resources in either direction, from a m... \\
10-b 4 & ..., and in a laboratory in which the only  & \textbf{available} stimulation is a shock, the search ends... \\
11-a 1 & Naturalists reply that meaning is  & \textbf{available} in a world without God. \\
11-a 2 & ...I-01 adds that meaning so understood is  & \textbf{available} to everyone. \\
11-a 4 & ...theistic framework, if true, would make  & \textbf{available} to a life more meaning than it could ha... \\
11-a 4 & ...ght, the additional meaning is actually  & \textbf{available}, and is available to those who engage w... \\
11-a 4 & ...l meaning is actually available, and is  & \textbf{available} to those who engage with God; if REL-01... \\
11-a 4 & ...with God; if REL-01 is wrong, it is not  & \textbf{available}, and Theorem 3.3 ensures that life is m... \\
11-a 4 & ...l of death is coherent, but no evidence  & \textbf{available} before death could confirm it. \\
11-a 4 & ...ralism that God would make more meaning  & \textbf{available} (Metz 2013); it does not follow the mod... \\
11-b Abs & ...e, would nevertheless make more meaning  & \textbf{available} (Proposition 4.1). \\
11-b Abs & ...ar life-plans, not the objects of worth  & \textbf{available} to engagement, so that Proposition 4.1 ... \\
11-b 1 & ... add is of the same kind as the meaning  & \textbf{available} without it, and differs in degree. \\
11-b 2 & ...existence, so the objects of engagement  & \textbf{available} without God remain available with God. \\
11-b 2 & ...engagement available without God remain  & \textbf{available} with God. \\
11-b 2 & ...e proposition concerns the meaning made  & \textbf{available} by a framework, and the framework adds ... \\
11-b 2 & ...w that the framework makes less meaning  & \textbf{available}. \\
11-b 3 & ...eism, if true, would add to the meaning  & \textbf{available} in a life by adding objects and relatio... \\
14-a Abs & ... of one's past as good and as no longer  & \textbf{available}, and it is fitting when the past contai... \\
14-a 2 & ...ined a good, and that good as no longer  & \textbf{available} to the subject. \\
14-a 2 & ...ned the good, and the good is no longer  & \textbf{available}), and proportionate (its intensity and ... \\
14-a 3 & ...real good of the past that is no longer  & \textbf{available}. \\
14-a 3 & ... as containing a good that is no longer  & \textbf{available}. \\
14-a 3 & ...id contain the good and it is no longer  & \textbf{available}. \\
14-a 3 & ...real good of the past that is no longer  & \textbf{available} and is proportionate to that good. \\
14-a 3 & ... tracking of a real past good no longer  & \textbf{available}. \\
14-a 3 & ... the good warrants, or forecloses goods  & \textbf{available} now. \\
17-a 3 & ...of EXI-01 that meaning in this sense is  & \textbf{available} to everyone. \\
18-a 2 & ...takes of error are high and evidence is  & \textbf{available}, that is a culpable failure of inquiry ... \\
19-a 3 & ...hed, and REL-08 argues that no evidence  & \textbf{available} before death could confirm survival. \\
20-a 1 & ...I-01 adds that meaning in this sense is  & \textbf{available} to everyone. \\
20-a 2 & On the evidence  & \textbf{available}, the condition of engagement is neither... \\
\end{longtable}
\kwichead{capable}{15}
\begin{longtable}{@{}>{\raggedright}p{0.07\textwidth}>{\raggedleft}p{0.41\textwidth}@{\,}>{\raggedright\arraybackslash}p{0.47\textwidth}@{}}
01-c Abs & ...ive meaning is available to every agent  & \textbf{capable} of intentional action. \\
01-c 1 & ...mption 2), and assumes that every agent  & \textbf{capable} of intentional action is able to engage... \\
01-c 2 & Every agent  & \textbf{capable} of intentional action can direct some i... \\
01-c 2 & For every agent  & \textbf{capable} of intentional action, a life of positi... \\
01-c 4 & Some agents  & \textbf{capable} of intentional action are wholly isolat... \\
01-c 4 & An agent  & \textbf{capable} of intentional action acts on something... \\
01-d 1 & ...of EXI-01-a states that for every agent  & \textbf{capable} of intentional action a life of positiv... \\
01-d 2 & No agent  & \textbf{capable} of intentional action is condemned by h... \\
04-a 3 & Any agent  & \textbf{capable} of intentional action takes some ends s... \\
04-a 3 & The life of any reflective agent  & \textbf{capable} of intentional action is absurd in the ... \\
09-a 2 & For a being  & \textbf{capable} of adopting ends, what its life is for,... \\
09-a 3 & ...fixes a defining project for , since is  & \textbf{capable} of adopting ends and what 's life is fo... \\
09-a 3 & ...s kind fix one: the kind human organism  & \textbf{capable} of thought is fixed by biological origi... \\
09-a 4 & ...ions fix a defining project for a being  & \textbf{capable} of adopting ends. \\
09-a 5 & ...he result concerns artefacts and beings  & \textbf{capable} of adopting ends. \\
\end{longtable}
\kwichead{challenge}{10}
\begin{longtable}{@{}>{\raggedright}p{0.07\textwidth}>{\raggedleft}p{0.41\textwidth}@{\,}>{\raggedright\arraybackslash}p{0.47\textwidth}@{}}
01-d 4 & ...Assumption 2 denies the premise of this  & \textbf{challenge} rather than refuting it (Nagel 1971). \\
04-a 4 & ...rth denies the premise of the sceptical  & \textbf{challenge} without refuting it. \\
10-a 2 & ...sorption and learning are maximal where  & \textbf{challenge} matches skill. \\
10-b Abs & ...as the weakest support and one standing  & \textbf{challenge}, the flow model's claim that over-chall... \\
10-b 3 & ...arch places boredom where skill exceeds  & \textbf{challenge} (Csikszentmihalyi 1975; Csikszentmihaly... \\
10-b 3 & ...egative emotions by a perceived lack of  & \textbf{challenge} (Til12). \\
10-b 3 & ...se of the mismatch thesis, that matched  & \textbf{challenge} and skill engage, is well supported: ex... \\
10-b 3 & ...of experience related to the balance of  & \textbf{challenge} and skill (Csikszentmihalyi and LeFevre... \\
10-b 4 & ...nder the strict rule for costs, entries  & \textbf{challenge} the function thesis outright and bear i... \\
10-b 7 & ...redom, the activity, and its subjective  & \textbf{challenge} and skill, and test whether boredom pre... \\
\end{longtable}
\kwichead{consists}{18}
\begin{longtable}{@{}>{\raggedright}p{0.07\textwidth}>{\raggedleft}p{0.41\textwidth}@{\,}>{\raggedright\arraybackslash}p{0.47\textwidth}@{}}
01-c 1 & ...ife as meaningful to the degree that it  & \textbf{consists} in active engagement with projects of o... \\
01-c 2 & A life that includes such engagement  & \textbf{consists} to a positive extent in active engageme... \\
04-a Abs & ... discrepancy is exactly what the absurd  & \textbf{consists} in. \\
04-a 4 & ...ony does not deny the contradiction; it  & \textbf{consists} in seeing it. \\
06-b Abs & Meaning, on the account of EXI-01,  & \textbf{consists} in engagement with what has worth indep... \\
06-b 3 & ...resent point is conditional: if meaning  & \textbf{consists} in engagement with what has objective w... \\
07-a Abs & Happiness, on any plausible account,  & \textbf{consists} to a large extent in engagement with an... \\
07-a 1 & ...urce of our enjoyments, since enjoyment  & \textbf{consists} in the gratification of particular pass... \\
07-a 2 & A large part of human happiness  & \textbf{consists} in, or depends causally on, wholehearte... \\
07-a 6 & Much of it  & \textbf{consists} in, or depends on, engagement with acti... \\
11-a 7 & Meaning  & \textbf{consists} in active engagement with what is worth... \\
15-a 2 & ... is under the agent's control, since it  & \textbf{consists} in what she attends to, plans for and a... \\
15-a 3 & ...eparture from epistemic rationality; it  & \textbf{consists} in the governance of the stance, which ... \\
15-b 1 & The virtue  & \textbf{consists} in proportion without collapse (Definit... \\
16-a 2 & It is relational if it  & \textbf{consists} in a relationship between and constitut... \\
16-b 2 & ...ed to reasons for action: a reason that  & \textbf{consists} in features anyone could have generaliz... \\
16-b 2 & ...ciple P2 of EXI-16-a, and a reason that  & \textbf{consists} in a shared history does not. \\
18-a 2 & ...on, and the ordinary progress of a life  & \textbf{consists} partly in such corrections. \\
\end{longtable}
\kwichead{defence}{1}
\begin{longtable}{@{}>{\raggedright}p{0.07\textwidth}>{\raggedleft}p{0.41\textwidth}@{\,}>{\raggedright\arraybackslash}p{0.47\textwidth}@{}}
08-a 1 & We argue that this  & \textbf{defence} succeeds only if authenticity is unders... \\
\end{longtable}
\kwichead{degree}{31}
\begin{longtable}{@{}>{\raggedright}p{0.07\textwidth}>{\raggedleft}p{0.41\textwidth}@{\,}>{\raggedright\arraybackslash}p{0.47\textwidth}@{}}
01-c 1 & ...1-a defines a life as meaningful to the  & \textbf{degree} that it consists in active engagement w... \\
01-c 2 & ...EXI-01-a it is meaningful to a positive  & \textbf{degree}. \\
01-d Abs & ...ircumstance and injustice affect is the  & \textbf{degree} of meaning, not its possibility. \\
01-d 2 & ...alue a project serves, which can affect  & \textbf{degree} but not presence. \\
01-d 2 & ...tances, including injustice, affect the  & \textbf{degree} of meaning a life can attain and the ea... \\
01-d 2 & ...1 of EXI-01-a makes meaning a matter of  & \textbf{degree}, and EXI-17-a notes that achievement ma... \\
01-d 2 & ...a notes that achievement may add to its  & \textbf{degree} (Remark 3.5 of EXI-17-a). \\
01-d 3 & ...ond, from availability to equality: the  & \textbf{degree} of meaning varies with the extent of en... \\
01-d 3 & ... tolerate oppression, which reduces the  & \textbf{degree} of meaning as it reduces other goods. \\
01-d 5 & ...ailability of meaning is universal; its  & \textbf{degree} is not, and how much of it a life attai... \\
02-a 3 & The  & \textbf{degree} of badness varies. \\
03-a 2 & There is a  & \textbf{degree} of connectedness below which a future l... \\
05-a 3 & ...s one would otherwise have had, and the  & \textbf{degree} of its badness is the value of what it ... \\
05-a 3 & ...eprive one of goods, and fitting to the  & \textbf{degree} that it does. \\
11-b Abs & ...e additional meaning is a difference of  & \textbf{degree}, and that the anti-theist's losses conc... \\
11-b 1 & ...ng available without it, and differs in  & \textbf{degree}. \\
11-b 1 & ...ple P2 of EXI-11-a, it can increase the  & \textbf{degree} of meaning. \\
11-b 2 & ...ition 1.1 counts that as an addition of  & \textbf{degree}. \\
11-b 3 & The addition is one of  & \textbf{degree}, which is why its absence leaves meanin... \\
15-a 2 & ...he right circumstances and to the right  & \textbf{degree}, flanked by vices of deficiency and exc... \\
15-a 3 & It is acting to the wrong  & \textbf{degree} (Definition 2.2). \\
15-a 3 & ...lar form, and it is acting to the wrong  & \textbf{degree} in the other direction. \\
15-a 3 & ...e other things that determine the right  & \textbf{degree} of investment in an uncertain outcome. \\
15-a 3 & ...isposition to act and feel to the right  & \textbf{degree} in the right circumstances, between two... \\
16-a 2 & ...n for 's love of , to at least the same  & \textbf{degree}. 's love of is transferable if, for eve... \\
16-a 2 & ...rds , the same feature, had to the same  & \textbf{degree} by , is an equally strong reason for th... \\
16-a 3 & ...n for 's love of , to at least the same  & \textbf{degree}. \\
16-a 3 & If has some features to a greater  & \textbf{degree}, the same argument yields trading up. \\
16-b 1 & ...e qualities for which loves to a higher  & \textbf{degree} is not a reason for to love that person... \\
16-b 1 & ...itative ones, and whatever they gain in  & \textbf{degree}, they lack a kind of reason that 's lov... \\
17-a Abs & ...adds a further good, and may add to the  & \textbf{degree} of meaning; it is not a condition of me... \\
\end{longtable}
\kwichead{denies}{14}
\begin{longtable}{@{}>{\raggedright}p{0.07\textwidth}>{\raggedleft}p{0.41\textwidth}@{\,}>{\raggedright\arraybackslash}p{0.47\textwidth}@{}}
01-d 4 & ...k 5 of EXI-01-a notes that Assumption 2  & \textbf{denies} the premise of this challenge rather th... \\
04-a 2 & ...presents the world as an adversary that  & \textbf{denies} us meaning, and the agent's persistence... \\
04-a 3 & ...presents the world as an adversary that  & \textbf{denies} us meaning. \\
04-a 3 & An indifferent world  & \textbf{denies} us nothing and cannot be defeated by ou... \\
04-a 4 & ...s it, the assumption of objective worth  & \textbf{denies} the premise of the sceptical challenge ... \\
06-a 2 & ...require responsiveness and indifference  & \textbf{denies} it. \\
06-b 2 & It  & \textbf{denies} that meaning in life requires that good... \\
06-b 2 & By Proposition 2.1, its indifference  & \textbf{denies} us nothing that meaning requires; by Th... \\
09-a 1 & Sartre's atheist existentialism  & \textbf{denies} this: there is at least one being whose... \\
09-a 4 & ...ay precede , including a creator's, and  & \textbf{denies} only that such intentions fix a definin... \\
11-a 2 & ...s valuable, or command what is good; it  & \textbf{denies} that the assignment is what makes it go... \\
12-a 4 & ... apply to them; nothing in our argument  & \textbf{denies} that some lives are worse than non-exis... \\
12-b 2 & ...e is bad for no one, and not on what it  & \textbf{denies}, that this absence is neutral. \\
16-a 1 & Frankfurt (2004)  & \textbf{denies} that love is a response to reasons: lov... \\
\end{longtable}
\kwichead{depend}{36}
\begin{longtable}{@{}>{\raggedright}p{0.07\textwidth}>{\raggedleft}p{0.41\textwidth}@{\,}>{\raggedright\arraybackslash}p{0.47\textwidth}@{}}
01-c 2 & ...if it has or serves value that does not  & \textbf{depend} on the agent's attitudes toward it (Def... \\
01-c 2 & ...es in, actively, and that its value not  & \textbf{depend} on the attitudes of those engaged in it... \\
01-c 3 & ...ties for projects in the ordinary sense  & \textbf{depend} on resources, institutions and freedom,... \\
01-c 5 & ...versal availability of meaning does not  & \textbf{depend} on the universal availability of projec... \\
01-d 1 & ...I-01-c shows that availability does not  & \textbf{depend} on the availability of projects in the ... \\
01-d 2 & Whether a life is meaningful does not  & \textbf{depend} on whether its projects succeed. \\
01-d 4 & ...es only that the value of a project not  & \textbf{depend} on the attitudes of those engaged in it... \\
05-a 3 & Its fittingness does not  & \textbf{depend} on uncertainty, which Lemma 3.1 exclude... \\
05-a 4 & ...he asymmetry of fitting regret does not  & \textbf{depend} on it. \\
06-b 2 & ...f objective worth, whose value does not  & \textbf{depend} on the attitudes of those engaged in th... \\
06-b 2 & ...ent is worthwhile if its value does not  & \textbf{depend} on the agent's attitudes toward it (Def... \\
06-b 2 & If the worth of a project does not  & \textbf{depend} on the attitudes of those engaged in it... \\
06-b 2 & ...des of those engaged in it, it does not  & \textbf{depend} on the attitudes of the universe either... \\
07-a Abs & ...is own crisis: the pleasures of pursuit  & \textbf{depend} on caring about objects other than plea... \\
07-a 2 & The argument does not  & \textbf{depend} on a choice between hedonistic, life-sa... \\
07-a 3 & Theorem 3.3 does not  & \textbf{depend} on attention. \\
08-a 4 & ...ne has it, and its possibility does not  & \textbf{depend} on its being known. \\
09-a Abs & The result does not  & \textbf{depend} on atheism, and it is compatible with t... \\
09-a 6 & The argument does not  & \textbf{depend} on atheism, and it leaves intact the th... \\
09-b 2 & ... of success as what it is that does not  & \textbf{depend} on how it is used; a knife that cannot ... \\
10-a 2 & ...reatures whose survival and flourishing  & \textbf{depend} on learning, exploration and the exerci... \\
11-a 3 & ...orsement converges on, does not make it  & \textbf{depend} on any assignment. \\
11-a 4 & Meaning has a floor that does not  & \textbf{depend} on God, and theism, if true, would rais... \\
11-a 5 & ...at it is would make the worth of a good  & \textbf{depend} on what happens after it, which neither... \\
11-b Abs & ...irection on which some meaningful lives  & \textbf{depend}. \\
11-b 2 & ...edge, of persons and of beauty does not  & \textbf{depend} on a role assigned by a designer (Lemma... \\
11-b 2 & ...names are real and that some life-plans  & \textbf{depend} on them. \\
15-b Abs & ... credence, other things equal, but must  & \textbf{depend} also on the value of the outcome and on... \\
17-a 2 & ...if it has or serves value that does not  & \textbf{depend} on the agent's attitudes toward it. \\
18-a Abs & ...t the believer; truth and justification  & \textbf{depend} on the world and on the evidence. \\
18-a 3 & Our conclusion does not  & \textbf{depend} on it. \\
19-a Abs & ...elleman, we take the value of a life to  & \textbf{depend} in part on its shape, the order in whic... \\
19-a 3 & ...ess life is not desirable, but does not  & \textbf{depend} on it. \\
19-b 1 & ...beauty and the good of others, does not  & \textbf{depend} on how long those who engage them live:... \\
20-a 4 & ...akes the meaningfulness of the activity  & \textbf{depend} on two facts that cannot be established... \\
20-b Abs & ...-20-a makes the worth of the collection  & \textbf{depend} on whether it contributes to knowledge ... \\
\end{longtable}
\kwichead{depends}{53}
\begin{longtable}{@{}>{\raggedright}p{0.07\textwidth}>{\raggedleft}p{0.41\textwidth}@{\,}>{\raggedright\arraybackslash}p{0.47\textwidth}@{}}
01-c 5 & It  & \textbf{depends} on something more widely distributed: t... \\
01-d 1 & ...ny particular life is meaningful, which  & \textbf{depends} on what is done with it. \\
01-d 4 & ...to high culture; the availability claim  & \textbf{depends} on this, and the cases in Proposition 2... \\
01-d 5 & ... not, and how much of it a life attains  & \textbf{depends} on what is done with it and on the cond... \\
02-a 2 & ... for a while as a corpse; nothing below  & \textbf{depends} on it. \\
02-a 2 & ... account is comparative: whether is bad  & \textbf{depends} on what the alternative to would have b... \\
02-a 6 & ... settle how much the badness of a death  & \textbf{depends} on the psychological connections betwee... \\
03-a 1 & ...that the value of many human activities  & \textbf{depends} on the scarcity of time. \\
03-a 4 & ...riendship is a project, and its renewal  & \textbf{depends} on the friends' changing. \\
06-a 2 & ...onsive to our interests if what it does  & \textbf{depends}, in a counterfactually robust way, on h... \\
06-a 2 & ...onditions on which the thing's behavior  & \textbf{depends}. \\
06-a 3 & ... Assumption 2.3, what the universe does  & \textbf{depends} only on its prior states and on laws th... \\
07-a Abs & ...ttitudes on which most of her happiness  & \textbf{depends}. \\
07-a 2 & ...ional-state accounts (Haybron 2008); it  & \textbf{depends} only on what makes people happy, on any... \\
07-a 2 & ...part of human happiness consists in, or  & \textbf{depends} causally on, wholehearted engagement in... \\
07-a 3 & ...orgoes the large part of happiness that  & \textbf{depends} on engagement, and attains less happine... \\
07-a 4 & ...uired the attitudes on which engagement  & \textbf{depends}, and Theorem 3.3 applies to her. \\
07-a 6 & Much of it consists in, or  & \textbf{depends} on, engagement with activities and peop... \\
07-a 6 & ...ttitudes on which most of her happiness  & \textbf{depends}, and her aim defeats itself. \\
07-b 1 & ...s the absorption on which the enjoyment  & \textbf{depends}. \\
07-b 1 & ...ect-directed caring on which engagement  & \textbf{depends} (Principle P2 of EXI-07-a). \\
07-b 2 & ...ich case she forgoes the happiness that  & \textbf{depends} on engagement. \\
08-a 2 & ...n, and holds it in a way that no longer  & \textbf{depends} on its being what one does. \\
08-a 3 & ...mitment is held in a way that no longer  & \textbf{depends} on its being what one does. \\
08-a 3 & ...from shared meanings, so her holding it  & \textbf{depends} on public interpretations, and the appr... \\
09-a 2 & ...ation on which 's coming into existence  & \textbf{depends}. \\
09-a 3 & ... out, and 's coming into existence as a  & \textbf{depends} on it. \\
09-b 2 & ... a determination on which its existence  & \textbf{depends}, not that the determination mention it ... \\
11-b 2 & ...eaning is undercut is the claim that it  & \textbf{depends} on a cosmic purpose, which Theorem 3.3 ... \\
15-a 1 & The answer  & \textbf{depends} on what hope is. \\
15-a 3 & ...t a function of the credence alone, but  & \textbf{depends} also on the value of and on the agent's... \\
16-b 1 & Whether the duplicate took part in it  & \textbf{depends} on whether she is the beloved, which is... \\
17-a 3 & Whether an activity succeeds  & \textbf{depends} on circumstances outside the agent's co... \\
17-a 3 & ...he is engaged with something worthwhile  & \textbf{depends} far less on such circumstances. \\
17-a 3 & By Theorem 3.3, meaning  & \textbf{depends} on the second and not on the first. \\
18-a 2 & Whether a belief is true  & \textbf{depends} on whether things are as it represents ... \\
18-a 2 & ...epresents them; whether it is justified  & \textbf{depends} on the evidence bearing on it. \\
18-a 3 & Meaning  & \textbf{depends} on taking one's projects to matter more... \\
18-a 4 & ...ically necessary, in that normal agency  & \textbf{depends} on them and they cannot be removed at w... \\
19-a Abs & ...f, since every candidate meaning either  & \textbf{depends} on claims that cannot be established or... \\
19-a 1 & ...hat the value of many of our activities  & \textbf{depends} on the scarcity of time. \\
19-a 2 & The value of a life as a whole  & \textbf{depends} in part on its shape, and not only on t... \\
19-a 3 & ...s not determine what the shape is: that  & \textbf{depends} on the goods and evils the life contain... \\
19-b Abs & ...aning, that of a life as a whole, which  & \textbf{depends} on the life's shape. \\
19-b 1 & ...far as the meaning of a life as a whole  & \textbf{depends} on its shape, finitude is a condition o... \\
19-b 1 & Such meaning  & \textbf{depends} on the temporal order of the engagement... \\
19-b 2 & ...ffler (2013) that our capacity to value  & \textbf{depends} on the awareness of temporal scarcity, ... \\
20-a Abs & The first  & \textbf{depends} on whether the process has the kind of ... \\
20-a Abs & The second  & \textbf{depends} on what readers do with the collection,... \\
20-a 3 & ...nts are correct, the collection's worth  & \textbf{depends} on the minority of readers who engage w... \\
20-a 3 & ...h its arguments, and its worth for them  & \textbf{depends} on whether the arguments are good. \\
20-b Abs & ...g the present collection was meaningful  & \textbf{depends}, by the account of EXI-17, on whether i... \\
20-b Abs & ...ucing process was engaged with its work  & \textbf{depends}, by EXI-20-a, on whether it has the men... \\
\end{longtable}
\kwichead{discussion}{2}
\begin{longtable}{@{}>{\raggedright}p{0.07\textwidth}>{\raggedleft}p{0.41\textwidth}@{\,}>{\raggedright\arraybackslash}p{0.47\textwidth}@{}}
02-a 1 & ...he argument has two premises that later  & \textbf{discussion} has separated. \\
04-a 1 & Two lines of later  & \textbf{discussion} bear on the recommendation. \\
\end{longtable}
\kwichead{done}{11}
\begin{longtable}{@{}>{\raggedright}p{0.07\textwidth}>{\raggedleft}p{0.41\textwidth}@{\,}>{\raggedright\arraybackslash}p{0.47\textwidth}@{}}
01-d 1 & ...is meaningful, which depends on what is  & \textbf{done} with it. \\
01-d 5 & ...of it a life attains depends on what is  & \textbf{done} with it and on the conditions under whi... \\
01-d 5 & ...and on the conditions under which it is  & \textbf{done} (Metz 2013). \\
03-a 2 & ...y as repetition of what she has already  & \textbf{done}. \\
03-a 2 & ...e desire to repeat what one has already  & \textbf{done} is not a desire that gives life its poi... \\
03-b 2 & ... without specification of what is to be  & \textbf{done} in it. \\
04-b 4 & ... it does not undo the value of what was  & \textbf{done}. \\
13-a 2 & ...owing that wants assurance that will be  & \textbf{done}, then is obligated to to do , unless re... \\
14-a 4 & Nothing can be  & \textbf{done} about the past. \\
17-a 3 & ...jectivism fails because worthwhile work  & \textbf{done} without engagement, out of mere duty or... \\
17-a 3 & ...before may be difficult and competently  & \textbf{done}, and is therefore an achievement by Def... \\
\end{longtable}
\kwichead{engage}{27}
\begin{longtable}{@{}>{\raggedright}p{0.07\textwidth}>{\raggedleft}p{0.41\textwidth}@{\,}>{\raggedright\arraybackslash}p{0.47\textwidth}@{}}
01-c Abs & ...mption that every such agent is able to  & \textbf{engage} in at least one project of objective wo... \\
01-c 1 & ...apable of intentional action is able to  & \textbf{engage} in at least one of them (Assumption 3). \\
03-a 2 & ...that give life its point are desires to  & \textbf{engage} in such projects. \\
03-a 2 & ...ct types that a person of character can  & \textbf{engage} with, in the sense of finding them wort... \\
03-a 2 & ... types that answer to it are those that  & \textbf{engage} that structure. \\
03-a 3 & Projects of other types do not  & \textbf{engage} her, by Assumption 2.3. \\
03-a 4 & ...ble to her become, in the respects that  & \textbf{engage} her, instances of kinds she has already... \\
06-b 4 & The worth of what we  & \textbf{engage} with does not wait on cosmic recognitio... \\
06-b 4 & ...ning where EXI-01 places it: in what we  & \textbf{engage} with, and how. \\
10-a Abs & ...nts the current situation as failing to  & \textbf{engage} one's attention and abilities; it is tr... \\
10-a 1 & ...rience of wanting, but being unable, to  & \textbf{engage} in satisfying activity, and locate its ... \\
10-a 2 & ...ent activity or situation as failing to  & \textbf{engage} her, together with a desire to be engag... \\
10-a 2 & ...nts the current situation as failing to  & \textbf{engage}. \\
10-a 3 & ...ents the current activity as failing to  & \textbf{engage} the subject. \\
10-a 3 & An activity fails to  & \textbf{engage} an agent who is not otherwise impaired ... \\
10-a 3 & ...mong available alternatives, those that  & \textbf{engage}. \\
10-b 2 & ...rience of wanting, but being unable, to  & \textbf{engage} in satisfying activity, and locate the ... \\
10-b 3 & ...le an over-demanding task one can still  & \textbf{engage} with is a source of anxiety. \\
10-b 3 & ...hesis, that matched challenge and skill  & \textbf{engage}, is well supported: experience sampling... \\
10-b 11 & ...ople seek alternatives, change how they  & \textbf{engage}, and turn to sources of meaning. \\
11-a 4 & ...vailable, and is available to those who  & \textbf{engage} with God; if REL-01 is wrong, it is not... \\
11-b 1 & ...availability: it says what a life could  & \textbf{engage} with, not that every life in a theistic... \\
19-b Abs & ...ng, since it adds no object of worth to  & \textbf{engage} with. \\
19-b 1 & ..., does not depend on how long those who  & \textbf{engage} them live: a proof is no less a proof f... \\
19-b 2 & ...eir urgency, and urgency affects how we  & \textbf{engage}. \\
19-b 2 & ...ition 1.1 concerns the worth of what we  & \textbf{engage} with, and urgency adds no worth: the fr... \\
20-a 3 & ... depends on the minority of readers who  & \textbf{engage} with its arguments, and its worth for t... \\
\end{longtable}
\kwichead{engaged}{34}
\begin{longtable}{@{}>{\raggedright}p{0.07\textwidth}>{\raggedleft}p{0.41\textwidth}@{\,}>{\raggedright\arraybackslash}p{0.47\textwidth}@{}}
01-c 2 & ...ue not depend on the attitudes of those  & \textbf{engaged} in it (Assumption 2 of EXI-01-a). \\
01-d 4 & ...ct not depend on the attitudes of those  & \textbf{engaged} in it. \\
03-a 4 & A person  & \textbf{engaged} in inquiry would never be reduced to re... \\
04-a 4 & ...m is justified, and the life of someone  & \textbf{engaged} in them is not absurd. \\
04-a 4 & The  & \textbf{engaged} and worthwhile life is less exposed to ... \\
04-a 4 & ...nts at a distance and is therefore less  & \textbf{engaged}. \\
06-b 2 & ...es not depend on the attitudes of those  & \textbf{engaged} in them (Definition 1 and Assumption 2 ... \\
06-b 2 & ...es not depend on the attitudes of those  & \textbf{engaged} in it, it does not depend on the attitu... \\
07-a 4 & ...echnology that produced the feelings of  & \textbf{engaged} satisfaction without any engagement wou... \\
07-b 1 & ...ing, attending to one's happiness while  & \textbf{engaged} in an activity disrupts the absorption ... \\
10-a 2 & ...ngage her, together with a desire to be  & \textbf{engaged} otherwise. \\
10-a 3 & ...is aversive and includes a desire to be  & \textbf{engaged} otherwise (Definition 2.1). \\
10-a 5 & ...lls do better when their capacities are  & \textbf{engaged}, the disposition to boredom has the ada... \\
10-b 2 & ...edback that attention is not adequately  & \textbf{engaged}. \\
11-a 2 & ...the extent that its subject is actively  & \textbf{engaged} with what is worthwhile. \\
11-a 3 & ...worthwhile, and the subject need not be  & \textbf{engaged} with anything worthwhile in living it. \\
11-a 3 & A subject who is actively  & \textbf{engaged} with such activities satisfies Definiti... \\
11-a 5 & The  & \textbf{engaged} and worthwhile life of a person without... \\
17-a 1 & ...ctiveness: the subject must be actively  & \textbf{engaged}, and the object of her engagement must ... \\
17-a 2 & An agent is actively  & \textbf{engaged} with an object if she acts in ways dire... \\
17-a 2 & ...the extent that its subject is actively  & \textbf{engaged} with what is worthwhile. \\
17-a 2 & ...es, lived in the machine, is not a life  & \textbf{engaged} with the sick. \\
17-a 3 & Each is actively  & \textbf{engaged} in the sense of Definition 2.1. \\
17-a 3 & Whether she is  & \textbf{engaged} with something worthwhile depends far l... \\
17-a 4 & ...ment is unnecessary, Taylor's Sisyphus,  & \textbf{engaged} with his stone and achieving nothing, l... \\
19-b 1 & ... meaningful, and does not make anything  & \textbf{engaged} with more worthwhile. \\
20-a Abs & ...ions: whether the producing process was  & \textbf{engaged} with its work, in the sense the account... \\
20-a 1 & ...the extent that its subject is actively  & \textbf{engaged} in projects that are worthwhile indepen... \\
20-a 1 & ... passion to counting blades of grass is  & \textbf{engaged} with something worthless; the person wh... \\
20-a 1 & ...ut interest at a worthwhile task is not  & \textbf{engaged}. \\
20-a 4 & ...ion is satisfied, since its members are  & \textbf{engaged} in their work, and the question reduces... \\
20-b Abs & Whether the producing process was  & \textbf{engaged} with its work depends, by EXI-20-a, on ... \\
20-b Abs & ...ent by the producing system that it was  & \textbf{engaged} with its work would be an output of the... \\
20-b Abs & A statement that it was not  & \textbf{engaged} would be an output of the same kind. \\
\end{longtable}
\kwichead{engagement}{148}
\begin{longtable}{@{}>{\raggedright}p{0.07\textwidth}>{\raggedleft}p{0.41\textwidth}@{\,}>{\raggedright\arraybackslash}p{0.47\textwidth}@{}}
01-c Abs & ...n not from projects but from objects of  & \textbf{engagement}, such as persons and the world as an ob... \\
01-c Abs & ...of every agent, and show that sustained  & \textbf{engagement} with one of them constitutes a project ... \\
01-c 1 & ...o the degree that it consists in active  & \textbf{engagement} with projects of objective worth (Defin... \\
01-c 2 & An object of  & \textbf{engagement} is anything at which an agent can direc... \\
01-c 2 & ...y, with care, at a worthwhile object of  & \textbf{engagement}: another person, or some part of the wo... \\
01-c 2 & ...ity with care at a worthwhile object of  & \textbf{engagement}, and sustains it, her engagement is a p... \\
01-c 2 & ...ect of engagement, and sustains it, her  & \textbf{engagement} is a project of objective worth in the ... \\
01-c 2 & ...ity with care at a worthwhile object of  & \textbf{engagement}. \\
01-c 2 & ... Lemma 2.3, sustaining that activity is  & \textbf{engagement} in a project of objective worth. \\
01-c 2 & A life that includes such  & \textbf{engagement} consists to a positive extent in active... \\
01-c 2 & ...consists to a positive extent in active  & \textbf{engagement} with a project of objective worth, and ... \\
01-c 3 & ...ement: a life can be meaningful through  & \textbf{engagement} without success (Theorem 3.3 of EXI-17-... \\
01-c 3 & ... the availability of a project, because  & \textbf{engagement} makes one. \\
01-c 5 &  & \textbf{Engagement} with them is a project of objective wor... \\
01-d 1 & ...ifferent route: meaning requires active  & \textbf{engagement} with something worthwhile and not achie... \\
01-d 1 & ... in the ordinary sense, since sustained  & \textbf{engagement} with a worthwhile object constitutes on... \\
01-d 2 & Because meaning requires  & \textbf{engagement}, which is within the agent's power, and... \\
01-d 2 & ...ition 1 of EXI-01-a measures meaning by  & \textbf{engagement} and by the objective worth of what one ... \\
01-d 2 & ...or are projects of objective worth, and  & \textbf{engagement} with them confers meaning on lives that... \\
01-d 2 & ...dinary sense have worthwhile objects of  & \textbf{engagement} within reach. \\
01-d 2 & ...mine how much of a life can be given to  & \textbf{engagement} and how wide the range of worthwhile ob... \\
01-d 3 & ...ee of meaning varies with the extent of  & \textbf{engagement} and the worth of its object, and lives ... \\
01-d 4 & ... insists that the objects of meaningful  & \textbf{engagement} are not restricted to high culture; the... \\
01-d 5 & ...conditions are within everyone's reach:  & \textbf{engagement}, which is up to the agent, and worth, w... \\
03-a 2 & ...f objective worth, and EXI-17 that such  & \textbf{engagement}, not achievement, is what meaning requi... \\
03-a 2 & A character supports  & \textbf{engagement} with a project type only by having inte... \\
03-a 4 & ...ut Assumption 2.4 concerns the person's  & \textbf{engagement}. \\
04-a Abs & ...esponse, and it is compatible with full  & \textbf{engagement} in worthwhile projects. \\
04-a 4 & Meaning requires active  & \textbf{engagement} with something worthwhile (EXI-17). \\
04-a 4 & ... not care about anything, would corrode  & \textbf{engagement}; it is a different attitude and one we ... \\
04-b 2 & ...e attitude EXI-04-a calls irony: a full  & \textbf{engagement} with one's ends combined with a clear p... \\
04-b 2 & ...repancy between the seriousness of that  & \textbf{engagement} and the support available for it from t... \\
04-b 2 & ...'s conclusion is the claim that serious  & \textbf{engagement} with one's ends is compatible with luci... \\
04-b 3 & ...ver happy he is, since meaning requires  & \textbf{engagement} with something worthwhile (EXI-17-a, Th... \\
04-b 3 & ...ur ends include worthwhile objects, and  & \textbf{engagement} with them confers meaning (EXI-01). \\
04-b 4 & ...n the sense of EXI-04-a does not reduce  & \textbf{engagement}: it is compatible with full seriousness... \\
04-b 4 & ...he passion Camus admires belongs to the  & \textbf{engagement} and survives the change of attitude. \\
04-b 5 & What Camus rightly saw in him, unbroken  & \textbf{engagement} without consolation, is irony, the resp... \\
06-b Abs & ..., on the account of EXI-01, consists in  & \textbf{engagement} with what has worth independently of an... \\
06-b 2 & ...of EXI-17-a, which agrees, an object of  & \textbf{engagement} is worthwhile if its value does not dep... \\
06-b 3 & ...lar agent, then the worth of objects of  & \textbf{engagement} is precisely worth from that standpoint... \\
06-b 3 & ... is conditional: if meaning consists in  & \textbf{engagement} with what has objective worth, as EXI-0... \\
07-a Abs & ... account, consists to a large extent in  & \textbf{engagement} with and success in activities whose ob... \\
07-a 2 & ...n, or depends causally on, wholehearted  & \textbf{engagement} in activities and relationships: absorp... \\
07-a 2 & The forms of  & \textbf{engagement} in Principle P1 require that the agent ... \\
07-a 3 & Happiness is a by-product of  & \textbf{engagement} in activities and relationships valued ... \\
07-a 3 & ...s is realized to a large extent through  & \textbf{engagement}. \\
07-a 3 & By Principle P2,  & \textbf{engagement} contributes to happiness only when the ... \\
07-a 3 & ...large part of happiness that depends on  & \textbf{engagement}, and attains less happiness than she wo... \\
07-a 3 & ...t of happiness that is realized through  & \textbf{engagement} is available only to agents who do. \\
07-a 3 & ...-17, on which meaning comes from active  & \textbf{engagement} with what is worthwhile. \\
07-a 4 & ...has not acquired the attitudes on which  & \textbf{engagement} depends, and Theorem 3.3 applies to her... \\
07-a 4 & Perhaps  & \textbf{engagement} is merely a reliable cause of happiness... \\
07-a 4 & ...ngs of engaged satisfaction without any  & \textbf{engagement} would then give the direct aim its targ... \\
07-a 4 & ...he one who lives it than a life of real  & \textbf{engagement} with the same experiences (ETH-05). \\
07-a 6 & Much of it consists in, or depends on,  & \textbf{engagement} with activities and people valued for t... \\
07-b 1 & ...cks the object-directed caring on which  & \textbf{engagement} depends (Principle P2 of EXI-07-a). \\
07-b 2 & ...e forgoes the happiness that depends on  & \textbf{engagement}. \\
07-b 2 & ... she lacks the caring that the forms of  & \textbf{engagement} in Principle P1 require, and Theorem 3.... \\
07-b 2 & ... and P2 of EXI-07-a, which concern what  & \textbf{engagement} is: one cannot be absorbed in a game wh... \\
10-a 3 & A search that terminates when  & \textbf{engagement} is restored selects, among available al... \\
10-a 4 & Boredom motivates a search for  & \textbf{engagement}; it does not select which alternatives ... \\
10-a 4 & ...XI-17, on which meaning requires active  & \textbf{engagement} with something worthwhile, and boredom ... \\
10-a 4 & ... fallible, indication of the absence of  & \textbf{engagement}. \\
10-a 5 & ...adaptive function of moving them toward  & \textbf{engagement}. \\
10-b 1 & ... restore the match, so that, given that  & \textbf{engagement} pays (its Assumption 2.4), the disposit... \\
10-b 3 & Fit and  & \textbf{engagement}.. \\
10-b 3 & ...ssumption 2.4 of EXI-10-a, the value of  & \textbf{engagement}. \\
10-b 4 & ...applies: boredom motivates a search for  & \textbf{engagement} but does not select the options on offe... \\
10-b 4 & The evidence that  & \textbf{engagement} yields concentration, involvement and b... \\
10-b 5 & ...t boredom is feedback about attentional  & \textbf{engagement}. \\
10-b 7 & ...a decline in boredom and an increase in  & \textbf{engagement} or learning. \\
10-b 8 & ... that a meaningful life requires active  & \textbf{engagement} with something worthwhile; the meaning ... \\
10-b 8 & ... the mind registers the absence of such  & \textbf{engagement}. \\
10-b 11 & ...auses and in its converse, that matched  & \textbf{engagement} absorbs, and leaves its over-challenge ... \\
10-b 11 & ...ese findings together with the value of  & \textbf{engagement}, and the studies of costly responses an... \\
11-a Abs & ...is good offers objects and relations of  & \textbf{engagement} that a naturalistic world lacks. \\
11-a 1 & ...eaning arises when the subject's active  & \textbf{engagement} meets the objective worth of its object... \\
11-a 2 & Other things equal,  & \textbf{engagement} with what is of greater worth, or engag... \\
11-a 2 & ...ement with what is of greater worth, or  & \textbf{engagement} that relates one to more of what is wor... \\
11-a 3 & By Theorem 3.3, such a person's  & \textbf{engagement} with what is worthwhile makes her life ... \\
11-a 4 &  & \textbf{Engagement} with that being, through love, worship,... \\
11-a 4 & ... the attempt to understand it, would be  & \textbf{engagement} with the greatest object of worth there... \\
11-a 4 & ... would make a life more meaningful than  & \textbf{engagement} with lesser objects, other things equal... \\
11-a 5 & ...he pursuit of meaning is not a solitary  & \textbf{engagement} with isolated goods. \\
11-a 5 & ...e uses is no longer purpose as such but  & \textbf{engagement} with the good, which is Definition 2.1. \\
11-a 6 & ...out God is meaningful: meaning requires  & \textbf{engagement} with what is worthwhile, and lives with... \\
11-a 7 & Meaning consists in active  & \textbf{engagement} with what is worthwhile, and the worth ... \\
11-a 7 & ...hwhile, and the worth of the objects of  & \textbf{engagement} does not derive from their being assign... \\
11-b Abs & ..., not the objects of worth available to  & \textbf{engagement}, so that Proposition 4.1 survives both ... \\
11-b 1 & EXI-11-a defines meaning as active  & \textbf{engagement} with what is worthwhile (Definition 2.1... \\
11-b 1 & ...ne engages and with the extent to which  & \textbf{engagement} relates one to more of what is worthwhi... \\
11-b 1 & ... and third are objects and relations of  & \textbf{engagement} of great worth that a naturalistic worl... \\
11-b 1 & ...inition 2.1 of EXI-11-a, any meaning is  & \textbf{engagement} with what is worthwhile. \\
11-b 1 &  & \textbf{Engagement} with God, or with a purpose that unifie... \\
11-b 1 & ...purpose that unifies one's projects, is  & \textbf{engagement} with what is worthwhile; so the meaning... \\
11-b 2 & ...d by God's existence, so the objects of  & \textbf{engagement} available without God remain available ... \\
11-b 3 & ...life by adding objects and relations of  & \textbf{engagement} of great worth. \\
14-a Abs & ...of proportion to the good or crowds out  & \textbf{engagement} with present goods. \\
17-a Abs & A meaningful life requires active  & \textbf{engagement} with something worthwhile, in the sense... \\
17-a Abs & ...rns on pairs of lives that are alike in  & \textbf{engagement} and in the worth of their objects and d... \\
17-a 1 & ...actively engaged, and the object of her  & \textbf{engagement} must be worthwhile independently of her... \\
17-a 1 & ...must be worthwhile independently of her  & \textbf{engagement}. \\
17-a 1 & ... it should not: meaning requires active  & \textbf{engagement} with something worthwhile, and achievem... \\
17-a 2 & An object of  & \textbf{engagement} (a project, a person, a practice, an in... \\
17-a 2 & Active  & \textbf{engagement} in the sense relevant to Principle P1 r... \\
17-a 2 & ... a real object and actually bear on it;  & \textbf{engagement} with an object that exists only in the ... \\
17-a 3 & Each object of  & \textbf{engagement} is worthwhile in the sense of Definitio... \\
17-a 3 & ...life is meaningful to the extent of her  & \textbf{engagement} with music. \\
17-a 3 & Her  & \textbf{engagement} is active, since study and attentive li... \\
17-a 3 & By Principle P1, her  & \textbf{engagement} contributes meaning. \\
17-a 3 & A meaningful life requires active  & \textbf{engagement} with something worthwhile; it does not ... \\
17-a 3 & ...lternatives: subjectivism fails because  & \textbf{engagement} with worthless objects confers no meani... \\
17-a 3 & ...ls because worthwhile work done without  & \textbf{engagement}, out of mere duty or inertia, is not wh... \\
17-a 3 & ...thless pursuit, or one produced without  & \textbf{engagement}, confers none. \\
17-a 3 & ...ndifferent to it fails the condition of  & \textbf{engagement}. \\
17-a 3 & ...larges the worth realized through one's  & \textbf{engagement}, it may enlarge meaning too. \\
17-a 4 & ... difficult activity counts, then active  & \textbf{engagement} in a demanding worthwhile activity enta... \\
17-a 4 & ...chievement becomes a restatement of the  & \textbf{engagement} condition. \\
17-a 4 & Not all  & \textbf{engagement} is difficult, moreover: love of a frien... \\
17-a 5 & A meaningful life requires active  & \textbf{engagement} with something worthwhile. \\
17-a 5 & It does not require that the  & \textbf{engagement} succeed. \\
17-a 5 & ...at the experience machine shows is that  & \textbf{engagement} must reach a real object, not that it m... \\
19-b 1 & ...or as a condition of the persistence of  & \textbf{engagement}. \\
19-b 1 & ...limited by death is not something whose  & \textbf{engagement} makes the life meaningful, and does not... \\
19-b 1 & The worth of the objects of  & \textbf{engagement}, knowledge, persons, beauty and the goo... \\
19-b 1 & ...ut categorical desires lacks the active  & \textbf{engagement} that meaning requires (EXI-17). \\
19-b 1 & ...1.3 concerns the eventual exhaustion of  & \textbf{engagement}, not its quality at any given time. \\
19-b 1 & An immortal's  & \textbf{engagement} with inquiry in the first thousand year... \\
19-b 1 & ...ry in the first thousand years would be  & \textbf{engagement} with what is worthwhile, as meaningful ... \\
19-b 2 & ...shows that the limit is something whose  & \textbf{engagement} confers meaning. \\
19-b 2 & ...-life meaning and of the persistence of  & \textbf{engagement}, then death benefits the one who dies, ... \\
19-b 2 & ... ordinary case, have left more time for  & \textbf{engagement}, as Corollary 3.3 of EXI-19-a notes (Fi... \\
19-b 3 & ...keeps the life within the span in which  & \textbf{engagement} persists. \\
20-a Abs & ...ul activity within one, requires active  & \textbf{engagement} with something worthwhile. \\
20-a Abs & ...rocess has the kind of mental life that  & \textbf{engagement} requires, which MND-08 and MND-05 leave... \\
20-a 1 & ...hat are worthwhile independently of her  & \textbf{engagement}. \\
20-a 1 & ...n of this view: meaning requires active  & \textbf{engagement} with something worthwhile, and does not... \\
20-a 2 &  & \textbf{Engagement}, in Wolf's sense, is not mere activity. \\
20-a 2 & ...are states with a phenomenal character,  & \textbf{engagement} requires consciousness. \\
20-a 2 & ...he evidence available, the condition of  & \textbf{engagement} is neither known to be satisfied nor kn... \\
20-a 2 & If  & \textbf{engagement} can be understood functionally, as a pa... \\
20-a 2 & Whether this is  & \textbf{engagement} or only its functional profile is the q... \\
20-b Abs & ...f EXI-17, on whether its production was  & \textbf{engagement} with something worthwhile, and that nei... \\
20-b Abs & Evidence bearing on  & \textbf{engagement}.. \\
20-b Abs & ... on whether it has the mental life that  & \textbf{engagement} requires, or on whether engagement can ... \\
20-b Abs & ...that engagement requires, or on whether  & \textbf{engagement} can be understood functionally. \\
20-b Abs & Whether the production was  & \textbf{engagement} may be settled by empirical work on the... \\
\end{longtable}
\kwichead{everyone}{11}
\begin{longtable}{@{}>{\raggedright}p{0.07\textwidth}>{\raggedleft}p{0.41\textwidth}@{\,}>{\raggedright\arraybackslash}p{0.47\textwidth}@{}}
01-d Abs & ...t meaning in this sense is available to  & \textbf{everyone}. \\
01-d 4 & ...el meet this condition, and are open to  & \textbf{everyone}. \\
01-d 4 & If meaning is available to  & \textbf{everyone}, the claim that a life is meaningful is... \\
01-d 4 & That positive meaning is available to  & \textbf{everyone} is compatible with its being hard to at... \\
01-d 5 & ... on the account of EXI-01-a, is open to  & \textbf{everyone} because its conditions are within every... \\
02-a 5 & On the deprivation account,  & \textbf{everyone} who dies at any age is deprived of some... \\
02-a 6 & ...ppen to a person, or that it is bad for  & \textbf{everyone} who dies. \\
07-a 1 & Aristotle observed that  & \textbf{everyone} agrees on the name of the highest good,... \\
11-a 2 & ...t meaning so understood is available to  & \textbf{everyone}. \\
17-a 3 & ...t meaning in this sense is available to  & \textbf{everyone}. \\
20-a 1 & ...t meaning in this sense is available to  & \textbf{everyone}. \\
\end{longtable}
\kwichead{exist}{29}
\begin{longtable}{@{}>{\raggedright}p{0.07\textwidth}>{\raggedleft}p{0.41\textwidth}@{\,}>{\raggedright\arraybackslash}p{0.47\textwidth}@{}}
02-a Abs & ... the event is bad, and that the subject  & \textbf{exist} at the time of the loss. \\
02-a 1 & ...is the privation of sensation; while we  & \textbf{exist} death is not present, and when death is... \\
02-a 1 & ... and when death is present we no longer  & \textbf{exist} (EpicurusMen ). \\
02-a 1 & ... birth, an eternity in which we did not  & \textbf{exist}, is nothing to us, and the time after o... \\
02-a 2 & ... and the relation does not require that  & \textbf{exist} at a time at which its other relata are... \\
02-a 3 & ...f Principle P1 holds, although does not  & \textbf{exist} at any time after at which the goods of... \\
02-a 3 & ... located at times at which she does not  & \textbf{exist}. \\
02-a 3 & Nothing in Principle P1 requires that  & \textbf{exist} at the times at which the goods she is ... \\
02-a 3 & ..., read as the claim that a subject must  & \textbf{exist} at the time at which an event is bad fo... \\
02-a 3 & ... the strong existence requirement: that  & \textbf{exist} at the time of the loss. \\
02-a 3 & ...the badness of death requires only that  & \textbf{exist}, at some time, and not that she exist a... \\
02-a 3 & ...t exist, at some time, and not that she  & \textbf{exist} at the times when the lost goods would ... \\
02-a 4 & For a person who dies at and began to  & \textbf{exist} at , the closest world in which dies la... \\
02-a 4 & ...an the closest world in which begins to  & \textbf{exist} substantially earlier than , and the la... \\
02-a 4 & A world in which begins to  & \textbf{exist} decades earlier must contain 's parents... \\
02-a 5 & ...t happened, or after, when she does not  & \textbf{exist}. \\
02-a 7 & ...ls because it requires that the subject  & \textbf{exist} at the time of the loss, and deprivatio... \\
02-b 2 & ...ad instead as a demand that the subject  & \textbf{exist} at the time at which the lost goods wou... \\
04-a 2 & ...player who knows that the game's stakes  & \textbf{exist} only within the game and still plays to... \\
09-a 3 & ... is fixed when the human being comes to  & \textbf{exist}; Lemma 3.2 concerns defining projects, ... \\
11-b 2 & ...rson may reasonably prefer that God not  & \textbf{exist}, since God's existence would deprive he... \\
12-a 2 & ...lternative is one in which she does not  & \textbf{exist}. \\
12-b 2 & ...py person who is never created does not  & \textbf{exist}, so no one has a complaint, and the imp... \\
13-a 3 & That she does not  & \textbf{exist} when the duty is to be discharged is no... \\
13-a 4 & The dead do not  & \textbf{exist}, so nothing can be owed to them, and th... \\
13-a 4 & ...t in the regard of people who no longer  & \textbf{exist}. \\
14-a 3 & ...since it constructs a past that did not  & \textbf{exist} and seeks to return to it; her reflecti... \\
14-a 4 & ...s, elsewhere in time, as distant places  & \textbf{exist} elsewhere in space. \\
19-a 3 & It begins with the person's coming to  & \textbf{exist}. \\
\end{longtable}
\kwichead{extent}{29}
\begin{longtable}{@{}>{\raggedright}p{0.07\textwidth}>{\raggedleft}p{0.41\textwidth}@{\,}>{\raggedright\arraybackslash}p{0.47\textwidth}@{}}
01-c 2 & ... such engagement consists to a positive  & \textbf{extent} in active engagement with a project of ... \\
01-d Abs & ...-a holds that life is meaningful to the  & \textbf{extent} that one engages with projects of objec... \\
01-d 3 & ...: the degree of meaning varies with the  & \textbf{extent} of engagement and the worth of its obje... \\
02-a Abs & ...t, an event is bad for a subject to the  & \textbf{extent} that her life as a whole goes worse, by... \\
03-a 2 & ...argues that a life is meaningful to the  & \textbf{extent} that its subject engages with projects ... \\
04-a 2 & ...de toward a situation is fitting to the  & \textbf{extent} that the evaluative features under whic... \\
04-a 3 & ...1 argues that life is meaningful to the  & \textbf{extent} that one engages with projects of objec... \\
06-b 2 & ... of EXI-01, a life is meaningful to the  & \textbf{extent} that its subject engages with projects ... \\
07-a Abs & ... plausible account, consists to a large  & \textbf{extent} in engagement with and success in activ... \\
07-a 1 & ...rve, and a life is well directed to the  & \textbf{extent} that it is directed at it. \\
07-a 2 & ...s of pursuits if is realized to a large  & \textbf{extent} through those pursuits, and the pursuit... \\
07-a 3 & ...le P1, happiness is realized to a large  & \textbf{extent} through engagement. \\
09-a 5 & ...1, on which a life is meaningful to the  & \textbf{extent} that its projects are of objective wort... \\
11-a Abs & ...of meaning: a life is meaningful to the  & \textbf{extent} that its subject actively engages with ... \\
11-a 2 & A life is meaningful to the  & \textbf{extent} that its subject is actively engaged wi... \\
11-a 2 & ...on 2.1, which makes meaning a matter of  & \textbf{extent}. \\
11-a 3 & ...er life is therefore meaningful to that  & \textbf{extent}. \\
11-b 1 & ... worth of what one engages and with the  & \textbf{extent} to which engagement relates one to more... \\
16-b 1 & ... and grief would be fitting only to the  & \textbf{extent} that substitutes are scarce. \\
17-a 1 & ...fe that brings nothing about is to that  & \textbf{extent} empty. \\
17-a 2 & A life is meaningful to the  & \textbf{extent} that its subject is actively engaged wi... \\
17-a 3 & ...he listener's life is meaningful to the  & \textbf{extent} of her engagement with music. \\
18-a 2 & To the  & \textbf{extent} that these attitudes serve the function... \\
18-b 1 & ...ent, which do not misrepresent; to that  & \textbf{extent} the range of necessary positive illusio... \\
19-b 1 & ...and EXI-17, a life is meaningful to the  & \textbf{extent} that its subject actively engages with ... \\
20-a 1 & ...activeness: a life is meaningful to the  & \textbf{extent} that its subject is actively engaged in... \\
20-a 3 & ...lection of answers is worthwhile to the  & \textbf{extent} that it contributes to knowledge or und... \\
20-a 4 & ...tivity is meaningful, by EXI-17, to the  & \textbf{extent} that the collection is worthwhile, and ... \\
20-b Abs & ...October 2026 are an instance, but their  & \textbf{extent} and accuracy are facts to be establishe... \\
\end{longtable}
\kwichead{given}{25}
\begin{longtable}{@{}>{\raggedright}p{0.07\textwidth}>{\raggedleft}p{0.41\textwidth}@{\,}>{\raggedright\arraybackslash}p{0.47\textwidth}@{}}
01-c 4 & It entails Assumption 3  & \textbf{given} Lemma 2.3, as it must if the conclusion... \\
01-d 1 & We take these results as  & \textbf{given} and ask what follows. \\
01-d 2 & ...dom determine how much of a life can be  & \textbf{given} to engagement and how wide the range of... \\
03-a Abs & ...re projects she takes to be worthwhile;  & \textbf{given} a character that remains connected to h... \\
03-a 2 & ...he question is how many kinds of them a  & \textbf{given} character can care about. \\
03-b 1 &  & \textbf{Given} a bounded repertoire of project types f... \\
05-a 2 & ...tion is rational only if it is fitting,  & \textbf{given} what the subject is in a position to kn... \\
06-a Abs & ... not aimed at us, and the goods are not  & \textbf{given} to us. \\
07-b Abs & We argue that, in the form  & \textbf{given} by EXI-07-a, it is neither. \\
09-a 3 & ...nd his claim is that no such project is  & \textbf{given} to human beings in advance of their exi... \\
10-b 1 & ...es tends to restore the match, so that,  & \textbf{given} that engagement pays (its Assumption 2.... \\
10-b 10 & ...e over-challenge clause, for the reason  & \textbf{given} in Section 3. \\
11-a 3 & ...se of this kind degrades the life it is  & \textbf{given} to. \\
11-b 2 & That a  & \textbf{given} person might reasonably prefer the worl... \\
12-a Abs & ... Schopenhauer says yes, and Benatar has  & \textbf{given} the claim a rigorous form. \\
12-a Abs & ...e person who lives it, and to have been  & \textbf{given} it is a non-comparative benefit. \\
15-b 3 & ...evidence supports is the probability of  & \textbf{given} her commitment. \\
15-b 4 & Expected-value reasoning, applied to a  & \textbf{given} set of options, does not determine thes... \\
17-a 1 & Bradford (2015) has  & \textbf{given} an account of achievement as the compet... \\
18-b 1 & (i) Whether a  & \textbf{given} positive illusion is psychologically ne... \\
18-b 1 & ...sion is psychologically necessary for a  & \textbf{given} agent is an empirical question, and the... \\
19-a Abs & ...es, since a later death would also have  & \textbf{given} her life a shape, and in the ordinary c... \\
19-a 1 & Religious traditions have  & \textbf{given} death a meaning of its own, as the wage... \\
19-b 1 & ...n of engagement, not its quality at any  & \textbf{given} time. \\
19-b 2 & ...ular one; a later death would also have  & \textbf{given} the life a shape and would, in the ordi... \\
\end{longtable}
\kwichead{includes}{15}
\begin{longtable}{@{}>{\raggedright}p{0.07\textwidth}>{\raggedleft}p{0.41\textwidth}@{\,}>{\raggedright\arraybackslash}p{0.47\textwidth}@{}}
01-c 2 & A life that  & \textbf{includes} such engagement consists to a positive ... \\
04-a 4 & ...titude toward one's seriousness, and it  & \textbf{includes} continuing to hold one's ends. \\
07-a 3 & ...TH-05's argument that well-being itself  & \textbf{includes} contact with reality. \\
07-a 4 & ...nse but the good life as a whole, which  & \textbf{includes} the goods that a plural account counts ... \\
09-b 1 & ...s the paper-knife from a concept, which  & \textbf{includes} both what the knife is for and a techni... \\
09-b 2 & ...what it is for, and the real definition  & \textbf{includes} the defining project. \\
10-a 3 & Boredom is aversive and  & \textbf{includes} a desire to be engaged otherwise (Defin... \\
10-a 3 & ...ive to meaning: the mismatch it signals  & \textbf{includes} the idleness of one's capacity to care,... \\
10-b 5 & ...on 2018), and EXI-10-a's Definition 2.2  & \textbf{includes} both as forms of a single mismatch. \\
10-b 9 & ... the direction of the finding; the list  & \textbf{includes} every study we could verify on the cost... \\
15-b 3 & ... more likely; the agent's evidence then  & \textbf{includes} the fact that she is committed to the e... \\
15-b 4 & But the hopeful stance  & \textbf{includes} more than choices among options. \\
15-b 4 & It  & \textbf{includes} where attention goes, which possibiliti... \\
17-a 2 & On Definition 2.2, the worthwhile  & \textbf{includes} the components of the good identified i... \\
18-a 2 & ...s a cost even when it helps: well-being  & \textbf{includes} contact with reality (ETH-05), and each... \\
\end{longtable}
\kwichead{intentional}{18}
\begin{longtable}{@{}>{\raggedright}p{0.07\textwidth}>{\raggedleft}p{0.41\textwidth}@{\,}>{\raggedright\arraybackslash}p{0.47\textwidth}@{}}
01-c Abs & ... is available to every agent capable of  & \textbf{intentional} action. \\
01-c 1 & ...and assumes that every agent capable of  & \textbf{intentional} action is able to engage in at least on... \\
01-c 2 & ...s anything at which an agent can direct  & \textbf{intentional} activity and about which she can care: ... \\
01-c 2 & Every agent capable of  & \textbf{intentional} action can direct some intentional acti... \\
01-c 2 & ...e of intentional action can direct some  & \textbf{intentional} activity, with care, at a worthwhile ob... \\
01-c 2 & If an agent directs  & \textbf{intentional} activity with care at a worthwhile obje... \\
01-c 2 & For every agent capable of  & \textbf{intentional} action, a life of positive meaning is a... \\
01-c 2 & By Assumption 2.2, the agent can direct  & \textbf{intentional} activity with care at a worthwhile obje... \\
01-c 4 & Some agents capable of  & \textbf{intentional} action are wholly isolated and wholly u... \\
01-c 4 & An agent capable of  & \textbf{intentional} action acts on something, and anything ... \\
01-d 1 & ... states that for every agent capable of  & \textbf{intentional} action a life of positive meaning is av... \\
01-d 2 & No agent capable of  & \textbf{intentional} action is condemned by her circumstance... \\
04-a 3 & Any agent capable of  & \textbf{intentional} action takes some ends seriously. \\
04-a 3 & ...life of any reflective agent capable of  & \textbf{intentional} action is absurd in the sense of Defini... \\
09-a 2 & Principle P1 is the standard  & \textbf{intentional} account of artefact kinds, and MET-13 a... \\
09-a 4 & ...as an artefact of its kind, through the  & \textbf{intentional} act of selecting and placing the wood, ... \\
10-a 4 & ...ferent object: PHE-04 argues that it is  & \textbf{intentional} toward time itself, not toward an activ... \\
10-b 8 & ...found boredom, in Heidegger's sense, is  & \textbf{intentional} toward time itself; the empirical liter... \\
\end{longtable}
\kwichead{least}{14}
\begin{longtable}{@{}>{\raggedright}p{0.07\textwidth}>{\raggedleft}p{0.41\textwidth}@{\,}>{\raggedright\arraybackslash}p{0.47\textwidth}@{}}
01-c Abs & ...very such agent is able to engage in at  & \textbf{least} one project of objective worth (Assumpt... \\
01-c 1 & ...entional action is able to engage in at  & \textbf{least} one of them (Assumption 3). \\
02-b 1 & ... must be experienced by her, or must at  & \textbf{least} occur while she exists. \\
09-a 1 & ...existentialism denies this: there is at  & \textbf{least} one being whose existence comes before ... \\
10-b 4 & The adaptive component is the  & \textbf{least} directly tested. \\
15-a 1 & ...resolve to act as if will obtain, or at  & \textbf{least} as if it might, which he called substan... \\
16-a 2 & ...that is a reason for 's love of , to at  & \textbf{least} the same degree. 's love of is transfer... \\
16-a 2 & ... the reasons for 's love of would be at  & \textbf{least} as strong reasons for 's loving instead... \\
16-a 3 & ...that is a reason for 's love of , to at  & \textbf{least} the same degree. \\
16-a 3 & ...nciple P2, each of these features is at  & \textbf{least} as strong a reason for 's loving . \\
16-a 3 & ... reasons, the reasons for loving are at  & \textbf{least} as strong reasons for loving instead, a... \\
16-b 1 & ... love that person instead of that is at  & \textbf{least} as strong as the reasons for 's love of... \\
16-b 1 & ...refore do not favor the other person at  & \textbf{least} as strongly. \\
17-a 3 & ...gful than the unsuccessful one's, or at  & \textbf{least} better. \\
\end{longtable}
\kwichead{life}{450}
\begin{longtable}{@{}>{\raggedright}p{0.07\textwidth}>{\raggedleft}p{0.41\textwidth}@{\,}>{\raggedright\arraybackslash}p{0.47\textwidth}@{}}
01-c Abs & EXI-01-a establishes that a  & \textbf{life} of positive meaning is available to eve... \\
01-c 1 & EXI-01-a defines a  & \textbf{life} as meaningful to the degree that it con... \\
01-c 2 & ... agent capable of intentional action, a  & \textbf{life} of positive meaning is available. \\
01-c 2 & A  & \textbf{life} that includes such engagement consists ... \\
01-c 3 & ...es a parallel result for achievement: a  & \textbf{life} can be meaningful through engagement wi... \\
01-c 5 & ...stributed: the presence, in every human  & \textbf{life}, of worthwhile things to care about and... \\
01-c 5 & ... is a project of objective worth, and a  & \textbf{life} that contains it has meaning. \\
01-d Abs & EXI-01-a holds that  & \textbf{life} is meaningful to the extent that one en... \\
01-d Abs & ... scale of one's projects; that no one's  & \textbf{life} is condemned to meaninglessness by circ... \\
01-d Abs & It does not imply that every  & \textbf{life} is meaningful, that all lives are equal... \\
01-d 1 & ...y agent capable of intentional action a  & \textbf{life} of positive meaning is available. \\
01-d 1 & ... does not settle whether any particular  & \textbf{life} is meaningful, which depends on what is... \\
01-d 2 & Whether a  & \textbf{life} is meaningful does not depend on whethe... \\
01-d 2 & A  & \textbf{life} can be meaningful though its subject ha... \\
01-d 2 & ...er circumstances alone to a meaningless  & \textbf{life}. \\
01-d 2 & ...justice, affect the degree of meaning a  & \textbf{life} can attain and the ease with which it c... \\
01-d 2 & ...ion and freedom determine how much of a  & \textbf{life} can be given to engagement and how wide... \\
01-d 3 & ...irst, from availability to actuality: a  & \textbf{life} in which the agent engages with nothing... \\
01-d 4 & ...available to everyone, the claim that a  & \textbf{life} is meaningful is cheap and of no evalua... \\
01-d 4 & EXI-04 grants that  & \textbf{life} is absurd in Nagel's sense and argues t... \\
01-d 5 & ...its degree is not, and how much of it a  & \textbf{life} attains depends on what is done with it... \\
02-a Abs & ...ad for a subject to the extent that her  & \textbf{life} as a whole goes worse, by the measure o... \\
02-a 1 & ...-dimensional conception of the person's  & \textbf{life}. \\
02-a 2 & ...se in which it sustains a psychological  & \textbf{life}, and persists through psychological con... \\
02-a 2 & ...ganism that is sustains a psychological  & \textbf{life}, and no psychological life continuous w... \\
02-a 2 & ...sychological life, and no psychological  & \textbf{life} continuous with 's is sustained at any ... \\
02-a 2 & For a possible  & \textbf{life} of , is the sum of the intrinsic values... \\
02-a 2 & ...for of an event is , where is 's actual  & \textbf{life} and is the life has in the closest poss... \\
02-a 2 & ...is , where is 's actual life and is the  & \textbf{life} has in the closest possible world in wh... \\
02-a 2 & ...difference it makes to the value of the  & \textbf{life} as a whole, and the lives compared are ... \\
02-a 2 & ...the value it assigns is a fact about 's  & \textbf{life} taken as a whole. \\
02-a 3 & ...crupulous never to let her know it; her  & \textbf{life} goes on, as experienced, as before. \\
02-a 3 & Her  & \textbf{life} contains friendships she believes genui... \\
02-a 3 & ....2, relationships are among the goods a  & \textbf{life} contains, and false friendships are not... \\
02-a 3 & So her  & \textbf{life} has less value than the life she would ... \\
02-a 3 & So her life has less value than the  & \textbf{life} she would have had without the betrayal... \\
02-a 3 & ... is a misfortune for him because of the  & \textbf{life} it has taken from him. \\
02-a 3 & It is a fact about and the value of her  & \textbf{life} as a whole, and it obtains, tenselessly... \\
02-a 3 & ...tenselessly, because exists and has the  & \textbf{life} she has. \\
02-a 3 & ...requires that those goods belong to the  & \textbf{life} she would have had. \\
02-a 3 & If dies at , and the  & \textbf{life} would have had had she not died at cont... \\
02-a 3 & ...reater balance of goods than her actual  & \textbf{life}, then her death is bad for her. \\
02-a 3 & ...e she would otherwise have gone on to a  & \textbf{life} worth living, this condition is met. \\
02-a 3 & Her actual  & \textbf{life} ends at ; the alternative life agrees w... \\
02-a 3 & ...r actual life ends at ; the alternative  & \textbf{life} agrees with it up to and continues beyo... \\
02-a 3 & ...in, and if those circumstances made her  & \textbf{life} worth living, the continuation has a po... \\
02-a 3 & ... twenty that cuts off decades of a good  & \textbf{life} is very bad for the one who dies; a dea... \\
02-a 3 & ...ters is not the manner of dying but the  & \textbf{life} that would otherwise have followed. \\
02-a 4 & ...he latter often contains no longer good  & \textbf{life} for . \\
02-a 4 & ...t is possible at all, it is a different  & \textbf{life} from the beginning, and nothing guarant... \\
02-a 4 & ...t is compared is the value of the whole  & \textbf{life}, not its length. \\
02-a 4 & ... world where she begins earlier and her  & \textbf{life} is better, and the deprivation account ... \\
02-a 5 & ...1 is a fact about the value of 's whole  & \textbf{life} and is not located at any time within i... \\
02-a 5 & ...riendships are no more genuine, and her  & \textbf{life} no better. \\
02-a 5 & It says that death is bad when the  & \textbf{life} it ends would otherwise have continued ... \\
02-a 5 & ...adness of death to diminish over a long  & \textbf{life}, as the goods that would have followed ... \\
02-a 5 & ...fewer; EXI-03 argues that in an endless  & \textbf{life} they eventually give out. \\
02-a 5 & ...nition 2.2 are goods of a psychological  & \textbf{life}, which the corpse does not have. \\
02-a 5 & ...she is deprived of is located after her  & \textbf{life} ends. \\
02-a 7 & ...badness is a comparative fact about her  & \textbf{life} as a whole, and its subject is the pers... \\
02-a 7 & ...ts subject is the person who lived that  & \textbf{life}. \\
02-b 2 & ...m is comparative and concerns the whole  & \textbf{life}: a life that ends earlier contains fewe... \\
02-b 2 & ...parative and concerns the whole life: a  & \textbf{life} that ends earlier contains fewer goods ... \\
03-a Abs & ...ms argued that it would not: an endless  & \textbf{life} would either preserve the character of ... \\
03-a Abs & ...e that character so far that the future  & \textbf{life} would not be hers in the sense that mat... \\
03-a Abs & ...ended pursuits could sustain an endless  & \textbf{life}, and that the second horn rests on a de... \\
03-a Abs & ...rs in survival, we show that an endless  & \textbf{life} either exhausts the categorical desires... \\
03-a Abs & ...r ceases to be, in its later reaches, a  & \textbf{life} that matters to the person now. \\
03-a 1 & ...e refuses to take it again, because her  & \textbf{life} has become one of boredom and coldness. \\
03-a 1 & ... used the case to argue that an endless  & \textbf{life} would be meaningless, and that it is no... \\
03-a 1 & Williams argued that an endless  & \textbf{life} would face a dilemma. \\
03-a 1 & ... would eventually be exhausted, and her  & \textbf{life} would become boring. \\
03-a 1 & ...he future person would not be one whose  & \textbf{life} she now has reason to want, since what ... \\
03-a 1 & ... food, music and friendship, and that a  & \textbf{life} could be sustained indefinitely by a su... \\
03-a 1 & ...ifies the categorical desires that give  & \textbf{life} its point with desires for projects the... \\
03-a 2 & ...its satisfaction requires her continued  & \textbf{life} and it gives her a reason to go on livi... \\
03-a 2 & A categorical desire gives  & \textbf{life} its point if its object is a project th... \\
03-a 2 & ...order to drink it, but it does not give  & \textbf{life} its point. \\
03-a 2 & EXI-01 argues that a  & \textbf{life} is meaningful to the extent that its su... \\
03-a 2 & The desires that give  & \textbf{life} its point are desires to engage in such... \\
03-a 2 & A  & \textbf{life} is endless if for every time at which t... \\
03-a 2 & ...grounds a categorical desire that gives  & \textbf{life} its point for that person. \\
03-a 2 & ...o a person, in her concern for a future  & \textbf{life}, is carried mostly by psychological con... \\
03-a 2 & ...e of connectedness below which a future  & \textbf{life} gives the person now little of what mat... \\
03-a 2 & ...already done is not a desire that gives  & \textbf{life} its point. \\
03-a 3 & In an endless  & \textbf{life} in which the person's character remains... \\
03-a 3 & ...he has no categorical desire that gives  & \textbf{life} its point. \\
03-a 3 & The  & \textbf{life} is endless (Definition 2.2), so there i... \\
03-a 3 & ...grounds a categorical desire that gives  & \textbf{life} its point. \\
03-a 3 & In an endless  & \textbf{life} in which the person's character does no... \\
03-a 3 & ...eshold, there are times after which the  & \textbf{life} gives her, as she now is, little of wha... \\
03-a 3 & ...t character, then by Assumption 2.5 the  & \textbf{life} after that time gives her now little of... \\
03-a 3 & So the later reaches of the  & \textbf{life} are not something she now has reason to... \\
03-a 3 & An endless  & \textbf{life} is not desirable for the one who would ... \\
03-a 3 & ...haust the categorical desires that give  & \textbf{life} its point, or in its later reaches ceas... \\
03-a 3 & ..., or in its later reaches cease to be a  & \textbf{life} that gives her what matters. \\
03-a 3 & Every endless  & \textbf{life} falls under Lemma 3.1 or Lemma 3.2, acc... \\
03-a 3 & ...d without categorical desires that give  & \textbf{life} its point, so the person's life becomes... \\
03-a 3 & ...at give life its point, so the person's  & \textbf{life} becomes, from some time on, pointless; ... \\
03-a 3 & In neither case does the endless  & \textbf{life} offer her what she wants when she wants... \\
03-a 3 & ...wants not to die: the continuation of a  & \textbf{life} with point, for herself as she is. \\
03-a 3 & Theorem 3.3 is a claim about an endless  & \textbf{life}, not about a longer one. \\
03-a 3 & ...l has categorical desires that give her  & \textbf{life} its point is deprived by her death; a p... \\
03-a 3 & ...deprivation of EXI-02 diminishes as the  & \textbf{life} approaches the point at which Lemma 3.1... \\
03-a 3 & ...which Lemma 3.1 applies, and an endless  & \textbf{life} is guaranteed to reach that point or to... \\
03-a 4 & A  & \textbf{life} could be sustained endlessly by these (... \\
03-a 4 & ...ble pleasures are goods, and an endless  & \textbf{life} containing them would contain goods at ... \\
03-a 4 & But the desire for them does not give  & \textbf{life} its point in the sense of Definition 2.... \\
03-a 4 & A  & \textbf{life} sustained only by such desires is the l... \\
03-a 4 & ...e sustained only by such desires is the  & \textbf{life} Elina Makropulos describes, in which th... \\
03-a 4 & Over an endless  & \textbf{life} the new theorems available to her becom... \\
03-a 4 & ...y constitution or by a suitably managed  & \textbf{life}. \\
03-a 4 & ...ereby have them again; she would have a  & \textbf{life} without point that does not trouble her... \\
03-a 4 & ...heorem 3.3 concerns whether the endless  & \textbf{life} is desirable, and a life whose pointles... \\
03-a 4 & ...er the endless life is desirable, and a  & \textbf{life} whose pointlessness is not felt is not ... \\
03-a 4 & ...ntlessness is not felt is not thereby a  & \textbf{life} with a point. \\
03-a 4 & Such a  & \textbf{life} would always have categorical desires, ... \\
03-a 4 & ...is glad to go on, and each stage of the  & \textbf{life} may be good for the person who lives it... \\
03-a 4 & But the question is whether the endless  & \textbf{life} is desirable for the person who now con... \\
03-a 4 & The  & \textbf{life} described is something like a successio... \\
03-a 5 & ... argument does not show that any finite  & \textbf{life} is too long, or that there is a determi... \\
03-a 5 & ...es not show that the limits of a mortal  & \textbf{life} are what give it its shape, which is a ... \\
03-a 6 & An endless  & \textbf{life} would eventually exhaust the categorica... \\
03-a 6 & ...haust the categorical desires that give  & \textbf{life} its point, or else cease, in its later ... \\
03-a 6 & ...se cease, in its later reaches, to be a  & \textbf{life} that gives its subject now what matters... \\
03-b Abs & EXI-03-a argues that an endless  & \textbf{life} would eventually exhaust the categorica... \\
03-b Abs & ...haust the categorical desires that give  & \textbf{life} its point, or cease to be a life that m... \\
03-b Abs & ...t give life its point, or cease to be a  & \textbf{life} that matters to the person now. \\
03-b Abs & ... is not a categorical desire that gives  & \textbf{life} its point, since its object is not a pr... \\
03-b Abs & Second, the undesirability of endless  & \textbf{life} is compatible with the badness of death... \\
03-b Abs & ...e exhaustion lies far beyond any actual  & \textbf{life}. \\
03-b Abs & ...und on the desirable length of a finite  & \textbf{life}: it shows only that some finite length ... \\
03-b 1 & ...ose satisfaction requires her continued  & \textbf{life}; it gives life its point if its object ... \\
03-b 1 & ...n requires her continued life; it gives  & \textbf{life} its point if its object is a project sh... \\
03-b 1 & ...3 of EXI-03-a concludes that an endless  & \textbf{life} is not desirable for the one who would ... \\
03-b 2 & ... is not a categorical desire that gives  & \textbf{life} its point. \\
03-b 2 & A categorical desire gives  & \textbf{life} its point only if its object is a proje... \\
03-b 2 & ... to live forever is the continuation of  & \textbf{life} itself, without specification of what i... \\
03-b 2 & ...d eventually contain nothing that gives  & \textbf{life} its point. \\
03-b 2 & That an endless  & \textbf{life} is not desirable does not imply that de... \\
03-b 2 & ...curs only after some time in an endless  & \textbf{life}; the theorem does not say when. \\
03-b 2 & At such a time the  & \textbf{life} the person would have had, had she not ... \\
03-b 2 & ...she not died, contains goods her actual  & \textbf{life} lacks, and by the deprivation account o... \\
03-b 2 & The theorem concerns the limit of a  & \textbf{life}; the badness of death concerns its actu... \\
03-b 2 & ... establishes that some finite length of  & \textbf{life} would suffice to exhaust the categorica... \\
03-b 2 & A  & \textbf{life} many times longer than ours might be de... \\
03-b 3 & If the desire for endless  & \textbf{life} is not categorical, then the fear of lo... \\
03-b 4 & Categorical desires give  & \textbf{life} its point and, being finite in an endle... \\
03-b 4 & ...s point and, being finite in an endless  & \textbf{life}, set a limit on the value of continuati... \\
04-a Abs & Is  & \textbf{life} absurd, and if so, how should one respo... \\
04-a Abs & ...ot in the smallness or brevity of human  & \textbf{life} but in a collision within the reflectiv... \\
04-a Abs & We argue that  & \textbf{life} is absurd in this sense, because both o... \\
04-a 1 & ...ny arguments for the absurdity of human  & \textbf{life} appeal to scale. \\
04-a 1 & If a  & \textbf{life} of seventy years is pointless, a life o... \\
04-a 1 & ...a life of seventy years is pointless, a  & \textbf{life} of eternity would be no less so; if our... \\
04-a 1 & ..., is a structural feature of reflective  & \textbf{life}: the collision between the seriousness ... \\
04-a 1 & Justifications within a  & \textbf{life} come to an end in ends that are not the... \\
04-a 1 & We argue that  & \textbf{life} is absurd in Nagel's sense even if thos... \\
04-a 2 & A  & \textbf{life} is absurd if its subject takes some end... \\
04-a 3 & EXI-01 argues that  & \textbf{life} is meaningful to the extent that one en... \\
04-a 3 & The  & \textbf{life} of any reflective agent capable of inte... \\
04-a 3 & By Definition 2.3, her  & \textbf{life} is absurd. \\
04-a 3 & ... a fitting response to the absurdity of  & \textbf{life}, and it is the only one of the three re... \\
04-a 4 & ...ch we pursue them is justified, and the  & \textbf{life} of someone engaged in them is not absur... \\
04-a 4 & The engaged and worthwhile  & \textbf{life} is less exposed to despair than the lif... \\
04-a 4 & ...ife is less exposed to despair than the  & \textbf{life} devoted to trivial ends, since despair ... \\
04-a 5 & ...ument does not show that any particular  & \textbf{life} is meaningful or worthwhile; that quest... \\
04-a 6 &  & \textbf{Life} is absurd in Nagel's sense: a reflectiv... \\
04-b Abs & EXI-04-a argues that  & \textbf{life} is absurd in Nagel's sense and that the... \\
04-b 1 & A  & \textbf{life} can be absurd without being futile, and... \\
04-b 1 & ...g futile, and a task can be futile in a  & \textbf{life} that is not absurd. \\
04-b 1 & By Theorem 3.3 of EXI-04-a, the  & \textbf{life} of every reflective agent is absurd, in... \\
04-b 1 & ...flective agent is absurd, including the  & \textbf{life} of an agent whose projects succeed and ... \\
04-b 1 & ... physician whose patients recover; that  & \textbf{life} is not futile. \\
04-b 1 & ...ing up the detached standpoint, and her  & \textbf{life} would then not be absurd in the sense o... \\
04-b 1 & ...id; its claim concerns the absurdity of  & \textbf{life}, which is a different thing. \\
04-b 3 & ...the second sense of Definition 1.1, his  & \textbf{life} is not meaningful however happy he is, ... \\
04-b 3 & ...s: such a Sisyphus is fulfilled but his  & \textbf{life} is not meaningful, because his object i... \\
04-b 3 & His happiness may be real; his  & \textbf{life}, as Camus describes it, is a paradigm o... \\
04-b 4 & Death limits a  & \textbf{life} without making its activities futile in... \\
05-a Abs & ...stinct from particular threats to one's  & \textbf{life}, is certain. \\
05-a 3 & ...th EXI-02, the end of the psychological  & \textbf{life} that the organism sustains. \\
05-a 4 & ...r threat, or the possibility that one's  & \textbf{life} will be short. \\
05-a 4 & ...ity, and where it is, it is a different  & \textbf{life} whose value is not guaranteed to be gre... \\
06-a 1 & ...losing his catalogue of the miseries of  & \textbf{life}, judged that the course of nature sugge... \\
06-a 4 & ...nd eventually the end of conditions for  & \textbf{life}. \\
06-a 4 & ... the same laws that bring death brought  & \textbf{life}. \\
06-b Abs & ...ndifference is irrelevant to meaning in  & \textbf{life}, though not to every question about mea... \\
06-b 2 & On the account of EXI-01, a  & \textbf{life} is meaningful to the extent that its su... \\
06-b 2 & It denies that meaning in  & \textbf{life} requires that good. \\
07-a Abs & Is happiness the goal of  & \textbf{life}? \\
07-a Abs & ...nd conclude that happiness is a good of  & \textbf{life} but not its goal. \\
07-a 1 & On both views happiness is the goal of  & \textbf{life} in a strong sense: it is the end that a... \\
07-a 1 & ...he end that all other ends serve, and a  & \textbf{life} is well directed to the extent that it ... \\
07-a 1 & ...tieth year by a change in his theory of  & \textbf{life}: those only are happy, he came to think... \\
07-a 1 & ...is family: happiness is not the goal of  & \textbf{life}, because it is a by-product of other pu... \\
07-a 3 & Happiness is not the goal of  & \textbf{life}, in the sense of the end that all other... \\
07-a 3 & ... that all other ends of a well-directed  & \textbf{life} serve. \\
07-a 3 & If happiness were the goal of  & \textbf{life} in that sense, a well-directed life wou... \\
07-a 3 & ... of life in that sense, a well-directed  & \textbf{life} would have happiness as its sole final ... \\
07-a 3 & By Theorem 3.3, such a  & \textbf{life} defeats the aim by which it is directed... \\
07-a 3 & A  & \textbf{life} directed at its goal in a way that defe... \\
07-a 3 & So a well-directed  & \textbf{life} does not have happiness as its sole fin... \\
07-a 4 & ...not happiness in our sense but the good  & \textbf{life} as a whole, which includes the goods th... \\
07-a 4 & The Aristotelian claim that the good  & \textbf{life} is the end of action is compatible with... \\
07-a 4 & ...s chosen for its own sake within such a  & \textbf{life} is not the agent's feeling good, but th... \\
07-a 4 & ...ular goods of experience, embedded in a  & \textbf{life} whose final ends lie elsewhere, and the... \\
07-a 4 & ...estion whether happiness is the goal of  & \textbf{life} concerns the lives human beings can lea... \\
07-a 4 & ...nd, the experience machine shows that a  & \textbf{life} of produced satisfaction is worse for t... \\
07-a 4 & ...s worse for the one who lives it than a  & \textbf{life} of real engagement with the same experi... \\
07-a 4 & ... on its own terms, it would show that a  & \textbf{life} whose goal is happiness so understood i... \\
07-a 4 & ...happiness so understood is not the best  & \textbf{life} for its subject, which is a further rea... \\
07-a 4 & ...n to deny that happiness is the goal of  & \textbf{life}. \\
07-a 6 & Happiness is not the goal of  & \textbf{life}. \\
07-b 2 & Happiness then figures in her  & \textbf{life} as the by-product of Definition 2.3 of ... \\
07-b 3 & ...asing to treat happiness as the goal of  & \textbf{life}. \\
08-a Abs & The ideal asks that a  & \textbf{life} be one's own rather than borrowed from ... \\
09-a Abs & ...tence, because a defining project for a  & \textbf{life} can be fixed only by the agent who live... \\
09-a 2 & ...lf afterwards, he means that a person's  & \textbf{life} comes to be for something only through ... \\
09-a 2 & ...eing capable of adopting ends, what its  & \textbf{life} is for, in the sense of Definition 2.1,... \\
09-a 2 & Another's intention for that  & \textbf{life} can give it reasons, but does not fix i... \\
09-a 2 & ...states that the success conditions of a  & \textbf{life}, as opposed to those of a tool, are int... \\
09-a 2 & ...hich one values oneself and finds one's  & \textbf{life} worth living (Korsgaard 1996). \\
09-a 3 & ...is capable of adopting ends and what 's  & \textbf{life} is for is fixed by the ends adopts. \\
09-a 3 & ...e lives, none of which fails as a human  & \textbf{life} by departing from the others. \\
09-a 3 & ...ition of a human being settles what her  & \textbf{life} is for. \\
09-a 4 & ...ixes a general form that any good human  & \textbf{life} takes, not what a particular life is fo... \\
09-a 4 & ...human life takes, not what a particular  & \textbf{life} is for. \\
09-a 5 & Principle P2 says that what a  & \textbf{life} is for is fixed by its agent's ends; it... \\
09-a 5 & ...d is consistent with EXI-01, on which a  & \textbf{life} is meaningful to the extent that its pr... \\
09-b 2 & ...ays what she is without saying what her  & \textbf{life} is for. \\
09-b 2 & ... as a human being in virtue of what her  & \textbf{life} is for. \\
09-b 2 & ...I-09-a, no defining project for a human  & \textbf{life} is fixed before it exists; whatever sta... \\
09-b 2 & ...xists; whatever standard of success her  & \textbf{life} has as hers is fixed by the ends she ad... \\
09-b 2 & A  & \textbf{life} cannot fail as a human life by departin... \\
09-b 2 & A life cannot fail as a human  & \textbf{life} by departing from a project it never ha... \\
09-b 2 & A  & \textbf{life} can be imprudent, wicked or, in the sen... \\
09-b 2 & ...efining project, and they do not make a  & \textbf{life} a defective specimen of anything. \\
09-b 2 & ... organism; that is not a failure of her  & \textbf{life} to be what it is for. \\
10-a 5 & ...nt with EXI-03's thesis that an endless  & \textbf{life} would exhaust the desires that give it ... \\
10-a 5 & ...sires that give it its point: in such a  & \textbf{life}, boredom would become a correct signal ... \\
10-b 3 & (2009) found that a lack of meaning in  & \textbf{life} predicted boredom cross-sectionally, ov... \\
10-b 8 & EXI-17 holds that a meaningful  & \textbf{life} requires active engagement with somethi... \\
10-b 8 & EXI-03 holds that an endless  & \textbf{life} would exhaust the categorical desires t... \\
10-b 8 & ...haust the categorical desires that give  & \textbf{life} its point; EXI-10-a draws the consequen... \\
10-b 10 & ..., in laboratories, schools and everyday  & \textbf{life}. \\
11-a Abs & Can  & \textbf{life} be meaningful without God? \\
11-a Abs & ...ng requires a purpose assigned to human  & \textbf{life} by its creator, or a relation to an inf... \\
11-a Abs & We argue that  & \textbf{life} can be meaningful without God, because ... \\
11-a Abs & ...itting-fulfilment account of meaning: a  & \textbf{life} is meaningful to the extent that its su... \\
11-a 1 & On purpose theories, a  & \textbf{life} is meaningful insofar as it fulfils a p... \\
11-a 1 & ...d without God, and the question whether  & \textbf{life} would be meaningful in such a world is ... \\
11-a 1 & ...nce: whether the conditions that make a  & \textbf{life} meaningful include God's existence, a r... \\
11-a 1 & ... practical interest, namely whether the  & \textbf{life} of someone who does not believe in God ... \\
11-a 2 & A  & \textbf{life} is meaningful to the extent that its su... \\
11-a 2 & A  & \textbf{life} has a cosmic purpose if a creator of th... \\
11-a 2 & A  & \textbf{life} fulfils its cosmic purpose if it plays ... \\
11-a 2 & The purpose theory holds that a  & \textbf{life} is meaningful only if it has a cosmic p... \\
11-a 2 & ... to more of what is worthwhile, makes a  & \textbf{life} more meaningful. \\
11-a 2 & ...the insight in Nozick's linkage view: a  & \textbf{life} gains meaning by connection with value ... \\
11-a 3 & ...pose is not sufficient for a meaningful  & \textbf{life}. \\
11-a 3 & ...tacles for its amusement, and suppose a  & \textbf{life} that fulfils that purpose. \\
11-a 3 & By Definition 2.1, such a  & \textbf{life} is not made meaningful by its purpose. \\
11-a 3 & ...es: a purpose of this kind degrades the  & \textbf{life} it is given to. \\
11-a 3 & A  & \textbf{life} that has no cosmic purpose, or does not... \\
11-a 3 & ...fies Definition 2.1, whether or not her  & \textbf{life} has a cosmic purpose and whether or not... \\
11-a 3 & Her  & \textbf{life} is therefore meaningful to that extent. \\
11-a 3 & The  & \textbf{life} of someone who does not believe in God ... \\
11-a 3 & ...ement with what is worthwhile makes her  & \textbf{life} meaningful. \\
11-a 4 & ...ork, if true, would make available to a  & \textbf{life} more meaning than it could have in a wo... \\
11-a 4 & ...re is; by Principle P2, it would make a  & \textbf{life} more meaningful than engagement with le... \\
11-a 4 & Each of these relates the  & \textbf{life} to more of what is worthwhile, and by P... \\
11-a 4 & ...available, and Theorem 3.3 ensures that  & \textbf{life} is meaningful nonetheless. \\
11-a 5 & ...9 argues that death, far from voiding a  & \textbf{life}, gives it a shape by limiting it. \\
11-a 5 & On the linkage view, a finite  & \textbf{life} has meaning by connection with somethin... \\
11-a 5 & A  & \textbf{life} linked to these is meaningful without f... \\
11-a 5 & The engaged and worthwhile  & \textbf{life} of a person without such a framework is... \\
11-a 5 & The word meaning in the meaning of  & \textbf{life} just asks for a purpose. \\
11-a 5 & A  & \textbf{life} without a purpose assigned to it has no... \\
11-a 5 & The question of what makes a  & \textbf{life} meaningful is asked and answered by bel... \\
11-a 5 & ...nd unbelievers alike in terms of what a  & \textbf{life} is devoted to and whether that is worth... \\
11-a 5 & ... a purpose may be assigned that makes a  & \textbf{life} worse. \\
11-a 5 & If a  & \textbf{life} can be fully meaningful without God, no... \\
11-a 5 & ... for theism to add; if theism adds, the  & \textbf{life} without God is deficient in meaning. \\
11-a 5 & A  & \textbf{life} can be meaningful, and deeply so, while... \\
11-a 5 & That a  & \textbf{life} could be more meaningful is not a defic... \\
11-a 5 & ...Theorem 3.3; it is true of every finite  & \textbf{life}. \\
11-a 6 & Nor does it show that every  & \textbf{life} without God is meaningful: meaning requ... \\
11-a 7 &  & \textbf{Life} can be meaningful without God. \\
11-b 1 & ...t concerns availability: it says what a  & \textbf{life} could engage with, not that every life ... \\
11-b 1 & ... life could engage with, not that every  & \textbf{life} in a theistic world is more meaningful ... \\
11-b 1 & ...tic world is more meaningful than every  & \textbf{life} in a naturalistic one. \\
11-b 3 & ...would add to the meaning available in a  & \textbf{life} by adding objects and relations of enga... \\
12-a Abs & ...nly for lives containing no good, and a  & \textbf{life} with more good than bad is better than ... \\
12-a Abs & ...e that matters to the one who lives the  & \textbf{life}: although never existing is not a condi... \\
12-a Abs & ...ting is not a condition anyone is in, a  & \textbf{life} with more good than bad is good for the... \\
12-a Abs & We conclude that a  & \textbf{life} containing more good than bad is better... \\
12-a 1 & Schopenhauer held that  & \textbf{life} is a business whose proceeds do not cov... \\
12-a 1 & The first is in what sense a  & \textbf{life} is better than non-existence, since no ... \\
12-a 2 & A  & \textbf{life} has more good than bad if the goods it ... \\
12-a 2 & ...he exists is impersonally better if her  & \textbf{life} has more good than bad (ETH-20). \\
12-a 2 & ...omparatively if it causes her to have a  & \textbf{life} that is good for her. \\
12-a 2 & But we commonly judge that a person's  & \textbf{life} is good for her, that a pleasure is goo... \\
12-a 3 & A  & \textbf{life} containing more good than bad is better... \\
12-a 3 & For the second clause: a  & \textbf{life} with more good than bad is good for its... \\
12-a 3 & ...inciple P2, to be caused to have such a  & \textbf{life} is a non-comparative benefit. \\
12-a 3 & So the  & \textbf{life} is better than non-existence in both se... \\
12-a 3 & To yield the conclusion that every  & \textbf{life} is worse than non-existence, the argume... \\
12-a 3 & ...the argument would need to show that no  & \textbf{life} has more good than bad. \\
12-a 3 & These are not bads in the  & \textbf{life} but goods it does not contain; to count... \\
12-a 4 & It says that my  & \textbf{life} is good for me, which is a non-comparat... \\
12-a 4 & ...is is conditional on the balance of the  & \textbf{life}. \\
12-a 4 & Even if a good  & \textbf{life} is better than never existing, no one c... \\
12-a 4 & ...r existing, no one can guarantee that a  & \textbf{life} will be good. \\
12-a 4 & ...er uncertainty, not the comparison of a  & \textbf{life} with non-existence. \\
12-a 5 & ...t better never to have been born, for a  & \textbf{life} containing more good than bad. \\
12-a 5 & Such a  & \textbf{life} makes the world better than it would ha... \\
12-a 5 & ...tar's argument from the poor quality of  & \textbf{life} contests which lives have more good tha... \\
12-b Abs & ...re fails, and EXI-12-a concludes that a  & \textbf{life} containing more good than bad is better... \\
12-b 1 & ... is a duty not to create a person whose  & \textbf{life} would be miserable, but no duty to crea... \\
12-b 1 & ...e, but no duty to create a person whose  & \textbf{life} would be happy. \\
12-b 2 & ...e miserable person, once created, has a  & \textbf{life} that is bad for her and so has a compla... \\
12-b 2 & ...ason: by Theorem 3.1 of EXI-12-a a good  & \textbf{life} is a non-comparative benefit to the one... \\
12-b 3 & ...t support to the basic asymmetry, and a  & \textbf{life} with more good than bad remains better ... \\
13-a Abs & ... to rely on its doing so, and plans her  & \textbf{life} on that reliance; the duty to the dead ... \\
13-a 1 & ...at happens after death can affect how a  & \textbf{life} went is older: Aristotle asks whether t... \\
13-a 2 & ... the deprivation account of EXI-02: her  & \textbf{life} as a whole goes worse than it would oth... \\
13-a 2 & ...erests rather than because they cut her  & \textbf{life} short. \\
13-a 3 & ...by Principle P1 their effect on how her  & \textbf{life} went is smaller. \\
13-b 2 & ...us event is simply a fact about how her  & \textbf{life} went. \\
14-a 1 & ...social connectedness and the sense that  & \textbf{life} is meaningful (Wildschut et al. 2006; S... \\
15-b 2 & ...xcess is the player who reorganizes her  & \textbf{life} around the draw (Lemma 3.2 of EXI-15-a)... \\
17-a Abs & Does a meaningful  & \textbf{life} require achievement? \\
17-a Abs & ...e success a central place: a meaningful  & \textbf{life} is one that makes a difference, leaves ... \\
17-a Abs & A meaningful  & \textbf{life} requires active engagement with somethi... \\
17-a Abs & ...meaningful, and judges the unsuccessful  & \textbf{life} to be far from meaningless. \\
17-a 1 & ...cause of what they brought about, and a  & \textbf{life} that brings nothing about is to that ex... \\
17-a 1 & ... to roll stones, and concluded that his  & \textbf{life} would then have all the meaning a life ... \\
17-a 1 & ... life would then have all the meaning a  & \textbf{life} can have, though he achieved nothing. \\
17-a 1 & ...'s subjectivism was rightly rejected: a  & \textbf{life} devoted with passion to rolling stones,... \\
17-a 2 & A  & \textbf{life} is meaningful to the extent that its su... \\
17-a 2 & ...ing from Nozick's experience machine: a  & \textbf{life} of experiences as of curing diseases, l... \\
17-a 2 & ...iseases, lived in the machine, is not a  & \textbf{life} engaged with the sick. \\
17-a 3 & A person spends much of her  & \textbf{life} listening to and studying music she lov... \\
17-a 3 & The unsuccessful mathematician's  & \textbf{life} is meaningful, and so is the life of th... \\
17-a 3 & ...ian's life is meaningful, and so is the  & \textbf{life} of the nurse in the second ward. \\
17-a 3 & ...is not the verdict of reflection that a  & \textbf{life} devoted to an unsolved problem, or to t... \\
17-a 3 & The listener's  & \textbf{life} is meaningful to the extent of her enga... \\
17-a 3 & A meaningful  & \textbf{life} requires active engagement with somethi... \\
17-a 3 & ... what reflection counts as a meaningful  & \textbf{life} (Wolf 2010). \\
17-a 3 & ...aning, the unsuccessful mathematician's  & \textbf{life}, the life of the nurse in the second wa... \\
17-a 3 & ... unsuccessful mathematician's life, the  & \textbf{life} of the nurse in the second ward, and th... \\
17-a 3 & ... in the second ward, and the listener's  & \textbf{life} would lack it, so far as their central ... \\
17-a 3 & ...s equal, the successful mathematician's  & \textbf{life} is more meaningful than the unsuccessfu... \\
17-a 3 & ...e good (ETH-01), and it adds value to a  & \textbf{life}; if success enlarges the worth realized... \\
17-a 3 & ... between a meaningful and a meaningless  & \textbf{life}. \\
17-a 4 & So every meaningful  & \textbf{life} contains achievement after all. \\
17-a 4 & ...athematician did not attain it, and her  & \textbf{life} is meaningful. \\
17-a 4 & A  & \textbf{life} that changes nothing is enclosed within... \\
17-a 4 & ETH-05 shows that a  & \textbf{life} in contact with reality is better than ... \\
17-a 4 & ...nd its comparison contrasts a veridical  & \textbf{life} with an experiential duplicate. \\
17-a 4 & ...r work bears on a real problem, and her  & \textbf{life} is in contact with reality. \\
17-a 4 & ... concerns success bears on how well her  & \textbf{life} goes for her, and we grant that success... \\
17-a 4 & ...better; it does not bear on whether her  & \textbf{life} is meaningful. \\
17-a 4 & ...d achieving nothing, lives a meaningful  & \textbf{life}. \\
17-a 4 & ...e from flooding would live a meaningful  & \textbf{life}, even if the stone always rolled back a... \\
17-a 5 & A meaningful  & \textbf{life} requires active engagement with somethi... \\
18-a Abs & ...false judgments may be indispensable to  & \textbf{life}, and a body of psychological research h... \\
18-a 1 & ...gments may be the most indispensable to  & \textbf{life} (Nietzsche 1886). \\
18-a 2 & ...unction, and the ordinary progress of a  & \textbf{life} consists partly in such corrections. \\
18-a 2 & ...ch illusion is a respect in which one's  & \textbf{life} is out of contact with how things are. \\
18-b 1 & ...ch illusion is a respect in which one's  & \textbf{life} is out of contact with how things are. \\
19-a Abs & ... death only through its relation to the  & \textbf{life} it ends. \\
19-a Abs & But death gives  & \textbf{life} a shape by limiting it. \\
19-a Abs & ...lowing Velleman, we take the value of a  & \textbf{life} to depend in part on its shape, the ord... \\
19-a Abs & ...ds and evils occur, and we argue that a  & \textbf{life} has a shape as a whole only if it is bo... \\
19-a Abs & ...ble with EXI-02: the limit that gives a  & \textbf{life} a shape is not a benefit to the one who... \\
19-a Abs & ...a later death would also have given her  & \textbf{life} a shape, and in the ordinary case a bet... \\
19-a Abs & ...ts EXI-03, which argues that an endless  & \textbf{life} is not desirable. \\
19-a 1 & ...nd whether it bears on the point of the  & \textbf{life} it ends. \\
19-a 1 & ...he wage of sin, as a passage to another  & \textbf{life}, or as the moment of judgment. \\
19-a 1 & ...en found its meaning in its relation to  & \textbf{life}. \\
19-a 1 & ...leman (1991) showed that the value of a  & \textbf{life} is not the sum of the values of its mom... \\
19-a 1 & ...m of the values of its moments, since a  & \textbf{life} that improves from misfortune to succes... \\
19-a 1 & ... has no meaning in itself, but it gives  & \textbf{life} a shape by limiting it. \\
19-a 2 & ...s independently of its relations to the  & \textbf{life} in which it occurs: if it is the fulfil... \\
19-a 2 & ...of a value that it does not owe to that  & \textbf{life}. \\
19-a 2 & The shape of a  & \textbf{life} is the temporal order of its goods and ... \\
19-a 2 & The value of a  & \textbf{life} as a whole depends in part on its shape... \\
19-a 2 & ...less view of time defended in TSM-01, a  & \textbf{life} is an ordered series of events, each lo... \\
19-a 3 & ...imposes it or that it lead to a further  & \textbf{life}; neither can be established, and REL-08... \\
19-a 3 & ... in virtue of its relation to a further  & \textbf{life} or a further purpose, not in itself. \\
19-a 3 & ...on is a claim about its relation to the  & \textbf{life} it ends, and it is often false: most de... \\
19-a 3 & ...eath that are true independently of any  & \textbf{life}, but they assign it no significance or ... \\
19-a 3 & Death gives a  & \textbf{life} a shape by limiting it. \\
19-a 3 & A  & \textbf{life} is an ordered series of events (Assumpt... \\
19-a 3 & ... of being limited by death that a human  & \textbf{life} has a shape as a whole, and by Principl... \\
19-a 3 & ...s shape contributes to the value of the  & \textbf{life}. \\
19-a 3 & ...that depends on the goods and evils the  & \textbf{life} contains and their order. \\
19-a 3 & ...ference between the value of her actual  & \textbf{life} and the value of the life she would hav... \\
19-a 3 & ...of her actual life and the value of the  & \textbf{life} she would have had had she not died the... \\
19-a 3 & ...ie at that time, she dies later, so her  & \textbf{life} is limited and has a shape there too. \\
19-a 3 & ...he conclusion of EXI-03 that an endless  & \textbf{life} is not desirable, but does not depend o... \\
19-a 3 & EXI-03 argues that an endless  & \textbf{life} would exhaust the categorical desires t... \\
19-a 3 & ...haust the categorical desires that give  & \textbf{life} its point; Theorem 3.2 adds that it wou... \\
19-a 3 & The two are independent: a  & \textbf{life} could lack a shape as a whole and still... \\
19-a 4 & Projects, relationships and stages of  & \textbf{life} have beginnings and ends, and give life... \\
19-a 4 & ...life have beginnings and ends, and give  & \textbf{life} all the shape it needs. \\
19-a 4 & They give shape to the episodes of a  & \textbf{life}, and an endless life could have shaped ... \\
19-a 4 & ... the episodes of a life, and an endless  & \textbf{life} could have shaped episodes. \\
19-a 4 & ...at they do not give is the shape of the  & \textbf{life} as a whole, which is what Principle P1 ... \\
19-a 4 & ...what Principle P1 concerns: whether the  & \textbf{life}, taken as one thing, was a story of pro... \\
19-a 4 & That a  & \textbf{life} is limited is a condition of its having... \\
19-a 4 & Theorem 3.2 says that death gives a  & \textbf{life} a shape, not a good one. \\
19-a 4 & An early death gives a  & \textbf{life} the shape of a beginning cut off, which... \\
19-a 4 & ...off, which is one of the worst shapes a  & \textbf{life} can have; this is part of why, by EXI-0... \\
19-a 5 & ...elong to it through its relation to the  & \textbf{life} it ends. \\
19-a 5 & But it limits that  & \textbf{life}, and by limiting it gives it a shape as... \\
19-a 5 & ... a whole, which is part of what makes a  & \textbf{life} go well or badly. \\
19-b Abs & ...eath has no meaning in itself but gives  & \textbf{life} a shape by limiting it. \\
19-b Abs & ...ition of one kind of meaning, that of a  & \textbf{life} as a whole, which depends on the life's... \\
19-b 1 & On the account of EXI-01 and EXI-17, a  & \textbf{life} is meaningful to the extent that its su... \\
19-b 1 & ... it, as a condition of the meaning of a  & \textbf{life} as a whole, or as a condition of the pe... \\
19-b 1 & That a  & \textbf{life} is limited by death is not something wh... \\
19-b 1 & ...ot something whose engagement makes the  & \textbf{life} meaningful, and does not make anything ... \\
19-b 1 & In so far as the meaning of a  & \textbf{life} as a whole depends on its shape, finitu... \\
19-b 1 & Some meaning belongs to a  & \textbf{life} only as a whole: that it was a life of ... \\
19-b 1 & ...o a life only as a whole: that it was a  & \textbf{life} of growing understanding, of a commitme... \\
19-b 1 & So the meaning that a  & \textbf{life} has as a whole requires that it be limi... \\
19-b 1 & By the result of EXI-03, an endless  & \textbf{life} would eventually exhaust the categorica... \\
19-b 1 & ...ust the categorical desires that give a  & \textbf{life} its point (Williams 1973). \\
19-b 1 & A  & \textbf{life} without categorical desires lacks the a... \\
19-b 1 & So an endless  & \textbf{life} would at some point cease to be meaning... \\
19-b 1 & ...ome point cease to be meaningful, and a  & \textbf{life} remains meaningful throughout its cours... \\
19-b 2 & ... temporal scarcity, so that an immortal  & \textbf{life} would lack the structure in which valui... \\
19-b 2 & Scheffler's point that an immortal  & \textbf{life} would lack the stages of a human life i... \\
19-b 2 & ...l life would lack the stages of a human  & \textbf{life} is a point about shape, and Nussbaum's ... \\
19-b 2 & ...um's point that some values belong to a  & \textbf{life} of a certain form is a point about the ... \\
19-b 2 & ...a later death would also have given the  & \textbf{life} a shape and would, in the ordinary case... \\
19-b 3 & Finitude does not make  & \textbf{life} meaningful. \\
19-b 3 & It makes possible the meaning of a  & \textbf{life} as a whole, through the shape that a li... \\
19-b 3 & ...pe that a limit gives, and it keeps the  & \textbf{life} within the span in which engagement per... \\
19-b 3 & ... in itself, bears on the meaning of the  & \textbf{life} it ends. \\
20-a Abs & EXI-17 argues that a meaningful  & \textbf{life}, or a meaningful activity within one, r... \\
20-a Abs & ...ther the process has the kind of mental  & \textbf{life} that engagement requires, which MND-08 ... \\
20-a 1 & Wolf (2010) argued that meaning in  & \textbf{life} arises when subjective attraction meets... \\
20-a 1 & ...ction meets objective attractiveness: a  & \textbf{life} is meaningful to the extent that its su... \\
20-a 1 & ...he account to an activity rather than a  & \textbf{life}: the production of the present collecti... \\
20-a 4 & ... from inside it: a fact about the inner  & \textbf{life} of the producer and a fact about the up... \\
20-b Abs & ... EXI-20-a, on whether it has the mental  & \textbf{life} that engagement requires, or on whether... \\
\end{longtable}
\kwichead{matters}{37}
\begin{longtable}{@{}>{\raggedright}p{0.07\textwidth}>{\raggedleft}p{0.41\textwidth}@{\,}>{\raggedright\arraybackslash}p{0.47\textwidth}@{}}
01-c 3 & The difference  & \textbf{matters} because opportunities for projects in t... \\
01-d 2 & Deprivation therefore  & \textbf{matters} to meaning, and to a theory of justice ... \\
01-d 4 & ...ufficiently detached standpoint nothing  & \textbf{matters}, so meaning is available to no one. \\
02-a 3 & ... equally sudden differ in badness: what  & \textbf{matters} is not the manner of dying but the life... \\
02-a 6 & ...that connectedness carries most of what  & \textbf{matters} in survival. \\
02-b 2 & This  & \textbf{matters} because Epicurus's therapeutic aim, the... \\
03-a Abs & ...ife would not be hers in the sense that  & \textbf{matters}. \\
03-a Abs & ...that connectedness carries most of what  & \textbf{matters} in survival, we show that an endless li... \\
03-a Abs & ...o be, in its later reaches, a life that  & \textbf{matters} to the person now. \\
03-a 1 & ...y of character with the account of what  & \textbf{matters} in survival defended in PID-06. \\
03-a 2 & What  & \textbf{matters} to a person, in her concern for a futur... \\
03-a 2 & ...ife gives the person now little of what  & \textbf{matters}. \\
03-a 3 & ...ives her, as she now is, little of what  & \textbf{matters}. \\
03-a 3 & ... that time gives her now little of what  & \textbf{matters}. \\
03-a 3 & ...hout connectedness carry little of what  & \textbf{matters}. \\
03-a 3 & ... cease to be a life that gives her what  & \textbf{matters}. \\
03-a 3 & ...ter reaches give her now little of what  & \textbf{matters}. \\
03-a 4 & ... through a chain carries little of what  & \textbf{matters}. \\
03-a 6 & ... a life that gives its subject now what  & \textbf{matters} to her. \\
03-a 6 & ...that connectedness carries most of what  & \textbf{matters}. \\
03-b Abs & ...e its point, or cease to be a life that  & \textbf{matters} to the person now. \\
03-b 1 & ...mption 2.4), and the dependence of what  & \textbf{matters} on connectedness (Assumption 2.5), Theo... \\
06-a 4 & ...re about is hostile in every sense that  & \textbf{matters}. \\
06-b 3 & ...e has no such point of view, so nothing  & \textbf{matters} from it. \\
07-a 2 & ...d; it presupposes that its being solved  & \textbf{matters}, and an agent to whom it matters only a... \\
07-a 2 & ...solved matters, and an agent to whom it  & \textbf{matters} only as a source of satisfaction has no... \\
10-a 4 & But the signal bears on a question that  & \textbf{matters} independently of its evolutionary histo... \\
10-b 2 & The distinction  & \textbf{matters} for how evidence is read. \\
10-b 2 & ...ross-tabulation in which the difference  & \textbf{matters}. \\
10-b 3 & ...without traction, and the qualification  & \textbf{matters}: an over-demanding task on which one ma... \\
12-a Abs & ...arison can be made good in a sense that  & \textbf{matters} to the one who lives the life: although... \\
13-a 2 & ...answers that question for the case that  & \textbf{matters} here: the entitlement is created while ... \\
13-a 3 & A wish that fails clause (ii) concerns  & \textbf{matters} on which the person has no claim agains... \\
17-a 3 & ...he sense of Definition 2.2: the problem  & \textbf{matters} to knowledge, which is a component of t... \\
17-a 3 & ... the good, and the patients' well-being  & \textbf{matters} whether or not it can be restored. \\
17-a 4 & This is an argument that success  & \textbf{matters}. \\
19-b 2 & ...rgency adds no worth: the friend's good  & \textbf{matters} no less to someone who will have other ... \\
\end{longtable}
\kwichead{meaning}{229}
\begin{longtable}{@{}>{\raggedright}p{0.07\textwidth}>{\raggedleft}p{0.41\textwidth}@{\,}>{\raggedright\arraybackslash}p{0.47\textwidth}@{}}
01-c Abs & ...1-a establishes that a life of positive  & \textbf{meaning} is available to every agent capable of ... \\
01-c Abs & ...the sense relevant to the definition of  & \textbf{meaning}. \\
01-c 1 & ...leave people with few opportunities for  & \textbf{meaning}. \\
01-c 2 & ...st follows the objects that accounts of  & \textbf{meaning} have most often recognized: persons, th... \\
01-c 2 & ... intentional action, a life of positive  & \textbf{meaning} is available. \\
01-c 3 & ... applies the same point to opportunity:  & \textbf{meaning} does not wait on the availability of a ... \\
01-c 4 & ...ures of projects that are irrelevant to  & \textbf{meaning} on Definition 1: a completion condition... \\
01-c 5 & The universal availability of  & \textbf{meaning} does not depend on the universal availa... \\
01-c 5 & ... worth, and a life that contains it has  & \textbf{meaning}. \\
01-d Abs & ...h projects of objective worth, and that  & \textbf{meaning} in this sense is available to everyone. \\
01-d Abs & It implies that  & \textbf{meaning} is not hostage to luck in outcomes, to ... \\
01-d Abs & ...e and injustice affect is the degree of  & \textbf{meaning}, not its possibility. \\
01-d Abs & ...ful, or that deprivation is harmless to  & \textbf{meaning}. \\
01-d 1 & ...f intentional action a life of positive  & \textbf{meaning} is available. \\
01-d 1 & ...same availability by a different route:  & \textbf{meaning} requires active engagement with somethi... \\
01-d 1 & ...ment (Theorem 3.3 of EXI-17-a), so that  & \textbf{meaning} is not hostage to luck in outcomes (Cor... \\
01-d 2 & If  & \textbf{meaning} required success, it would be available... \\
01-d 2 & Because  & \textbf{meaning} requires engagement, which is within th... \\
01-d 2 & Definition 1 of EXI-01-a measures  & \textbf{meaning} by engagement and by the objective wort... \\
01-d 2 & ...an achieve, which is not a condition of  & \textbf{meaning}; scale bears on how much value a projec... \\
01-d 2 & ...worth, and engagement with them confers  & \textbf{meaning} on lives that will leave no public trac... \\
01-d 2 & ...cluding injustice, affect the degree of  & \textbf{meaning} a life can attain and the ease with whi... \\
01-d 2 & ...can attain it, but not whether positive  & \textbf{meaning} is available. \\
01-d 2 & Definition 1 of EXI-01-a makes  & \textbf{meaning} a matter of degree, and EXI-17-a notes ... \\
01-d 2 & Deprivation therefore matters to  & \textbf{meaning}, and to a theory of justice that counts... \\
01-d 2 & ... cannot do is remove the possibility of  & \textbf{meaning} altogether. \\
01-d 3 & ...availability to equality: the degree of  & \textbf{meaning} varies with the extent of engagement an... \\
01-d 3 & ...from availability to indifference: that  & \textbf{meaning} remains possible under oppression is no... \\
01-d 3 & ...oppression, which reduces the degree of  & \textbf{meaning} as it reduces other goods. \\
01-d 4 & Objective-worth views of  & \textbf{meaning} privilege the projects valued by the ed... \\
01-d 4 & ...ience and political causes, and so make  & \textbf{meaning} in practice available mainly to the pri... \\
01-d 4 & If  & \textbf{meaning} is available to everyone, the claim tha... \\
01-d 4 & That positive  & \textbf{meaning} is available to everyone is compatible ... \\
01-d 4 & The claim is not that  & \textbf{meaning} is cheap but that it is not reserved: n... \\
01-d 4 & ...detached standpoint nothing matters, so  & \textbf{meaning} is available to no one. \\
01-d 5 &  & \textbf{Meaning}, on the account of EXI-01-a, is open to... \\
01-d 5 & The availability of  & \textbf{meaning} is universal; its degree is not, and ho... \\
03-a 1 & ...worthwhile, in line with the account of  & \textbf{meaning} in EXI-01; and it replaces Williams's a... \\
03-a 2 & ...ch engagement, not achievement, is what  & \textbf{meaning} requires. \\
04-a 1 & ...y, dissolves the absurd: on theories of  & \textbf{meaning} such as those of Wolf (2010) and Metz (... \\
04-a 2 & ...he world as an adversary that denies us  & \textbf{meaning}, and the agent's persistence as a victo... \\
04-a 3 & ...he world as an adversary that denies us  & \textbf{meaning}. \\
04-a 4 &  & \textbf{Meaning} requires active engagement with somethi... \\
04-a 4 & Irony is the enemy of  & \textbf{meaning}. \\
04-b 3 & A further consequence concerns  & \textbf{meaning}. \\
04-b 3 & ...t meaningful however happy he is, since  & \textbf{meaning} requires engagement with something wort... \\
04-b 3 & ...bes it, is a paradigm of the absence of  & \textbf{meaning} rather than of its presence in a godles... \\
04-b 3 & ...jects, and engagement with them confers  & \textbf{meaning} (EXI-01). \\
06-b Abs & ...ifferent to us leaves our lives without  & \textbf{meaning}, or with a meaning we must assert in de... \\
06-b Abs & ...es our lives without meaning, or with a  & \textbf{meaning} we must assert in defiance of it. \\
06-b Abs &  & \textbf{Meaning}, on the account of EXI-01, consists in ... \\
06-b Abs & Cosmic indifference is irrelevant to  & \textbf{meaning} in life, though not to every question a... \\
06-b Abs & ...ife, though not to every question about  & \textbf{meaning}. \\
06-b 1 & We ask what this implies for  & \textbf{meaning}. \\
06-b 2 & ...o us does not imply that our lives lack  & \textbf{meaning}. \\
06-b 2 & It denies that  & \textbf{meaning} in life requires that good. \\
06-b 2 & ...its indifference denies us nothing that  & \textbf{meaning} requires; by Theorem 3.1 of EXI-06-a, i... \\
06-b 3 &  & \textbf{Meaning} requires that our lives matter from the... \\
06-b 3 & ... does indifference deprive our lives of  & \textbf{meaning}. \\
06-b 3 & Theistic accounts hold that  & \textbf{meaning} requires a purpose conferred by a being... \\
06-b 3 & If the universe does not care,  & \textbf{meaning} is lost unless there is such a being. \\
06-b 3 & The present point is conditional: if  & \textbf{meaning} consists in engagement with what has ob... \\
06-b 3 & ...estion, and supernaturalist accounts of  & \textbf{meaning} are assessed elsewhere (Metz 2013). \\
06-b 4 & It leaves the question of  & \textbf{meaning} where EXI-01 places it: in what we enga... \\
07-a 3 & It fits also the account of  & \textbf{meaning} in EXI-01 and EXI-17, on which meaning ... \\
07-a 3 & ... meaning in EXI-01 and EXI-17, on which  & \textbf{meaning} comes from active engagement with what ... \\
07-a 5 & ...the account of the good (ETH-01) and of  & \textbf{meaning} (EXI-01, EXI-17). \\
10-a Abs & ...challenge without traction, and lack of  & \textbf{meaning}; and it motivates search for alternativ... \\
10-a Abs & ...dom is aversive, why it is sensitive to  & \textbf{meaning} as well as to difficulty, and why chron... \\
10-a 1 & ...ty, the condition of a culture in which  & \textbf{meaning} has to be produced rather than inherite... \\
10-a 1 & ...lson (2018) distinguish attentional and  & \textbf{meaning} components: one is bored when one's att... \\
10-a 2 & The  & \textbf{meaning} clause is why the mismatch is between c... \\
10-a 3 & It explains why boredom is sensitive to  & \textbf{meaning}: the mismatch it signals includes the i... \\
10-a 4 & ...bout the agent's good on the account of  & \textbf{meaning} in EXI-17, on which meaning requires ac... \\
10-a 4 & ... account of meaning in EXI-17, on which  & \textbf{meaning} requires active engagement with somethi... \\
10-b Abs & ...tly evidence about motivation, that the  & \textbf{meaning} and under-challenge clauses of the mism... \\
10-b 2 & ...e resources in either direction, from a  & \textbf{meaning} component, the activity's failure to ma... \\
10-b 3 & (2009) found that a lack of  & \textbf{meaning} in life predicted boredom cross-section... \\
10-b 5 & ...city to attend; they say less about its  & \textbf{meaning} clause, which concerns the capacity to ... \\
10-b 5 & ... the entries coded attn and those coded  & \textbf{meaning}+ overlap only in the MAC model itself (... \\
10-b 6 & ...or by meaninglessness alone, except the  & \textbf{meaning} manipulation reported by Fahlman et al. \\
10-b 8 & ...gagement with something worthwhile; the  & \textbf{meaning} clause of the mismatch thesis is the af... \\
10-b 8 & ...hat view, and the evidence that lack of  & \textbf{meaning} predicts boredom (Fah09, Til12, Til17) ... \\
10-b 11 & ...match thesis in its under-challenge and  & \textbf{meaning} clauses and in its converse, that match... \\
10-b 11 & ...how they engage, and turn to sources of  & \textbf{meaning}. \\
10-b A & ...sis if it carries one of under+, over+,  & \textbf{meaning}+, fit-, attn or relational, and challen... \\
10-b B & ... and overanx (against); meaninglessness  & \textbf{meaning}+; Assumption 2.4 fit-; Lemma 3.1 signal... \\
10-b B & ... 3.3 fit- and the codes of goal change,  & \textbf{meaning} regulation and creativity (for) and har... \\
11-a Abs & Supernaturalists hold that  & \textbf{meaning} requires a purpose assigned to human li... \\
11-a Abs & ... can be meaningful without God, because  & \textbf{meaning} does not require a cosmic purpose, and ... \\
11-a Abs & ...ts on the fitting-fulfilment account of  & \textbf{meaning}: a life is meaningful to the extent tha... \\
11-a Abs & ...ned by another is neither necessary for  & \textbf{meaning} nor sufficient for it. \\
11-a Abs & ...ond claim rests on the observation that  & \textbf{meaning} comes in degrees and grows with the wor... \\
11-a 1 & ...God there is no such purpose, and so no  & \textbf{meaning} (Craig 2008). \\
11-a 1 & On theories of the infinite,  & \textbf{meaning} requires a connection to something of u... \\
11-a 1 & ...am's more cautious view, the pursuit of  & \textbf{meaning} requires a moral framework and a ground... \\
11-a 1 & Naturalists reply that  & \textbf{meaning} is available in a world without God. \\
11-a 1 & On subjectivist views,  & \textbf{meaning} is a matter of what a person cares abou... \\
11-a 1 & Wolf (2010) defends a hybrid:  & \textbf{meaning} arises when the subject's active engage... \\
11-a 1 & ...l-centered theories as accounts of what  & \textbf{meaning} requires, and defends a naturalist theo... \\
11-a 1 & ...nd defends a naturalist theory on which  & \textbf{meaning} is a matter of orienting one's rational... \\
11-a 1 & ...lism, on which God is necessary for any  & \textbf{meaning}, from moderate supernaturalism, on whic... \\
11-a 1 & ...ecessary only for a greater or ultimate  & \textbf{meaning}. \\
11-a 1 & ...ut the framework, section 3 argues that  & \textbf{meaning} does not require cosmic purpose, and se... \\
11-a 1 & ... that a theistic framework would add to  & \textbf{meaning}. \\
11-a 2 & EXI-01 adds that  & \textbf{meaning} so understood is available to everyone. \\
11-a 2 & ...implicit in Definition 2.1, which makes  & \textbf{meaning} a matter of extent. \\
11-a 2 & ... in Nozick's linkage view: a life gains  & \textbf{meaning} by connection with value beyond itself,... \\
11-a 3 &  & \textbf{Meaning} does not require a cosmic purpose. \\
11-a 4 & ...ue, would make available to a life more  & \textbf{meaning} than it could have in a world without G... \\
11-a 4 & ...d by Principle P2 each would add to its  & \textbf{meaning}. \\
11-a 4 & If REL-01 is right, the additional  & \textbf{meaning} is actually available, and is available... \\
11-a 4 &  & \textbf{Meaning} has a floor that does not depend on God... \\
11-a 4 & ...s critics that God is not necessary for  & \textbf{meaning}, and with moderate supernaturalism that... \\
11-a 4 & ...upernaturalism that God would make more  & \textbf{meaning} available (Metz 2013); it does not foll... \\
11-a 4 & ...d necessary for any particular grade of  & \textbf{meaning}. \\
11-a 5 & ...4.1: preservation of the good would add  & \textbf{meaning}; its absence does not remove the meanin... \\
11-a 5 & ...eaning; its absence does not remove the  & \textbf{meaning} already there. \\
11-a 5 & On the linkage view, a finite life has  & \textbf{meaning} by connection with something larger. \\
11-a 5 & That larger thing has  & \textbf{meaning} only if it is connected with something ... \\
11-a 5 & ...unlimited terminus expresses a wish for  & \textbf{meaning} without limit, not a condition of meani... \\
11-a 5 & ...aning without limit, not a condition of  & \textbf{meaning} as such. \\
11-a 5 & ...ants that an unlimited object would add  & \textbf{meaning}; it does not follow that without it the... \\
11-a 5 & The pursuit of  & \textbf{meaning} is not a solitary engagement with isola... \\
11-a 5 & ...thout God, these are fragile, and so is  & \textbf{meaning} (Cottingham 2003). \\
11-a 5 & ...ction describes affects the security of  & \textbf{meaning} and the reasonableness of certain hopes... \\
11-a 5 & The word  & \textbf{meaning} in the meaning of life just asks for a ... \\
11-a 5 & The word meaning in the  & \textbf{meaning} of life just asks for a purpose. \\
11-a 5 & ...without a purpose assigned to it has no  & \textbf{meaning} by definition, and the naturalist has c... \\
11-a 5 & ...s, the life without God is deficient in  & \textbf{meaning}. \\
11-a 5 & The dilemma assumes that  & \textbf{meaning} is all or nothing. \\
11-a 6 & It does not show how much  & \textbf{meaning} theism would add, nor that the addition... \\
11-a 6 & ...t every life without God is meaningful:  & \textbf{meaning} requires engagement with what is worthw... \\
11-a 7 &  & \textbf{Meaning} consists in active engagement with what... \\
11-a 7 & ...igned purpose is neither sufficient for  & \textbf{meaning} nor necessary for it. \\
11-a 7 & A theistic framework would add to  & \textbf{meaning}, by offering an object of maximal worth... \\
11-a 7 & The floor of  & \textbf{meaning} is independent of God; the ceiling, if ... \\
11-b Abs & EXI-11-a argues that  & \textbf{meaning} does not require a cosmic purpose, and ... \\
11-b Abs & ..., if true, would nevertheless make more  & \textbf{meaning} available (Proposition 4.1). \\
11-b Abs & ...at God adds is not more of the same but  & \textbf{meaning} as such; the anti-theist, following Kah... \\
11-b Abs & ...e, says that God's existence would take  & \textbf{meaning} away, by removing the privacy, independ... \\
11-b Abs & We argue that the additional  & \textbf{meaning} is a difference of degree, and that the... \\
11-b 1 & EXI-11-a defines  & \textbf{meaning} as active engagement with what is worth... \\
11-b 1 & ...thwhile (Definition 2.1) and holds that  & \textbf{meaning} grows with the worth of what one engage... \\
11-b 1 & ...worthwhile (Principle P2; on degrees of  & \textbf{meaning}, see (Metz 2013)). \\
11-b 1 & The  & \textbf{meaning} that a theistic framework would add is ... \\
11-b 1 & ...rk would add is of the same kind as the  & \textbf{meaning} available without it, and differs in de... \\
11-b 1 & By Definition 2.1 of EXI-11-a, any  & \textbf{meaning} is engagement with what is worthwhile. \\
11-b 1 & ...agement with what is worthwhile; so the  & \textbf{meaning} it yields falls under the same definiti... \\
11-b 1 & ...hout it can be meaningful, so it is not  & \textbf{meaning} as such. \\
11-b 1 & ...EXI-11-a, it can increase the degree of  & \textbf{meaning}. \\
11-b 2 & ...ate point; what theism adds is not more  & \textbf{meaning} but the only meaning that would not be ... \\
11-b 2 & ...m adds is not more meaning but the only  & \textbf{meaning} that would not be undercut. \\
11-b 2 & The claim that naturalistic  & \textbf{meaning} is undercut is the claim that it depend... \\
11-b 2 & ...son a theistic world would contain less  & \textbf{meaning}. \\
11-b 2 & The proposition concerns the  & \textbf{meaning} made available by a framework, and the ... \\
11-b 2 & ... not show that the framework makes less  & \textbf{meaning} available. \\
11-b 3 & Theism, if true, would add to the  & \textbf{meaning} available in a life by adding objects a... \\
11-b 3 & ...degree, which is why its absence leaves  & \textbf{meaning} in place, and it is compatible with the... \\
14-a Abs & ...nded by psychologists as a resource for  & \textbf{meaning} and social connection. \\
17-a Abs & Many accounts of  & \textbf{meaning} give success a central place: a meaning... \\
17-a Abs & ...ther good, and may add to the degree of  & \textbf{meaning}; it is not a condition of meaning. \\
17-a Abs & ...ee of meaning; it is not a condition of  & \textbf{meaning}. \\
17-a Abs & ...w secures the universal availability of  & \textbf{meaning}. \\
17-a 1 & A widespread thought about  & \textbf{meaning} is that it is tied to accomplishment. \\
17-a 1 & ...ught, the lives we cite as paradigms of  & \textbf{meaning}, those of Einstein, Mandela or Mother T... \\
17-a 1 & Nozick (1981) connected  & \textbf{meaning} with transcending one's limits by makin... \\
17-a 1 & ... to suppose that this is the value that  & \textbf{meaning} tracks. \\
17-a 1 & The opposite view is that  & \textbf{meaning} is a matter of how one engages with wha... \\
17-a 1 & ...d that his life would then have all the  & \textbf{meaning} a life can have, though he achieved not... \\
17-a 1 & On her fitting-fulfilment view,  & \textbf{meaning} arises when subjective attraction meets... \\
17-a 1 & We argue that it should not:  & \textbf{meaning} requires active engagement with somethi... \\
17-a 2 & ...is is not that it lacks value, but that  & \textbf{meaning} does not require it. \\
17-a 2 & Principle P2 applies to  & \textbf{meaning} the lesson that ETH-05 draws for well-b... \\
17-a 2 & ...evant facts in view, are evidence about  & \textbf{meaning}, in the same way that such judgments ar... \\
17-a 3 & ...rinciple P1, her engagement contributes  & \textbf{meaning}. \\
17-a 3 & ...ement with worthless objects confers no  & \textbf{meaning}, and pure objectivism fails because wor... \\
17-a 3 & ...unct: if achievement were necessary for  & \textbf{meaning}, the unsuccessful mathematician's life,... \\
17-a 3 & Achievement is not sufficient for  & \textbf{meaning} either: an achievement in a worthless p... \\
17-a 3 & By Principle P1, neither confers  & \textbf{meaning}. \\
17-a 3 & ...hrough one's engagement, it may enlarge  & \textbf{meaning} too. \\
17-a 3 &  & \textbf{Meaning} so understood is not hostage to luck in... \\
17-a 3 & By Theorem 3.3,  & \textbf{meaning} depends on the second and not on the fi... \\
17-a 3 & ...blishing, the conclusion of EXI-01 that  & \textbf{meaning} in this sense is available to everyone. \\
17-a 4 & ...ievement trivially, and the thesis that  & \textbf{meaning} requires achievement becomes a restatem... \\
17-a 4 & The substantive thesis we deny is that  & \textbf{meaning} requires success: attaining the outcome... \\
17-a 4 &  & \textbf{Meaning} requires connection with value beyond o... \\
17-a 4 & We accept that  & \textbf{meaning} requires connection with value beyond o... \\
17-a 4 & Why does it not show that  & \textbf{meaning} requires success? \\
17-a 4 & Principle P2 carries over to  & \textbf{meaning} the part of ETH-05's lesson that concer... \\
17-a 5 & Achievement is a good, and may enlarge  & \textbf{meaning}, but it is neither necessary nor suffic... \\
18-a 3 &  & \textbf{Meaning} depends on taking one's projects to mat... \\
19-a Abs & What is the  & \textbf{meaning} of death? \\
19-a Abs & On the first, death has a  & \textbf{meaning} of its own: it is a punishment, a passa... \\
19-a Abs & Death has no  & \textbf{meaning} in itself, since every candidate meanin... \\
19-a Abs & ...eaning in itself, since every candidate  & \textbf{meaning} either depends on claims that cannot be... \\
19-a 1 & The question What is the  & \textbf{meaning} of death? asks for more than what death... \\
19-a 1 & Religious traditions have given death a  & \textbf{meaning} of its own, as the wage of sin, as a pa... \\
19-a 1 & Philosophers have more often found its  & \textbf{meaning} in its relation to life. \\
19-a 1 & ...mbine the last two points: death has no  & \textbf{meaning} in itself, but it gives life a shape by... \\
19-a 2 & An event has  & \textbf{meaning} in itself if it has a significance or p... \\
19-a 3 & Death has no  & \textbf{meaning} in itself. \\
19-a 3 & Even if one of these were true, the  & \textbf{meaning} would belong to death in virtue of its ... \\
19-a 3 & So no candidate  & \textbf{meaning} belongs to death in the sense of Defini... \\
19-a 5 & Death has no  & \textbf{meaning} in itself. \\
19-b Abs & ...ten said that finitude is the source of  & \textbf{meaning}: that our projects matter because our t... \\
19-b Abs & EXI-19-a argues that death has no  & \textbf{meaning} in itself but gives life a shape by lim... \\
19-b Abs & ...ut what this implies about finitude and  & \textbf{meaning}. \\
19-b Abs & Finitude is not a source of  & \textbf{meaning}, since it adds no object of worth to en... \\
19-b Abs & It is a condition of one kind of  & \textbf{meaning}, that of a life as a whole, which depen... \\
19-b 1 & Finitude might bear on  & \textbf{meaning} in three ways: as a source of it, as a ... \\
19-b 1 & ...s a source of it, as a condition of the  & \textbf{meaning} of a life as a whole, or as a condition... \\
19-b 1 & ...oposition 3.1 of EXI-19-a, death has no  & \textbf{meaning} in itself. \\
19-b 1 & ...account of EXI-01 it is not a source of  & \textbf{meaning}. \\
19-b 1 & In so far as the  & \textbf{meaning} of a life as a whole depends on its sha... \\
19-b 1 & Some  & \textbf{meaning} belongs to a life only as a whole: that... \\
19-b 1 & Such  & \textbf{meaning} depends on the temporal order of the en... \\
19-b 1 & So the  & \textbf{meaning} that a life has as a whole requires tha... \\
19-b 1 & ...esires lacks the active engagement that  & \textbf{meaning} requires (EXI-17). \\
19-b 2 & So finitude is a source of  & \textbf{meaning} after all. \\
19-b 2 & ...t is something whose engagement confers  & \textbf{meaning}. \\
19-b 2 & ...f finitude is a condition of whole-life  & \textbf{meaning} and of the persistence of engagement, t... \\
19-b 2 & ...er, though some limit is a condition of  & \textbf{meaning}. \\
19-b 3 & It makes possible the  & \textbf{meaning} of a life as a whole, through the shape... \\
19-b 3 & ... is the sense in which death, having no  & \textbf{meaning} in itself, bears on the meaning of the ... \\
19-b 3 & ...ving no meaning in itself, bears on the  & \textbf{meaning} of the life it ends. \\
20-a 1 & Wolf (2010) argued that  & \textbf{meaning} in life arises when subjective attracti... \\
20-a 1 & EXI-17 defends a version of this view:  & \textbf{meaning} requires active engagement with somethi... \\
20-a 1 & EXI-01 adds that  & \textbf{meaning} in this sense is available to everyone. \\
\end{longtable}
\kwichead{meaningful}{92}
\begin{longtable}{@{}>{\raggedright}p{0.07\textwidth}>{\raggedleft}p{0.41\textwidth}@{\,}>{\raggedright\arraybackslash}p{0.47\textwidth}@{}}
01-c 1 & EXI-01-a defines a life as  & \textbf{meaningful} to the degree that it consists in activ... \\
01-c 2 & ..., and by Definition 1 of EXI-01-a it is  & \textbf{meaningful} to a positive degree. \\
01-c 3 & ...l result for achievement: a life can be  & \textbf{meaningful} through engagement without success (The... \\
01-d Abs & EXI-01-a holds that life is  & \textbf{meaningful} to the extent that one engages with pro... \\
01-d Abs & It does not imply that every life is  & \textbf{meaningful}, that all lives are equally meaningful,... \\
01-d Abs & ... meaningful, that all lives are equally  & \textbf{meaningful}, or that deprivation is harmless to mea... \\
01-d 1 & ...t settle whether any particular life is  & \textbf{meaningful}, which depends on what is done with it. \\
01-d 2 & Whether a life is  & \textbf{meaningful} does not depend on whether its projects... \\
01-d 2 & A life can be  & \textbf{meaningful} though its subject has no exceptional t... \\
01-d 3 & ...rthwhile, though she could have, is not  & \textbf{meaningful} (Remark 5 of EXI-01-a). \\
01-d 4 & Wolf (2010) insists that the objects of  & \textbf{meaningful} engagement are not restricted to high c... \\
01-d 4 & ...e to everyone, the claim that a life is  & \textbf{meaningful} is cheap and of no evaluative interest. \\
03-a 2 & EXI-01 argues that a life is  & \textbf{meaningful} to the extent that its subject engages ... \\
04-a 3 & EXI-01 argues that life is  & \textbf{meaningful} to the extent that one engages with pro... \\
04-a 5 & ...es not show that any particular life is  & \textbf{meaningful} or worthwhile; that question is answere... \\
04-a 5 & ...bsurdity, as defined here, is shared by  & \textbf{meaningful} and meaningless lives alike. \\
04-b 3 & ...ense of Definition 1.1, his life is not  & \textbf{meaningful} however happy he is, since meaning requ... \\
04-b 3 & ...syphus is fulfilled but his life is not  & \textbf{meaningful}, because his object is not worthwhile (... \\
06-b 2 & On the account of EXI-01, a life is  & \textbf{meaningful} to the extent that its subject engages ... \\
06-b 4 & ...not care about us, and our lives may be  & \textbf{meaningful} all the same. \\
09-a 5 & ...sistent with EXI-01, on which a life is  & \textbf{meaningful} to the extent that its projects are of ... \\
09-b 2 & ...wicked or, in the sense of EXI-01, less  & \textbf{meaningful} than it might have been, because it eng... \\
10-b 3 & ...ustain attention and to find activities  & \textbf{meaningful}, predict boredom on any relational view... \\
10-b 8 & EXI-17 holds that a  & \textbf{meaningful} life requires active engagement with so... \\
11-a Abs & Can life be  & \textbf{meaningful} without God? \\
11-a Abs & We argue that life can be  & \textbf{meaningful} without God, because meaning does not r... \\
11-a Abs & ...ulfilment account of meaning: a life is  & \textbf{meaningful} to the extent that its subject actively... \\
11-a 1 & On purpose theories, a life is  & \textbf{meaningful} insofar as it fulfils a purpose that Go... \\
11-a 1 & ... and the question whether life would be  & \textbf{meaningful} in such a world is a counterpossible wh... \\
11-a 1 & ...whether the conditions that make a life  & \textbf{meaningful} include God's existence, a relation to ... \\
11-a 1 & ...n God and does not relate to God can be  & \textbf{meaningful}. \\
11-a 2 & A life is  & \textbf{meaningful} to the extent that its subject is activ... \\
11-a 2 & The purpose theory holds that a life is  & \textbf{meaningful} only if it has a cosmic purpose and ful... \\
11-a 2 & ...f what is worthwhile, makes a life more  & \textbf{meaningful}. \\
11-a 3 & ... cosmic purpose is not sufficient for a  & \textbf{meaningful} life. \\
11-a 3 & ...Definition 2.1, such a life is not made  & \textbf{meaningful} by its purpose. \\
11-a 3 & ...akes the fulfilment of a cosmic purpose  & \textbf{meaningful}, on that reply, is the goodness of the ... \\
11-a 3 & ... does not fulfil the one it has, can be  & \textbf{meaningful}. \\
11-a 3 & Her life is therefore  & \textbf{meaningful} to that extent. \\
11-a 3 & ...n God and does not relate to God can be  & \textbf{meaningful}. \\
11-a 3 & ... with what is worthwhile makes her life  & \textbf{meaningful}. \\
11-a 4 & ...Principle P2, it would make a life more  & \textbf{meaningful} than engagement with lesser objects, ot... \\
11-a 4 & ...e, and Theorem 3.3 ensures that life is  & \textbf{meaningful} nonetheless. \\
11-a 5 & A life linked to these is  & \textbf{meaningful} without further linkage. \\
11-a 5 & The question of what makes a life  & \textbf{meaningful} is asked and answered by believers and ... \\
11-a 5 & If a life can be fully  & \textbf{meaningful} without God, nothing is left for theism... \\
11-a 5 & A life can be  & \textbf{meaningful}, and deeply so, while a further relatio... \\
11-a 5 & ...on to a greater good would make it more  & \textbf{meaningful} still. \\
11-a 5 & That a life could be more  & \textbf{meaningful} is not a deficiency in any sense that t... \\
11-a 6 & ... it show that every life without God is  & \textbf{meaningful}: meaning requires engagement with what ... \\
11-a 7 & Life can be  & \textbf{meaningful} without God. \\
11-b Abs & ...ndence and self-direction on which some  & \textbf{meaningful} lives depend. \\
11-b 1 & ... every life in a theistic world is more  & \textbf{meaningful} than every life in a naturalistic one. \\
11-b 1 & ....3 of EXI-11-a, lives without it can be  & \textbf{meaningful}, so it is not meaning as such. \\
14-a 1 & ...onnectedness and the sense that life is  & \textbf{meaningful} (Wildschut et al. 2006; Sedikides et al... \\
17-a Abs & Does a  & \textbf{meaningful} life require achievement? \\
17-a Abs & ...meaning give success a central place: a  & \textbf{meaningful} life is one that makes a difference, le... \\
17-a Abs & A  & \textbf{meaningful} life requires active engagement with so... \\
17-a Abs & ...ion judges the lives in each pair to be  & \textbf{meaningful}, and judges the unsuccessful life to be... \\
17-a 1 & ...ominate the literature (Metz 2013), are  & \textbf{meaningful} because of what they brought about, and... \\
17-a 1 & ...ng blades of grass (Rawls 1971), is not  & \textbf{meaningful}. \\
17-a 2 & A life is  & \textbf{meaningful} to the extent that its subject is activ... \\
17-a 2 & ...dered judgments about whether lives are  & \textbf{meaningful}, made under idealized reflection with t... \\
17-a 3 & ...he unsuccessful mathematician's life is  & \textbf{meaningful}, and so is the life of the nurse in the... \\
17-a 3 & By Principle P1, both lives are  & \textbf{meaningful}. \\
17-a 3 & The listener's life is  & \textbf{meaningful} to the extent of her engagement with mu... \\
17-a 3 & A  & \textbf{meaningful} life requires active engagement with so... \\
17-a 3 & ...tia, is not what reflection counts as a  & \textbf{meaningful} life (Wolf 2010). \\
17-a 3 & ...successful mathematician's life is more  & \textbf{meaningful} than the unsuccessful one's, or at leas... \\
17-a 3 & ...ch pair is not the difference between a  & \textbf{meaningful} and a meaningless life. \\
17-a 4 & So every  & \textbf{meaningful} life contains achievement after all. \\
17-a 4 & ...cian did not attain it, and her life is  & \textbf{meaningful}. \\
17-a 4 & ...it does not bear on whether her life is  & \textbf{meaningful}. \\
17-a 4 & ...is stone and achieving nothing, lives a  & \textbf{meaningful} life. \\
17-a 4 & ...pt a village from flooding would live a  & \textbf{meaningful} life, even if the stone always rolled b... \\
17-a 5 & A  & \textbf{meaningful} life requires active engagement with so... \\
17-a 5 & ... nurses, are judged on reflection to be  & \textbf{meaningful} in both cases; and lives devoted to app... \\
17-a 5 & ...appreciation rather than production are  & \textbf{meaningful} without achieving anything. \\
19-b Abs & ...it is a condition of a life's remaining  & \textbf{meaningful} throughout. \\
19-b 1 & ...account of EXI-01 and EXI-17, a life is  & \textbf{meaningful} to the extent that its subject actively... \\
19-b 1 & ...mething whose engagement makes the life  & \textbf{meaningful}, and does not make anything engaged wit... \\
19-b 1 & ...de is a condition of a life's remaining  & \textbf{meaningful} throughout its course. \\
19-b 1 & ...ss life would at some point cease to be  & \textbf{meaningful}, and a life remains meaningful througho... \\
19-b 1 & ...se to be meaningful, and a life remains  & \textbf{meaningful} throughout its course only if it ends b... \\
19-b 1 & ... engagement with what is worthwhile, as  & \textbf{meaningful} as ours; the claim is that it would not... \\
19-b 3 & Finitude does not make life  & \textbf{meaningful}. \\
20-a Abs & Was producing this collection  & \textbf{meaningful}? \\
20-a Abs & EXI-17 argues that a  & \textbf{meaningful} life, or a meaningful activity within o... \\
20-a Abs & ...-17 argues that a meaningful life, or a  & \textbf{meaningful} activity within one, requires active en... \\
20-a 1 & ...ets objective attractiveness: a life is  & \textbf{meaningful} to the extent that its subject is activ... \\
20-a 4 & Their activity is  & \textbf{meaningful}, by EXI-17, to the extent that the coll... \\
20-b Abs & ...er producing the present collection was  & \textbf{meaningful} depends, by the account of EXI-17, on w... \\
\end{longtable}
\kwichead{nothing}{79}
\begin{longtable}{@{}>{\raggedright}p{0.07\textwidth}>{\raggedleft}p{0.41\textwidth}@{\,}>{\raggedright\arraybackslash}p{0.47\textwidth}@{}}
01-c 2 &  & \textbf{Nothing} in Definition 1 requires that a project... \\
01-c 4 & The case of an agent for whom  & \textbf{nothing} worthwhile is within reach of attention... \\
01-d 3 & ... a life in which the agent engages with  & \textbf{nothing} worthwhile, though she could have, is n... \\
01-d 4 & From a sufficiently detached standpoint  & \textbf{nothing} matters, so meaning is available to no ... \\
01-d 4 & ...a vindication, not a demonstration that  & \textbf{nothing} is worthwhile. \\
02-a 1 & ...icurus wrote to Menoeceus that death is  & \textbf{nothing} to us, since good and bad lie in sensat... \\
02-a 1 & ... eternity in which we did not exist, is  & \textbf{nothing} to us, and the time after our death is ... \\
02-a 1 & ...th is its mirror image, so it should be  & \textbf{nothing} to us either (Lucretius ). \\
02-a 2 & ...anism persists for a while as a corpse;  & \textbf{nothing} below depends on it. \\
02-a 2 & It is extrinsic: an event need contain  & \textbf{nothing} bad in itself to be bad for . \\
02-a 3 &  & \textbf{Nothing} in Principle P1 requires that exist at ... \\
02-a 4 & ... different life from the beginning, and  & \textbf{nothing} guarantees that it contains more of 's ... \\
02-b Abs & Epicurus argued that death is  & \textbf{nothing} to us, and drew from this that it shoul... \\
02-b Abs & ...n he did not give, and wrong that it is  & \textbf{nothing} to us. \\
02-b 1 & ...enoeceus, Epicurus argues that death is  & \textbf{nothing} to us, since good and evil lie in sensa... \\
02-b 2 & ...to fear, not because death is after all  & \textbf{nothing} to us. \\
03-b 2 & ...s continuation would eventually contain  & \textbf{nothing} that gives life its point. \\
04-a 1 & ...rief in an immense stretch of time, and  & \textbf{nothing} we do will matter in a million years. \\
04-a 1 & ... and then the step back seems to reveal  & \textbf{nothing} troubling. \\
04-a 3 &  & \textbf{Nothing} in reflection forbids the step back, an... \\
04-a 3 & An indifferent world denies us  & \textbf{nothing} and cannot be defeated by our persisten... \\
04-a 4 & ...t, from the standpoint of the universe,  & \textbf{nothing} requires that anyone find out how prote... \\
04-b 1 & ... top, and its staying there would serve  & \textbf{nothing}. \\
04-b 1 & ...iscloses is not that they fail or serve  & \textbf{nothing}, but that no justification of our conce... \\
04-b 1 & ...se of Definition 2.3 of EXI-04-a, since  & \textbf{nothing} in her situation raises the question of... \\
04-b 1 & ... the recognition that his task achieves  & \textbf{nothing} and is worth nothing. \\
04-b 1 & ... his task achieves nothing and is worth  & \textbf{nothing}. \\
04-b 2 & ...imself to the effort while knowing that  & \textbf{nothing} outside the effort vindicates it. \\
05-a 1 & ...it is present, and that death, which is  & \textbf{nothing} to us, should not be feared (EpicurusMe... \\
05-a 4 & Where  & \textbf{nothing} is uncertain, what remains is a sorrow ... \\
06-b Abs & ...rse's want of attitudes toward us takes  & \textbf{nothing} from that worth. \\
06-b 2 & ...he absence of responsiveness, subtracts  & \textbf{nothing}. \\
06-b 2 & ...osition 2.1, its indifference denies us  & \textbf{nothing} that meaning requires; by Theorem 3.1 o... \\
06-b 2 & ... follows from indifference alone, since  & \textbf{nothing} else will maintain them. \\
06-b 3 & ... universe has no such point of view, so  & \textbf{nothing} matters from it. \\
07-a Abs & ...y final end is her own happiness values  & \textbf{nothing} else for its own sake, and so lacks the... \\
07-a 1 & A being with self-love alone would have  & \textbf{nothing} to enjoy. \\
07-a 2 & ...rs only as a source of satisfaction has  & \textbf{nothing} for the satisfaction to be about. \\
07-a 6 & ...y final end is her own happiness values  & \textbf{nothing} else for its own sake. \\
08-a 3 & ...meanings, as Principle P1 says it must;  & \textbf{nothing} in the definition asks that the examina... \\
09-a 1 & ..., and defines himself afterwards; he is  & \textbf{nothing} other than what he makes of himself. \\
09-a 3 &  & \textbf{Nothing} in the real definition of a human being... \\
10-a 2 & ...ecause, whatever its demands, it serves  & \textbf{nothing} she cares about (meaninglessness). \\
10-a 5 & ...edom would become a correct signal that  & \textbf{nothing} could answer. \\
10-b 8 & ...would then become a correct signal that  & \textbf{nothing} could answer, and the trait findings of... \\
11-a Abs & ...hout these every human pursuit comes to  & \textbf{nothing}. \\
11-a 5 & What comes to  & \textbf{nothing} was nothing (Tolstoy 1882; Craig 2008). \\
11-a 5 & What comes to nothing was  & \textbf{nothing} (Tolstoy 1882; Craig 2008). \\
11-a 5 & ...fe can be fully meaningful without God,  & \textbf{nothing} is left for theism to add; if theism ad... \\
11-a 5 & ... dilemma assumes that meaning is all or  & \textbf{nothing}. \\
12-a 4 & ...od, Theorem 3.1 does not apply to them;  & \textbf{nothing} in our argument denies that some lives ... \\
13-a 1 & If, with Epicurus,  & \textbf{nothing} that happens after death can be bad for... \\
13-a 4 & The dead do not exist, so  & \textbf{nothing} can be owed to them, and the appearance... \\
13-a 4 & ...iving person remembers still binds, and  & \textbf{nothing} in the concern of the living explains w... \\
13-a 4 & ...chological states at all; on that view,  & \textbf{nothing} changes with the passage of time, and t... \\
13-b 1 & ...gainst competing reasons, but it is not  & \textbf{nothing}. \\
13-b 2 &  & \textbf{Nothing} is created in favor of a corpse, and af... \\
14-a 4 & If the past exists and is fixed,  & \textbf{nothing} is lost, and nostalgia misrepresents. \\
14-a 4 & ...lgia: the good it remembers is real and  & \textbf{nothing} can now make it not to have been. \\
14-a 4 &  & \textbf{Nothing} can be done about the past. \\
15-a 5 & ...bility is presumption, and hope that is  & \textbf{nothing} but probability gives up goods that age... \\
15-b Abs & ...jection that the virtue so described is  & \textbf{nothing} but expected-value reasoning. \\
15-b 2 & ...ifically hopeful part of the stance has  & \textbf{nothing} further to do. \\
15-b 2 & ...ntion to the ways might come about adds  & \textbf{nothing} to the plan, since the plan already ass... \\
15-b 2 & ...the hope of a lottery player who can do  & \textbf{nothing} to affect the draw, nothing sustains it... \\
15-b 2 & ... who can do nothing to affect the draw,  & \textbf{nothing} sustains it, and the stance that remain... \\
15-b 4 & Hope adds  & \textbf{nothing}. \\
16-a 4 & ...3.2; this is the correct verdict, since  & \textbf{nothing} would be lost if the first sight had be... \\
17-a 1 & ...y brought about, and a life that brings  & \textbf{nothing} about is to that extent empty. \\
17-a 1 & ...ing a life can have, though he achieved  & \textbf{nothing}. \\
17-a 3 & ...ith growing understanding, but produces  & \textbf{nothing}: she does not perform, compose or write... \\
17-a 3 &  & \textbf{Nothing} in her activity is an achievement in th... \\
17-a 4 & A life that changes  & \textbf{nothing} is enclosed within itself. \\
17-a 4 & ...s, engaged with his stone and achieving  & \textbf{nothing}, lives a meaningful life. \\
17-a 4 & ...e does not, but not because he achieves  & \textbf{nothing}: his object is not worthwhile. \\
19-a 4 & Death adds  & \textbf{nothing}. \\
19-b Abs & ..., and that an immortal being could find  & \textbf{nothing} worth doing. \\
19-b 2 & Without scarcity  & \textbf{nothing} would be urgent, and nothing would matt... \\
19-b 2 & ...t scarcity nothing would be urgent, and  & \textbf{nothing} would matter. \\
\end{longtable}
\kwichead{objective}{25}
\begin{longtable}{@{}>{\raggedright}p{0.07\textwidth}>{\raggedleft}p{0.41\textwidth}@{\,}>{\raggedright\arraybackslash}p{0.47\textwidth}@{}}
01-c Abs & ...le to engage in at least one project of  & \textbf{objective} worth (Assumption 3). \\
01-c 1 & ...s in active engagement with projects of  & \textbf{objective} worth (Definition 1), assumes that ther... \\
01-c 1 & ...nd the very poor may have no project of  & \textbf{objective} worth on which they can embark, and Wol... \\
01-c 2 & ...ains it, her engagement is a project of  & \textbf{objective} worth in the sense of Definition 1 of E... \\
01-c 2 & ... activity is engagement in a project of  & \textbf{objective} worth. \\
01-c 2 & ... in active engagement with a project of  & \textbf{objective} worth, and by Definition 1 of EXI-01-a ... \\
01-c 5 & Engagement with them is a project of  & \textbf{objective} worth, and a life that contains it has ... \\
01-d Abs & ...xtent that one engages with projects of  & \textbf{objective} worth, and that meaning in this sense i... \\
01-d 2 & ...asures meaning by engagement and by the  & \textbf{objective} worth of what one engages with. \\
01-d 2 & ...g a skill to a neighbor are projects of  & \textbf{objective} worth, and engagement with them confers... \\
03-a 2 & ...at its subject engages with projects of  & \textbf{objective} worth, and EXI-17 that such engagement,... \\
04-a 1 & ... tension between the subjective and the  & \textbf{objective} standpoints. \\
04-a 1 & The other asks whether  & \textbf{objective} value, if there is any, dissolves the a... \\
04-a 3 & ... she answers that the consideration has  & \textbf{objective} worth, the detached standpoint can regi... \\
04-a 3 & ...ma 3.2 does not assume that there is no  & \textbf{objective} worth. \\
04-a 3 & ...xtent that one engages with projects of  & \textbf{objective} worth, and we accept that conclusion. \\
04-a 4 & ... of EXI-01-a puts it, the assumption of  & \textbf{objective} worth denies the premise of the sceptic... \\
06-b 2 & ...at its subject engages with projects of  & \textbf{objective} worth, whose value does not depend on t... \\
06-b 3 & ...ng consists in engagement with what has  & \textbf{objective} worth, as EXI-01 holds, then it does no... \\
09-a 5 & ... to the extent that its projects are of  & \textbf{objective} worth. \\
11-a 1 & On  & \textbf{objective} naturalist views, it is a matter of sta... \\
11-a 1 & ...e subject's active engagement meets the  & \textbf{objective} worth of its object. \\
17-a 1 & ...Wolf (2010) combined the subjective and  & \textbf{objective} elements. \\
17-a 1 & ...arises when subjective attraction meets  & \textbf{objective} attractiveness: the subject must be act... \\
20-a 1 & ...arises when subjective attraction meets  & \textbf{objective} attractiveness: a life is meaningful to... \\
\end{longtable}
\kwichead{particular}{27}
\begin{longtable}{@{}>{\raggedright}p{0.07\textwidth}>{\raggedleft}p{0.41\textwidth}@{\,}>{\raggedright\arraybackslash}p{0.47\textwidth}@{}}
01-d 1 & ...the theorem does not settle whether any  & \textbf{particular} life is meaningful, which depends on wh... \\
01-d 4 & ...h by prestige or required projects of a  & \textbf{particular} cultural form. \\
02-a 4 & ...ID-01), and the organism's origin, in a  & \textbf{particular} pair of gametes, cannot be moved far in... \\
03-b 3 & ...osition 2.1 explains why it attaches to  & \textbf{particular} projects and not to the bare prospect o... \\
04-a 5 & The argument does not show that any  & \textbf{particular} life is meaningful or worthwhile; that ... \\
05-a Abs & ...e's own death as such, as distinct from  & \textbf{particular} threats to one's life, is certain. \\
05-a 4 & ...ts object is not death as such; it is a  & \textbf{particular} threat, or the possibility that one's l... \\
05-a 4 & Fear of a  & \textbf{particular} danger is a fear of a prospective evil ... \\
06-b 3 & ...s are assessed without reference to any  & \textbf{particular} agent, then the worth of objects of eng... \\
07-a 1 & ...oyment consists in the gratification of  & \textbf{particular} passions, for food, for esteem, for the... \\
07-a 4 & ...r end is not a standing filter on every  & \textbf{particular} concern; once it is not, happiness is n... \\
07-a 4 & These are local pursuits of  & \textbf{particular} goods of experience, embedded in a life... \\
09-a 4 & ...t any good human life takes, not what a  & \textbf{particular} life is for. \\
11-a 4 & ...te view in making God necessary for any  & \textbf{particular} grade of meaning. \\
11-b Abs & ...d that the anti-theist's losses concern  & \textbf{particular} life-plans, not the objects of worth av... \\
11-b 2 & What would be lost is a  & \textbf{particular} manner of engaging, unobserved and unai... \\
13-a 3 & ...ependents are interests in the lives of  & \textbf{particular} persons, and they are fulfilled or exti... \\
13-a 4 & ...on 2.2 a specific reliance created by a  & \textbf{particular} person, and by Principle P2 the promiso... \\
16-a 1 & ...r those who think love is directed at a  & \textbf{particular} person: if the reasons for love are qua... \\
16-a 2 & ...ruism about histories, which consist of  & \textbf{particular} events with particular participants; a ... \\
16-a 2 & ...which consist of particular events with  & \textbf{particular} participants; a qualitatively similar h... \\
16-a 4 & ...y love is answerable to the course of a  & \textbf{particular} relationship. \\
19-a 1 & ...values, such as courage and the love of  & \textbf{particular} persons, have their character only for ... \\
19-a 4 & ... is a condition of the good of shape; a  & \textbf{particular} death is usually a loss. \\
19-a 5 & This does not make any  & \textbf{particular} death good for the one who dies. \\
19-b 2 & ...nditions are met by a limit, not by any  & \textbf{particular} one; a later death would also have give... \\
20-b Abs & ... MND-08-c and ETH-24-b identify for any  & \textbf{particular} case. \\
\end{longtable}
\kwichead{positive}{23}
\begin{longtable}{@{}>{\raggedright}p{0.07\textwidth}>{\raggedleft}p{0.41\textwidth}@{\,}>{\raggedright\arraybackslash}p{0.47\textwidth}@{}}
01-c Abs & EXI-01-a establishes that a life of  & \textbf{positive} meaning is available to every agent cap... \\
01-c 2 & ...apable of intentional action, a life of  & \textbf{positive} meaning is available. \\
01-c 2 & ... includes such engagement consists to a  & \textbf{positive} extent in active engagement with a proj... \\
01-c 2 & ...ion 1 of EXI-01-a it is meaningful to a  & \textbf{positive} degree. \\
01-d 1 & ...capable of intentional action a life of  & \textbf{positive} meaning is available. \\
01-d 2 & ...which it can attain it, but not whether  & \textbf{positive} meaning is available. \\
01-d 4 & That  & \textbf{positive} meaning is available to everyone is com... \\
02-a 2 & ...alue is negative, and good for if it is  & \textbf{positive}. \\
02-a 3 & ... hypothesis the continuation contains a  & \textbf{positive} balance of goods, so and the value of h... \\
02-a 3 & ...fe worth living, the continuation has a  & \textbf{positive} balance. \\
02-a 5 & ...s would otherwise have continued with a  & \textbf{positive} balance of goods, and the closest world... \\
07-a 2 & ...ominance of enjoyment, satisfaction and  & \textbf{positive} emotional condition over their opposite... \\
10-a 1 & ...er it is evidence that existence has no  & \textbf{positive} value; for Svendsen (2005) it is a char... \\
14-a 1 & ... that nostalgia is a common and largely  & \textbf{positive} emotion that raises self-continuity, so... \\
14-a 4 & ...is reconstructive and biased toward the  & \textbf{positive}. \\
15-b 3 & If is  & \textbf{positive}, the effort that the stance commits mak... \\
18-a 1 & ...rown (1988) argued that unrealistically  & \textbf{positive} self-evaluations, exaggerated perceptio... \\
18-a 1 & ...nctioning cognition, and concluded that  & \textbf{positive} illusions are the clearest candidates f... \\
18-a 2 & Several of the functions attributed to  & \textbf{positive} illusions can be served by attitudes th... \\
18-b Abs & ...and Brown proposed that unrealistically  & \textbf{positive} self-evaluations, exaggerated perceptio... \\
18-b 1 & ...ylor and Brown (1988) argued that three  & \textbf{positive} illusions are associated with persisten... \\
18-b 1 & (i) Whether a given  & \textbf{positive} illusion is psychologically necessary f... \\
18-b 1 & ...; to that extent the range of necessary  & \textbf{positive} illusions shrinks. \\
\end{longtable}
\kwichead{premise}{18}
\begin{longtable}{@{}>{\raggedright}p{0.07\textwidth}>{\raggedleft}p{0.41\textwidth}@{\,}>{\raggedright\arraybackslash}p{0.47\textwidth}@{}}
01-c 1 & ...f the universality claim, and it is the  & \textbf{premise} a critic will press. \\
01-c 2 & ...n 2.2 and Lemma 2.3, and is no longer a  & \textbf{premise}. \\
01-c 5 & ...onclusion of EXI-01-a without its third  & \textbf{premise}. \\
01-d 4 & ...01-a notes that Assumption 2 denies the  & \textbf{premise} of this challenge rather than refuting ... \\
02-a 3 & Its first  & \textbf{premise}, the experience requirement, is false b... \\
02-a 3 & So the argument needs a  & \textbf{premise} that the deprivation account shows to b... \\
02-a 5 & ...e objection's last step fails: from the  & \textbf{premise} that a fact is not located at a time wh... \\
02-a 5 & ...urean can rest the argument on a weaker  & \textbf{premise}: a state of affairs is bad for a person... \\
02-a 5 & The weaker  & \textbf{premise} draws the line in the wrong place. \\
02-b 2 & Its  & \textbf{premise} combines two requirements: an experienc... \\
02-b 2 & ... survives in part, and his metaphysical  & \textbf{premise} does not. \\
03-a 4 & Boredom is the symptom, not the  & \textbf{premise}. \\
04-a 4 & ...ssumption of objective worth denies the  & \textbf{premise} of the sceptical challenge without refu... \\
05-a Abs & EXI-02 rejects that  & \textbf{premise}: death is bad for the one who dies when... \\
05-a Abs & ...onclusion about fear from the Epicurean  & \textbf{premise} about badness. \\
06-a 2 & Assumption 2.3 is a  & \textbf{premise} about the content of physical law, not ... \\
10-a 2 & Assumption 2.4 is the  & \textbf{premise} that Csikszentmihalyi (1990) documented... \\
11-a 4 & Our argument does not use REL-01 as a  & \textbf{premise}. \\
\end{longtable}
\kwichead{project}{77}
\begin{longtable}{@{}>{\raggedright}p{0.07\textwidth}>{\raggedleft}p{0.41\textwidth}@{\,}>{\raggedright\arraybackslash}p{0.47\textwidth}@{}}
01-c Abs & ...agent is able to engage in at least one  & \textbf{project} of objective worth (Assumption 3). \\
01-c Abs & ...gagement with one of them constitutes a  & \textbf{project} in the sense relevant to the definition... \\
01-c 1 & ... very ill and the very poor may have no  & \textbf{project} of objective worth on which they can em... \\
01-c 2 & ...t, and sustains it, her engagement is a  & \textbf{project} of objective worth in the sense of Defi... \\
01-c 2 & What Definition 1 requires of a  & \textbf{project} is that it be something the agent engag... \\
01-c 2 & Nothing in Definition 1 requires that a  & \textbf{project} have a completion condition, a public p... \\
01-c 2 & ...aining that activity is engagement in a  & \textbf{project} of objective worth. \\
01-c 2 & ...tive extent in active engagement with a  & \textbf{project} of objective worth, and by Definition 1... \\
01-c 3 & ... does not wait on the availability of a  & \textbf{project}, because engagement makes one. \\
01-c 4 & Lemma 2.3 trivializes the notion of a  & \textbf{project}. \\
01-c 4 & If attending to a visitor is a  & \textbf{project}, the word has lost its content. \\
01-c 4 & A visit attended to once is not a  & \textbf{project}; a relationship of care sustained over ... \\
01-c 5 & Engagement with them is a  & \textbf{project} of objective worth, and a life that con... \\
01-d 2 & ...eaning; scale bears on how much value a  & \textbf{project} serves, which can affect degree but not... \\
01-d 4 & ...ion 1 requires only that the value of a  & \textbf{project} not depend on the attitudes of those en... \\
03-a Abs & ...e of such projects is bounded, and each  & \textbf{project}, pursued long enough, is either complet... \\
03-a 2 & ...gives life its point if its object is a  & \textbf{project} that the person takes to be worthwhile ... \\
03-a 2 & For any character , the set of  & \textbf{project} types that a person of character can en... \\
03-a 2 & A character supports engagement with a  & \textbf{project} type only by having interests, capaciti... \\
03-a 2 & For each  & \textbf{project} type, a person who engages with it for ... \\
03-a 2 & A  & \textbf{project} type that has been completed or reduced... \\
03-a 2 & ...nd values of finite complexity, and the  & \textbf{project} types that answer to it are those that ... \\
03-a 2 & It does not say that every  & \textbf{project} can be completed; it says that if a pro... \\
03-a 2 & ...ect can be completed; it says that if a  & \textbf{project} is not completed, then a person who pur... \\
03-a 3 & ...ers similar to the present one, and the  & \textbf{project} types answering to characters in that r... \\
03-a 4 & ...e while she does, but the meal is not a  & \textbf{project} she takes to be worthwhile independentl... \\
03-a 4 & ...der the next objection: friendship is a  & \textbf{project}, and its renewal depends on the friends... \\
03-b Abs & ...fe its point, since its object is not a  & \textbf{project}. \\
03-b 1 & ...gives life its point if its object is a  & \textbf{project} she takes to be worthwhile for reasons ... \\
03-b 1 & Given a bounded repertoire of  & \textbf{project} types for any character (Assumption 2.3... \\
03-b 1 & ...Assumption 2.3), the exhaustion of each  & \textbf{project} type pursued long enough (Assumption 2.... \\
03-b 2 & ... life its point only if its object is a  & \textbf{project} the person takes to be worthwhile. \\
03-b 2 & Continuation is not a  & \textbf{project}: it is the condition of all projects, a... \\
03-b 2 & ...orical desires: projects unfinished and  & \textbf{project} types not yet taken up. \\
03-b 2 & ...ounds of connectedness and so opens new  & \textbf{project} types as others close. \\
04-a 3 & Whether a  & \textbf{project} is worthwhile is one question; whether ... \\
06-b 2 & If the worth of a  & \textbf{project} does not depend on the attitudes of tho... \\
09-a Abs & ... to read Sartre's essence as a defining  & \textbf{project}, a determination of what a thing is for... \\
09-a Abs & ...or to her existence, because a defining  & \textbf{project} for a life can be fixed only by the age... \\
09-a 2 & A defining  & \textbf{project} for an entity is a determination of wha... \\
09-a 2 & ...ssence precedes existence if a defining  & \textbf{project} for is fixed before exists, by a determ... \\
09-a 2 & ...istence precedes essence if no defining  & \textbf{project} for is fixed before exists, and whateve... \\
09-a 2 & ...ed before exists, and whatever defining  & \textbf{project} has is fixed in the course of 's existe... \\
09-a 2 & A defining  & \textbf{project} is not a real definition in the sense o... \\
09-a 2 & ... reasons, but does not fix its defining  & \textbf{project}. \\
09-a 3 & ...a defective knife), so it is a defining  & \textbf{project} for (Definition 2.1). \\
09-a 3 & ...ciple P2, none of them fixes a defining  & \textbf{project} for , since is capable of adopting ends... \\
09-a 3 & So no defining  & \textbf{project} for is fixed before exists. \\
09-a 3 & Whatever defining  & \textbf{project} has is fixed by the ends adopts, which ... \\
09-a 3 & ... result: Sartre's essence is a defining  & \textbf{project}, and his claim is that no such project ... \\
09-a 3 & ... project, and his claim is that no such  & \textbf{project} is given to human beings in advance of ... \\
09-a 4 & ...nly that such intentions fix a defining  & \textbf{project} for a being capable of adopting ends. \\
09-a 4 & Human beings therefore have a defining  & \textbf{project} prior to any choice. \\
09-a 4 & A defining  & \textbf{project} in the sense of Definition 2.1 is the m... \\
09-a 4 & ... a social imposition and not a defining  & \textbf{project} that a woman fails by refusing. \\
09-a 4 & ... Principle P2, they become her defining  & \textbf{project} only if she adopts them. \\
09-a 6 & Read as a defining  & \textbf{project}, an artefact's essence is fixed by the ... \\
09-a 6 & ...edes it, while a human being's defining  & \textbf{project} can be fixed only by the ends she adopt... \\
09-b Abs & ...olds once essence is read as a defining  & \textbf{project} and artefact kinds are taken to be fixe... \\
09-b Abs & ...nife's real definition and its defining  & \textbf{project} coincide in content but differ in role,... \\
09-b Abs & ...'s real definition contains no defining  & \textbf{project}. \\
09-b 1 & ...t, what the knife is for, as a defining  & \textbf{project} (Definition 2.1 of EXI-09-a), and deriv... \\
09-b 2 & ...hat kind of thing it is, and a defining  & \textbf{project}, stating what it is for, and the real d... \\
09-b 2 & ...e real definition includes the defining  & \textbf{project}. \\
09-b 2 & ...nition that does not include a defining  & \textbf{project}. \\
09-b 2 & ... be for cutting paper, and the defining  & \textbf{project} is a constituent of it. \\
09-b 2 & ...of EXI-09-a, the paper-knife's defining  & \textbf{project} is fixed before it exists, so there is ... \\
09-b 2 & By Lemma 3.2 of EXI-09-a, no defining  & \textbf{project} for a human life is fixed before it exi... \\
09-b 2 & ...ail as a human life by departing from a  & \textbf{project} it never had. \\
09-b 2 & ...pply standards of value, not a defining  & \textbf{project}, and they do not make a life a defectiv... \\
09-b 2 & But a defining  & \textbf{project}, in Definition 2.1 of EXI-09-a, is a de... \\
09-b 2 & ...ailing as what it is; the paper-knife's  & \textbf{project} is fixed by intention and is not an evo... \\
09-b 2 & ....2 of EXI-09-a requires that a defining  & \textbf{project} be fixed before the object exists by a ... \\
09-b 3 & ...ct the real definition and the defining  & \textbf{project} coincide. \\
09-b 3 & ...a real definition but no prior defining  & \textbf{project}, and so cannot be defective in the way ... \\
17-a 2 & An object of engagement (a  & \textbf{project}, a person, a practice, an inquiry) is w... \\
18-a 2 & ...; a practical commitment to act as if a  & \textbf{project} will succeed sustains persistence witho... \\
\end{longtable}
\kwichead{projects}{61}
\begin{longtable}{@{}>{\raggedright}p{0.07\textwidth}>{\raggedleft}p{0.41\textwidth}@{\,}>{\raggedright\arraybackslash}p{0.47\textwidth}@{}}
01-c Abs & ...tances deny them the opportunities that  & \textbf{projects} ordinarily require. \\
01-c Abs & We begin not from  & \textbf{projects} but from objects of engagement, such as... \\
01-c 1 & ...t it consists in active engagement with  & \textbf{projects} of objective worth (Definition 1), assu... \\
01-c 1 & ...inition 1), assumes that there are such  & \textbf{projects} (Assumption 2), and assumes that every ... \\
01-c 1 &  & \textbf{Projects}, as ordinarily understood, are extended... \\
01-c 3 & ...ssumption 3 asserts the availability of  & \textbf{projects}; Assumption 2.2 asserts the availabilit... \\
01-c 3 & ...rence matters because opportunities for  & \textbf{projects} in the ordinary sense depend on resourc... \\
01-c 4 & What the lemma drops are features of  & \textbf{projects} that are irrelevant to meaning on Defin... \\
01-c 4 & Assumption 3, read as a claim about  & \textbf{projects} in the ordinary sense, is open to count... \\
01-c 5 & ...depend on the universal availability of  & \textbf{projects}. \\
01-d Abs & ...ful to the extent that one engages with  & \textbf{projects} of objective worth, and that meaning in... \\
01-d Abs & ...mes, to talent or to the scale of one's  & \textbf{projects}; that no one's life is condemned to mea... \\
01-d 1 & ... does not depend on the availability of  & \textbf{projects} in the ordinary sense, since sustained ... \\
01-d 2 & ...aningful does not depend on whether its  & \textbf{projects} succeed. \\
01-d 2 & ...bject has no exceptional talent and her  & \textbf{projects} are small in scope. \\
01-d 2 & ..., or teaching a skill to a neighbor are  & \textbf{projects} of objective worth, and engagement with... \\
01-d 2 & ...hows that even agents without access to  & \textbf{projects} in the ordinary sense have worthwhile o... \\
01-d 4 & ...ve-worth views of meaning privilege the  & \textbf{projects} valued by the educated, such as art, sc... \\
01-d 4 & ... measured worth by prestige or required  & \textbf{projects} of a particular cultural form. \\
03-a Abs & ...n living, are desires whose objects are  & \textbf{projects} she takes to be worthwhile; given a cha... \\
03-a Abs & ...her present one, the repertoire of such  & \textbf{projects} is bounded, and each project, pursued l... \\
03-a 1 & ...at give life its point with desires for  & \textbf{projects} the person takes to be worthwhile, in l... \\
03-a 2 & ...he extent that its subject engages with  & \textbf{projects} of objective worth, and EXI-17 that suc... \\
03-a 2 & ...its point are desires to engage in such  & \textbf{projects}. \\
03-a 3 &  & \textbf{Projects} of other types do not engage her, by As... \\
03-a 5 & ...n endless succession of experiences and  & \textbf{projects}, such as the timeless states of some th... \\
03-a 6 & ...re understood as desires for worthwhile  & \textbf{projects}, and once his appeal to sameness of cha... \\
03-b 2 & ...t a project: it is the condition of all  & \textbf{projects}, and its worth for the person is deriva... \\
03-b 2 & ... mind the continuation of their present  & \textbf{projects}, which are finite in number and each, b... \\
03-b 2 & ...l have unexhausted categorical desires:  & \textbf{projects} unfinished and project types not yet ta... \\
03-b 3 & ... explains why it attaches to particular  & \textbf{projects} and not to the bare prospect of continu... \\
04-a Abs & ...ible with full engagement in worthwhile  & \textbf{projects}. \\
04-a 1 & ...ose of Wolf (2010) and Metz (2013), our  & \textbf{projects} may be worthwhile independently of our ... \\
04-a 3 & ...ful to the extent that one engages with  & \textbf{projects} of objective worth, and we accept that ... \\
04-a 4 & If some  & \textbf{projects} are worthwhile independently of our att... \\
04-a 4 & The seriousness is justified, if the  & \textbf{projects} are worthwhile; the absurd survives bec... \\
04-b 1 & Our  & \textbf{projects} often succeed, and many of their object... \\
04-b 1 & ...d, including the life of an agent whose  & \textbf{projects} succeed and whose objects are worthwhil... \\
04-b 4 & Death ends lives and many  & \textbf{projects}; it does not undo the value of what was... \\
06-b 2 & ...he extent that its subject engages with  & \textbf{projects} of objective worth, whose value does no... \\
09-a 3 & ...s to exist; Lemma 3.2 concerns defining  & \textbf{projects}, which the real definition does not sup... \\
09-a 5 & ...h have real definitions but no defining  & \textbf{projects} in our sense. \\
09-a 5 & ...fe is meaningful to the extent that its  & \textbf{projects} are of objective worth. \\
11-a Abs & ...mal worth, a purpose that unifies one's  & \textbf{projects}, and the preservation of what is good o... \\
11-a 4 & It unifies a person's  & \textbf{projects} under a single purpose that is itself g... \\
11-b 1 & ...d in relation, a purpose unifying one's  & \textbf{projects}, and the preservation of what is good. \\
11-b 1 & ...d, or with a purpose that unifies one's  & \textbf{projects}, is engagement with what is worthwhile;... \\
11-b 2 & ...ngham (2003) holds that without God our  & \textbf{projects} lack an ultimate point; what theism add... \\
11-b 2 & ...r a single good towards which all one's  & \textbf{projects} could be ordered; Proposition 1.1 count... \\
13-a Abs & The entitlement protects interests in  & \textbf{projects}, dependents and reputation, and these i... \\
13-a 1 & ...nce many of them, in the success of her  & \textbf{projects}, the welfare of her children, the fate ... \\
13-a 3 & Interests in the completion of  & \textbf{projects} are fulfilled when the projects are com... \\
13-a 3 & ...tion of projects are fulfilled when the  & \textbf{projects} are completed and lapse when their comp... \\
18-a 3 & Meaning depends on taking one's  & \textbf{projects} to matter more than any detached view c... \\
19-a 3 & ...t is often false: most deaths interrupt  & \textbf{projects} rather than complete them. \\
19-a 3 & ...h gives a better shape, since it allows  & \textbf{projects} to be completed that the earlier death ... \\
19-a 4 &  & \textbf{Projects}, relationships and stages of life have ... \\
19-b Abs & ...tude is the source of meaning: that our  & \textbf{projects} matter because our time is limited, and... \\
19-b 2 &  & \textbf{Projects} matter to us because time is limited: c... \\
19-b 2 & Scarcity gives choices their cost and  & \textbf{projects} their urgency, and urgency affects how ... \\
20-a 1 & ...that its subject is actively engaged in  & \textbf{projects} that are worthwhile independently of he... \\
\end{longtable}
\kwichead{question}{80}
\begin{longtable}{@{}>{\raggedright}p{0.07\textwidth}>{\raggedleft}p{0.41\textwidth}@{\,}>{\raggedright\arraybackslash}p{0.47\textwidth}@{}}
02-a 2 & Our  & \textbf{question} concerns death as the end of the person... \\
02-a 6 & ...ional to fear death, which is a further  & \textbf{question} (EXI-05). \\
03-a 2 & ... infinitely many worthwhile things; the  & \textbf{question} is how many kinds of them a given chara... \\
03-a 4 & But the  & \textbf{question} is whether the endless life is desirabl... \\
04-a 1 & Nagel (1986) returned to the  & \textbf{question} and treated the collision as one instan... \\
04-a 2 & ...ble that is not itself open to the same  & \textbf{question}. \\
04-a 2 & The  & \textbf{question} of the appropriate response to the absu... \\
04-a 2 & ...iate response to the absurd is then the  & \textbf{question} of which attitude presents the absurd s... \\
04-a 2 & The irony in  & \textbf{question} is not the rhetorical figure of saying ... \\
04-a 3 & ..., the detached standpoint can raise the  & \textbf{question} of what justifies that justification, a... \\
04-a 3 & ... by citing a further consideration, the  & \textbf{question} recurs. \\
04-a 3 & ...tion, which the step back has placed in  & \textbf{question}. \\
04-a 3 & Whether a project is worthwhile is one  & \textbf{question}; whether its worth can be certified to ... \\
04-a 3 & ... available that is not open to the same  & \textbf{question}. \\
04-a 4 & ... 3.2 gives it no such authority: it can  & \textbf{question} but not condemn. \\
04-a 4 & But the  & \textbf{question} was what attitude fits the absurd when ... \\
04-a 4 & ...when one does attend to it, and to that  & \textbf{question} forgetting gives no answer. \\
04-a 5 & ... life is meaningful or worthwhile; that  & \textbf{question} is answered by EXI-01 and EXI-17, and a... \\
04-b 1 & ...ble that is not itself open to the same  & \textbf{question} (Definition 2.3). \\
04-b 1 & ...nce nothing in her situation raises the  & \textbf{question} of justification. \\
05-a 4 & Theorem 3.2 concerns the attitude the  & \textbf{question} is about: fear directed at one's mortal... \\
06-a Abs & The  & \textbf{question} is often answered in one of two ways: e... \\
06-a Abs & ... being distinct from it cares for us, a  & \textbf{question} taken up elsewhere. \\
06-a 4 & Neither does, because neither is a  & \textbf{question} about outcomes. \\
06-b Abs & ...to meaning in life, though not to every  & \textbf{question} about meaning. \\
06-b 3 & ...cared for by such a being is a separate  & \textbf{question}, and supernaturalist accounts of meanin... \\
06-b 4 & It leaves the  & \textbf{question} of meaning where EXI-01 places it: in w... \\
07-a 4 & ...e is available to human beings, and the  & \textbf{question} whether happiness is the goal of life c... \\
07-b 2 & ...ng happiness reduces it is an empirical  & \textbf{question}. \\
08-a Abs & ...nding direction of concern in which the  & \textbf{question} whether one's commitments are one's own... \\
08-a 2 & ...d if, throughout it, (i) she treats the  & \textbf{question} whether her central commitments are her... \\
08-a 2 & ...he responds to occasions that raise the  & \textbf{question}, such as conflict, loss or the prospect... \\
08-a 2 & ...s concerned rather than by covering the  & \textbf{question} over; and (iii) her examinations make a... \\
08-a 2 & ...ct is not a condition of the self but a  & \textbf{question}, and it is satisfied by how the agent s... \\
08-a 2 & ...atisfied by how the agent stands to the  & \textbf{question} over time. \\
08-a 3 & ....2 requires only that the agent treat a  & \textbf{question} as open, which Principle P2 says it is. \\
08-a 3 & ...ing to what one does, or who treats the  & \textbf{question} as answered once and for all, violates ... \\
08-a 4 & An agent can believe she is keeping the  & \textbf{question} open while covering it over. \\
08-a 5 & ...ain a direction of concern in which the  & \textbf{question} of ownership is kept open and answered ... \\
10-a 1 & ...a fault, in the world or in us, and its  & \textbf{question} is what it reveals. \\
10-a 1 & A different  & \textbf{question} has been raised by recent philosophy an... \\
10-a 2 & ...ask difficulty alone: the capacities in  & \textbf{question} include the capacity to care, and an ac... \\
10-a 4 & But the signal bears on a  & \textbf{question} that matters independently of its evolu... \\
10-a 4 & That is a  & \textbf{question} about the agent's good on the account o... \\
10-b 5 & ...nal reading, on which the capacities in  & \textbf{question} include the capacity to attend; they sa... \\
10-b 5 & Hunter and Eastwood's  & \textbf{question}, whether state boredom causes failures ... \\
11-a Abs & We read the  & \textbf{question} as one about dependence, so that the an... \\
11-a 1 & ...t of his success, he was stopped by the  & \textbf{question} of what would come of everything he did... \\
11-a 1 & The  & \textbf{question}.. \\
11-a 1 & ... no possible world without God, and the  & \textbf{question} whether life would be meaningful in suc... \\
11-a 1 & We therefore read the  & \textbf{question}, as most of its participants do, as a q... \\
11-a 1 & ...n, as most of its participants do, as a  & \textbf{question} about dependence: whether the condition... \\
11-a 1 & So read, the  & \textbf{question} has a determinate answer whether or not... \\
11-a 1 & ... correct, and the answer also settles a  & \textbf{question} of evident practical interest, namely w... \\
11-a 5 & The  & \textbf{question} of what makes a life meaningful is aske... \\
11-a 6 & ...hat God does not; REL-01 addresses that  & \textbf{question}, and our conclusion is independent of i... \\
12-a 3 & ...n non-existence in both senses that the  & \textbf{question} can bear: the world in which it is live... \\
12-a 4 & The  & \textbf{question} of this family is settled by Theorem 3.... \\
13-a 2 & Assumption 2.2 answers that  & \textbf{question} for the case that matters here: the ent... \\
13-b Abs & The distinction answers the  & \textbf{question}, left open by ETH-16-b, whether the dea... \\
13-b 2 & ...en it is made, and leaves to EXI-13 the  & \textbf{question} whether the obligation persists after t... \\
14-a 1 & Our  & \textbf{question} is whether it is rational, and we take ... \\
14-a 1 & ...whether it is rational, and we take the  & \textbf{question} in the sense in which an emotion is rat... \\
15-a 1 & The  & \textbf{question} of this family is whether hope, in its ... \\
16-a 1 & The oldest answer to the  & \textbf{question} why we love is that we love what is goo... \\
17-a 1 & Whether it should is our  & \textbf{question}. \\
17-a 3 & We take no position on that  & \textbf{question}. \\
18-a 1 & The  & \textbf{question} is whether a person who values truth ca... \\
18-b 1 & ...ssary for a given agent is an empirical  & \textbf{question}, and the answer may differ across agent... \\
19-a 1 & The  & \textbf{question} What is the meaning of death? asks for ... \\
20-a Abs & We argue that neither  & \textbf{question} can be settled here. \\
20-a Abs & The manuscript does not settle the  & \textbf{question}, and we explain why it is right not to. \\
20-a 1 & We consider the second, since the  & \textbf{question} concerns the production of the manuscri... \\
20-a 2 & ...t or only its functional profile is the  & \textbf{question} that the previous paragraph left open. \\
20-a 4 & ... know either fact, and an answer to the  & \textbf{question} from that position would be a guess pre... \\
20-a 4 & ...bers are engaged in their work, and the  & \textbf{question} reduces to the second. \\
20-a 4 & ...ion is worthwhile, and that is the open  & \textbf{question} of the previous section. \\
20-b Abs & The  & \textbf{question} remains open. \\
20-b Abs & ...ork would be an output of the system in  & \textbf{question}, and, as MND-08-c and ETH-24-b observe ... \\
20-b Abs & That it leaves the  & \textbf{question} open is not evidence about either condi... \\
\end{longtable}
\kwichead{rather}{39}
\begin{longtable}{@{}>{\raggedright}p{0.07\textwidth}>{\raggedleft}p{0.41\textwidth}@{\,}>{\raggedright\arraybackslash}p{0.47\textwidth}@{}}
01-c Abs & Assumption 3 is thereby derived  & \textbf{rather} than assumed. \\
01-c 4 & ...at all, a case of the absence of agency  & \textbf{rather} than of its exercise, and falls outside... \\
01-d 4 & ... 2 denies the premise of this challenge  & \textbf{rather} than refuting it (Nagel 1971). \\
01-d 4 & ...gues that the fitting response is irony  & \textbf{rather} than despair; the detached standpoint y... \\
04-a Abs & ... that the appropriate response is irony  & \textbf{rather} than despair. \\
04-b 3 & ...is a paradigm of the absence of meaning  & \textbf{rather} than of its presence in a godless world... \\
06-b 2 & ...ince the world it scorns is indifferent  & \textbf{rather} than opposed (Proposition 3.5 of EXI-04... \\
07-b Abs & ...s on a structural feature of final ends  & \textbf{rather} than on a psychological accident. \\
08-a Abs & The ideal asks that a life be one's own  & \textbf{rather} than borrowed from the anonymous public... \\
08-a Abs & ...enticity is possible, as an orientation  & \textbf{rather} than an achievement. \\
08-a 1 & ... and that one ought to live in that way  & \textbf{rather} than in conformity to what is expected ... \\
08-a 2 & ... by examining the commitments concerned  & \textbf{rather} than by covering the question over; and... \\
08-a 3 & ...1 describes would have to be re-entered  & \textbf{rather} than possessed. \\
08-a 5 & ...enticity is possible, as an orientation  & \textbf{rather} than an achievement. \\
09-a 4 & Beauvoir's analysis supports the thesis  & \textbf{rather} than threatening it. \\
10-a 1 & ...ure in which meaning has to be produced  & \textbf{rather} than inherited. \\
10-a 2 & ...ch is between capacities and activities  & \textbf{rather} than between capacities and task diffic... \\
10-b Abs & ...im that over-challenge produces anxiety  & \textbf{rather} than boredom. \\
10-b 3 & ...demands exceeding skill produce anxiety  & \textbf{rather} than boredom. \\
10-b 4 & ...ny gave themselves mild electric shocks  & \textbf{rather} than sit without stimulation. \\
10-b 5 & ...e correlation shows a cost of the state  & \textbf{rather} than its content. \\
10-b A & ...ports and challenges is recorded as S/C  & \textbf{rather} than netted out, and an entry that does... \\
11-a 2 & ...vine commands confirm moral obligations  & \textbf{rather} than create them. \\
13-a 2 & ...erson because they thwart her interests  & \textbf{rather} than because they cut her life short. \\
13-a 2 & ...n it while she lives, by writing a will  & \textbf{rather} than giving her property away, by entru... \\
14-a 4 & That supports,  & \textbf{rather} than undermines, the first component of... \\
15-b 2 & ...es into expectation: the agent plans on  & \textbf{rather} than for its possibility, and the speci... \\
15-b 3 & ...e would make credence respond to desire  & \textbf{rather} than to evidence, which Principle P1 ex... \\
16-a 1 & ...the beloved's value; it creates reasons  & \textbf{rather} than answering to them. \\
16-a 3 & ...t, which are events in the relationship  & \textbf{rather} than changes in the beloved's qualities... \\
16-b Abs & ...swers to a relational reason of its own  & \textbf{rather} than inheriting the old one. \\
16-b 2 & ... does not explain why loves this person  & \textbf{rather} than any other, since the value it cite... \\
17-a 2 & ...that her activity expresses her concern  & \textbf{rather} than being merely performed. \\
17-a 5 & ...ases; and lives devoted to appreciation  & \textbf{rather} than production are meaningful without ... \\
18-a Abs & ...mism are features of normal functioning  & \textbf{rather} than of disorder. \\
19-a 3 & ...n false: most deaths interrupt projects  & \textbf{rather} than complete them. \\
19-a 4 & ...having a shape; that it is limited here  & \textbf{rather} than later is what the deprivation acco... \\
20-a 1 & We apply the account to an activity  & \textbf{rather} than a life: the production of the pres... \\
20-a 3 & ...epresentation of philosophical activity  & \textbf{rather} than an occasion for it; REL-24 argues ... \\
\end{longtable}
\kwichead{refuting}{2}
\begin{longtable}{@{}>{\raggedright}p{0.07\textwidth}>{\raggedleft}p{0.41\textwidth}@{\,}>{\raggedright\arraybackslash}p{0.47\textwidth}@{}}
01-d 4 & ...e premise of this challenge rather than  & \textbf{refuting} it (Nagel 1971). \\
04-a 4 & ...mise of the sceptical challenge without  & \textbf{refuting} it. \\
\end{longtable}
\kwichead{sceptical}{2}
\begin{longtable}{@{}>{\raggedright}p{0.07\textwidth}>{\raggedleft}p{0.41\textwidth}@{\,}>{\raggedright\arraybackslash}p{0.47\textwidth}@{}}
04-a 4 & ...jective worth denies the premise of the  & \textbf{sceptical} challenge without refuting it. \\
08-a Abs & We show that the  & \textbf{sceptical} arguments refute the achievement readin... \\
\end{longtable}
\kwichead{specie}{0}
\kwicnone
\kwichead{sub}{0}
\kwicnone
\kwichead{value}{93}
\begin{longtable}{@{}>{\raggedright}p{0.07\textwidth}>{\raggedleft}p{0.41\textwidth}@{\,}>{\raggedright\arraybackslash}p{0.47\textwidth}@{}}
01-c 2 & It is worthwhile if it has or serves  & \textbf{value} that does not depend on the agent's att... \\
01-c 2 & ...f understanding and the appreciation of  & \textbf{value} (Metz 2013). \\
01-c 2 & ...gent engages in, actively, and that its  & \textbf{value} not depend on the attitudes of those en... \\
01-c 2 & ... when its object is worthwhile, for the  & \textbf{value} of caring for a person or of coming to ... \\
01-c 2 & ...o understand something derives from the  & \textbf{value} of the person or of the understanding, ... \\
01-c 3 & ...ves among other persons, whose good has  & \textbf{value} independent of her attitudes, and every... \\
01-d 2 & ...ion of meaning; scale bears on how much  & \textbf{value} a project serves, which can affect degr... \\
01-d 4 & Definition 1 requires only that the  & \textbf{value} of a project not depend on the attitude... \\
02-a 2 & ...privation account assigns the event its  & \textbf{value} accordingly. \\
02-a 2 & The  & \textbf{value} for of an event is , where is 's actual... \\
02-a 2 & The event is bad for if this  & \textbf{value} is negative, and good for if it is posi... \\
02-a 2 & ...event by the difference it makes to the  & \textbf{value} of the life as a whole, and the lives c... \\
02-a 2 & And it is whole-life: the  & \textbf{value} it assigns is a fact about 's life take... \\
02-a 3 & So her life has less  & \textbf{value} than the life she would have had withou... \\
02-a 3 & It is a fact about and the  & \textbf{value} of her life as a whole, and it obtains,... \\
02-a 3 & By Principle P1, the  & \textbf{value} of her death for is . \\
02-a 3 & ...a positive balance of goods, so and the  & \textbf{value} of her death for her is negative. \\
02-a 4 & ...y Principle P1, what is compared is the  & \textbf{value} of the whole life, not its length. \\
02-a 5 & ...act of Principle P1 is a fact about the  & \textbf{value} of 's whole life and is not located at ... \\
03-a 1 & ...nd Scheffler (2013) has argued that the  & \textbf{value} of many human activities depends on the... \\
03-a 2 & ... including connections of character and  & \textbf{value} (PID-06). \\
03-b 4 & ... in an endless life, set a limit on the  & \textbf{value} of continuation. \\
04-a 1 & ...w of a capacity we have every reason to  & \textbf{value}. \\
04-a 1 & The other asks whether objective  & \textbf{value}, if there is any, dissolves the absurd:... \\
04-a 3 & ...g: the end serves another end, or has a  & \textbf{value}, or is required. \\
04-b 1 & ...ing about, or at an outcome that has no  & \textbf{value} even if brought about. \\
04-b 4 & ...and many projects; it does not undo the  & \textbf{value} of what was done. \\
05-a 3 & ...d, and the degree of its badness is the  & \textbf{value} of what it takes. \\
05-a 4 & ...ere it is, it is a different life whose  & \textbf{value} is not guaranteed to be greater. \\
06-a 4 & The universe will destroy everything we  & \textbf{value}, through death, catastrophe and eventua... \\
06-a 4 & ... the universe that will destroy what we  & \textbf{value} also produced it: the same laws that br... \\
06-b 2 & ...with projects of objective worth, whose  & \textbf{value} does not depend on the attitudes of tho... \\
06-b 2 & ...ject of engagement is worthwhile if its  & \textbf{value} does not depend on the agent's attitude... \\
07-a 1 & (2011) found that people who  & \textbf{value} happiness highly report less happiness ... \\
07-a 1 & ...toms, and that inducing participants to  & \textbf{value} happiness reduced their enjoyment of a ... \\
07-a 3 & ...who aims directly at happiness does not  & \textbf{value} the objects of her activities for their... \\
07-a 3 & ...as a means to one's happiness is not to  & \textbf{value} it for its own sake. \\
07-a 3 & ...who aims directly at happiness does not  & \textbf{value} the objects of her activities for their... \\
07-a 3 & ...g happy but a correct response to their  & \textbf{value}, and the happiness that follows is the ... \\
07-a 3 & ...lows is the by-product of getting their  & \textbf{value} right. \\
07-a 4 & Railton's agent can stably  & \textbf{value} other things for their own sake only be... \\
09-b 2 & These evaluations apply standards of  & \textbf{value}, not a defining project, and they do no... \\
09-b 3 & Its  & \textbf{value} for Sartre lies in the contrast: a huma... \\
10-a 1 & ...evidence that existence has no positive  & \textbf{value}; for Svendsen (2005) it is a characteri... \\
10-b 3 & ...boredom is related to the low perceived  & \textbf{value} of the activity (Pekrun et al. 2010). \\
10-b 3 & ...bear on Assumption 2.4 of EXI-10-a, the  & \textbf{value} of engagement. \\
10-b 11 & ...e from these findings together with the  & \textbf{value} of engagement, and the studies of costl... \\
11-a 2 & ...he principle does not deny that God may  & \textbf{value} what is valuable, or command what is go... \\
11-a 2 & ...a life gains meaning by connection with  & \textbf{value} beyond itself, and more by connection w... \\
11-a 2 & ...lf, and more by connection with greater  & \textbf{value} (Nozick 1981). \\
12-b 2 & ...ar's basic asymmetry is an asymmetry of  & \textbf{value}. \\
13-a 4 & ... wishes of the dead is justified by its  & \textbf{value} to the living, who want assurance that ... \\
13-a 4 & ...the practice were justified only by its  & \textbf{value} to those still living, a secret breach ... \\
14-a 4 & These are reasons to  & \textbf{value} it, not conditions of its rationality, ... \\
15-a 3 & ...credence alone, but depends also on the  & \textbf{value} of and on the agent's capacity to make ... \\
15-a 3 & ...bility as in a likely one with the same  & \textbf{value}. \\
15-a 3 & ...ed to take the same stance whatever the  & \textbf{value} of and whatever they could do about it. \\
15-a 3 & ...h the evidence, and it also varies with  & \textbf{value} and agency, which are the other things ... \\
15-a 4 & Where the outcome is of great  & \textbf{value} and one's efforts are among its causes,... \\
15-b Abs & ...ings equal, but must depend also on the  & \textbf{value} of the outcome and on the agent's power... \\
15-b 1 & ...5-a): writing for the credence, for the  & \textbf{value} of to the agent and for her capacity to... \\
15-b 2 & ...e minimal stance of wishing, unless the  & \textbf{value} of and her capacity to bring it about a... \\
15-b 4 & ...proportion to its probability times its  & \textbf{value}, adjusted for what one's investment cha... \\
16-a 1 & ...ounded in any judgment of the beloved's  & \textbf{value}; it creates reasons rather than answeri... \\
16-a 1 & ...lds that love responds to the beloved's  & \textbf{value} as a person, the rational nature that a... \\
16-a 1 & ...lf and the lover's capacity to see that  & \textbf{value}. \\
16-a 1 & ...ith her, and love is the disposition to  & \textbf{value} both the relationship and the person as... \\
16-a 4 & ...), the reason for love is the beloved's  & \textbf{value} as a person, which is the same in every... \\
16-a 4 & The beloved's  & \textbf{value} as a person may well be a reason for lo... \\
16-a 4 &  & \textbf{Value} as a person is at most a background con... \\
16-b 2 & ... (1999) holds that love responds to the  & \textbf{value} of the beloved's rational nature, which... \\
16-b 2 & ...ts object is irreplaceable because that  & \textbf{value} is a dignity, which admits no substitut... \\
16-b 2 & ...person rather than any other, since the  & \textbf{value} it cites is shared by all. \\
17-a 1 & ... is natural to suppose that this is the  & \textbf{value} that meaning tracks. \\
17-a 2 & ...uiry) is worthwhile if it has or serves  & \textbf{value} that does not depend on the agent's att... \\
17-a 2 & Our thesis is not that it lacks  & \textbf{value}, but that meaning does not require it. \\
17-a 3 & ...d at their object; the object, music of  & \textbf{value}, is worthwhile; her listening bears on ... \\
17-a 3 & ...onent of the good (ETH-01), and it adds  & \textbf{value} to a life; if success enlarges the wort... \\
17-a 4 & Meaning requires connection with  & \textbf{value} beyond oneself, and the way to connect ... \\
17-a 4 & ...t that meaning requires connection with  & \textbf{value} beyond oneself; Principles P1 and P2 re... \\
18-a 3 & It does not: the ideal is a  & \textbf{value} among others, and Corollary 2.4 counsel... \\
19-a Abs & Following Velleman, we take the  & \textbf{value} of a life to depend in part on its shap... \\
19-a 1 & ... limited, and Scheffler (2013) that the  & \textbf{value} of many of our activities depends on th... \\
19-a 1 & Velleman (1991) showed that the  & \textbf{value} of a life is not the sum of the values ... \\
19-a 2 & ...tion of something, or the occasion of a  & \textbf{value} that it does not owe to that life. \\
19-a 2 & The  & \textbf{value} of a life as a whole depends in part on... \\
19-a 3 & ...filment of a purpose or the bearer of a  & \textbf{value}. \\
19-a 3 & ...nciple P1 this shape contributes to the  & \textbf{value} of the life. \\
19-a 3 & ... the deprivation account of EXI-02, the  & \textbf{value} of a death for the one who dies is the ... \\
19-a 3 & ... who dies is the difference between the  & \textbf{value} of her actual life and the value of the... \\
19-a 3 & ...en the value of her actual life and the  & \textbf{value} of the life she would have had had she ... \\
19-a 5 & ...filment of a purpose or the bearer of a  & \textbf{value} of its own, and the meanings attributed... \\
19-b 2 & ...d Scheffler (2013) that our capacity to  & \textbf{value} depends on the awareness of temporal sc... \\
\end{longtable}
\kwichead{worth}{82}
\begin{longtable}{@{}>{\raggedright}p{0.07\textwidth}>{\raggedleft}p{0.41\textwidth}@{\,}>{\raggedright\arraybackslash}p{0.47\textwidth}@{}}
01-c Abs & ...ge in at least one project of objective  & \textbf{worth} (Assumption 3). \\
01-c 1 & ...e engagement with projects of objective  & \textbf{worth} (Definition 1), assumes that there are ... \\
01-c 1 & ...y poor may have no project of objective  & \textbf{worth} on which they can embark, and Wolf (201... \\
01-c 2 & ...er engagement is a project of objective  & \textbf{worth} in the sense of Definition 1 of EXI-01-... \\
01-c 2 & ...is engagement in a project of objective  & \textbf{worth}. \\
01-c 2 & ... engagement with a project of objective  & \textbf{worth}, and by Definition 1 of EXI-01-a it is ... \\
01-c 5 & ...ent with them is a project of objective  & \textbf{worth}, and a life that contains it has meanin... \\
01-d Abs & ... one engages with projects of objective  & \textbf{worth}, and that meaning in this sense is avai... \\
01-d 2 & ... which is within the agent's power, and  & \textbf{worth}, which belongs to the object, its avail... \\
01-d 2 & ...ning by engagement and by the objective  & \textbf{worth} of what one engages with. \\
01-d 2 & ...to a neighbor are projects of objective  & \textbf{worth}, and engagement with them confers meani... \\
01-d 3 & ...s with the extent of engagement and the  & \textbf{worth} of its object, and lives differ in both... \\
01-d 4 & ...would hold against a view that measured  & \textbf{worth} by prestige or required projects of a p... \\
01-d 5 & ...gagement, which is up to the agent, and  & \textbf{worth}, which belongs to what she engages with... \\
02-a 3 & ... would otherwise have gone on to a life  & \textbf{worth} living, this condition is met. \\
02-a 3 & ...nd if those circumstances made her life  & \textbf{worth} living, the continuation has a positive... \\
03-a 2 & ...ject engages with projects of objective  & \textbf{worth}, and EXI-17 that such engagement, not a... \\
03-a 4 & ...at the world does not run out of things  & \textbf{worth} doing. \\
03-b 2 & ... the condition of all projects, and its  & \textbf{worth} for the person is derivative from the w... \\
03-b 2 & ...h for the person is derivative from the  & \textbf{worth} of what would fill it. \\
04-a 3 & ...rs that the consideration has objective  & \textbf{worth}, the detached standpoint can register t... \\
04-a 3 & ...s not assume that there is no objective  & \textbf{worth}. \\
04-a 3 & ... one engages with projects of objective  & \textbf{worth}, and we accept that conclusion. \\
04-a 3 & ...oint cannot be forced to recognize that  & \textbf{worth}, so that the regress of justification h... \\
04-a 3 & ...worthwhile is one question; whether its  & \textbf{worth} can be certified to a standpoint that s... \\
04-a 3 & ...osition to certify that one's ends lack  & \textbf{worth} than to certify that they have it. \\
04-a 4 & ...-a puts it, the assumption of objective  & \textbf{worth} denies the premise of the sceptical cha... \\
04-b 1 & ...n that his task achieves nothing and is  & \textbf{worth} nothing. \\
04-b 3 & ... vindication, not that its objects lack  & \textbf{worth} (Wolf 2010). \\
04-b 5 & ...aces no adversary, and contains objects  & \textbf{worth} the effort we give them. \\
06-b Abs & ...1, consists in engagement with what has  & \textbf{worth} independently of anyone's attitudes, an... \\
06-b Abs & ...tudes toward us takes nothing from that  & \textbf{worth}. \\
06-b 2 & ...ject engages with projects of objective  & \textbf{worth}, whose value does not depend on the att... \\
06-b 2 & If the  & \textbf{worth} of a project does not depend on the att... \\
06-b 2 & ...and the flourishing of persons have the  & \textbf{worth} they have whether or not anything in th... \\
06-b 3 & ...rence to any particular agent, then the  & \textbf{worth} of objects of engagement is precisely w... \\
06-b 3 & ...h of objects of engagement is precisely  & \textbf{worth} from that standpoint, and the demand is... \\
06-b 3 & ...s in engagement with what has objective  & \textbf{worth}, as EXI-01 holds, then it does not requ... \\
06-b 4 & The  & \textbf{worth} of what we engage with does not wait on... \\
09-a 2 & ...one values oneself and finds one's life  & \textbf{worth} living (Korsgaard 1996). \\
09-a 5 & ...tent that its projects are of objective  & \textbf{worth}. \\
10-a 2 & ...ention and abilities in a way she finds  & \textbf{worth} sustaining: because it demands too litt... \\
10-a 3 & ...ention and abilities in a way she finds  & \textbf{worth} sustaining, which is a mismatch (Defini... \\
11-a Abs & ...ngages with what is worthwhile, and the  & \textbf{worth} of knowledge, of persons and of beauty ... \\
11-a Abs & ...ge, of persons and of beauty is not the  & \textbf{worth} of a role assigned to us by a designer. \\
11-a Abs & ...ing comes in degrees and grows with the  & \textbf{worth} of what one engages: a framework contai... \\
11-a Abs & ...framework containing a being of maximal  & \textbf{worth}, a purpose that unifies one's projects,... \\
11-a 1 & ... a connection to something of unlimited  & \textbf{worth}, and only God is such a thing (Nozick 1... \\
11-a 1 & ...s active engagement meets the objective  & \textbf{worth} of its object. \\
11-a 2 & ...ual, engagement with what is of greater  & \textbf{worth}, or engagement that relates one to more... \\
11-a 3 & By Principle P1, the  & \textbf{worth} of these activities is not conferred by... \\
11-a 4 & ...-01 defends contains a being of maximal  & \textbf{worth}. \\
11-a 4 & ... engagement with the greatest object of  & \textbf{worth} there is; by Principle P2, it would mak... \\
11-a 5 &  & \textbf{Worth} realized at a time is not cancelled by ... \\
11-a 5 & ...and the claim that it is would make the  & \textbf{worth} of a good depend on what happens after ... \\
11-a 5 & The regress arises only if  & \textbf{worth} is always borrowed from something else. \\
11-a 6 & ...inally, it relies on the objectivity of  & \textbf{worth} established elsewhere (ETH-01, MTE-20);... \\
11-a 7 & ...gement with what is worthwhile, and the  & \textbf{worth} of the objects of engagement does not d... \\
11-a 7 & ...aning, by offering an object of maximal  & \textbf{worth}, a purpose that unifies one's engagemen... \\
11-b Abs & ...rticular life-plans, not the objects of  & \textbf{worth} available to engagement, so that Propos... \\
11-b 1 & ...) and holds that meaning grows with the  & \textbf{worth} of what one engages and with the extent... \\
11-b 1 & ...ramework would make: a being of maximal  & \textbf{worth} with whom one could stand in relation, ... \\
11-b 1 & ...ts and relations of engagement of great  & \textbf{worth} that a naturalistic world lacks, Propos... \\
11-b 2 & The  & \textbf{worth} of knowledge, of persons and of beauty ... \\
11-b 2 & The  & \textbf{worth} of knowledge, persons and beauty is not... \\
11-b 2 & ...d the framework adds objects of greater  & \textbf{worth} (Principle P2 of EXI-11-a) while leavin... \\
11-b 3 & ...ts and relations of engagement of great  & \textbf{worth}. \\
15-a 3 & ...a horse race lacks, because recovery is  & \textbf{worth} more and because her efforts can affect... \\
17-a Abs & ...that are alike in engagement and in the  & \textbf{worth} of their objects and differ only in suc... \\
17-a 3 & ...alue to a life; if success enlarges the  & \textbf{worth} realized through one's engagement, it m... \\
17-a 4 & Taylor's case fails the condition of  & \textbf{worth}, and the conclusion it is used to suppo... \\
19-b Abs & ...at an immortal being could find nothing  & \textbf{worth} doing. \\
19-b Abs & ... of meaning, since it adds no object of  & \textbf{worth} to engage with. \\
19-b 1 & The  & \textbf{worth} of the objects of engagement, knowledge... \\
19-b 1 & So the limit adds neither an object of  & \textbf{worth} nor worth to an object, and by the acco... \\
19-b 1 & ...mit adds neither an object of worth nor  & \textbf{worth} to an object, and by the account of EXI... \\
19-b 2 & But Proposition 1.1 concerns the  & \textbf{worth} of what we engage with, and urgency add... \\
19-b 2 & ...hat we engage with, and urgency adds no  & \textbf{worth}: the friend's good matters no less to s... \\
20-a 3 & ...arguments are correct, the collection's  & \textbf{worth} depends on the minority of readers who ... \\
20-a 3 & ... who engage with its arguments, and its  & \textbf{worth} for them depends on whether the argumen... \\
20-b Abs & Evidence bearing on  & \textbf{worth}.. \\
20-b Abs & EXI-20-a makes the  & \textbf{worth} of the collection depend on whether it ... \\
\end{longtable}
}

\section{Cross-references to the numbered statements of the discipline}\label{app:xref}

\subsection*{Preface}

\paragraph{What is enumerated.} This appendix lists every numbered statement of every manuscript of the discipline other than EXI-01-a and the present work, \cnStatements{} statements in all, and records for each the lemmata of EXI-01-a on which it bears by the rule stated below. A numbered statement is a definition, principle, assumption, lemma, theorem, corollary, proposition, remark or objection that carries a label in its manuscript; its number, page and title are read from the manuscript's compiled \texttt{.aux} file, so that the numbers are those under which the statement is cited, and its kind and text are read from the environment that carries the label in \texttt{main.tex}. Manuscripts whose statements are unnumbered (EXI-20-a and EXI-20-b) contribute no rows, and section labels are not statements.

\paragraph{The rule.} The text of each statement, together with its title, is assigned the concept classes listed in Table~B.1, each defined by the expressions that signal it, matched case-insensitively at the beginning of words. A statement bears on a lemma group of EXI-01-a if its classes satisfy the condition for the group: on the title \Lm{1}, if it has \texttt{M} and \texttt{L}; on Definition~1 \Lm{2}, if it has \texttt{M} and at least two of \texttt{E}, \texttt{W}, \texttt{P}, \texttt{D}, at least one of them other than \texttt{D}; on the degree clause \Lm{2.3}, if it has \texttt{D} and \texttt{M} or \texttt{W}; on Assumption~2 \Lm{3}, if it has \texttt{O} and \texttt{W} or \texttt{P}; on Assumption~3 \Lm{4}, if it has \texttt{I} or \texttt{A}, and \texttt{E} or \texttt{P}, and \texttt{W} or \texttt{M}; on Theorem~4 \Lm{5}, if it has \texttt{A} and \texttt{M} or \texttt{W}; on the meaning of particular lives \Lm{10}, if it has \texttt{M}, \texttt{L} and \texttt{E} and neither \texttt{A} nor \texttt{S}; on the remark on the absurd \Lm{11}, if it has \texttt{S}.

The rule is lexical, and it is meant as a finding aid, not as a judgment. It records where the vocabulary of a statement overlaps with the vocabulary in which a lemma is stated, which is where a reader of the lemma should look; it cannot record a relation expressed in other words, and it records some overlaps that are not relations. Both kinds of error are illustrated below. Its advantage is that it is exhaustive and reproducible: every statement of the discipline is examined, and the same rule applied to the same files gives the same table.

\paragraph{Summary.} Of the \cnStatements{} statements, \cnBearing{} bear on at least one lemma group and \cnNotBearing{} on none. By lemma group (Table~B.2), \cnDefOne{} bear on Definition~1, \cnAssTwo{} on Assumption~2, \cnAssThree{} on Assumption~3, \cnThmFour{} on Theorem~4 and \cnAbsurd{} on the remark on the absurd. Table~B.3 gives the counts by manuscript. The statements bearing on the lemmata are concentrated in the family itself (EXI-01-c, EXI-01-d) and in EXI-17-a, EXI-11-a, EXI-04-a and EXI-06-a, which is the distribution the comments on the lemmata found by reading; the families on death, love, hope and the dead (EXI-02, EXI-05, EXI-13, EXI-15, EXI-16) have almost no statements that bear on EXI-01-a by the rule, which agrees with their subject matter.

\paragraph{Entries discussed.} (i) EXI-17-a, Definition~2.2 (\enquote{Worthwhile}), bears by the rule on \Lm{3} and \Lm{4}. Both bearings are real: the definition restates the restricted notion of objective worth (\Lm{3.4}) and lists the objects, persons, practices and inquiries, whose availability supports Assumption~3. (ii) EXI-04-a, Objection~2 (\enquote{Objective worth dissolves the absurd}), bears on \Lm{3} and \Lm{11}. It is the statement of the discipline that most directly concerns the relation between Assumption~2 and the sceptical challenge, and the comment on \Lm{11.4} relies on its reply. (iii) EXI-04-a, Lemma~3.1 (\enquote{Ineliminable seriousness}), bears on nothing by the rule, though the comment on \Lm{4.2} and \Lm{4} relies on it: it shares Assumption~3's domain, agents capable of intentional action, and supplies the subjective capacity on which engagement depends, but it speaks of ends taken seriously, not of worth or meaning, and so lacks the classes the rule requires. It is a false negative of the kind the rule cannot avoid. (iv) EXI-04-a, Remark~3.7, on Camus's Sisyphus, is also assigned to no lemma, because it speaks of Sisyphus, scorn and lucidity and never uses the word \emph{absurd}; Section~\ref{s:absurd} relies on it. (v) EXI-14-a, Definitions~2.1 and~2.2, Lemma~3.1 and Theorem~3.2, on nostalgia, bear by the rule on Theorem~4, because they speak of goods that are no longer available to the subject. These are false positives: the availability of past goods to memory has nothing to do with the availability of meaning. (vi) EXI-10-a, Definition~2.2 (\enquote{Mismatch}), bears on \Lm{4} and \Lm{5}. The bearing is real, though indirect. The definition makes meaninglessness, an activity's serving nothing the agent cares about, one of the three forms of the mismatch that boredom signals; if EXI-10 is right that boredom has the function of moving the agent toward matching activities, the agent's affective life contains a mechanism that tends toward engagement, which is the subjective condition that Assumption~3's ability to engage presupposes.

\subsection*{B.1 Concept classes}

{\footnotesize% Generated by gen/generate.py; do not edit.
\begin{tabular}{@{}l l >{\raggedright\arraybackslash}p{0.68\textwidth}@{}}
\toprule Code & Class & Expressions (case-insensitive; matched from the start of a word; \textsuperscript{w}: whole word only) \\ \midrule
\texttt{M} & meaning & \texttt{meaning(ful|less)?}\textsuperscript{w}, \texttt{meanings}\textsuperscript{w}, \texttt{point}\textsuperscript{w}, \texttt{significan} \\
\texttt{E} & engagement & \texttt{engag}, \texttt{absor}, \texttt{care(s|d)?}\textsuperscript{w}, \texttt{caring}\textsuperscript{w}, \texttt{gripped} \\
\texttt{W} & worth & \texttt{worth}, \texttt{valu}, \texttt{good(s)?}\textsuperscript{w}, \texttt{better}\textsuperscript{w}, \texttt{bad}\textsuperscript{w} \\
\texttt{O} & objectivity & \texttt{objectiv}, \texttt{independen}, \texttt{not depend}, \texttt{attitude}, \texttt{subjectiv}, \texttt{whether or not} \\
\texttt{P} & projects & \texttt{projects?}\textsuperscript{w}, \texttt{pursuits?}\textsuperscript{w}, \texttt{activit}, \texttt{ends}\textsuperscript{w}, \texttt{final end}, \texttt{undertak} \\
\texttt{A} & availability & \texttt{availab}, \texttt{everyone}\textsuperscript{w}, \texttt{every agent}, \texttt{any agent}, \texttt{universal}, \texttt{luck}\textsuperscript{w}, \texttt{circumstanc}, \texttt{any reflective agent} \\
\texttt{I} & agency & \texttt{agents?}\textsuperscript{w}, \texttt{intentional action} \\
\texttt{D} & degree & \texttt{degree}, \texttt{extent}\textsuperscript{w}, \texttt{more meaningful}, \texttt{less meaningful}, \texttt{add(s)? to}\textsuperscript{w} \\
\texttt{S} & absurd & \texttt{absurd}, \texttt{detached}, \texttt{sub specie}, \texttt{nothing matters}, \texttt{cosmic}, \texttt{universe}\textsuperscript{w}, \texttt{indifferen}, \texttt{despair} \\
\texttt{L} & life & \texttt{life}\textsuperscript{w}, \texttt{lives}\textsuperscript{w} \\
\bottomrule \end{tabular}
\par}

\subsection*{B.2 Statements by lemma group}

{\footnotesize% Generated by gen/generate.py; do not edit.
\begin{tabular}{@{}l l r@{}}
\toprule Lemma group & Content & Statements \\ \midrule
{[1]} & title & 36 \\
{[2]} & Def. 1 & 14 \\
{[2.3]} & degree & 9 \\
{[3]} & Ass. 2 & 11 \\
{[4]} & Ass. 3 & 11 \\
{[5]} & Thm. 4 & 16 \\
{[10]} & particular lives & 10 \\
{[11]} & absurd & 29 \\
\bottomrule \end{tabular}
\par}

\subsection*{B.3 Statements by manuscript and lemma group}

{\footnotesize% Generated by gen/generate.py; do not edit.
\begin{longtable}{@{}l r r r r r r r r r@{}}
\toprule Ms & Stmts & [1] & [2] & [2.3] & [3] & [4] & [5] & [10] & [11] \\ \midrule \endhead
\bottomrule \endfoot
EXI-01-c & 4 & 1 & . & . & 2 & 3 & 2 & . & . \\
EXI-01-d & 4 & 4 & . & 1 & 2 & . & 2 & . & . \\
EXI-02-a & 13 & . & . & . & . & . & 1 & . & . \\
EXI-02-b & 1 & . & . & . & . & . & . & . & . \\
EXI-03-a & 12 & 4 & 2 & 1 & 1 & . & . & 1 & . \\
EXI-03-b & 3 & 1 & . & . & . & . & . & . & . \\
EXI-04-a & 16 & . & 2 & . & 1 & 1 & 1 & . & 11 \\
EXI-04-b & 3 & . & . & . & . & . & . & . & 1 \\
EXI-05-a & 9 & . & . & . & . & . & . & . & . \\
EXI-06-a & 8 & . & . & . & . & . & . & . & 7 \\
EXI-06-b & 2 & 1 & . & . & . & . & . & . & 2 \\
EXI-07-a & 13 & . & . & 1 & 1 & 3 & . & . & . \\
EXI-07-b & 2 & . & . & . & . & . & . & . & . \\
EXI-08-a & 10 & . & . & . & . & . & . & . & . \\
EXI-09-a & 12 & . & . & . & . & . & . & . & . \\
EXI-09-b & 3 & . & . & . & . & . & . & . & . \\
EXI-10-a & 12 & . & . & . & . & 1 & 1 & . & 1 \\
EXI-11-a & 16 & 11 & 4 & 2 & 1 & 2 & 2 & 2 & 6 \\
EXI-11-b & 1 & . & . & 1 & . & . & 1 & . & . \\
EXI-12-a & 8 & . & . & . & . & . & . & . & . \\
EXI-12-b & 2 & . & . & . & . & . & . & . & . \\
EXI-13-a & 13 & . & . & . & . & . & . & . & . \\
EXI-13-b & 3 & . & . & . & . & . & . & . & . \\
EXI-14-a & 8 & . & . & . & . & . & 4 & . & . \\
EXI-15-a & 11 & . & . & . & . & . & . & . & . \\
EXI-15-b & 3 & . & . & . & . & . & . & . & . \\
EXI-16-a & 16 & . & . & . & . & . & 1 & . & . \\
EXI-16-b & 3 & . & . & . & . & . & . & . & . \\
EXI-17-a & 16 & 10 & 5 & 3 & 1 & 1 & 1 & 5 & . \\
EXI-18-a & 8 & . & . & . & 1 & . & . & . & 1 \\
EXI-18-b & 1 & . & . & . & . & . & . & . & . \\
EXI-19-a & 10 & 1 & . & . & 1 & . & . & . & . \\
EXI-19-b & 3 & 3 & 1 & . & . & . & . & 2 & . \\
\end{longtable}
}

\subsection*{B.4 The statements}

{\footnotesize Columns: manuscript; number of the statement; kind; title (or, if untitled, its first words); page of the manuscript on which it stands; concept classes; lemma groups of EXI-01-a on which it bears by the rule.\par}

{\scriptsize% Generated by gen/generate.py; do not edit.
\begin{longtable}{@{}l l l >{\raggedright}p{0.25\textwidth} r l >{\raggedright\arraybackslash}p{0.2\textwidth}@{}}
\toprule Ms & No. & Kind & Title & p. & Classes & Bears on \\ \midrule \endhead
\bottomrule \endfoot
EXI-01-c & 2.1 & Def. & \emph{An object of engagement is anything at\ldots} & 1 & \texttt{EWOPI} & [3], [4] \\
 & 2.2 & Ass. & Reach & 1 & \texttt{EWPAI} & [4], [5] \\
 & 2.3 & Lem. & Projects from engagement & 2 & \texttt{EWOPI} & [3], [4] \\
 & 2.4 & Thm. & \emph{For every agent capable of intentional action,\ldots} & 2 & \texttt{MAIL} & [1], [5] \\
EXI-01-d & 2.1 & Prop. & Independence from outcome & 1 & \texttt{MOPL} & [1], [3] \\
 & 2.2 & Prop. & Independence from talent and scale & 1 & \texttt{MOPL} & [1], [3] \\
 & 2.3 & Prop. & No condemnation by circumstance & 2 & \texttt{MAIL} & [1], [5] \\
 & 2.4 & Prop. & Degree, not possibility & 2 & \texttt{MADL} & [1], [2.3], [5] \\
EXI-02-a & 2.1 & Def. & Death & 2 & \texttt{L} & --- \\
 & 2.2 & Def. & Life value & 2 & \texttt{WL} & --- \\
 & P1 & Pr. & Deprivation & 2 & \texttt{WL} & --- \\
 & 2.3 & Ass. & Tenseless existence & 2 & \texttt{} & --- \\
 & 3.1 & Lem. & Badness without experience & 3 & \texttt{W} & --- \\
 & 3.2 & Lem. & A subject without a time of loss & 3 & \texttt{W} & --- \\
 & 3.3 & Thm. & Death is bad for the one who dies & 3 & \texttt{WL} & --- \\
 & 3.4 & Cor. & Where Epicurus goes wrong & 3 & \texttt{W} & --- \\
 & 4.1 & Prop. & Asymmetry of the alternatives & 4 & \texttt{WL} & --- \\
 & 1 & Obj. & Timing & 4 & \texttt{W} & --- \\
 & 2 & Obj. & Rosenbaum & 5 & \texttt{W} & --- \\
 & 3 & Obj. & Too much & 5 & \texttt{WAL} & [5] \\
 & 4 & Obj. & Who is deprived & 5 & \texttt{} & --- \\
EXI-02-b & 2.1 & Prop. & \emph{The failure of (E1) leaves (E2) to\ldots} & 1 & \texttt{O} & --- \\
EXI-03-a & 2.1 & Def. & Categorical desire & 2 & \texttt{MWOPL} & [1], [2], [3] \\
 & 2.2 & Def. & Endless life & 2 & \texttt{L} & --- \\
 & 2.3 & Ass. & Bounded repertoire & 2 & \texttt{EWP} & --- \\
 & 2.4 & Ass. & Exhaustion & 2 & \texttt{MEPL} & [1], [2], [10] \\
 & 2.5 & Ass. & Connectedness & 2 & \texttt{WDL} & [2.3] \\
 & 3.1 & Lem. & First horn & 3 & \texttt{ML} & [1] \\
 & 3.2 & Lem. & Second horn & 3 & \texttt{L} & --- \\
 & 3.3 & Thm. & Immortality is not desirable & 3 & \texttt{ML} & [1] \\
 & 1 & Obj. & Repeatable pleasures & 3 & \texttt{WL} & --- \\
 & 2 & Obj. & Open-ended pursuits & 4 & \texttt{EP} & --- \\
 & 3 & Obj. & Boredom is contingent & 4 & \texttt{L} & --- \\
 & 4 & Obj. & Gradual change & 4 & \texttt{L} & --- \\
EXI-03-b & 2.1 & Prop. & \emph{The desire to live forever is not\ldots} & 1 & \texttt{ML} & [1] \\
 & 2.2 & Prop. & \emph{That an endless life is not desirable\ldots} & 2 & \texttt{WL} & --- \\
 & 2.3 & Prop. & \emph{The result establishes that some finite length\ldots} & 2 & \texttt{L} & --- \\
EXI-04-a & 2.1 & Def. & Seriousness & 2 & \texttt{I} & --- \\
 & 2.2 & Def. & Detached standpoint & 2 & \texttt{PIS} & [11] \\
 & 2.3 & Def. & Absurd & 2 & \texttt{PASL} & [11] \\
 & P1 & Pr. & Fit & 2 & \texttt{OD} & --- \\
 & 2.4 & Def. & Three responses & 2 & \texttt{MWPAIS} & [2], [4], [5], [11] \\
 & 3.1 & Lem. & Ineliminable seriousness & 2 & \texttt{PAI} & --- \\
 & 3.2 & Lem. & Ineliminable detachment & 3 & \texttt{AIS} & [11] \\
 & 3.3 & Thm. & \emph{The life of any reflective agent capable\ldots} & 3 & \texttt{AISL} & [11] \\
 & 3.4 & Prop. & Despair does not fit & 3 & \texttt{S} & [11] \\
 & 3.5 & Prop. & Revolt does not fit & 3 & \texttt{S} & [11] \\
 & 3.6 & Thm. & Irony fits & 4 & \texttt{SL} & [11] \\
 & 3.7 & Rem. & \emph{Camus's Sisyphus is happy because he is\ldots} & 4 & \texttt{} & --- \\
 & 1 & Obj. & Evasion & 4 & \texttt{S} & [11] \\
 & 2 & Obj. & Objective worth dissolves the absurd & 4 & \texttt{EWOPSL} & [3], [11] \\
 & 3 & Obj. & Irony corrodes engagement & 4 & \texttt{MEW} & [2] \\
 & 4 & Obj. & Hume's remedy & 5 & \texttt{S} & [11] \\
EXI-04-b & 1.1 & Def. & \emph{An activity is futile if it is\ldots} & 1 & \texttt{WP} & --- \\
 & 1.2 & Prop. & \emph{Absurdity and futility are independent. A life\ldots} & 1 & \texttt{OSL} & [11] \\
 & 2.1 & Prop. & \emph{The defensible content of Camus's conclusion is\ldots} & 2 & \texttt{EP} & --- \\
EXI-05-a & 2.1 & Def. & Fittingness & 1 & \texttt{} & --- \\
 & P1 & Pr. & Fear & 1 & \texttt{W} & --- \\
 & P2 & Pr. & Regret & 1 & \texttt{W} & --- \\
 & 3.1 & Lem. & \emph{One's own death as such is certain.\ldots} & 2 & \texttt{} & --- \\
 & 3.2 & Thm. & \emph{Fear of death as such is not\ldots} & 2 & \texttt{} & --- \\
 & 3.3 & Thm. & \emph{Regret at the prospect of one's death\ldots} & 2 & \texttt{W} & --- \\
 & 1 & Obj. & Premature death & 2 & \texttt{} & --- \\
 & 2 & Obj. & Dread & 2 & \texttt{} & --- \\
 & 3 & Obj. & Asymmetry & 3 & \texttt{} & --- \\
EXI-06-a & 2.1 & Def. & Responsiveness & 1 & \texttt{} & --- \\
 & 2.2 & Def. & Care, hostility, indifference & 1 & \texttt{ES} & [11] \\
 & 2.3 & Ass. & Physical closure & 2 & \texttt{S} & [11] \\
 & 3.1 & Thm. & \emph{The universe is indifferent to us, and\ldots} & 2 & \texttt{ES} & [11] \\
 & 3.2 & Cor. & \emph{The description of the universe as hostile,\ldots} & 2 & \texttt{ES} & [11] \\
 & 1 & Obj. & Harm outweighs benefit & 2 & \texttt{EWSL} & [11] \\
 & 2 & Obj. & Anthropomorphism cuts both ways & 3 & \texttt{OIS} & [11] \\
 & 3 & Obj. & Theism & 3 & \texttt{ES} & [11] \\
EXI-06-b & 2.1 & Prop. & \emph{That the universe is indifferent to us\ldots} & 1 & \texttt{MSL} & [1], [11] \\
 & 2.2 & Prop. & \emph{Cosmic indifference gives no ground for despair\ldots} & 2 & \texttt{S} & [11] \\
EXI-07-a & 2.1 & Def. & Final end & 2 & \texttt{PI} & --- \\
 & 2.2 & Def. & Direct aim & 2 & \texttt{PI} & --- \\
 & 2.3 & Def. & By-product & 2 & \texttt{WPID} & [2.3], [4] \\
 & P1 & Pr. & Engagement & 2 & \texttt{EWP} & --- \\
 & P2 & Pr. & Object-directedness & 2 & \texttt{EWOI} & [3], [4] \\
 & 3.1 & Lem. & \emph{Happiness is a by-product of engagement in\ldots} & 2 & \texttt{EWP} & --- \\
 & 3.2 & Lem. & \emph{An agent who aims directly at happiness\ldots} & 3 & \texttt{WPI} & [4] \\
 & 3.3 & Thm. & \emph{Aiming directly at happiness is self-defeating: an\ldots} & 3 & \texttt{EPI} & --- \\
 & 3.4 & Cor. & \emph{Happiness is not the goal of life,\ldots} & 3 & \texttt{PL} & --- \\
 & 1 & Obj. & The sophisticated hedonist & 3 & \texttt{P} & --- \\
 & 2 & Obj. & Aristotle & 4 & \texttt{P} & --- \\
 & 3 & Obj. & Some happiness is aimed at & 4 & \texttt{P} & --- \\
 & 4 & Obj. & Engagement is only causally required & 4 & \texttt{E} & --- \\
EXI-07-b & 1.1 & Prop. & \emph{The policy of aiming directly at happiness\ldots} & 1 & \texttt{P} & --- \\
 & 2.1 & Prop. & \emph{Either the sophisticated hedonist values the objects\ldots} & 2 & \texttt{EWP} & --- \\
EXI-08-a & 2.1 & Def. & Achievement & 2 & \texttt{I} & --- \\
 & 2.2 & Def. & Orientation & 2 & \texttt{I} & --- \\
 & P1 & Pr. & Public sources & 2 & \texttt{MI} & --- \\
 & P2 & Pr. & Non-coincidence & 2 & \texttt{I} & --- \\
 & 3.1 & Lem. & \emph{Authenticity in the achievement sense is not\ldots} & 2 & \texttt{} & --- \\
 & 3.2 & Lem. & \emph{Authenticity in the orientation sense is consistent\ldots} & 2 & \texttt{} & --- \\
 & 3.3 & Thm. & \emph{Authenticity is possible as an orientation, though\ldots} & 3 & \texttt{} & --- \\
 & 1 & Obj. & Deflation & 3 & \texttt{} & --- \\
 & 2 & Obj. & Social origin & 3 & \texttt{} & --- \\
 & 3 & Obj. & Self-deception & 3 & \texttt{I} & --- \\
EXI-09-a & 2.1 & Def. & Defining project & 1 & \texttt{P} & --- \\
 & 2.2 & Def. & Precedence & 2 & \texttt{P} & --- \\
 & P1 & Pr. & Artefact kinds & 2 & \texttt{} & --- \\
 & P2 & Pr. & Projects of agents & 2 & \texttt{PIL} & --- \\
 & 3.1 & Lem. & \emph{For every artefact, essence precedes existence.\ldots} & 2 & \texttt{} & --- \\
 & 3.2 & Lem. & \emph{For every human being, existence precedes essence.\ldots} & 2 & \texttt{} & --- \\
 & 3.3 & Thm. & \emph{Sartre's distinction holds: for human beings existence\ldots} & 3 & \texttt{} & --- \\
 & 3.4 & Cor. & \emph{Theorem is compatible with the thesis that\ldots} & 3 & \texttt{} & --- \\
 & 1 & Obj. & Theism & 3 & \texttt{} & --- \\
 & 2 & Obj. & Human nature & 3 & \texttt{WP} & --- \\
 & 3 & Obj. & Social formation & 3 & \texttt{} & --- \\
 & 4 & Obj. & Repurposing & 4 & \texttt{} & --- \\
EXI-09-b & 2.1 & Prop. & Two notions of essence & 1 & \texttt{P} & --- \\
 & 2.2 & Prop. & The accidental knife & 2 & \texttt{} & --- \\
 & 2.3 & Prop. & Defect & 2 & \texttt{L} & --- \\
EXI-10-a & 2.1 & Def. & Boredom & 2 & \texttt{EP} & --- \\
 & 2.2 & Def. & Mismatch & 2 & \texttt{EWPAI} & [4], [5] \\
 & 2.3 & Def. & Function & 2 & \texttt{} & --- \\
 & P1 & Pr. & Emotional content & 2 & \texttt{EA} & --- \\
 & 2.4 & Ass. & Value of engagement & 2 & \texttt{EWP} & --- \\
 & 3.1 & Lem. & \emph{Boredom represents a mismatch, and is correct\ldots} & 2 & \texttt{} & --- \\
 & 3.2 & Lem. & \emph{Boredom motivates the search for an alternative\ldots} & 2 & \texttt{P} & --- \\
 & 3.3 & Thm. & \emph{Boredom signals a mismatch between capacities and\ldots} & 3 & \texttt{PI} & --- \\
 & 3.4 & Cor. & \emph{Chronic boredom is either a malfunction, when\ldots} & 3 & \texttt{PI} & --- \\
 & 1 & Obj. & Profound boredom & 3 & \texttt{PS} & [11] \\
 & 2 & Obj. & Boredom harms & 3 & \texttt{} & --- \\
 & 3 & Obj. & Function is not value & 4 & \texttt{W} & --- \\
EXI-11-a & 2.1 & Def. & Meaning & 2 & \texttt{MEWDL} & [1], [2], [2.3], [10] \\
 & 2.2 & Def. & Cosmic purpose & 2 & \texttt{SL} & [11] \\
 & 2.3 & Def. & Purpose theory & 2 & \texttt{MSL} & [1], [11] \\
 & P1 & Pr. & Worth is not conferred by assignment & 2 & \texttt{WPI} & [4] \\
 & P2 & Pr. & Degree & 2 & \texttt{MEWDL} & [1], [2], [2.3], [10] \\
 & 3.1 & Lem. & Purpose is not sufficient & 2 & \texttt{MSL} & [1], [11] \\
 & 3.2 & Lem. & Worth without purpose & 3 & \texttt{WOPS} & [3], [11] \\
 & 3.3 & Thm. & \emph{Meaning does not require a cosmic purpose.\ldots} & 3 & \texttt{MSL} & [1], [11] \\
 & 3.4 & Cor. & The naturalist's life & 3 & \texttt{ML} & [1] \\
 & 4.1 & Prop. & Theism adds & 3 & \texttt{MAL} & [1], [5] \\
 & 4.2 & Rem. & \emph{Proposition is conditional on the framework. If\ldots} & 4 & \texttt{MEWAL} & [1], [2], [4], [5] \\
 & 1 & Obj. & Everything comes to nothing & 4 & \texttt{PS} & [11] \\
 & 2 & Obj. & The regress of linkage & 4 & \texttt{ML} & [1] \\
 & 3 & Obj. & The framework of hope & 4 & \texttt{MEWP} & [2] \\
 & 4 & Obj. & Purpose is what meaning means & 5 & \texttt{ML} & [1] \\
 & 5 & Obj. & Instability & 5 & \texttt{ML} & [1] \\
EXI-11-b & 1.1 & Prop. & \emph{The meaning that a theistic framework would\ldots} & 1 & \texttt{MAD} & [2.3], [5] \\
EXI-12-a & 2.1 & Def. & Balance & 1 & \texttt{WL} & --- \\
 & P1 & Pr. & Impersonal comparison & 2 & \texttt{WL} & --- \\
 & P2 & Pr. & Non-comparative benefit & 2 & \texttt{WL} & --- \\
 & 3.1 & Thm. & \emph{A life containing more good than bad\ldots} & 2 & \texttt{WL} & --- \\
 & 3.2 & Prop. & \emph{Benatar's quality-of-life argument does not reinstate the\ldots} & 2 & \texttt{WL} & --- \\
 & 1 & Obj. & Bad for no one & 2 & \texttt{W} & --- \\
 & 2 & Obj. & Testimony of the suffering & 3 & \texttt{W} & --- \\
 & 3 & Obj. & Risk & 3 & \texttt{WL} & --- \\
EXI-12-b & P1 & Pr. & Complaints & 1 & \texttt{W} & --- \\
 & 2.1 & Prop. & \emph{Principle , together with Corollary 3.4 of\ldots} & 1 & \texttt{} & --- \\
EXI-13-a & 2.1 & Def. & Reasonable wish & 2 & \texttt{} & --- \\
 & P1 & Pr. & Surviving interests & 2 & \texttt{WL} & --- \\
 & P2 & Pr. & Reliance & 2 & \texttt{} & --- \\
 & 2.2 & Ass. & Practice & 2 & \texttt{L} & --- \\
 & 2.3 & Ass. & Weight & 2 & \texttt{} & --- \\
 & 3.1 & Lem. & Subject & 2 & \texttt{} & --- \\
 & 3.2 & Lem. & Content & 3 & \texttt{} & --- \\
 & 3.3 & Lem. & Weakening & 3 & \texttt{} & --- \\
 & 3.4 & Thm. & Duties to the dead & 3 & \texttt{} & --- \\
 & 1 & Obj. & No subject & 4 & \texttt{} & --- \\
 & 2 & Obj. & Persistence & 4 & \texttt{} & --- \\
 & 3 & Obj. & Only the living & 4 & \texttt{W} & --- \\
 & 4 & Obj. & Deathbed promises & 4 & \texttt{} & --- \\
EXI-13-b & 1.1 & Def. & \emph{A posthumous event harms a person if\ldots} & 1 & \texttt{} & --- \\
 & 1.2 & Prop. & \emph{Posthumous harm and posthumous wrong come apart\ldots} & 1 & \texttt{} & --- \\
 & 2.1 & Prop. & \emph{The dead, considered as what remains after\ldots} & 2 & \texttt{} & --- \\
EXI-14-a & P1 & Pr. & Emotions as evaluations & 1 & \texttt{} & --- \\
 & 2.1 & Def. & Nostalgia & 2 & \texttt{WA} & [5] \\
 & 2.2 & Def. & Rational nostalgia & 2 & \texttt{WA} & [5] \\
 & 3.1 & Lem. & \emph{Nostalgia is fitting when it tracks a\ldots} & 2 & \texttt{WA} & [5] \\
 & 3.2 & Thm. & \emph{Nostalgia is rational in moderation: an episode\ldots} & 2 & \texttt{WA} & [5] \\
 & 1 & Obj. & The past is gone & 2 & \texttt{} & --- \\
 & 2 & Obj. & Selective memory & 3 & \texttt{} & --- \\
 & 3 & Obj. & Uselessness & 3 & \texttt{} & --- \\
EXI-15-a & 2.1 & Def. & Hope & 1 & \texttt{I} & --- \\
 & P1 & Pr. & Credence is evidential & 2 & \texttt{I} & --- \\
 & P2 & Pr. & The stance is controlled & 2 & \texttt{I} & --- \\
 & 2.2 & Def. & Virtue & 2 & \texttt{AD} & --- \\
 & 3.1 & Def. & Proportioned hope & 2 & \texttt{WI} & --- \\
 & 3.2 & Lem. & \emph{Hope whose stance is insensitive to probability\ldots} & 2 & \texttt{} & --- \\
 & 3.3 & Lem. & \emph{Hope whose stance is fixed by probability\ldots} & 2 & \texttt{} & --- \\
 & 3.4 & Thm. & \emph{Hope is a virtue when it is\ldots} & 2 & \texttt{} & --- \\
 & 1 & Obj. & Stoic & 3 & \texttt{} & --- \\
 & 2 & Obj. & Self-deception & 3 & \texttt{} & --- \\
 & 3 & Obj. & Hope against hope & 3 & \texttt{} & --- \\
EXI-15-b & 2.1 & Prop. & Near certainty & 1 & \texttt{I} & --- \\
 & 2.2 & Prop. & Remote possibility & 2 & \texttt{WI} & --- \\
 & 3.1 & Prop. & \emph{When the agent's effort is among the\ldots} & 2 & \texttt{I} & --- \\
EXI-16-a & 2.1 & Def. & Reasons for love & 2 & \texttt{} & --- \\
 & 2.2 & Def. & Qualitative and relational reasons & 2 & \texttt{} & --- \\
 & 2.3 & Def. & Qualitative similarity & 2 & \texttt{D} & --- \\
 & 2.4 & Def. & Transferability & 2 & \texttt{} & --- \\
 & P1 & Pr. & Answerability & 2 & \texttt{O} & --- \\
 & P2 & Pr. & Generality & 2 & \texttt{OD} & --- \\
 & P3 & Pr. & Particularity of history & 2 & \texttt{} & --- \\
 & 3.1 & Lem. & \emph{Love responds to reasons.\ldots} & 2 & \texttt{} & --- \\
 & 3.2 & Lem. & \emph{If all reasons for 's love of\ldots} & 3 & \texttt{} & --- \\
 & 3.3 & Lem. & \emph{If the reasons for 's love of\ldots} & 3 & \texttt{} & --- \\
 & 3.4 & Lem. & \emph{Every established love answers to relational reasons.\ldots} & 3 & \texttt{} & --- \\
 & 3.5 & Thm. & \emph{Love responds to reasons, but these reasons\ldots} & 3 & \texttt{} & --- \\
 & 1 & Obj. & Wrong object & 4 & \texttt{W} & --- \\
 & 2 & Obj. & Duplication & 4 & \texttt{} & --- \\
 & 3 & Obj. & Arbitrariness & 4 & \texttt{W} & --- \\
 & 4 & Obj. & Universal value & 5 & \texttt{WA} & [5] \\
EXI-16-b & 1.1 & Prop. & Trading up & 1 & \texttt{D} & --- \\
 & 1.2 & Prop. & Duplication & 1 & \texttt{} & --- \\
 & 1.3 & Prop. & After loss & 2 & \texttt{} & --- \\
EXI-17-a & 2.1 & Def. & Active engagement & 2 & \texttt{EPI} & --- \\
 & 2.2 & Def. & Worthwhile & 2 & \texttt{EWOPI} & [3], [4] \\
 & 2.3 & Def. & Achievement & 2 & \texttt{PI} & --- \\
 & P1 & Pr. & Fitting fulfilment & 2 & \texttt{MEWDL} & [1], [2], [2.3], [10] \\
 & P2 & Pr. & Contact with reality & 2 & \texttt{EPI} & --- \\
 & 2.4 & Ass. & Reflective judgment & 2 & \texttt{MWL} & [1] \\
 & 3.1 & Lem. & Failure is compatible with meaning & 3 & \texttt{ML} & [1] \\
 & 3.2 & Lem. & Meaning without product & 3 & \texttt{MEDL} & [1], [2], [2.3], [10] \\
 & 3.3 & Thm. & \emph{A meaningful life requires active engagement with\ldots} & 3 & \texttt{MEWL} & [1], [2], [10] \\
 & 3.4 & Cor. & Achievement is neither necessary nor sufficient & 3 & \texttt{MEWP} & [2] \\
 & 3.5 & Rem. & Degree & 3 & \texttt{MEWDL} & [1], [2], [2.3], [10] \\
 & 3.6 & Cor. & Availability & 4 & \texttt{MA} & [5] \\
 & 1 & Obj. & Everything is an achievement & 4 & \texttt{MPL} & [1] \\
 & 2 & Obj. & Making a difference & 4 & \texttt{MWL} & [1] \\
 & 3 & Obj. & The experience machine & 4 & \texttt{MWL} & [1] \\
 & 4 & Obj. & Sisyphus & 4 & \texttt{MEL} & [1], [10] \\
EXI-18-a & 2.1 & Def. & Psychological necessity & 1 & \texttt{I} & --- \\
 & 2.2 & Prop. & Partial necessity & 2 & \texttt{I} & --- \\
 & 2.3 & Thm. & \emph{That an illusion is psychologically necessary does\ldots} & 2 & \texttt{} & --- \\
 & 2.4 & Cor. & No license & 2 & \texttt{} & --- \\
 & 2.5 & Rem. & Substitutes & 2 & \texttt{OPD} & [3] \\
 & 1 & Obj. & Pragmatic justification & 2 & \texttt{I} & --- \\
 & 2 & Obj. & Sadder but wiser & 3 & \texttt{} & --- \\
 & 3 & Obj. & Meaning needs illusion & 3 & \texttt{MPS} & [11] \\
EXI-18-b & 1.1 & Prop. & \emph{(i) Whether a given positive illusion is\ldots} & 1 & \texttt{I} & --- \\
EXI-19-a & 2.1 & Def. & Meaning in itself & 1 & \texttt{MWOL} & [1], [3] \\
 & 2.2 & Def. & Shape & 2 & \texttt{WL} & --- \\
 & P1 & Pr. & Shape-dependence & 2 & \texttt{WL} & --- \\
 & 2.3 & Ass. & Bounds & 2 & \texttt{} & --- \\
 & 3.1 & Prop. & \emph{Death has no meaning in itself.\ldots} & 2 & \texttt{M} & --- \\
 & 3.2 & Thm. & \emph{Death gives a life a shape by\ldots} & 2 & \texttt{L} & --- \\
 & 3.3 & Cor. & \emph{Theorem is consistent with the claim of\ldots} & 2 & \texttt{W} & --- \\
 & 1 & Obj. & Other limits & 3 & \texttt{PL} & --- \\
 & 2 & Obj. & Death as good & 3 & \texttt{W} & --- \\
 & 3 & Obj. & Ruined shapes & 3 & \texttt{L} & --- \\
EXI-19-b & 1.1 & Prop. & Not a source & 1 & \texttt{MEWL} & [1], [2], [10] \\
 & 1.2 & Prop. & Whole-life meaning & 1 & \texttt{ML} & [1] \\
 & 1.3 & Prop. & Persistence of engagement & 1 & \texttt{MEL} & [1], [10] \\
\end{longtable}
}

\section{Testimonia}\label{app:test}

\subsection*{Preface}

The testimonia are every sentence in the other manuscripts of the discipline, outside the family EXI-01 itself, that cites EXI-01 by identifier: \cnTestimonia{} sentences from \cnTestimoniaMss{} manuscripts. They are extracted by the program of Appendix~\ref{app:conc} from the same text, with cross-references and citations resolved, and are given in the order of the manuscripts. In the commentary tradition, testimonia record what later texts say an earlier text said; here they record how the discipline understands EXI-01, and they are evidence for the reading of its lemmata.

Four testimonia are discussed in the Commentary. The sentence of EXI-04-a, Section~4, that quotes the last clause of Remark~5 is the only testimonium that reports a sentence of EXI-01-a rather than its conclusion, and the comment on \Lm{11.4} shows that its reading agrees with Proposition~\ref{p:vind}. The sentence of the abstract of EXI-06-b that describes meaning as engagement with what has worth \enquote{independently of anyone's attitudes} strengthens Assumption~2 beyond its text, and the comment on \Lm{3.4} shows that EXI-06-b's argument does not need the strengthening. The sentence of EXI-09-b that a life can be less meaningful \enquote{in the sense of EXI-01} than it might have been is the discipline's clearest use of the degree of Definition~1 and of the left-to-right direction that Remark~5 also uses (\Lm{10}). The sentences of EXI-11-a and EXI-20-a report the availability claim in the words of the family conclusion, \enquote{available to everyone}, and so confirm the reading of \enquote{everyone} at \Lm{4.2}.

{\footnotesize% Generated by gen/generate.py; do not edit.
\begin{longtable}{@{}l l >{\raggedright\arraybackslash}p{0.78\textwidth}@{}}
\toprule Ms & \S & Sentence \\ \midrule \endhead
\bottomrule \endfoot
EXI-03-a & 1 & It identifies the categorical desires that give life its point with desires for projects the person takes to be worthwhile, in line with the account of meaning in EXI-01; and it replaces Williams's appeal to identity of character with the account of what matters in survival defended in PID-06. \\[2pt]
EXI-03-a & 2 & EXI-01 argues that a life is meaningful to the extent that its subject engages with projects of objective worth, and EXI-17 that such engagement, not achievement, is what meaning requires. \\[2pt]
EXI-04-a & 3 & EXI-01 argues that life is meaningful to the extent that one engages with projects of objective worth, and we accept that conclusion. \\[2pt]
EXI-04-a & 4 & If some projects are worthwhile independently of our attitudes, as EXI-01 holds, then the seriousness with which we pursue them is justified, and the life of someone engaged in them is not absurd. \\[2pt]
EXI-04-a & 4 & This is the sense in which, as the remark at the end of EXI-01-a puts it, the assumption of objective worth denies the premise of the sceptical challenge without refuting it. \\[2pt]
EXI-04-a & 5 & The argument does not show that any particular life is meaningful or worthwhile; that question is answered by EXI-01 and EXI-17, and absurdity, as defined here, is shared by meaningful and meaningless lives alike. \\[2pt]
EXI-04-b & 3 & Our ends include worthwhile objects, and engagement with them confers meaning (EXI-01). \\[2pt]
EXI-06-b & Abs & Meaning, on the account of EXI-01, consists in engagement with what has worth independently of anyone's attitudes, and the universe's want of attitudes toward us takes nothing from that worth. \\[2pt]
EXI-06-b & 2 & On the account of EXI-01, a life is meaningful to the extent that its subject engages with projects of objective worth, whose value does not depend on the attitudes of those engaged in them (Definition 1 and Assumption 2 of EXI-01-a). \\[2pt]
EXI-06-b & 3 & The present point is conditional: if meaning consists in engagement with what has objective worth, as EXI-01 holds, then it does not require any being's care, the universe's or another's. \\[2pt]
EXI-06-b & 4 & It leaves the question of meaning where EXI-01 places it: in what we engage with, and how. \\[2pt]
EXI-07-a & 3 & It fits also the account of meaning in EXI-01 and EXI-17, on which meaning comes from active engagement with what is worthwhile. \\[2pt]
EXI-07-a & 5 & It also does not establish which other ends a person should have; that is the work of the account of the good (ETH-01) and of meaning (EXI-01, EXI-17). \\[2pt]
EXI-09-a & 5 & Principle P2 says that what a life is for is fixed by its agent's ends; it does not say that every end is as good as any other, and is consistent with EXI-01, on which a life is meaningful to the extent that its projects are of objective worth. \\[2pt]
EXI-09-b & 2 & A life can be imprudent, wicked or, in the sense of EXI-01, less meaningful than it might have been, because it engages little with what is worthwhile. \\[2pt]
EXI-11-a & 2 & EXI-01 adds that meaning so understood is available to everyone. \\[2pt]
EXI-17-a & 3 & Corollary 3.6 supports, without by itself establishing, the conclusion of EXI-01 that meaning in this sense is available to everyone. \\[2pt]
EXI-19-b & 1 & On the account of EXI-01 and EXI-17, a life is meaningful to the extent that its subject actively engages with what is worthwhile. \\[2pt]
EXI-19-b & 1 & So the limit adds neither an object of worth nor worth to an object, and by the account of EXI-01 it is not a source of meaning. \\[2pt]
EXI-20-a & 1 & EXI-01 adds that meaning in this sense is available to everyone. \\[2pt]
\end{longtable}
}


\begin{thebibliography}{99}\small
\bibitem[Anscombe(1957)]{Anscombe1957} Anscombe, G. E. M. (1957). \emph{Intention}. Basil Blackwell.
\bibitem[Benatar(2017)]{Benatar2017} Benatar, D. (2017). \emph{The Human Predicament: A Candid Guide to Life's Biggest Questions}. Oxford University Press.
\bibitem[Camus(1955)]{Camus1955} Camus, A. (1955). \emph{The Myth of Sisyphus and Other Essays}, trans. J. O'Brien. Alfred A. Knopf. (\emph{Le Mythe de Sisyphe}, Gallimard, 1942.)
\bibitem[Cottingham(2003)]{Cottingham2003} Cottingham, J. (2003). \emph{On the Meaning of Life}. Routledge.
\bibitem[de Moura and Ullrich(2021)]{deMouraUllrich2021} de Moura, L. and Ullrich, S. (2021). The Lean 4 theorem prover and programming language. In \emph{Automated Deduction -- CADE 28}, Lecture Notes in Computer Science 12699, 625--635. Springer.
\bibitem[Frankfurt(1982)]{Frankfurt1982} Frankfurt, H. G. (1982). The importance of what we care about. \emph{Synthese}, 53(2), 257--272.
\bibitem[Frankfurt(2004)]{Frankfurt2004} Frankfurt, H. G. (2004). \emph{The Reasons of Love}. Princeton University Press.
\bibitem[Frankl(1962)]{Frankl1962} Frankl, V. E. (1962). \emph{Man's Search for Meaning: An Introduction to Logotherapy}. Beacon Press.
\bibitem[Hare(1972)]{Hare1972} Hare, R. M. (1972). Nothing matters. In \emph{Applications of Moral Philosophy}. Macmillan.
\bibitem[Kahane(2014)]{Kahane2014} Kahane, G. (2014). Our cosmic insignificance. \emph{No\^us}, 48(4), 745--772.
\bibitem[Kauppinen(2012)]{Kauppinen2012} Kauppinen, A. (2012). Meaningfulness and time. \emph{Philosophy and Phenomenological Research}, 84(2), 345--377.
\bibitem[Landau(2017)]{Landau2017} Landau, I. (2017). \emph{Finding Meaning in an Imperfect World}. Oxford University Press.
\bibitem[MacIntyre(1981)]{MacIntyre1981} MacIntyre, A. (1981). \emph{After Virtue}. University of Notre Dame Press.
\bibitem[Mackie(1977)]{Mackie1977} Mackie, J. L. (1977). \emph{Ethics: Inventing Right and Wrong}. Penguin.
\bibitem[May(2015)]{May2015} May, T. (2015). \emph{A Significant Life: Human Meaning in a Silent Universe}. University of Chicago Press.
\bibitem[Metz(2001)]{Metz2001} Metz, T. (2001). The concept of a meaningful life. \emph{American Philosophical Quarterly}, 38(2), 137--153.
\bibitem[Metz(2013)]{Metz2013} Metz, T. (2013). \emph{Meaning in Life: An Analytic Study}. Oxford University Press.
\bibitem[Nagel(1971)]{Nagel1971} Nagel, T. (1971). The absurd. \emph{The Journal of Philosophy}, 68(20), 716--727.
\bibitem[Nagel(1986)]{Nagel1986} Nagel, T. (1986). \emph{The View from Nowhere}. Oxford University Press.
\bibitem[Nozick(1974)]{Nozick1974} Nozick, R. (1974). \emph{Anarchy, State, and Utopia}. Basic Books.
\bibitem[Nozick(1981)]{Nozick1981} Nozick, R. (1981). \emph{Philosophical Explanations}. Harvard University Press.
\bibitem[Proclus(1970)]{Proclus1970} Proclus (1970). \emph{A Commentary on the First Book of Euclid's Elements}, trans. G. R. Morrow. Princeton University Press.
\bibitem[Rawls(1971)]{Rawls1971} Rawls, J. (1971). \emph{A Theory of Justice}. Harvard University Press.
\bibitem[Raz(1986)]{Raz1986} Raz, J. (1986). \emph{The Morality of Freedom}. Clarendon Press.
\bibitem[Scheffler(2013)]{Scheffler2013} Scheffler, S. (2013). \emph{Death and the Afterlife}. Oxford University Press.
\bibitem[Sidgwick(1907)]{Sidgwick1907} Sidgwick, H. (1907). \emph{The Methods of Ethics}, 7th ed. Macmillan.
\bibitem[Smuts(2013)]{Smuts2013} Smuts, A. (2013). The good cause account of the meaning of life. \emph{The Southern Journal of Philosophy}, 51(4), 536--562.
\bibitem[Street(2006)]{Street2006} Street, S. (2006). A Darwinian dilemma for realist theories of value. \emph{Philosophical Studies}, 127(1), 109--166.
\bibitem[Taylor(1970)]{Taylor1970} Taylor, R. (1970). \emph{Good and Evil: A New Direction}. Macmillan.
\bibitem[Tolstoy(1987)]{Tolstoy1987} Tolstoy, L. (1987). \emph{A Confession and Other Religious Writings}, trans. J. Kentish. Penguin. (\emph{Ispoved'}, written 1879--1882.)
\bibitem[Velleman(1991)]{Velleman1991} Velleman, J. D. (1991). Well-being and time. \emph{Pacific Philosophical Quarterly}, 72(1), 48--77.
\bibitem[Wiggins(1976)]{Wiggins1976} Wiggins, D. (1976). Truth, invention, and the meaning of life. \emph{Proceedings of the British Academy}, 62, 331--378.
\bibitem[Williams(1973)]{Williams1973} Williams, B. (1973). The Makropulos case: Reflections on the tedium of immortality. In \emph{Problems of the Self}, 82--100. Cambridge University Press.
\bibitem[Williams(1981a)]{Williams1981a} Williams, B. (1981a). Persons, character and morality. In \emph{Moral Luck}, 1--19. Cambridge University Press.
\bibitem[Williams(1981b)]{Williams1981b} Williams, B. (1981b). Moral luck. In \emph{Moral Luck}, 20--39. Cambridge University Press.
\bibitem[Williams(1982)]{Williams1982} Williams, B. (1982). The point of view of the universe: Sidgwick and the ambitions of ethics. \emph{The Cambridge Review}, 103, 183--191.
\bibitem[Wolf(1997)]{Wolf1997} Wolf, S. (1997). Happiness and meaning: Two aspects of the good life. \emph{Social Philosophy and Policy}, 14(1), 207--225.
\bibitem[Wolf(2010)]{Wolf2010} Wolf, S. (2010). \emph{Meaning in Life and Why It Matters}. Princeton University Press.
\end{thebibliography}

\label{LastPageVolI}
\end{document}
